% \iffalse ltxdoc klootch
% AASTeX.dtx: document class for AAS journal.
% Copyright(C) 2003 American Astronomical Society.
%
% Disclaimer:
%   This file is distributed WITHOUT ANY WARRANTY;
%   without even the implied warranty of
%   MERCHANTABILITY or FITNESS FOR A PARTICULAR PURPOSE.
%
% ReadMe:
%   For the documentation and more detailed instructions for 
%   installation, typeset this document with \LaTeX.
%   
% \fi
% \CheckSum{5961}
%% \CharacterTable
%%  {Upper-case    \A\B\C\D\E\F\G\H\I\J\K\L\M\N\O\P\Q\R\S\T\U\V\W\X\Y\Z
%%   Lower-case    \a\b\c\d\e\f\g\h\i\j\k\l\m\n\o\p\q\r\s\t\u\v\w\x\y\z
%%   Digits        \0\1\2\3\4\5\6\7\8\9
%%   Exclamation   \!     Double quote  \"     Hash (number) \#
%%   Dollar        \$     Percent       \%     Ampersand     \&
%%   Acute accent  \'     Left paren    \(     Right paren   \)
%%   Asterisk      \*     Plus          \+     Comma         \,
%%   Minus         \-     Point         \.     Solidus       \/
%%   Colon         \:     Semicolon     \;     Less than     \<
%%   Equals        \=     Greater than  \>     Question mark \?
%%   Commercial at \@     Left bracket  \[     Backslash     \\
%%   Right bracket \]     Circumflex    \^     Underscore    \_
%%   Grave accent  \`     Left brace    \{     Vertical bar  \|
%%   Right brace   \}     Tilde         \~}%
%
% \iffalse ltxdoc klootch
%%%  @LaTeX-file{
%%%     filename        = "aastex.dtx",
%%%     version         = "5.2",
%%%     date            = "2005/06/22",
%%%     time            = "16:23:00 GMT",
%%%     checksum        = "5964",
%%%     author          = "Arthur Ogawa (mailto:ogawa@teleport.com)",
%%%     revised by      = "SR Nova Private Ltd."
%%%     copyright       = "Copyright (C) 2003 American Astronomical Society
%%%
%%%                        This work may be distributed and/or modified under the 
%%%                        conditions of the LaTeX Project Public License, either version 1.3
%%%                        of this license or (at your option) any later version.
%%%                        The latest version of this license is in
%%%                        http://www.latex-project.org/lppl.txt
%%%                        and version 1.3 or later is part of all distributions of LaTeX
%%%                        version 2003/12/01 or later.
%%%
%%%                        This work has the LPPL maintenance status "maintained".
%%% 
%%%                        The Current Maintainer of this work is the American Astronomical Society.
%%%
%%%                        This work consists of all files listed in the document README. 
%%%
%%%     address         = "American Astronomical Society
%%%                        USA",
%%%     telephone       = "+1 ???",
%%%     FAX             = "",
%%%     email           = "aastex-help@aas.org",
%%%     codetable       = "ISO/ASCII",
%%%     keywords        = "latex, AAS, journal",
%%%     supported       = "yes",
%%%     abstract        = "formatter for AAS journal submissions",
%%%     docstring       = "The checksum field above generated by ltxdoc",
%%%  }
% \fi
%
% \iffalse ltxdoc klootch
% The following references the \file{Read.Me} file,
% which contains basic information about this document class.
% The contents of this file are generated when 
% you typeset the programmer's documentation.
% Search on "{filecontents*}{Read.Me}" to locate it.
% \fi\input{Read.Me}%
%
% \subsection{Bill of Materials}
%
% Following is a list of the files in this distribution arranged
% according to provenance.
%
% \subsubsection{Primary Source}%
% One single file generates all macro files and documentation.
% \begin{verbatim}
%aasclass.dtx
% \end{verbatim}
%
% \subsubsection{Generated by \texttt{latex aasclass.dtx}}%
% Typesetting the source file under \LaTeX\ 
% generates the readme, the installer, and a configuration file for 
% this run.
% \begin{verbatim}
%Read.Me         aasclass.ins     ltxdoc.cfg
% \end{verbatim}
%
% \subsubsection{Generated by \texttt{tex aasclass.ins}}%
% Typesetting the installer generates
% the class file, style file, user guide, and sample documents,
% \begin{verbatim}
%aastex.cls    aastex.sty      aasclass.tex  aasguide.tex
%sample.tex   table.tex
%
% \end{verbatim}
%
% \subsubsection{Documentation}%
% The following are the online documentation:
% \begin{verbatim}
%aasclass.dvi   aasguide.dvi
%sample.dvi   table.dvi
% \end{verbatim}
%
% \subsubsection{Auxiliary}%
% The following are auxiliary files generated 
% in the course of running \LaTeX:
% \begin{verbatim}
%aasclass.aux    aasclass.toc    
%aasclass.idx    aasclass.ilg    aasclass.ind
%aasguide.aux    aasguide.toc
%aasguide.idx    aasguide.ilg    aasguide.ind
% \end{verbatim}
%
% \subsubsection{Font files}%
% In support of the |astro| font, certain files are provided.
%
% For those running Metafont, the primary source for the |astro| font is
% at |CTAN://tex-archive/fonts/astro| and consists of:
% \begin{verbatim}
%astrosym.doc
%astrosym.mf
%astrosym.uni
%astrosym.mac
%astrosym.com
%astrosym.tex
% \end{verbatim}
% with the first-mentioned file representing the documentation.
%
% Metrics for the font are provided as:
% \begin{verbatim}
%astrosym.tfm
%astrosym.pl
% \end{verbatim}
%
% For those with PostScript-capable previewers and printers,
% we provide the font in an ATM-compatible PostScript Type~1 format.
%
% For MacOS with ATM, 
% \begin{verbatim}
%Astrosym.sit.bin
%  metrics
%    astrosym metrics % For Textures
%  font
%    Astro
%    Astrosym.bmap
% \end{verbatim}
%
% For MacOS with ATM,
% \begin{verbatim}
%ASTROSYM.72
%ASTROSYM.FON
%ASTROSYM.INF
%ASTROSYM.PFB
%ASTROSYM.PFM
% \end{verbatim}
%
% For UNIX,
% \begin{verbatim}
%72.Astrosym
%Astrosym.afm
%Astrosym.pfa
% \end{verbatim}
%
%
% \section{The {\sc docstrip} modules}
%
% \subsection{Code common to all modules}%
% The following may look a bit klootchy, but we
% want to require only one place in this file
% where the version number is stated,
% and we also want to ensure that the version
% number is embedded into every generated file.
%
% Here is some code for the document style,
% which can be run under both \LaTeX2.09 or 
% \LaTeXe.
% If the latter, we terminate reading of the 
% style file and 
% simply load the document class instead.
% If the former, we provide dummy
% definitions for the \LaTeXe\ commands
% appearing in the style file.
%
%    \begin{macrocode}
%<style>\ifx\ProvidesFile\undefined
%<style> \def\next{}%
%<style>\else
%<style> \def\next{%
%<style>  \endinput
%<style>  \@obsoletefile{aastex.cls}{aastex.sty}\LoadClass{aastex}%
%<style> }%
%<style>\fi\next
%<style>\def\ProvidesFile#1[#2]{\message{File: #1 #2}}%
%    \end{macrocode}
%
% To start with, we declare that
% this file can only be used with \LaTeXe.
% An appropriate message is displayed if 
% a different \TeX{} format is used.
%    \begin{macrocode}
%<*doc|class|guide|sample|table>
\NeedsTeXFormat{LaTeX2e}[1995/12/01]%
%</doc|class|guide|sample|table>
%    \end{macrocode}
% As desired, the following modules all
% take common version information:
%    \begin{macrocode}
%<class>\ProvidesClass{aastex}%
%<style>\ProvidesFile{aastex.sty}%
%<guide>\ProvidesFile{aasguide.tex}%
%<symbols>\ProvidesFile{aassymbols.tex}%
%<sample>\ProvidesFile{sample.tex}%
%<table>\ProvidesFile{table1.tex}%
%<*doc>
\expandafter\ProvidesFile\expandafter{\jobname.dtx}%
%</doc>
%    \end{macrocode}
%
% \subsection{Version and Date}%
%
% The following line contains, for once and for all,
% the version and date information.
% By various means, this information is reproduced
% consistently in all generated files and in the 
% typeset documentation.
%    \begin{macrocode}
 [2005/06/22 5.2/AAS markup document class]%
%    \end{macrocode}
%
% \subsection{The driver module \texttt{doc}}
%
% Because the only uncommented-out lines of code at the beginning of
% this file constitute the \file{doc} module itself,
% we can simply typeset the |.dtx| file directly,
% and there is thus rarely any need to
% generate the ``doc'' {\sc docstrip} module.
% Module delimiters are nonetheless required so that
% this code does not find its way into the other modules.
% 
% The |\end{document}| command concludes the typesetting run.
%
%    \begin{macrocode}
%<*doc>
%    \end{macrocode}
%
% \subsection{The ltxdoc configuration file}
%
% This {\sc docstrip} module generates a local file \file{ltxdoc.cfg},
% which contains customizing commands for
% when the programmer documentation \file{aasclass.dtx} is typeset.
%
% In particular, we request the production of an index of macros.
%
% For more information on such proccessing options, 
% see the standard \LaTeX\ file \file{ltxdoc.dtx}.
%
%    \begin{macrocode}
\begin{filecontents}{ltxdoc.cfg}
% ltxdoc.cfg - local copy for aasclass directory
\AtBeginDocument{\CodelineIndex\EnableCrossrefs}
\AtEndDocument{\PrintIndex}
\end{filecontents}
%    \end{macrocode}
%
% \subsection{The Document Class Statement}
% The programmers documentation is formatted 
% with the |ltxdoc| class.
%    \begin{macrocode}
\documentclass[draft]{ltxdoc}
%    \end{macrocode}
%
% \subsubsection{Indexing directives}
%    We don't want everything to appear in the index
%    \begin{macrocode}
\DoNotIndex{\',\.,\@M,\@@input,\@Alph,\@alph,\@addtoreset,\@arabic}
\DoNotIndex{\@badmath,\@centercr,\@cite}
\DoNotIndex{\@dotsep,\@empty,\@float,\@gobble,\@gobbletwo,\@ignoretrue}
\DoNotIndex{\@input,\@ixpt,\@m,\@minus,\@mkboth}
\DoNotIndex{\@ne,\@nil,\@nomath,\@plus,\roman,\@set@topoint}
\DoNotIndex{\@tempboxa,\@tempcnta,\@tempdima,\@tempdimb}
\DoNotIndex{\@tempswafalse,\@tempswatrue,\@viipt,\@viiipt,\@vipt}
\DoNotIndex{\@vpt,\@warning,\@xiipt,\@xipt,\@xivpt,\@xpt,\@xviipt}
\DoNotIndex{\@xxpt,\@xxvpt,\\,\ ,\addpenalty,\addtolength,\addvspace}
\DoNotIndex{\advance,\ast,\begin,\begingroup,\bfseries,\bgroup,\box}
\DoNotIndex{\bullet}
\DoNotIndex{\cdot,\cite,\CodelineIndex,\cr,\day,\DeclareOption}
\DoNotIndex{\def,\DisableCrossrefs,\divide,\DocInput,\documentclass}
\DoNotIndex{\DoNotIndex,\egroup,\ifdim,\else,\fi,\em,\endtrivlist}
\DoNotIndex{\EnableCrossrefs,\end,\end@dblfloat,\end@float,\endgroup}
\DoNotIndex{\endlist,\everycr,\everypar,\ExecuteOptions,\expandafter}
\DoNotIndex{\fbox}
\DoNotIndex{\filedate,\filename,\fileversion,\fontsize,\framebox,\gdef}
\DoNotIndex{\global,\halign,\hangindent,\hbox,\hfil,\hfill,\hrule}
\DoNotIndex{\hsize,\hskip,\hspace,\hss,\if@tempswa,\ifcase,\or,\fi,\fi}
\DoNotIndex{\ifhmode,\ifvmode,\ifnum,\iftrue,\ifx,\fi,\fi,\fi,\fi,\fi}
\DoNotIndex{\input}
\DoNotIndex{\jobname,\kern,\leavevmode,\let,\leftmark}
\DoNotIndex{\list,\llap,\long,\m@ne,\m@th,\mark,\markboth,\markright}
\DoNotIndex{\month,\newcommand,\newcounter,\newenvironment}
\DoNotIndex{\NeedsTeXFormat,\newdimen}
\DoNotIndex{\newlength,\newpage,\nobreak,\noindent,\null,\number}
\DoNotIndex{\numberline,\OldMakeindex,\OnlyDescription,\p@}
\DoNotIndex{\pagestyle,\par,\paragraph,\paragraphmark,\parfillskip}
\DoNotIndex{\penalty,\PrintChanges,\PrintIndex,\ProcessOptions}
\DoNotIndex{\protect,\ProvidesClass,\raggedbottom,\raggedright}
\DoNotIndex{\refstepcounter,\relax,\renewcommand}
\DoNotIndex{\rightmargin,\rightmark,\rightskip,\rlap,\rmfamily}
\DoNotIndex{\secdef,\selectfont,\setbox,\setcounter,\setlength}
\DoNotIndex{\settowidth,\sfcode,\skip,\sloppy,\slshape,\space}
\DoNotIndex{\symbol,\the,\trivlist,\typeout,\tw@,\undefined,\uppercase}
\DoNotIndex{\usecounter,\usefont,\usepackage,\vfil,\vfill,\viiipt}
\DoNotIndex{\viipt,\vipt,\vskip,\vspace}
\DoNotIndex{\wd,\xiipt,\year,\z@}
\DoNotIndex{\next}
%    \end{macrocode}
%    We \emph{do} want an index, using linenumbers
%    \begin{macrocode}
\EnableCrossrefs
%    \end{macrocode}
%
% \subsubsection{Docstrip and info directives}
%    We use so many {\sc docstrip} modules that we set the
%    \texttt{StandardModuleDepth} counter to 1.
%    \begin{macrocode}
\setcounter{StandardModuleDepth}{1}
%    \end{macrocode}
%    The following command retrieves the date and version information
%    from this file.
%    \begin{macrocode}
\expandafter\GetFileInfo\expandafter{\jobname.dtx}%
%    \end{macrocode}
%    Direct \file{ltxdoc} to produce an index and cross references.
%    \begin{macrocode}
\CodelineIndex\EnableCrossrefs
%    \end{macrocode}
%
% \subsubsection{Macros}
%    Some commonly used abbreviations
%    \begin{macrocode}
\newcommand*{\Lopt}[1]{\textsf {#1}}%
%\newcommand*{\file}[1]{\texttt {#1}}%
\newcommand*{\Lcount}[1]{\textsl {\small#1}}%
\newcommand*{\pstyle}[1]{\textsl {#1}}%
\newcommand\aastex{AAS\TeX}%
\setlength\arraycolsep{0pt}%
\let\oarg\undefined
%    \end{macrocode}
% Modify |theindex| environment so that it produces a TOC entry
%    \begin{macrocode}
\makeatletter
\renewenvironment{theindex}
               {\if@twocolumn
                  \@restonecolfalse
                \else
                  \@restonecoltrue
                \fi
                \columnseprule \z@
                \columnsep 35\p@
                \twocolumn[\section{\indexname}]%
%               \@mkboth{\MakeUppercase\indexname}%
%                       {\MakeUppercase\indexname}%
                \thispagestyle{plain}\parindent\z@
                \parskip\z@ \@plus .3\p@\relax
                \let\item\@idxitem}
               {\if@restonecol\onecolumn\else\clearpage\fi}
\makeatother
%    \end{macrocode}
% Look out! We are about to reach deep.
%    \begin{macrocode}
\makeatletter
%    \end{macrocode}
% We paste in a bit of very useful code extracted from 
% \file{ltxguide.cls}, by Alan Jeffrey.
%    \begin{macrocode}
 % This code stolen from ltxguide.cls:
% Some hacks with verbatim... NB: this would be better done with the
% verbatim package, but this document has to run on any LaTeX
% installation.
 
\let\o@verbatim\verbatim
\def\verbatim{%
  \ifhmode\unskip\par\fi
% \nopagebreak              % Overridden by list penalty
  \ifx\@currsize\normalsize
     \small
  \fi
  \o@verbatim
}
 
% Here we extend the font-setting command to include making <> active
% (ie adjusting the input encoding).
\renewcommand \verbatim@font {%
  \normalfont \ttfamily
  \catcode`\<=\active
  \catcode`\>=\active
}
 
% Make |...| a synonym for \verb|...|.
\RequirePackage{shortvrb}
\MakeShortVerb{\|}

 
% Make active <...> produce italics surrounded by angle brackets
% (used in verbatim and \verb).
% << produces a less-than, and >> produces a greater-than.
 
\begingroup
  \catcode`\<=\active
  \catcode`\>=\active
  \gdef<{\@ifnextchar<\@lt\@meta}
  \gdef>{\@ifnextchar>\@gt\@gtr@err}
  \gdef\@meta#1>{\m{#1}}
  \gdef\@lt<{\char`\<}
  \gdef\@gt>{\char`\>}
\endgroup
\def\@gtr@err{%
   \ClassError{ltxguide}{%
      Isolated \protect>%
   }{%
      In this document class, \protect<...\protect>
      is used to indicate a parameter.\MessageBreak
      I've just found a \protect> on its own.
      Perhaps you meant to type \protect>\protect>?
   }%
}
\def\verbatim@nolig@list{\do\`\do\,\do\'\do\-}

% Various forms of argument:
 
\newcommand{\m}[1]{\mbox{$\langle$\it #1\/$\rangle$}}
\renewcommand{\arg}[1]{{\tt\string{}\m{#1}{\tt\string}}}
\newcommand{\oarg}[1]{{\tt[%]
}\m{#1}{\tt%[
]}}
\def\cmd#1{\cs{\expandafter\cmd@to@cs\string#1}}
\def\cmd@to@cs#1#2{\char\number`#2\relax}
\DeclareRobustCommand\cs[1]{\texttt{\char`\\#1}}
\newcommand*{\file}[1]{\texttt{#1}}%
 % end of code stolen from ltxguide.cls. Thanks, Alan.
%    \end{macrocode}
%
% Define sectioning command below |\subsubsection|
%    \begin{macrocode}
\newcommand\subsubsubsection{%
 \@startsection{subsubsection}{4}%
  {\z@}{-15\p@\@plus-5\p@\@minus-2\p@}%
  {5\p@}{\normalfont\normalsize\itshape}%
}%
%    \end{macrocode}
%
% Coax the version number into |filecontents|-generated files.
%    \begin{macrocode}
\def\endfilecontents{%
 \immediate\write\reserved@c{%
  \string\iffalse\space ltxdoc klootch^^J%
  This file has version number \fileversion, last revised \filedate.%
  \string\fi
 }%
 \immediate\closeout\reserved@c
 \def\T##1##2##3{%
  \ifx##1\@undefined\else
    \@latex@warning@no@line{##2 has been converted to Blank ##3e}%
  \fi
 }%
 \T\L{Form Feed}{Lin}%
 \T\I{Tab}{Spac}%
 \immediate\write\@unused{}%
}%
\expandafter\let\csname endfilecontents*\endcsname\endfilecontents
%    \end{macrocode}
%
% OK, an end to the hacking!
%    \begin{macrocode}
\makeatother
%    \end{macrocode}
%
% \subsubsection{The installer file}
%
% The installer \file{aasclass.ins} appears here.
% If you have retrieved the standard distribution of this class,
% the installer file is already on your filesystem.
% If you are bootstrapping,
% the first typesetting of the \file{.dtx} file
% will cause the installer to be generated.
%
% The following modules are used to direct
% {\sc docstrip} in generating the external files:
% \begin{center}
% \begin{tabular}{lll}
%   \textbf{Module}&\textbf{File}&\textbf{Description}\\
%   class    &\file{aastex.cls}  &document class\\
%   style    &\file{aastex.sty}  &document style\\
%   doc      &\file{aasclass.tex}&driver for programmer's documantation\\
%   guide    &\file{aasguide.tex}&user guide\\
%   symbols  &\file{aassymbols.tex}&AAS and AMS symbols crib sheet\\
%   sample  &\file{sample.tex} &sample\\
%   table &\file{table.tex}&complex table\\
% \end{tabular}
% \end{center}
%
%    \begin{macrocode}
\begin{filecontents}{aasclass.ins}
%% This file will generate documentation and runtime files
%% from aasclass.dtx when run through LaTeX or TeX.
\input docstrip
\preamble

This is a generated file;
altering it directly is inadvisable;
instead, modify the original source file.
%See the URL in the file Read.Me.

Copyright notice.

   These files are distributed
   WITHOUT ANY WARRANTY; without even the implied warranty of
   MERCHANTABILITY or FITNESS FOR A PARTICULAR PURPOSE.

\endpreamble
\keepsilent
\generate{%
 \file{aastex.cls}{\from{aasclass.dtx}{class}}%
 \file{aastex.sty}{\from{aasclass.dtx}{style}}%
 \file{aasclass.tex}{\from{aasclass.dtx}{doc}}%
 \file{aasguide.tex}{\from{aasclass.dtx}{guide}}%
 \file{aassymbols.tex}{\from{aasclass.dtx}{symbols}}%
 \file{sample.tex}{\from{aasclass.dtx}{sample}}%
 \file{table.tex}{\from{aasclass.dtx}{table}}%
}%
\ifToplevel{
\Msg{***********************************************************}
\Msg{*}
\Msg{* To finish the installation, please move the following}
\Msg{* files into a directory searched by TeX:}
\Msg{*}
\Msg{* \space\space aastex.cls}
\Msg{*}
\Msg{* To produce the user documentation,
       run aasguide.tex through LaTeX.}
\Msg{*}
\Msg{* To produce the AAS and AMS symbols crib sheet,
       run aassymbols.tex through LaTeX.}
\Msg{*}
\Msg{* To typeset the aastex sample document,
       run sample.tex and table.tex 
       through LaTeX.}
\Msg{*}
\Msg{* To produce the programmer documentation,
       run aasclass.dtx through LaTeX.}
\Msg{*}
\Msg{* Happy TeXing}
\Msg{***********************************************************}
}
\endbatchfile
\end{filecontents}
%    \end{macrocode}
% Note that, because all of the files generated by the installer
% are part of the standard distribution, it will
% be necessary to run the installer only when bootstrapping
% (or, of course, during development).
% Note, too, that it is rare to generate the |doc|
% module because it suffices to simply typeset the |.dtx| file itself.
%
% \subsubsection{The \file{Read.Me} File}
% As promised above, here is the contents of the 
% readme file. That file serves a double purpose,
% since it also constitutes the beginining of the
% programmer's documentation. What better thing, after 
% all, to have appear at the beginning of the 
% typeset documentation?
%
% A good discussion of how to write a ReadMe file can be found in
% Engst, Tonya, ``Writing a ReadMe File? Read This''
% \emph{MacTech} October 1998, p. 58.
%
% Note the appearance of the 
% |\StopEventually| command, which marks the 
% dividing line between the user documentation
% and the programmer documentation.
% The former simply summarize the installation
% instructions and refers the reader to 
% the user guide \file{awljavgd}.
%
% The usual user will not be asked to 
% do a full build, not to speak
% of the bootstrap. 
% Instructions for carrying these processes
% begin the programmer's manual.
%
%    \begin{macrocode}
\begin{filecontents*}{Read.Me}
\title{%
 The \aastex\ Document Class For \LaTeXe\\
 Programmer's Manual%
 \thanks{%
  This file has version number \fileversion, last revised \filedate.%
  % For version number and date, search on "\fileversion" in the .dtx file,
  % or see the end of the Read.Me file.
 }%
}%

\author{%
Chris Biemesderfer (\texttt{mailto:chris@seagoat.com}), v4\\
Arthur Ogawa (\texttt{mailto:ogawa@teleport.com}), v\fileversion\\
Copyright (C) 1998 AAS
}%
\maketitle

\aastex\ provides a markup system and formatter for authoring 
papers to be submitted electronically to a number of
astonomy journals.
This file embodies the document class,
the implementation, and the user documentation.

The distribution point for \aastex\ is
\file{ftp://ftp.aas.org/pubs/aastex50.zip}, which contains
fully unpacked, prebuilt runtime files and documentation.

To use this document class, you must have a working
\TeX\ installation equipped with \LaTeXe.
Certain features of \aastex\ (such as landscape tables)
require a PostScript printer.

To install, retrieve the distribution,
unpack it into a directory on the target computer, 
and move the file \file{aastex.cls}
into a location in your filesystem where it will be found by \LaTeX.

If you will be using the |astro| font in your \aastex\ documents,
put the file \file{Astrosym.tfm} where \TeX\ can find it
(Textures users will put \file{Astrosym metrics} into their \file{TeX fonts} folder).
Now either install the (ATM-compatible) PostScript Type~1 fonts
appropriate to your operating system and \TeX\ implementation,
or install the MetaFont source (both provided with \aastex).
In any case, you are responsible for properly installing the 
|astro| font on your system.

Then print out the user guide \file{aasguide.dvi} and 
the two sample documents, \file{sample.dvi} and \file{table.dvi}. 
If need be, print the programmer's manual \file{aasclass.dvi}.

\StopEventually{}\tableofcontents

\section{Processing Instructions}

The document class \file{aastex.cls} (plus other macro files),
the user documentation \file{aasguide.tex},
and the sample documents
are generated from this file, \file{aasclass.dtx},
via the {\sc docstrip} facility of \LaTeX.
The documentation that you are now reading is generated from
the same file by simply typesetting it with \LaTeX.

\subsection{Build Instructions for V5.0}

You may bootstrap \aastex\ solely from \file{aasclass.dtx}.
Prepare by installing \LaTeXe\ on your computer,
then carry out the following steps:
\begin{enumerate}

\item
Within an otherwise empty directory,
typeset \file{aasclass.dtx} with \LaTeX;
you will obtain the typeset documentation you are now reading,
along with the installer \file{aasclass.ins} and the file \file{Read.Me}.

Note that you will have to run MakeIndex (using the \LaTeX\ standard
\file{gind.ist}) and retypeset with \LaTeX\ as appropriate.

\item
Now typeset \file{aasclass.ins}, thereby generating
macro files \file{aastex.cls}, and \file{aastex.sty}, user documentation \file{aasguide.tex}, \file{aassymbols.tex}, 
the sample file \file{sample.tex} and the sample table
\file{table.tex}.

\item
Install \file{aastex.cls} and \file{aastex.sty}:
move them to a location
in your file system where they will be found by \LaTeX.

\item
Typeset \file{aasguide.tex} to obtain the user documentation
(you will have to run MakeIndex and retypeset as needed),

\item
Typeset the sample file \file{sample.tex}.
\end{enumerate}

For your reference,
there are six \TeX\ log files appended to this file
(locate them by searching on |preloaded format=|).
They result respectively from typesetting
\file{aasclass.dtx}, \file{aasclass.ins}, \file{aasguide.tex},
\file{aassymbols.tex}, 
\file{sample.tex}, \file{table.tex}.

\subsection{Build Instructions for V4.1}

When you bootstrap V5.0, you automatically create the files rquired for the 
V4.1 distribution, which consists of a number of the documents for V4.0
along with style files to patch it, plus revised sample documents.
The user guide for V5.0 is also the guide for 4.1; relevant differences
appear in an appendix.

The files for version~4.1 are located in a separate part of the distribution tree,
a directory called \file{./V4.1/}.

To build, simply copy the file
\file{aastex.sty} into the 4.1 directory.
Then copy
the sample files \file{sample.tex}, \file{sample2.tex}, and \file{samptbls.tex},
and the sample tables
\file{table1.tex}, \file{table2.tex}, \file{table3.tex}, and \file{complext.tex}
into the same directory.
Finally, copy 
\file{aasguide.dvi} and \file{aasguide.ps} into the same directory.

\end{filecontents*}
%    \end{macrocode}
%
% \subsubsection{The Body of the Programmers Documentation}
% Here is the document body, containing only a
% |\DocInput| directive---referring to this very file.
% This very cute self-reference is a common |ltxdoc| idiom.
%    \begin{macrocode}
\begin{document}%
\expandafter\DocInput\expandafter{\jobname.dtx}%
\setcounter{secnumdepth}{0}%
\PrintIndex
% ^^A\PrintChanges
\end{document}
%    \end{macrocode}
%
% End of the |doc| module:
%    \begin{macrocode}
%</doc>
%    \end{macrocode}
%
%
%
% \section{Common to Document Style and Class}%
%    \begin{macrocode}
%<*class|style>
%    \end{macrocode}
%
% \begin{macro}{\substitute@option}
%    \begin{macrocode}
\def\substitute@option#1#2{%
 \ClassWarning{aastex}{%
  Substyle #1 is deprecated in aastex.
  Using #2 instead (please fix your document).
 }\@nameuse{ds@#2}%
}%
%    \end{macrocode}
% \end{macro}
%
%    \begin{macrocode}
%</class|style>
%    \end{macrocode}
%
% \section{Document Style Module}%
%
% The document style is provided for compatibility sake.
% If the user's document has
% \begin{verbatim}
%\documentstyle[<options>]{aastex}%
% \end{verbatim}
% and it is run under \LaTeXe, then this style file will be called,
% but the document class will be loaded instead.
% A warning message is given admonishing the user to switch to
% using |\documentclass|.
%
% If the same document is run under \LaTeX 2.09, \aastex\ will run 
% version 4.1 (which is an updated version~4.0).
%
% Here begins the style file:
%    \begin{macrocode}
%<*style>
%    \end{macrocode}
%
% \subsection{Document Style Options}%
%
% \begin{macro}{\ClassWarning}
%    \begin{macrocode}
\def\ClassWarning#1#2{\@warning{Style #1: #2}}%
%    \end{macrocode}
% \end{macro}
%
% \begin{macro}{ds@manuscript}
% \begin{macro}{ds@preprint}
% \begin{macro}{ds@preprint2}
% The substyle options are defined here.
% As in the document class, exactly one substyle will be in effect;
% the default being the |manuscript| substyle.
%    \begin{macrocode}
\def\ds@manuscript{\def\aastex@substyle{aasms4}}%
\def\aastex@substyle{aasms4}%
\def\ds@preprint{\def\aastex@substyle{aaspp4}}%
\@namedef{ds@preprint2}{\def\aastex@substyle{aas2pp4}}%
%    \end{macrocode}
% \end{macro}
% \end{macro}
% \end{macro}
%
% 
% \begin{macro}{ds@aasms4}
% \begin{macro}{ds@aaspp4}
% \begin{macro}{ds@aas2pp4}
% \begin{macro}{ds@ajpt4}
% \begin{macro}{ds@apjpt4}
% For backward compatibility, we declare options corresponding to V4.0
%    \begin{macrocode}
\@namedef{ds@aasms4}{\substitute@option{aasms4}{manuscript}}%
\@namedef{ds@aaspp4}{\substitute@option{aaspp4}{preprint}}%
\@namedef{ds@aas2pp4}{\substitute@option{aas2pp4}{preprint2}}%
\@namedef{ds@aj_pt4}{\substitute@option{aj_pt4}{preprint}}%
\@namedef{ds@apjpt4}{\substitute@option{apjpt4}{preprint}}%
%    \end{macrocode}
% \end{macro}
% \end{macro}
% \end{macro}
% \end{macro}
% \end{macro}
%
% 
% \begin{macro}{ds@eqsecnum}
% \begin{macro}{ds@tighten}
% \begin{macro}{ds@flushrt}
% The remaining document style options are defined here.
% They have the same significance as those of the 
% document class.
%    \begin{macrocode}
\def\ds@eqsecnum{\AtEndOfClass{\eqsecnum}}%
\def\ds@tighten{\AtEndOfClass{\tightenlines}}%
\def\ds@flushrt{%
 \AtEndOfClass{%
  \@rightskip=\z@\rightskip\@rightskip % Flush right
  \pretolerance=100\relax              % Enable hyphenation
 }%
}%
%    \end{macrocode}
% \end{macro}
% \end{macro}
% \end{macro}
% 
% \subsection{\LaTeXe\ Equipment Needed}%
%
% \begin{macro}{\@testopt}
% The following is snitched from \LaTeX2e.
%    \begin{macrocode}
\def\@testopt#1#2{%
  \@ifnextchar[{#1}{#1[#2]}}
%    \end{macrocode}
% \end{macro}
%
% The remaining definitions merely provide for
% a bare emulation of \LaTeXe.
%
% \begin{macro}{\AtEndOfClass}
% \begin{macro}{\AtBeginDocument}
% 
%    \begin{macrocode}
\def\appdef#1#2{%
 \toks@\expandafter{#1}\toks\@ne{#2}\edef#1{\the\toks@\the\toks\@ne}%
}%
\def\AtEndOfClass{\appdef\@styleh@@k}%
\def\@styleh@@k{}%
\def\AtBeginDocument{\appdef\@begindocumenthook}%
\def\@begindocumenthook{}%
\appdef\document{\@begindocumenthook}%
\def\@latex@warning{\@warning}%
%    \end{macrocode}
% \end{macro}
% \end{macro}
%
% Some low-level macros.
%    \begin{macrocode}
\def\protected@edef{\edef}%
\def\@finalstrut#1{\strut}%
\let\color@begingroup\begingroup
\let\color@endgroup\endgroup
\def\G@refundefinedtrue{}%
\long\def\@firstofone#1{#1}%
%    \end{macrocode}
%
% Font-oriented macros.
%    \begin{macrocode}
\def\rmfamily{\rm}%
\def\itshape{\it}%
%    \end{macrocode}
%
% \subsection{Base Document Style and Packages}%
%
% Read in the base document style, the \LaTeX\ standard
% article style.
% It will execute the document style options.
%    \begin{macrocode}
\@@input article.sty\relax
%    \end{macrocode}
% 
% Add in the |natbib| package for \LaTeX2.09
%    \begin{macrocode}
\@@input natbib209.sty\relax
%    \end{macrocode}
% 
% Most of the remainder of the document style 
% is drawn from the code going into \file{aastex.cls}.
% 
% We first read in the kernel modifications,
% then read in the selected V4.0 substyle,
% insert our V4.1 overrides,
% then act on the document style options.
%    \begin{macrocode}
%</style>
%    \end{macrocode}
%
%
% \section{Document Class Module}
% The document class module comprises this and the next
% five sections.
%    \begin{macrocode}
%<*class>
\ClassInfo{aastex}{%
 Original \protect\LaTeX2.09 style 
 by Chris Biemesderfer (chris@seagoat.com).
 Adapted to \protect\LaTeXe\space
 by A. Ogawa (ogawa@teleport.com)%
}%
%    \end{macrocode}
%
% \subsection{Type Size}
% \begin{macro}{\aas@ptsize}
% The macro |\aas@ptsize| works very much like 
% the |\@ptsize| macro of \file{classes.dtx},
% but in this class.
%    \begin{macrocode}
\newif\iflong@abstract
\global\long@abstractfalse
\DeclareOption{10pt}{\newcommand\aas@ptsize{0}}%
\DeclareOption{11pt}{\newcommand\aas@ptsize{1}}%
\DeclareOption{12pt}{\newcommand\aas@ptsize{2}}%
%    \end{macrocode}
% \end{macro}
%
% \subsection{Landscape Orientation}
%    The option \Lopt{landscape} is disabled in this document class.
%    \begin{macrocode}
\DeclareOption{landscape}{\OptionNotUsed}%
%    \end{macrocode}
%
% \subsection{Equations Numbered by Section}
%
%    The option \Lopt{eqsecnum} invokes |\eqsecnum|, q.v.
%    \begin{macrocode}
\DeclareOption{eqsecnum}{\AtEndOfClass{\eqsecnum}}
%    \end{macrocode}
%
% \subsection{Longabstract}
%    The option \Lopt{longabstract} cause the long abstract text to break into 
% next page within preprint2. Using this option has no effect with other styles.
%
%    \begin{macrocode}
\DeclareOption{longabstract}{\global\long@abstracttrue}
% use when abstract exceed title page
%    \end{macrocode}
%
% \subsection{Set Text Justified Rather Than Ragged Right}
%
%    The option \Lopt{flushrt} causes the text to set 
% justified, hyphenated rather than ragged right, nohyphens
% as is normal in the |manuscript| substyle.
%
% Asserting this option presently has no effect in
% the other substyles because they make no changes 
% to the justification. If this situation changes,
% be careful to supply the appropriate hook as needed.
%
%    \begin{macrocode}
\DeclareOption{flushrt}{\let\opt@just\just@just}%
\let\opt@just\@empty
\def\just@just{%
 \@rightskip\z@skip
 \rightskip\@rightskip
 \pretolerance=100\relax
}%
%    \end{macrocode}
%
% \subsection{Use Tight Leading Rather Than Doublespace}
%
%    The option \Lopt{tighten} causes the text to set 
% with tight leading rather than double spacing
% as is normal in the |manuscript| substyle.
%
% Asserting this option presently has no effect in
% the other substyles because they make no changes 
% to the justification. If this situation changes,
% be careful to supply the appropriate hook as needed.
%
% Not! In the manuscript substyle, leading is doubled and
% the |\tightenlines| command is effectively disabled.
% So, what gives?
%
%    \begin{macrocode}
\DeclareOption{tighten}{\def\opt@tighten{\tightenlines}}%
\let\opt@tighten\@empty
%    \end{macrocode}

% \subsection{Summary}%
% The following are local options to this document class.
%
% \begin{center}
% \begin{tabular}{ll}
% |eqsecnum|&equations numbered by section\\
% |flushrt|&set text justified\\
% |tighten|&use tight leading\\
% \end{tabular}
% \end{center}
%
% \subsection{Substyle Options}%
% An \aastex\ document must have exactly one class option selected from 
% the following list; the default is |manuscript|.
% This option determines the substyle, signifying the 
% basic formatting choice for the document.
% \begin{center}
% \begin{tabular}{ll}
% \bfseries Class Option&\bfseries Purpose\\
% |manuscript|&Manuscript (the default)\\
% |preprint|&Preprint\\
% |preprint2|&Preprint in two columns
% \end{tabular}
% \end{center}
% Note that in the case of |preprint2|, we immediately signal
% the base document class that the two-column grid will be used.
%    \begin{macrocode}
\DeclareOption{manuscript}{\def\aas@substyle{manuscript}\long@abstractfalse}
\DeclareOption{preprint}{\def\aas@substyle{preprint}\long@abstractfalse}
\DeclareOption{preprint2}{\def\aas@substyle{preprint2}\@twocolumntrue}
%    \end{macrocode}
%
% For backward compatibility, we declare options corresponding to V4.0
%    \begin{macrocode}
\DeclareOption{aasms4}{\substitute@option{aasms4}{manuscript}}%
\DeclareOption{aaspp4}{\substitute@option{aaspp4}{preprint}}%
\DeclareOption{aas2pp4}{\substitute@option{aas2pp4}{preprint2}}%
\DeclareOption{aj_pt4}{\substitute@option{aj_pt4}{preprint}}%
\DeclareOption{apjpt4}{\substitute@option{apjpt4}{preprint}}%
%    \end{macrocode}
%
% \begin{macro}{\aas@substyle}
% At the end of package processing, we invoke the selected 
% substyle.
%    \begin{macrocode}
\def\aas@substyle{manuscript}%
\AtEndOfPackage{%
 \@nameuse{aas@\aas@substyle}%
}%
%    \end{macrocode}
% \end{macro}
%
% \subsection{Font Complement}
% \begin{macro}{astro}
% Invoking the |astro| class option causes the 
% astro font to be loaded at the appropriate size.
%    \begin{macrocode}
\DeclareOption{astro}{\AtBeginDocument{\load@astro}}%
%    \end{macrocode}
% \end{macro}
%
% \subsection{Default Option}%
%    \begin{macrocode}
\DeclareOption*{\PassOptionsToClass{\CurrentOption}{article}}%
%    \end{macrocode}
%
% \subsection{Executing Options}
%
%    Here we execute the default options to initialize certain
%    variables. 
%    \begin{macrocode}
\ExecuteOptions{letterpaper,oneside,onecolumn,final}%
%    \end{macrocode}
%
%    The |\ProcessOptions| command causes the execution of the code
%    for every option \Lopt{FOO}
%    which is declared and for which the user typed
%    the \Lopt{FOO} option in his
%    |\documentclass| command.  For every option \Lopt{BAR} he typed,
%    which is not declared, the option is assumed to be a global option.
%    All options will be passed as document options to any
%    |\usepackage| command in the document preamble.
%
% Note that |\ProcessOptions| is followed by |\relax| in order to 
% signal conclusively that the *-form has not been used.
%    \begin{macrocode}
\ProcessOptions\relax
%    \end{macrocode}
%
% \begin{macro}{aas@size@manuscript}
% \begin{macro}{aas@size@preprint}
% \begin{macro}{aas@size@preprint2}
% Now that local options have been processed, we know what substyle to use.
% We execute a procedure to pass the proper type size option to 
% the base class.
% Note that the macro |\aas@ptsize| acquires a value as a side effect of this 
% processing.
%
% As should be evident, the manuscript substyle mandates 12point, the default 
% size for preprint is 11 point, and that for two-column preprint is 10 point.
%
% Implementation note: here we must use \LaTeXe's cumbersome methods because we
% have not yet defined the Boolean control macros.
%    \begin{macrocode}
\@namedef{aas@size@manuscript}{%
 \def\aas@ptsize{2}%
 \pass@aassize
}%
\@namedef{aas@size@preprint}{%
 \ifx\undefined\aas@ptsize\expandafter\@firstoftwo\else\expandafter\@secondoftwo\fi
 {\def\aas@ptsize{1}}{}%
 \pass@aassize
}%
\@namedef{aas@size@preprint2}{%
 \ifx\undefined\aas@ptsize\expandafter\@firstoftwo\else\expandafter\@secondoftwo\fi
 {\def\aas@ptsize{0}}{}%
 \pass@aassize
}%
%    \end{macrocode}
% \end{macro}
% \end{macro}
% \end{macro}
%
% \begin{macro}{\pass@aassize}
% The procedure |\pass@aassize| is the bottleproc for passing the size option
% down to the base class, which is taken to be the |article| class.
%    \begin{macrocode}
\def\pass@aassize{%
 \expandafter\PassOptionsToClass\expandafter{\expandafter1\aas@ptsize pt}{article}%
}%
%    \end{macrocode}
% \end{macro}
%
% Now exercise the selected substyle option.
%    \begin{macrocode}
\@nameuse{aas@size@\aas@substyle}%
%    \end{macrocode}
%
%
%
% \section{Loading the Base Document Class and Packages}
%\aastex\ is a variant of the standard \LaTeX\ article class:
% |CTAN://macros/latex/unpacked|.
%    \begin{macrocode}
\LoadClass{article}[1995/12/01]%
%    \end{macrocode}
% 
% We always load the standard \LaTeX\ symbol package:
% |CTAN://macros/latex/base|.
%    \begin{macrocode}
\RequirePackage{latexsym}%
%    \end{macrocode}
% 
% We always load the standard \LaTeX\ extended graphics package:
% |CTAN://macros/latex/packages/graphics|.
%    \begin{macrocode}
\RequirePackage{graphicx}%
%    \end{macrocode}
% 
% The AMS |amssym| package provides access to symbols in the AMS fonts:
% |CTAN://macros/latex/packages/amslatex|.
%    \begin{macrocode}
\RequirePackage{amssymb}%
%    \end{macrocode}
% 
% The contributed \LaTeX\ |natbib| package:
% |CTAN://macros/latex/contrib/supported/natbib|.
%    \begin{macrocode}
\RequirePackage{natbib}%
\newcommand{\bibstyle@aas}{\bibpunct{(}{)}{;}{a}{}{,}}
\@citestyle{aas}
%    \end{macrocode}
% 
% The standard \LaTeX\ verbatim package may not be really necessary, but
% this code may serve as a model for loading other classes:
% |CTAN://macros/latex/packages/tools|.
%    \begin{macrocode}
\RequirePackage{verbatim}%
%    \end{macrocode}
%
% Here ends the code for the document class.
% The next section provides code that is common to 
% the document class and document style.
%    \begin{macrocode}
%</class>
%    \end{macrocode}
%
%
% \section{Extensions to the \LaTeX\ Kernel}%
% We extend and modify \LaTeX.
% Most of these tools are to help make the job of the 
% class file writer easier; some correct errors introduced 
% in \LaTeXe.
%
% All of the kernel extensions are incorporated into 
% the document style.
%    \begin{macrocode}
%<*class|style>
%    \end{macrocode}
%
% \subsection{New Tools}%
%
% \begin{macro}{\@to}
% \begin{macro}{\ulap}
% \begin{macro}{\dlap}
%
% The macro |\@to| is simple shorthand for |to|;
% the user commands |\ulap| and |\dlap| augment |\llap| and |\rlap|.
% \begin{verbatim}
%  \ulap{<vertical mode material>} % Will protrude upward
%  \dlap{<vertical mode material>} % Will protrude downward
% \end{verbatim}
%
%    \begin{macrocode}
\def\@to{to}%
\newcommand\ulap[1]{\vbox\@to\z@{{\vss#1}}}%
\newcommand\dlap[1]{\vbox\@to\z@{{#1\vss}}}%
%    \end{macrocode}
% \end{macro}
% \end{macro}
% \end{macro}
%
% \begin{macro}{\let@environment}
% The directive |\let@environment| takes care of a common programming 
% idiom whereby one environment is made a synonym for another.
%    \begin{macrocode}
\newcommand\let@environment[2]{%
 \expandafter\let
 \csname#1\expandafter\endcsname\csname#2\endcsname
 \expandafter\let
 \csname end#1\expandafter\endcsname\csname end#2\endcsname
 }%
%    \end{macrocode}
% \end{macro}
%
% \begin{macro}{\say}
% \begin{macro}{\sayythe}
% The commands |\say| and |\saythe| cause diagnostic messages in the 
% \TeX\ log that give the value of a control sequence name or a register 
% respectively.
%    \begin{macrocode}
\newcommand\say[1]{\typeout{<\noexpand#1=\meaning#1>}}%
\newcommand\saythe[1]{\typeout{<\noexpand#1=\the#1>}}%
%    \end{macrocode}
% \end{macro}
% \end{macro}
%
% \begin{macro}{\tracingplain}
% The command |\tracingplain| causes \TeX's tracing parameters to 
% return to the values set by default. This command is sometimes
% useful when you have said |\tracingall| somewhere and want to
% restore.
% The |\traceoutput| command causes |\tracingoutput| diagnostics
% upon |\shipout|.
%    \begin{macrocode}
\newcommand\tracingplain{%
 \tracingonline\z@\tracingcommands\z@\tracingstats\z@
 \tracingpages\z@\tracingoutput\z@\tracinglostchars\@ne
 \tracingmacros\z@\tracingparagraphs\z@\tracingrestores\z@
 \showboxbreadth5\showboxdepth3\relax %\errorstopmode
 }%
\newcommand\traceoutput{%
 \appdef\@resetactivechars{\showoutput}%
}%
%    \end{macrocode}
% \end{macro}
%
% \begin{macro}{\nohyphenation}
% \begin{macro}{\offhyphenation}
% The command |\nohyphenation| immediately preceding a word will
% prevent \TeX\ from hyphenating that word;
% the |\offhyphenation| command causes \TeX\ to not hyphenate within its
% scope.
%    \begin{macrocode}
\newcommand\nohyphenation{\kern\z@}%
\newcommand\offhyphenation{%
 \lefthyphenmin\maxdimen\emergencystretch8\p@
}%
%    \end{macrocode}
% \end{macro}
% \end{macro}
%
% \begin{macro}{\substitute@command}
% \begin{macro}{\remove@command}
%    \begin{macrocode}
\def\substitute@command#1#2{%
 \ClassWarning{aastex}{%
  Command \protect#1\space is deprecated in aastex.
  Using \protect#2\space instead (please fix your document).
 }%
 #2%
}%
\def\remove@command#1{%
 \ClassWarning{aastex}{%
  Command \protect#1\space is deprecated in aastex.
  Please remove it from your document%
 }%
}%
%    \end{macrocode}
% \end{macro}
% \end{macro}
%
%
% \subsection{Boolean Control}%
% \begin{macro}{\prepdef}
% \begin{macro}{\appdef}
%
% Provide the capability of performing head- and tail patches.
% The procedure |\prepdef| prepends to the given macro 
% the tokens specified in its second argument.
% Likewise for |\appdef|, except that it appends.
% Note that the first 10 toks registers are utility registers,
% and we simply make a control sequence name, |\toks@ii|, for one of 
% them.
%    \begin{macrocode}
\def\@provide#1{%Implicit #2
 \@ifxundefined#1{\def#1}{\def\@tempa}% 
}% 
\def\prepdef#1#2{%
 \@ifxundefined#1{\toks@{}}{\toks@\expandafter{#1}}%
 \toks@ii{#2}%
 \edef#1{\the\toks@ii\the\toks@}%
}%
\def\appdef#1#2{%
 \@ifxundefined#1{\toks@{}}{\toks@\expandafter{#1}}%
 \toks@ii{#2}%
 \edef#1{\the\toks@\the\toks@ii}%
}%
\def\appgdef#1#2{%
 \@ifxundefined#1{\toks@{}}{\toks@\expandafter{#1}}%
 \toks@ii{#2}%
 \xdef#1{\the\toks@\the\toks@ii}%
}% 
\toksdef\toks@ii=\tw@
%    \end{macrocode}
% \end{macro}
% \end{macro}
%
% \begin{macro}{\@ifxundefined}
% We introduce just enough of the Boolean calculus for \TeX.
% Alan Jeffrey was the pioneer here, with an article in TUGboat
% (Vol. 11, No. 2, page 237).
% This implementation owes a debt to
% William Baxter (web@superscript.com).
% See articles by Baxter and Ogawa in the proceedings of the 
% 1994 TUG meeting.
% 
% Other utility procedures used |\@ifxundefined|,
% which is defined here in terms of |\@ifx|.
%
%    \begin{macrocode}
\long\def\@ifxundefined#1{\@ifx{\undefined#1}}%
%    \end{macrocode}
% \end{macro}
%
% \begin{macro}{\@boolean}
% \begin{macro}{\@boole@def}
% In order to define |\@ifx|, we first must create the
% ``defining word'' (term taken form our Forth vocabulary)
% |\@boole@def|, which employs |\@boolean| to do its job.
%    \begin{macrocode}
\def\@boolean#1#2{%
  \long\def#1{%
    #2% \if<something>
      \expandafter\boolean@true
    \else
      \expandafter\boolean@false
    \fi
  }%
}%
\def\@boole@def#1#{\@boolean{#1}}% Implicit #2
%    \end{macrocode}
% \end{macro}
% \end{macro}
%
% \begin{macro}{\@ifx}
% \begin{macro}{\@ifx@empty}
% We can now invoke the defining word to create
% the procedure |\@ifx|.
%
%    \begin{macrocode}
\@boole@def\@ifx#1{\ifx#1}%
\@boole@def\@ifx@empty#1{\ifx\@empty#1}%
\@boole@def\@ifnum#1{\ifnum#1}%
\@boole@def\@ifdim#1{\ifdim#1}%
\@boole@def\@ifvoid#1{\ifvoid#1}%
%    \end{macrocode}
% \end{macro}
% \end{macro}
%
% \begin{macro}{\boolean@true}
% \begin{macro}{\boolean@false}
%
% Note that when a Boolean operator expands, it
% employs two macros that act as selectors, defined here.
%
%    \begin{macrocode}
\long\def\boolean@true#1#2{#1}%
\long\def\boolean@false#1#2{#2}%
%    \end{macrocode}
% \end{macro}
% \end{macro}
%
% \subsection{Damage Control}%
%
% \subsubsection{Vertical Stuff}%
% \begin{macro}{\addpenalty}
% \begin{macro}{\addvspace}
% Repair some \LaTeX\ damage.
%
%    \begin{macrocode}
\renewcommand\addpenalty[1]{%
  \ifvmode
    \if@minipage
    \else
      \if@nobreak
      \else
        \ifinner
          \@tempskipa\lastskip \unskip
          \@tempcnta\lastpenalty \unpenalty
          \ifnum#1<\@tempcnta
            \penalty#1\relax
          \else
            \penalty\@tempcnta
          \fi
          \vskip\@tempskipa
        \else
          \ifdim\lastskip=\z@
            \penalty#1\relax
          \else
            \@tempskipb\lastskip
            \vskip-\lastskip
            \penalty#1%
            \vskip\@tempskipb
          \fi
        \fi
      \fi
    \fi
  \else
    \@noitemerr
  \fi
}%
%
\renewcommand\addvspace[1]{%
  \ifvmode
    \if@minipage
    \else
      \ifinner
        \@tempskipa\lastskip \unskip
        \ifdim\@tempskipa=\z@
          \vskip#1\relax
        \else
          \@tempskipb#1\relax
          \ifdim\@tempskipa<\@tempskipb
            \vskip\@tempskipb
          \else
            \ifdim\@tempskipb<\z@
              \ifdim\@tempskipa<\z@
                \vskip\@tempskipa
              \else
                \advance\@tempskipb\@tempskipa
                \vskip\@tempskipb
              \fi
            \else
              \vskip\@tempskipa
            \fi
          \fi
        \fi
      \else
        \ifdim\lastskip=\z@
          \vskip#1\relax
        \else
          \@tempskipb#1\relax
          \@xaddvskip
        \fi
      \fi
    \fi
  \else
    \@noitemerr
  \fi
}%
%    \end{macrocode}
% \end{macro}
% \end{macro}
%
% \subsubsection{Title Page Formatting}%
% \begin{macro}{\maketitle@page}
% The |\maketitle| (for a title page) of \file{classes.dtx}.
%    \begin{macrocode}
\newcommand\maketitle@page{%
 \maketitle@disarm
 \begin{titlepage}%
  \let\footnotesize\small
  \let\footnoterule\relax
  \let \footnote \thanks
%    \end{macrocode}
%
%    \begin{macrocode}
  \@maketitle
%    \end{macrocode}
%
%    \begin{macrocode}
  \@thanks
  \vfil\null
 \end{titlepage}%
%    \end{macrocode}
%
%    \begin{macrocode}
 \after@maketitle
}%
%    \end{macrocode}
% \end{macro}
% \begin{macro}{\maketitle@runin}
% The |\maketitle| (for a run-in title page) of \file{classes.dtx}
% appears to be broken. 
%    \begin{macrocode}
\newcommand\maketitle@runin{%
  \par
  \maketitle@disarm
  \begingroup
    %%%%\renewcommand\thefootnote{\@fnsymbol\c@footnote}%
    \def\@makefnmark{\rlap{\@textsuperscript{\normalfont\@thefnmark}}}%
    \long\def\@makefntext##1{%
     \parindent 1em\noindent
     \hb@xt@1.8em{%
      \hss\@textsuperscript{\normalfont\@thefnmark}%
     }%
     ##1%
    }%
%    \end{macrocode}
% Here, we swap the two predicates:
%    \begin{macrocode}
    \if@twocolumn
      \@ifnum{\col@number=\@ne}{%
        \global\let\@title@footnote@text\@empty%
        \twocolumn[\let\@footnotetext\@title@footnotetext\@maketitle]%
        \@title@footnote@text%
      }{%
        \@maketitle
      }%
    \else
%    \end{macrocode}
%
%    \begin{macrocode}
      \newpage
      \global\@topnum\z@   % Prevents figures from going at top of page.
      \@maketitle
    \fi
%    \end{macrocode}
%
%    \begin{macrocode}
    %\thispagestyle{plain}%
    \@thanks
%    \end{macrocode}
%
%    \begin{macrocode}
  \endgroup
  \after@maketitle
}%
%    \end{macrocode}
% \end{macro}
% \begin{macro}{\after@maketitle}
% Common processing for title whether run-in or page. And switching from one-column 
% to two-column mode if \Lopt{longabstract} option is used.
%    \begin{macrocode}
\def\after@maketitle{%
 \iflong@abstract %
   \let\@makecol\@makecol@pptt
   \let\@outputdblcol\@outputdblcol@pptt
   \global\columnwidth\textwidth%
   \global\advance\columnwidth-\columnsep
   \global\divide\columnwidth\tw@
   \global\hsize\columnwidth
   \global\linewidth\columnwidth
   \global\@twocolumntrue
   \global\@firstcolumntrue
   \col@number \tw@
   \@ifnextchar [\@topnewpage\@floatplacement
   \begin{strip}% separate the abstract text from two-column text 
   \end{strip}
 \fi
 %%%\setcounter{footnote}{0}%
 \global\let\thanks\@gobble
 \global\let\maketitle\relax
 \global\let\@maketitle\relax
 \global\let\@thanks\@empty
 \global\let\@author\@empty
 \global\let\@date\@empty
 \global\let\@title\@empty
 \global\let\title\relax
 \global\let\author\relax
 \global\let\date\relax
 \global\let\and\relax
 \global\let\email\@text@email
 \global\firsttabskip\z@
}%
%    \end{macrocode}
% \end{macro}
% Defer until the end of the processing for the document class
% the decision between the two formatting procedures.
%    \begin{macrocode}
\AtEndOfClass{%
 \if@titlepage
  \let\maketitle\maketitle@page
 \else
  \let\maketitle\maketitle@runin
 \fi
}%
%    \end{macrocode}
%
%
% \subsubsection{Repairing Footnotes}%
%
% \begin{macro}{\@mpfootnotetext}
% \begin{macro}{\@footnotetext}
% Reparametrize so that |\footnotesize| (a user-level command)
% is not invoked---use |\footnote@size| instead.
% In general, user-level commands should not be invoked by
% low-level formatting procedures.
%
% Also, factor out the shared code into a procedure called |\set@footnotetext|.
%    \begin{macrocode}
%
\def\open@footnotetext{% 
  \insert\footins\bgroup % \egroup in \close@footnotetext 
  \interlinepenalty\interfootnotelinepenalty 
  \splittopskip\footnotesep 
  \splitmaxdepth\dp\strutbox 
  \floatingpenalty \@MM 
  \reset@font\footnote@size 
  \hsize\columnwidth 
  \@parboxrestore 
}% 
\def\close@footnotetext#1{% 
  \color@begingroup 
   \@makefntext{% 
    \rule\z@\footnotesep\ignorespaces 
    #1\@finalstrut\strutbox 
   }% 
  \color@endgroup 
 \egroup 
}% 
\def\@footnotetext{% Implict #1 
 \open@footnotetext 
  \set@footnotetext{p@footnote}% 
}%
% 
\def\@mpfootnotetext{% Implict #1 
 \global\setbox\@mpfootins\vbox\bgroup % \egroup in \set@footnotetext 
  \unvbox\@mpfootins 
  \set@footnotetext{p@mpfootnote}% 
}% 
\long\def\set@footnotetext#1{%Implicit #2 
  \protected@edef\@currentlabel{% 
   \csname #1\endcsname\@thefnmark 
  }% 
 \close@footnotetext 
}% 
\def\@title@footnotetext#1{% 
 \protected@edef\@tempa{\csname p@footnote\endcsname\@thefnmark}% 
 \expandafter\appgdef 
 \expandafter\@title@footnote@text 
 \expandafter{% 
 \expandafter\insert@footnotetext 
 \expandafter{\@tempa}{#1}}% 
}% 
\def\insert@footnotetext#1{%Implicit #2 
 \open@footnotetext 
  \def\@thefnmark{#1}% 
  \def\@currentlabel{#1}% 
 \close@footnotetext 
}% 
%
\AtBeginDocument{%
 \ifx\footnote@size\undefined
  \let\footnote@size\footnotesize
 \fi
}%
%    \end{macrocode}
% \end{macro}
% \end{macro}
%
%
% \subsection{Eqnarray Correction}%
%
% \begin{macro}{\eqnarray@broken}
% \begin{macro}{\eqnarray@fleqn@fixed}
% Team \LaTeX\ still hasn't repaired Leslie's broken 
% definition of |eqnarray|.
% Let us be bold\dots.
%
%    \begin{macrocode}
\long\def\eqnarray@broken{%
 \stepcounter{equation}\def\@currentlabel{\p@equation\theequation}%
 \global\@eqnswtrue\m@th\global\@eqcnt\z@
 \tabskip\mathindent
 \let\\=\@eqncr
 \setlength\abovedisplayskip{\topsep}%
 \ifvmode\addtolength\abovedisplayskip{\partopsep}\fi
 \addtolength\abovedisplayskip{\parskip}%
 \setlength\belowdisplayskip{\abovedisplayskip}%
 \setlength\belowdisplayshortskip{\abovedisplayskip}%
 \setlength\abovedisplayshortskip{\abovedisplayskip}%
 $$%
 \everycr{}%
 \halign to\linewidth\bgroup
  \hskip\@centering$\displaystyle\tabskip\z@skip{##}$\@eqnsel
  &\global\@eqcnt\@ne
   \hskip\tw@\arraycolsep
   \hfil${##}$\hfil
  &\global\@eqcnt\tw@
   \hskip\tw@\arraycolsep
   $\displaystyle{##}$\hfil\tabskip\@centering
  &\global\@eqcnt\thr@@\hb@xt@\z@\bgroup\hss##\egroup
   \tabskip\z@skip
  \cr
}%
\long\def\eqnarray@fleqn@fixed{%
 \stepcounter{equation}\def\@currentlabel{\p@equation\theequation}%
 \global\@eqnswtrue\m@th\global\@eqcnt\z@
 \tabskip\mathindent
 \let\\=\@eqncr
 \setlength\abovedisplayskip{\topsep}%
 \ifvmode\addtolength\abovedisplayskip{\partopsep}\fi
 \addtolength\abovedisplayskip{\parskip}%
 \setlength\belowdisplayskip{\abovedisplayskip}%
 \setlength\belowdisplayshortskip{\abovedisplayskip}%
 \setlength\abovedisplayshortskip{\abovedisplayskip}%
 $$%
 \everycr{}%
 \halign\@to\linewidth\bgroup
  \hskip\@centering$\displaystyle\tabskip\z@skip{##}$\@eqnsel
  &\global\@eqcnt\@ne
   \hskip\tw@\eqncolsep
   \hfil${{}##{}}$\hfil
  &\global\@eqcnt\tw@
   \hskip\tw@\eqncolsep 
   $\displaystyle{##}$\hfil\tabskip\@centering
  &\global\@eqcnt\thr@@\hb@xt@\z@\bgroup\hss##\egroup
   \tabskip\z@skip
  \cr
}%
\ifx\eqnarray\eqnarray@broken
 \ClassWarningNoLine{aastex}{Repairing broken {eqnarray} definition}%
 \let\eqnarray\eqnarray@fleqn@fixed
\fi
\let\eqnarray@broken\relax
\let\eqnarray@fleqn@fixed\relax
%    \end{macrocode}
% \end{macro}
% \end{macro}
%
% \subsection{Customization of Sections}%
% \begin{macro}{\@startsection}
% Patch a head hook into the basic sectioning command.
%    \begin{macrocode}
\def\@startsection#1#2#3#4#5#6{%
 \@startsection@hook
  \if@noskipsec \leavevmode \fi
  \par
  \@tempskipa #4\relax
  \@afterindenttrue
  \ifdim \@tempskipa <\z@
    \@tempskipa -\@tempskipa \@afterindentfalse
  \fi
  \if@nobreak
    \everypar{}%
  \else
    \addpenalty\@secpenalty\addvspace\@tempskipa
  \fi
  \@ifstar
    {\@ssect{#3}{#4}{#5}{#6}}%
    {\@dblarg{\@sect{#1}{#2}{#3}{#4}{#5}{#6}}}%
}%
\def\@startsection@hook{}%
%    \end{macrocode}
% \end{macro}
%
% \begin{macro}{\@sect}
% \begin{macro}{\@ssect}
% Adjust section heading typography slightly.  The manuscript editors
% seem to like the idea that the heads will be centered, and that there
% will be a |.| trailing the section number.  This requires hacking 
% below the normal \LaTeX\ programming interface (|\@startsection|), 
% since the centering has to be turned on and off again before the 
% section starts.
%
% FIXME: cf \LaTeX\ standard definitions. Hooks for |\sec@upcase| and 
% centering.
%    \begin{macrocode}
\def\@sect#1#2#3#4#5#6[#7]#8{%
 \ifnum #2>\c@secnumdepth
  \let\@svsec\@empty
 \else
  \refstepcounter{#1}\protected@edef\@svsec{\@seccntformat{#1}\relax}
 \fi
  \@tempskipa #5\relax
  \ifdim \@tempskipa>\z@
  \begingroup \section@centering#6{%
    \@hangfrom{\hskip #3\relax\@svsec}\interlinepenalty \@M\sec@upcase{#8}\@@par}%
  \endgroup 
  \csname #1mark\endcsname{#7}\addcontentsline
  {toc}{#1}{\ifnum #2>\c@secnumdepth \else
               \protect\numberline{\csname the#1\endcsname}\fi
    #7}\else
  \def\@svsechd{%
   #6{\hskip#3\relax\@svsec 
    \sec@upcase{#8}}% 
   \csname #1mark\endcsname
   {#7}\addcontentsline{toc}{#1}{% 
    \ifnum #2>\c@secnumdepth 
    \else 
     \protect\numberline{\csname the#1\endcsname}% 
    \fi 
    #7}}\fi 
 \@xsect{#5}% 
}% 
%
\def\@ssect#1#2#3#4#5{\@tempskipa #3\relax
 \ifdim \@tempskipa>\z@ 
  \begingroup 
   \section@centering
   #4{%
    \@hangfrom{\hskip #1}\interlinepenalty \@M
    \sec@upcase{#5}\@@par}\endgroup
   \else 
    \def\@svsechd{#4{\hskip #1\relax \sec@upcase{#5}}}\fi
    \@xsect{#3}}%
\def\@seccntformat#1{\csname the#1\endcsname.\quad}
\def\section@centering{\centering}%
\def\sec@upcase#1{\relax{#1}}%
%    \end{macrocode}
% \end{macro}
% \end{macro}
%
% \subsection{Customization of Citations}%
% \begin{macro}{\@citex}
% Modify |\@citex| so that individual citation texts are {\em not} 
% |\hbox|ed.
%
% This allows \TeX\ to break lines at spaces between words in citations.
% Remove the obnoxious |{\bf ?}| that was inserted in the text when
% undefined citations are encountered; here, the warning message is
% printed only on the terminal.
%
% FIXME: doublecheck compatibility with natbib!
%
%    \begin{macrocode}
\def\@citex[#1]#2{%
 \let\@citea\@empty
 \@cite{%
  \@for\@citeb:=#2\do{%
   \@citea
   \def\@citea{,\penalty\@m\ }%
   \edef\@citeb{\expandafter\@firstofone\@citeb\@empty}%
   \if@filesw\immediate\write\@auxout{\string\citation{\@citeb}}\fi
   \@ifundefined{b@\@citeb}{%
%   \mbox{\reset@font\bfseries ?}%
    \G@refundefinedtrue
    \@latex@warning{%
     Citation `\@citeb' on page \thepage \space undefined%
    }%
   }{%
%   \hbox{%
     \csname b@\@citeb\endcsname
%   }%
   }%
  }%
 }{#1}%
}%
%    \end{macrocode}
% \end{macro}
%
% All of the kernel extensions have been incorporated into 
% the document class and document style.
%    \begin{macrocode}
%</class|style>
%    \end{macrocode}
%
%
% \section{More of the Document Style Module}%
%    \begin{macrocode}
%<*style>
%    \end{macrocode}
%
% Now that kernel modifications are complete, read in the selected substyle.
%    \begin{macrocode}
\expandafter\@@input\aastex@substyle.sty\relax
%    \end{macrocode}
%
% And begin patching the code.
%    \begin{macrocode}
\def\revtex@ver{4.1}%
\def\revtex@date{25 Jan 99}%
%    \end{macrocode}
%
% Undefine commands that we are about to override.
%    \begin{macrocode}
\let\lefthead\undefined
\let\righthead\undefined
\let\references\undefined
\let\reference\undefined
\let\markcite\undefined
\let\colhead\undefined
\let\tablecolumns\undefined
\let\tablecaption\undefined
\let\tablehead\undefined
\let\tabletail\undefined
\let\tablewidth\undefined
\let\tableheadfrac\undefined
\let\deluxetable\undefined
\let\startdata\undefined
\let\tablevspace\undefined
\let\tablebreak\undefined
\let\cutinhead\undefined
\let\sidehead\undefined
\let\tablerefs\undefined
\let\tablecomments\undefined
\let\phn\undefined
\let\phd\undefined
\let\phs\undefined
\let\phm\undefined
\let\sq\undefined
\let\arcdeg\undefined
\let\arcmin\undefined
\let\arcsec\undefined
\let\fd\undefined
\let\fh\undefined
\let\fm\undefined
\let\fs\undefined
\let\fdg\undefined
\let\fp\undefined
\let\micron\undefined
\let\case\undefined
\let\slantfrac\undefined
\let\onehalf\undefined
\let\onethird\undefined
\let\twothirds\undefined
\let\onequarter\undefined
\let\threequarters\undefined
\let\ubvr\undefined
\let\ub\undefined
\let\bv\undefined
\let\vr\undefined
\let\ur\undefined
\let\ion\undefined
\let\nodata\undefined
\let\sun\undefined
\let\earth\undefined
%    \end{macrocode}
%
% The overrides and additions to the V4.0
% document style are drawn from the same code that goes into
% the document class.
%
% Once those emendments have been read in,
% we will return to the document style to finish things off.
%    \begin{macrocode}
%</style>
%    \end{macrocode}
%
%
% \section{\aastex\ Additions and Amendments}
%  This section presents commands that are specific to the \aastex\
%  macro package.
%
% All of the \aastex\ emendments are incorporated into the document class.
%    \begin{macrocode}
%<*class>
%    \end{macrocode}
%
% \subsection{A Special Note on the Two-Column Preprint Style}%
%
% It is not the objective of the \aastex\ project to develop \LaTeX\ 
% styles that produce pages that mimic the appearance of specific AAS 
% journals.
% However, it would be nice to offer format options within the preprint
% styles so that output format can be changed.
% (Within the ``manuscript''
% style, |\tighten| is available, and this effectively turns off the
% doublespacing.)  The primary motivation behind this idea is to provide 
% some definitions that permit institutions to distribute ``pretty'' 
% preprints with this package.
% (Of course, I make the glaring assumption that I
% can determine what someone else will think is pretty.)
%
% This style has the principle function of setting up two-column output.
% Although it should be painfully obvious, it is important to remember
% that text lines are considerably shorter when two of them are typeset
% side by side on a page.  Long equations, wide tables and figures, and
% the like, may not slide into this format without some adjustments.
% It is our opinion that great effort to adapt copy and markup for
% two-column pages is counterproductive; remember that the main goals
% of this package are to enable comprehensive tagging of the article
% and, secondarily, to produce ``correct'' draft (or referee) format 
% pages.
%
% I elect not to impose a format for the front matter, although I feel
% it looks best if the title, author, abstract, and keyword material is
% typeset on its own page at full text width.  The author may supply the
% |\twocolumn| command wherever desired, and, for that matter, is free
% to switch between |\onecolumn| and |\twocolumn| modes.
% (However, this is bad form.)
% \begin{verbatim}
%     \twocolumn % Enter two-column page grid
%     \onecolumn % Enter one-column page grid
% \end{verbatim}
% Note that the two-column page grid begins at the point |\twocolumn| 
% appears in the text, and if that point is before the frontmatter, 
% the latter will get typeset in two-column mode
% along with the rest of it.
% For purposes of producing ``pretty'' output,
% it is probably desirable to 
% put this command after the abstract and keywords, just before the body 
% of the paper.
% If the command is not specified explicitly, 
% the document class will automatically switch between the one-column
% and two-column page grid as needed.
%
% \subsection{Configuration Information}
%
% \begin{macro}{\revtex@ver}
% \begin{macro}{\revtex@date}
% \begin{macro}{\revtex@org}
% \begin{macro}{\revtex@jnl}
% \begin{macro}{\revtex@genre}
% \begin{macro}{\revtex@pageid}
% My convention for identifying configuration information is for styles
% to define the following five parameters:
% |\revtex@ver| is the version number of this file;
% |\revtex@date| signifies the revision date;
% |\revtex@org| is the sponsoring organization;
% and
% |\revtex@genre| gives the applicable journal, if specified.
%    \begin{macrocode}
\def\revtex@ver{5.2}%
\def\revtex@date{12 Dec 2003}%
\def\revtex@org{AAS}%
\def\revtex@jnl{AAS}%
\def\revtex@genre{manuscript}%
%\def\revtex@genre{preprint}%
%    \end{macrocode}
%
% Display identifying message on screen and in logfile.
%
%    \begin{macrocode}
\def\class@typeout{%
 \typeout{%
  \revtex@org\space
% \revtex@jnl\space
  \revtex@genre\space substyle, v\revtex@ver\space <\revtex@date>.%
 }%
}%
\def\class@typeout@jnl{%
 \typeout{%
  \revtex@org\space
  \revtex@jnl\space
  \revtex@genre\space substyle, v\revtex@ver\space <\revtex@date>.%
 }%
}%
\AtBeginDocument{%
 \class@typeout
}%
%    \end{macrocode}
%
% The following mini-colophon identifies the article as having been
% prepared with this macro package.  The |\revtex@pageid| command only
% defines the formatting of this colophon; its existence and placement
% within the printed article is controlled by the use of this command
% in the style file.
%
%    \begin{macrocode}
\def\revtex@pageid{%
 \xdef\@thefnmark{\null}%
 \@footnotetext{%
  This \revtex@genre\space was prepared with the
  \revtex@org\space \LaTeX\ macros v\revtex@ver.%
 }%
}%
%    \end{macrocode}
% \end{macro}
% \end{macro}
% \end{macro}
% \end{macro}
% \end{macro}
% \end{macro}
%
% \subsection{Page Style}
%
% \begin{macro}{\ps@plaintop}
% Create |plaintop| page style that has folios at top of page,
% and use |plaintop| pagestyle by default.
%
%    \begin{macrocode}
\def\ps@plaintop{%
 \let\@mkboth\@gobbletwo
 \def\@oddfoot{}%
 \def\@oddhead{\rmfamily\hfil--\space\thepage\space--\hfil}%
 \def\@evenfoot{}%
 \let\@evenhead\@oddhead
}%
\def\ps@msstitle{% mss ID on title page
 \let\@mkboth\@gobbletwo
 \def\@oddfoot{}%
 \def\@oddhead{\rmfamily\hfil\ms@id}%
 \def\@evenfoot{}%
 \let\@evenhead\@oddhead
}%
%    \end{macrocode}
% \end{macro}
%
% \subsection{Page Layout}
%
% \subsubsection{The pica area}
%
% \begin{macro}{\textwidth}
% Set the text width
%    \begin{macrocode}
\setlength\textwidth{6.5in}%
%    \end{macrocode}
% \end{macro}
%
% \begin{macro}{\textheight}
% The |\textheight| is the height of text
%    (including footnotes and figures, excluding running head and
%    foot).
%
%    \begin{macrocode}
\setlength\textheight{8.4in}%
%    \end{macrocode}
% \end{macro}
%
% \subsubsection{Vertical spacing}
%
% \begin{macro}{\headheight}
% \begin{macro}{\headsep}
% \begin{macro}{\topskip}
%    The |\headheight| is the height of the box that will contain the
%    running head. The |\headsep| is the distance between the bottom
%    of the running head and the top of the text. The |\topskip| is
%    the |\baselineskip| for the first line on a page; \LaTeX's output
%    routine will not work properly if it has the value 0pt, so do not
%    do that!
%    \begin{macrocode}
\setlength\headheight{.15in}%
%\setlength\headsep   {.35in}%
%\setlength\topskip   {10\p@}%
%    \end{macrocode}
% \end{macro}
% \end{macro}
% \end{macro}
%
% \subsubsection{Margins}
%
% \begin{macro}{\marginparsep}
% \begin{macro}{\marginparpush}
%    The horizontal space between the main text and marginal notes is
%    determined by |\marginparsep|, the minimum vertical separation
%    between two marginal notes is controlled by |\marginparpush|.
%    \begin{macrocode}
 \setlength\marginparsep {10\p@}%
 \setlength\marginparpush{5\p@}%
%    \end{macrocode}
% \end{macro}
% \end{macro}
%
%    Now we can give the values for the other margin parameters. For
%    native \LaTeXe, these are calculated.
% \begin{macro}{\oddsidemargin}
% \begin{macro}{\evensidemargin}
% \begin{macro}{\marginparwidth}
%
%    Values for two-sided printing:
%    \begin{macrocode}
 \setlength\oddsidemargin   {0in}%
 \setlength\evensidemargin  {0in}%
 \setlength\marginparwidth  {0in}%
%    \end{macrocode}
% \end{macro}
% \end{macro}
% \end{macro}
%
% \begin{macro}{\topmargin}
%    The |\topmargin| is the distance between the top of `the
%    printable area'---which is 1 inch below the top of the
%    paper--and the top of the box which contains the running head.
%
%    \begin{macrocode}
 \setlength\topmargin{0in}%
%    \end{macrocode}
% \end{macro}
%
% \subsection{Paragraphing}
%
% \begin{macro}{\lineskip}
% \begin{macro}{\normallineskip}
%    These parameters control \TeX's behavior when two lines tend to
%    come too close together.
%    \begin{macrocode}
\setlength\lineskip{1\p@}%
\setlength\normallineskip{1\p@}%
%    \end{macrocode}
% \end{macro}
% \end{macro}
%
% \begin{macro}{\@doubleleading}
% \begin{macro}{\@tightleading}
% The two alternatives for leading are defined here.
%    \begin{macrocode}
\def\@doubleleading{1.6}%
\def\@tightleading{1.1}%
%    \end{macrocode}
% \end{macro}
% \end{macro}
%
% \begin{macro}{\baselinestretch}
%    This is used as a multiplier for |\baselineskip|.
%    Note that if this command
%    doesn't resolve to ``empty'' any \texttt{plus} or \texttt{minus}
%    part in the specification of |\baselineskip| is ignored.
%
% Double spacing is the default: the style uses
% the \@doubleleading factor, defined above.
%
%    \begin{macrocode}
\renewcommand\baselinestretch{\@tightleading}%
%    \end{macrocode}
% \end{macro}
%
% \begin{macro}{\tighten}
% The |\tighten| command merely resets the |\baselinestretch| to a
% smaller factor, resulting in single-column text that is more closely
% spaced.  (This will always work.)  Tightening the vertical spacing
% in the text results in output that is more attractive for distribution
% to colleagues who are primarily interested in reading the paper, as
% opposed to copy editing it.  Double-spaced output for referees and
% copy editors is, however, the main objective of this style, hence
% double-spacing occurs by default.  The same effect as |\tighten|
% can be achieved by using the |tighten| document style option.
%
% The |\tighten| command is intended for use in the preamble of the 
% article.
%
%    \begin{macrocode}
\newcommand\tighten{\def\baselinestretch{\@tightleading}}%
\newcommand\tightenlines{\tighten}%
%    \end{macrocode}
% \end{macro}
%
% \begin{macro}{\singlespace}
% \begin{macro}{\doublespace}
% For authors who feel compelled to change between single and double
% spacing in the manuscript, |\singlespace| and |\doublespace| commands
% are provided.  (Following recommendations of AAS Publications Board
% Subcommittee on Preprints.)  These are no-ops in manuscript style.
%
%    \begin{macrocode}
\newcommand\singlespace{%
 \def\baselinestretch{\@tightleading}%
 \normalsize
}%
\newcommand\doublespace{%
 \def\baselinestretch{\@doubleleading}%
 \normalsize
}%
%    \end{macrocode}
%
% Here are the corresponding command definitions for the ApJ style:
%    \begin{macrocode}
\def\singlespace@apj{%
 \def\baselinestretch{\@tightleading}%
 \def\arraystretch{\@tightleading}%
 \normalsize
}%
\def\doublespace@apj{%
 \def\baselinestretch@apj{\@doubleleading}%
 \def\arraystretch{\@doubleleading}%
 \normalsize
}%
%    \end{macrocode}
% \end{macro}
% \end{macro}
%
% \begin{macro}{\sec@upcase}
% Upper case for sections (optional upper case items).  These are 
% created so that some headings can be toggled between mixed case and 
% upper case readily.
% Headings that might be changed can be wrapped in the style file in 
% |\sec@upcase{|\m{text}|}| constructs;
% the expansion of |\sec@upcase| is
% controlled here.  It is |\relax| by default (mixed case heads), and 
% can easily be changed to |\uppercase| if desired.
% If mixed case headings are wanted by the editor, authors {\em must} 
% supply mixed case text, although this is what authors should be doing 
% anyway.
% (Mixed can be converted to upper,
% but the reverse transformation cannot be automated.)
%
%    \begin{macrocode}
%    \end{macrocode}
% \end{macro}
%
% \begin{macro}{\parskip}
% \begin{macro}{\parindent}
%    |\parskip| gives extra vertical space between paragraphs and
%    |\parindent| is the width of the paragraph indentation. The value
%    of |\parindent| depends on whether we are in two column mode.
%    \begin{macrocode}
\setlength\parskip{1.5ex}%
\setlength\parindent{2em}%
%    \end{macrocode}
% \end{macro}
% \end{macro}
%
%  \begin{macro}{\smallskipamount}
%  \begin{macro}{\medskipamount}
%  \begin{macro}{\bigskipamount}
%    The values for these three parameters are set in the \LaTeX\
%    kernel.
%    \begin{macrocode}
\setlength\smallskipamount{3\p@\@plus1\p@\@minus1\p@}%
\setlength\medskipamount{6\p@\@plus2\p@\@minus2\p@}%
\setlength\bigskipamount{12\p@\@plus4\p@\@minus4\p@}%
%    \end{macrocode}
% \end{macro}
% \end{macro}
% \end{macro}
%
% \begin{macro}{\@lowpenalty}
% \begin{macro}{\@medpenalty}
% \begin{macro}{\@highpenalty}%
%    The commands |\nopagebreak| and |\nolinebreak| put in penalties
%    to discourage these breaks at the point they are put in.
%    They use |\@lowpenalty|, |\@medpenalty| or |\@highpenalty|,
%    dependent on their argument.
%    \begin{macrocode}
\@lowpenalty   51
\@medpenalty  151
\@highpenalty 301
%    \end{macrocode}
% \end{macro}
% \end{macro}
% \end{macro}
%
% \begin{macro}{\clubpenalty}
% \begin{macro}{\widowpenalty}
%    These penalties are use to discourage club and widow lines.
%    Because we use their default values we only show them here,
%    commented out.
%    \begin{macrocode}
% \clubpenalty  150
% \widowpenalty 150
%    \end{macrocode}
% \end{macro}
% \end{macro}
%
% \begin{macro}{\displaywidowpenalty}
% \begin{macro}{\predisplaypenalty}
% \begin{macro}{\postdisplaypenalty}
%    Discourage (but not so much) widows in front of a math display
%    and forbid breaking directly in front of a display. Allow break
%    after a display without a penalty. Again the default values are
%    used, therefore we only show them here.
%    \begin{macrocode}
% \displaywidowpenalty 50
% \predisplaypenalty   10000
% \postdisplaypenalty  0
%    \end{macrocode}
% \end{macro}
% \end{macro}
% \end{macro}
%
% \begin{macro}{\interlinepenalty}
%    Allow the breaking of a page in the middle of a paragraph.
%    \begin{macrocode}
% \interlinepenalty 0
%    \end{macrocode}
% \end{macro}
%
%
% \begin{macro}{\brokenpenalty}
%    We allow the breaking of a page after a hyphenated line.
%    \begin{macrocode}
% \brokenpenalty 100
%    \end{macrocode}
% \end{macro}
%
% \subsection{Layout for Landscape Tables}%
%
% \begin{macro}{\ptlandscape}
% This command switches to a page layout appropriate to
% a table in landscape orientation.
%    \begin{macrocode}
\newcommand\ptlandscape{%
 \setlength\textwidth{9in}%
 \setlength\textheight{6.5in}%
 \setlength\pt@width{\textwidth}%
 \setlength\hoffset{0.0in}%
 \setlength\voffset{-0.25in}%
}%
%    \end{macrocode}
% \end{macro}
%
%
% \subsection{Formatting Options}
%
% \begin{macro}{\eqsecnum}
% Sometimes authors want to number equations by section,
% rather than sequentially through the whole paper.
%
% To invoke this processing option, the user may 
% give the following command in the document preamble:
% \begin{verbatim}
% \eqsecnum
% \end{verbatim}
% The preferred way is to call |eqsecnum| as a 
% document class option:
% \begin{verbatim}
% \documentclass[eqsecnum]{aastex}
% \end{verbatim}
%    \begin{macrocode}
\def\eqsecnum{%
 \@addtoreset{equation}{section}%
 \def\theequation{%
 %\mbox{\normalsize\arabic{section}-\arabic{equation}}% Wrong!
  \arabic{section}-\arabic{equation}%
 }%
}%
%    \end{macrocode}
% \end{macro}
%
% \subsection{Data Commands}%
%
% There are some obvious style-dependent parameters and formatting
% instructions in this file.  The object has been to keep the author
% commands (the user interface, if you will) society-, journal-, and
% style-independent, while the style-specific parameters and commands
% are kept below that level.
%
% \subsubsection{Article Attributes}%
%
% \begin{macro}{\@journalname}
% \begin{macro}{\cpr@holder}
% Set the journal name and copyright holder.
%    \begin{macrocode}
\def\@journalname{The Astrophysical Journal}%
\def\cpr@holder{American Astronomical Society}%
%    \end{macrocode}
% \end{macro}
% \end{macro}
%
% Journal and article identification information is established by the
% editorial staff.  The following macros are used by personnel at the
% editorial office to annotate/record slug-line data and should be 
% included in the manuscript preamble, i.e., before the \LaTeX\ 
% |\begin{document}|.
% All they do is globally define other \TeX\ control sequences that are 
% used by other macros at a later time.
% \begin{verbatim}
%    \received{<receipt date>}
%    \revised{<revision date>}
%    \accepted{<acceptance date>}
%    \journalid{<vol>}{<journal date>}
%    \articleid{<start page>}{<end page>}
%    \paperid{<id>}
%    \ccc{<code>}
% \end{verbatim}
%
% \begin{macro}{\received}
% \begin{macro}{\accepted}
% \begin{macro}{\revised}
% \begin{macro}{\journalid}
% \begin{macro}{\articleid}
% \begin{macro}{\paperid}
% \begin{macro}{\ccc}
% For preprints and submitted manuscripts in draft/referee format, etc.,
% the slug-line information is irrelevant and in those kinds of style,
% the data are never used.  In my estimation, it is perfectly acceptable
% to have these commands in the file (for future use), so they are 
% defined in the usual way.
%    \begin{macrocode}
\newcommand\received[1]{\gdef\@recvdate{#1}}%
\newcommand\accepted[1]{\gdef\@accptdate{#1}}%
\newcommand\revised[1]{\gdef\@revisedate{#1}}%
\newcommand\journalid[2]{\gdef\@jourvol{#1}\gdef\@jourdate{#2}}%
\newcommand\articleid[2]{\gdef\@startpage{#1}\gdef\@finishpage{#2}}%
\newcommand\paperid[1]{\gdef\@paperid{#1}}%
\newcommand\ccc[1]{\gdef\CCC@code{#1}}%
%    \end{macrocode}
% We provide initialization code. If
% at some future date someone manages to give these 
% macros a value before the end of the class,
% |\providecommand| will nicely bow out.
%    \begin{macrocode}
\AtEndOfClass{%
 \@provide\@recvdate{}%
 \@provide\@accptdate{}%
 \@provide\@revisedate{}%
 \@provide\@jourvol{}%
 \@provide\@jourdate{}%
 \@provide\@startpage{}%
 \@provide\@finishpage{}%
 \@provide\@paperid{MS-0001-SAMP}%
 \@provide\CCC@code{000-00\$75.95-CDB}%
}% 
%    \end{macrocode}
% \end{macro}
% \end{macro}
% \end{macro}
% \end{macro}
% \end{macro}
% \end{macro}
% \end{macro}
%
% \subsubsection{Copyright data}%
%
% \begin{macro}{\cpright}
%
% We permit the author/editor to supply a copyright ``type'' as well as 
% the year.
% The formatting of the copyright notices is defined by |\@slug|.
% FIXME: |\@slug| is never used.
%
%    \begin{macrocode}
\newcommand\cpright[2]{%
 \global\expandafter\let\expandafter\cpr@type\csname cpr@#1\endcsname
 \gdef\cpr@year{#2}%
 \typeout{`#1' copyright \cpr@year.}%
}%
\def\cpr@AAS{Copyright \cpr@year\space by the \cpr@holder.}%
\def\cpr@PD{This article is in the public domain.}%
\def\cpr@Crown{Crown copyright \cpr@year\space by the \cpr@holder.}%
\def\cpr@none{No copyright is claimed for this article.}%
\def\cpr@ASP{Copyright \cpr@year\space by the ASP.}%
%    \end{macrocode}
% Default copyright year and type:
%    \begin{macrocode}
\def\cpr@year{\number\year}%
%    \end{macrocode}
% \end{macro}
%
% \begin{macro}{\@slug}
% |\@slug| determines formatting of slugline information.
%    \begin{macrocode}
\def\@slug{%
 \par
 \noindent\cpr@type\par
 \noindent Manuscript number \@paperid.\par
 \noindent\CCC@code\par
}%
%    \end{macrocode}
% \end{macro}
%
% Here ends the code for the document class.
% Hereinafter appears code that is common to 
% the document class and document style.
%    \begin{macrocode}
%</class>
%    \end{macrocode}
%
% \subsubsection{Running heads}%
%
% The change in the user-level markup is incorporated into the V4.1
% document style.
%    \begin{macrocode}
%<*class|style>
%    \end{macrocode}
% 
% \begin{macro}{\lefthead}
% \begin{macro}{\righthead}
%
% The |\lefthead| and |\righthead| commands are now deprecated and will give
% a diagnostic message in the \TeX\ log (NB: without halting \TeX).
%    \begin{macrocode}
\newcommand\lefthead{\substitute@command\lefthead\shorttitle}%
\newcommand\righthead{\substitute@command\righthead\shortauthors}%
%    \end{macrocode}
% \end{macro}
% \end{macro}
%
% \begin{macro}{\shorttitle}
% \begin{macro}{\shortauthors}
% Authors may be asked to supply running head information.  We use the
% language of the publisher here to distinguish between two pieces of
% information: the left head frequently contains an author list (last
% names, possibly truncated as et al.), while the right head frequently
% contains an abbreviated form of the paper title.  The |\shorttitle| and
% |\shortauthors| commands merely collect the information and buffer the 
% text.
% \begin{verbatim}
%    \shorttitle{<text>}
%    \shortauthors{<text>}
% \end{verbatim}
% For the ApJ, right heads are the short titles, and the publisher 
% prefers that the author specify a 44-character or less abbreviated 
% title to use as the running head on right pages.
% 
%    \begin{macrocode}
\newcommand\shorttitle[1]{\gdef\@versohead{#1}}%
\newcommand\shortauthors[1]{\gdef\@rectohead{#1}}%
%    \end{macrocode}
%
% Default values:
%    \begin{macrocode}
\shorttitle{\relax}%
\shortauthors{\relax}%
%    \end{macrocode}
% \end{macro}
% \end{macro}
%
% Here ends the portion of code for the document style.
%    \begin{macrocode}
%</class|style>
%    \end{macrocode}
%
% Hereinafter appears code for the document class.
%    \begin{macrocode}
%<*class>
%    \end{macrocode}
%
% \begin{macro}{\@runheads}
% The |\@runheads| procedure determines the formatting of running head 
% information.
% It produces a set of running head lines, one for each page of the 
% article.
%
%    \begin{macrocode}
\def\@runheads{%
 \@tempcnta\c@page
 \@whilenum\@tempcnta>\z@\do{%
  \vskip 3ex
  \hbox\@to30pc{%
   \small
   \expandafter\uppercase\expandafter{\@versohead}:
   \expandafter\uppercase\expandafter{\@rectohead}\hfil
  }%
  \advance\@tempcnta\m@ne
 }%
}%
%    \end{macrocode}
% \end{macro}
%
% \begin{macro}{\slugcomment}
% The |\slugcomment| command is provided so authors or editors may 
% include a remark about the manuscript that may be desired in print for 
% some styles.
% For instance, authors like to include information about the
% journal and publication date of accepted papers ("To appear in AJ,
% November 1992").
% \begin{verbatim}
%    \slugcomment{<text>}
% \end{verbatim}
% The formatting of the |\slugcomment| information is controlled by the
% |\@makeslugcmmnt| command, and its placement on the page is dictated
% by the presence of |\@makeslugcmmnt| in |\sluginfo| (or whatever).
%
%    \begin{macrocode}
\newcommand\msid[1]{\gdef\ms@id{#1}}% mss ID
\def\ms@id{}%
\newcommand\slugcomment[1]{\gdef\slug@comment{#1}}%
%    \end{macrocode}
% Default value:
%    \begin{macrocode}
\slugcomment{}%
%    \end{macrocode}
% \end{macro}
%
% \begin{macro}{\@slugcmmntwidth}
% \begin{macro}{\@makeslugcmmnt}
% Low-level formatting procedure:
%    \begin{macrocode}
\newdimen\@slugcmmntwidth
\long\def\@makeslugcmmnt@mss{%
 \@ifx@empty\slug@comment{\relax}{%
  \vskip 5ex
  \sbox\@tempboxa{\slug@comment}%
  \@ifdim{\wd\@tempboxa >\@slugcmmntwidth}{%
   \hbox\@to\textwidth{%
    \hss
    \parbox\@slugcmmntwidth\slug@comment
   }%
  }{%
   \hbox\@to\textwidth{\hfil\box\@tempboxa\hfil}%
  }%
 }%
}%
\long\def\@makeslugcmmnt@ppt{%
 \@ifx@empty\slug@comment{\relax}{%
  \sbox\@tempboxa{\slug@comment}%
  \@ifdim{\wd\@tempboxa >\@slugcmmntwidth}{%
   \hbox\@to\textwidth{%
    \hss
    \parbox\@slugcmmntwidth\slug@comment
   }%
  }{%
   \hbox\@to\textwidth{\hfil\box\@tempboxa}%
  }%
  \vskip 2ex
 }%
}%
%    \end{macrocode}
% \end{macro}
%
%
% \subsubsection{Receipt and Acceptance Dates}%
% Editorial offices like to have receipt and acceptance dates on the
% front page for tracking purposes.
% The current practice is to require the author
% to type Received \dots\ and Accepted \dots\ lines
% for the staff to type on.
%
% The |\@dates| command will typeset either rules
% or the dates specified in |\received| and/or |\accepted| commands.
% Note that authors should not include |\received| and |\accepted|.
%
% \begin{macro}{\@rcvaccrule}
% This is the rule following the date.
%    \begin{macrocode}
\def\@rcvaccrule{\vrule\@width1.75in\@height0.5pt\@depth\z@}%
%    \end{macrocode}
% \end{macro}
%
% \begin{macro}{\@dates}
% |\@dates| determines formatting of manuscript date information.
%
%    \begin{macrocode}
\def\@dates{\ifx\@title\@empty\else {Received}\space%
 \ifx\@empty\@recvdate\@rcvaccrule\else\@recvdate\fi;%
 \hspace{1.5em}{accepted}\space% 
 \ifx\@empty\@accptdate\@rcvaccrule\else\@accptdate\fi%
\fi}%                                                 %        
\def\@dates@pptt{% 
 \@ifx@empty\@recvdate{}{% 
  \@ifx@empty\@accptdate{}{% 
   {\center 
     \@ifnum{\col@number=\tw@}{\small}{\normalsize}% 
     {\itshape Received}\space 
     \@recvdate 
     \@ifnum{\col@number=\tw@}{\\[.5ex]}{\hspace{1.5em}}% 
     {\itshape accepted}\space 
     \@accptdate 
    \endcenter 
   }% 
  }% 
 }% 
}% 
%    \end{macrocode}
% \end{macro}
%
% \subsubsection{Formatting the Slug Information}%
%
% \begin{macro}{\sluginfo}
% The |\sluginfo| command is executed automatically when a 
% |\begin{abstract}| is
% encountered, so it is not necessary (or appropriate) for the author to
% include this line explicitly.  The received/accepted business is a
% concern only for the editorial staff, hence we have arranged for the
% author not to have to fuss with it.
%
%    \begin{macrocode}
\def\sluginfo@mss{% 
 {% 
  \addpenalty{\@M}% 
  \addvspace{6\p@}% 
  \centering 
  \@dates 
  \@makeslugcmmnt 
  \par 
  \addvspace{6\p@}% 
 }% 
}% 
\newcommand\sluginfo{% 
 {% 
  \addpenalty{\@M}% 
  \addvspace{6\p@}% 
  \centering 
  \@dates 
  \par 
  \addvspace{6\p@}% 
 }% 
}% 
%    \end{macrocode}
% For backwards compatibility:
%    \begin{macrocode}
\let\dates=\sluginfo
%    \end{macrocode}
% \end{macro}
%
%
% \subsection{Front Matter}%
%
% \subsubsection{Abstract}%
%
% \begin{macro}{\abstract}
% The abstract environment simply formats the
% material into a certain box for later use.
% Note that the assignment of the measure
% (|\hsize|) is taken from the code
% for |\begin{document}|,
% as is the call to |\process@table|, which
% turns into a no-op after |\begin{document}|.
%    \begin{macrocode}
\renewenvironment{abstract}{%
  \iflong@abstract% Abstract in one-column mode
    \onecolumn
  \fi
 \global\setbox\abstract@box\vtop\bgroup
% \let\@nodocument\@empty
  \columnwidth\textwidth\hsize\columnwidth\linewidth\hsize
% \process@table
  \normalsize
  \color@begingroup
  \@parboxrestore
  \@setminipage
  \quotation
}{%
  \par
  \endquotation
  \removelastskip
  \vskip-\prevdepth
  \color@endgroup
 \egroup
}%
\newbox\abstract@box
%    \end{macrocode}
% Make adjustments to \LaTeX\ abstract style: optional upper case 
% heading, retain current type size (should be |\normalsize|),
% ignore twocolumn option in the |manuscript| substyle,
% add |\sluginfo|, and start abstract on new page.
%    \begin{macrocode}
\def\abstractname{ABSTRACT}%
\def\@abstract@mss{%
 \sluginfo
 \clearpage
 \@ifvoid\abstract@box{}{%
  \begin{center}\bfseries\abstractname\end{center}%
  \contribute@box\abstract@box
 }%
}%
\def\contribute@box#1{%
 \dimen@\ht#1%
 \leavevmode\hbox{\vrule\@height\ht#1\@width\z@}\par
 \nointerlineskip\nobreak\kern-\ht#1\unvbox#1\prevdepth\z@
}%
\def\@abstract@ppt{%
%\sluginfo
%\clearpage
 \@ifvoid\abstract@box{}{%
  \begin{center}\bfseries\abstractname\end{center}%
  \contribute@box\abstract@box
 }%
}%
\def\@abstract@pptt{%
 \@ifvoid\abstract@box{}{%
  \@ifnum{\col@number=\tw@}{%
   \section*{Abstract}%
  }{%
   \vspace{3em}%
   \begin{center}%
    \large
    \bfseries\abstractname\vspace{-.5em}\vspace{\z@}%
   \end{center}%
  }%
  \contribute@box\abstract@box
 }%
}%
%    \end{macrocode}
% \end{macro}
%
% \subsubsection{Title, Author, etc.}\label{sec-preamble}%
% Title and author identification is by way of standard \LaTeX\ commands
% |\title| and |\author|.
% The authors' principal affiliation is specified with a separate macro |\affil|.
% Each |\author| command should be followed by its own |\affil| (or address).
%
% Authors often have affiliations above and beyond their main employer,
% and these are specified with the |\altaffilmark| and |\altaffiltext| 
% commands.
% These commands behave like the |\footnotemark| and |\footnotetext| 
% commands
% of \LaTeX.
% |\altaffilmark| is appended to author's names in the |\author| lists,
% and generates the superscript identification numbers.
% The text for the individual alternate affiliations is generated by the
% |\altaffiltext| command.  It is up to the author to make sure that the
% |\altaffilmark| numbers attached to authors' names correspond to the 
% correct alternate affiliation.
%
% The following commands may appear in the title block,
% with the |\title| command appearing exactly once, and 
% all others appearing as many times as needed:
% \begin{verbatim}
%     \title{<lucid text>}
%     \author{<name(s)>}
%     \authoraddr{<author's postal address>}% Deprecated (V5.0) in favor of \affil
%     \email{<author's e-mail address>}
%     \affil{<affiliation>}
%     \altaffilmark{<tag number(s)>}
%     \altaffiltext{<numerical tag>}{<text>}
%     \and
% \end{verbatim}
%
% \begin{macro}{\title}
% \begin{macro}{\author}
% \begin{macro}{\authoraddr}
% \begin{macro}{\email}
% \begin{macro}{\affil}
% \begin{macro}{\altaffilmark}
% \begin{macro}{\altaffiltext}
% \begin{macro}{\and}
% The title block commands simply queue up the 
% given tokens into the appropriate control sequence.
%
% Implementation note: 
% The |\email| command has one meaning until the titlepage
% is set, then acquires a different meaning, appropriate for text use.
%    \begin{macrocode}
\def\title#1{\appdef\@title{\make@title{#1}}}%
\def\author#1{\appdef\@author{\make@author{#1}}}%
%%% \def\authoraddr{\substitute@command\authoraddr\affil}%
\def\email#1{\appdef\@author{\make@authoremail{#1}}}%
\def\affil#1{\appdef\@author{\make@affil{#1}}}%
\def\altaffilmark#1{\appdef\@author{\make@altaffilmark{#1}}}%
\def\altaffiltext#1#2{\appdef\@author{\make@altaffiltext{#1}{#2}}}%
\def\and{\appdef\@author{\make@and}}%
\let\authoraddr\@gobble
\let\@title\@empty
\let\@author\@empty
\let\@date\@empty
%    \end{macrocode}
% \end{macro}
% \end{macro}
% \end{macro}
% \end{macro}
% \end{macro}
% \end{macro}
% \end{macro}
% \end{macro}
%
% \begin{macro}{\@maketitle}
% To format the title page, simply
% expand the tokens that were saved up.
% Remember that some users incorporate
% title block commands (notably |\altaffilmark|)
% into the argument of the |\author| command itself.
% Therefore, we might encounter
% certain of the above commands while formatting
% the title page itself.
%    \begin{macrocode}
\def\@maketitle{%
 \newpage
 \begingroup
  \let\footnote\thanks
  \let\email\make@authoremail
  \let\affil\make@affil
  \let\altaffilmark\make@altaffilmark
  \let\altaffiltext\make@altaffiltext
  \let\and\make@and
  \@title
  \@author
  \@date
  \par
  \@abstract
  \@ifxundefined\keyword@list{}{%
   \expandafter\@keywords
   \expandafter{\keyword@list}%
  }%
 \endgroup
}%
%    \end{macrocode}
%
% \end{macro}
%
% The typesetting performed by the titlepage commands is appropriate
% for the top of the first text page of the manuscript.
% They could as well be used on a title cover page, but vertical 
% placement would have to be controlled by the author.
% This is not to be encouraged, since I want to discourage
% the introduction of any extra horizontal or vertical space.
%
% \begin{macro}{\make@title}
% \begin{macro}{\make@author}
% \begin{macro}{\make@authoraddr}
% \begin{macro}{\make@authoremail}
% \begin{macro}{\make@affil}
% \begin{macro}{\make@altaffilmark}
% \begin{macro}{\make@altaffiltext}
% \begin{macro}{\make@and}
% Here are the formatting procedures for the title page.
%    \begin{macrocode}
\def\make@title@mss#1{%
 {%
  \def\baselinestretch{\@tightleading}%
  \center\large\bfseries{#1}\endcenter %%
 }%
 \thispagestyle{msstitle}% mss ID
}%
\def\make@title@ppt#1{%
 \@makeslugcmmnt
 {\center\large\bfseries{#1}\endcenter}%
 \thispagestyle{empty}%
}%
\def\make@title@pptt#1{%
 {%
  \center
  \@ifnum{\col@number=\tw@}{\large}{\Large}%
  \bfseries
  #1%
  \endcenter
 }%
}%
%    \end{macrocode}
% Here are the formatting procedures for the author list.
%    \begin{macrocode}
\def\make@author#1{%
 {\topsep\z@\center\normalsize#1\endcenter}%
}%
\def\make@author@pptt#1{%
 {\topsep\z@
  \center
  \@ifnum{\col@number=\tw@}{\normalsize}{\vspace{4ex}\large}%
  #1%
  \endcenter
 }%
}%
% FIXME: produce text when in preprint substyle
\let\make@authoraddr=\@gobble
\def\make@authoremail#1{%
 {\topsep\z@\center\normalsize\tt#1\endcenter}%
}% 
\newcommand\make@affil[1]{%
 \vspace*{-2.5ex}% 
{\topsep\z@\center#1\endcenter}%
}%
\def\make@affil@ppt#1{%
 \vspace*{-0.8ex}% 
 {%
  \topsep\z@ 
  \center 
   \@ifnum{\col@number=\tw@}{\small}{\normalsize}% 
   \itshape 
   #1% 
  \endcenter 
 }%
}%
\def\make@affil@pptt#1{%
 \vspace*{-0.8ex}% 
 {%
  \topsep\z@ 
  \center 
   \@ifnum{\col@number=\tw@}{\small}{\normalsize}% 
   \itshape 
   #1% 
  \endcenter 
 }%
}%
% 
\def\thefootnote{\@arabic\c@footnote}% initial style
%
\newcommand\make@altaffilmark[1]{$^{#1}$}%
%
\newcommand\make@altaffiltext[2]{%          
\iflong@abstract% 
   \footnotetext@ass[#1]{#2}
\else
   \footnotetext@ass[#1]{\hsize\columnwidth #2}
\fi
%\stepcounter{footnote}
}%
%% 
\def\thanks#1{\footnotemark
    \protected@xdef\@thanks{\@thanks
    \iflong@abstract% 
        \protect\footnotetext[\the\c@footnote]{\hsize18pc #1}
    \else
        \protect\footnotetext[\the\c@footnote]{#1}
    \fi
    }%
}
%%
\def\footnotetext@ass{%
     \@ifnextchar [\@xfootnotenext@ass
       {%\protected@xdef\@thefnmark{\thempfn}%
    \@footnotetext}}
\def\@xfootnotenext@ass[#1]{%
  \begingroup
     %\csname c@\@mpfn\endcsname #1\relax
     \unrestored@protected@xdef\@thefnmark{#1}%
  \endgroup
  \@footnotetext}
%    \end{macrocode}
%
% Redefine the \LaTeX\ |\and| command for this style.
%    \begin{macrocode}
\def\make@and{\vspace*{-0.5ex}{\topsep\z@\center and\endcenter}}%
%    \end{macrocode}
% \end{macro}
% \end{macro}
% \end{macro}
% \end{macro}
% \end{macro}
% \end{macro}
% \end{macro}
% \end{macro}
%
% \subsubsection{Keywords}%
%
% \begin{macro}{\keywords}
% \begin{macro}{\keyword@list}
% \begin{macro}{\@keywords}
% \begin{macro}{\@keywords@pptt}
% \begin{macro}{\@keywordtext}
% \begin{macro}{\@keyworddelim}
% \begin{macro}{\subjectheadings}
% Keywords, subject headings, etc., are accommodated
% only as a piece of text; the |\keywords| command indents if necessary 
% and supplies the proper lead-in text
% (``Keywords:'', ``Subject headings:'', etc.)
% but does not apply formatting to the user's keyword text.  The author
% must supply the correct punctuation according to the journal style.
%
% Note that the editors of A\&A, ApJ, and MNRAS have discussed the issue
% of keywords, and there is disagreement about whether
% ``subject headings'' or
% ``keywords'' should appear after abstracts in astronomical papers.
% In any case, the user-level command is |\keywords|.
% 
%    \begin{macrocode}
\newcommand\keywords{\appdef\keyword@list}%
\let\keyword@list\undefined
\def\@keywords#1{%
 \vspace*{-.7ex}%
 \@ifnum{\col@number=\tw@}{%
  \noindent
  {{\itshape\@keywordtext:\/}\space\@kwds{#1}}%
 }{%
  {%
   \quote
   {\itshape\@keywordtext:\/}\space
   \@kwds{#1}%
   \endquote
  }%
 }%
}%
\def\@keywords@pptt#1{% 
 \vspace*{-.7ex}% 
 \@ifnum{\col@number=\tw@}{% 
  \noindent 
  {% 
   \small 
   {\itshape\@keywordtext:\/}\space 
   \@kwds{#1}% 
  }% 
 }{% 
  {% 
   \quote
   \small
   {\itshape\@keywordtext:\/}\space
   \@kwds{#1}%
   \endquote
  }%     
 }% 
}% 
%
\def\@keywords@mss#1{%%%   
 \vspace*{-.7ex}% 
 \@ifnum{\col@number=\tw@}{% 
  \noindent 
  {% 
   \small 
   {\itshape\@keywordtext:\/}\space 
   \@kwds{#1}% 
  }% 
 }{% 
  {% 
   \quote 
   \small 
   {\itshape\@keywordtext:\/}\space 
   \@kwds{#1}% 
   \endquote%
   \clearpage 
  }% 
 }% 
}% 
%
\let\subjectheadings=\keywords
\def\@keywordtext{Subject headings}%
\def\@keyworddelim{---}%
%    \end{macrocode}
% \end{macro}
% \end{macro}
% \end{macro}
% \end{macro}
% \end{macro}
% \end{macro}
% \end{macro}
%
% \begin{macro}{\@kwds}
% It might be possible to handle keywords and their punctuation more
% automatically at some stage, but at the moment using this looping
% construct is serious overkill.  It also causes problems because the
% loop delimiter is a comma, and there are keywords that legitimately
% contain commas.  It also seem ponderous to me for keywords in the
% |\keyword| list to be delimited by \verb+|+, so I make the author 
% responsible for keyword list punctuation.
% 
%    \begin{macrocode}
\def\@kwds@jnl#1{%
 \def\@kwddlm{}%
 \@for\@kwd:=#1\do{%
  \@kwddlm
  \def\@kwddlm{\space\@keyworddelim\penalty\@m\space}%
  {\@kwd}%
 }%
}%
\def\@kwds#1{#1\relax}%
%    \end{macrocode}
% \end{macro}
%
% \subsubsection{Automatically Triggering the Title}%
% In order to accomodate the many different formatting options
% for the title page, it is best for the user to 
% present the title block and then to 
% invoke the |\maketitle| command explicitly.
% 
% The custom in V4.0 documents has been to give 
% the commands of the title block
% as if they formatted the given material immediately---which they did.
% But this defies the \LaTeX\ markup scheme in which those commands
% are expected to act like ``data'' commands: they are legal 
% in the preamble (before |\begin{document}|).
%
% Coupled with the need to automatically switch between
% one- and two-column page grids, it would be best if the
% title block commands were to return to their nature 
% as ``data'' commands.
%
% An \aastex\ document will therefore begin with its 
% |\begin{document}| declaration and its title block
% (which could legally appear in the preamble), followed
% by an optional |\maketitle| command, succeeded by 
% optional |\tableofcontents| and similar commands, 
% succeeded by the body of the article, which 
% may begin with a sectioning command like |\section|, 
% or may even begin with paragraph text.
%
% In order to trigger at the start of paragraph text, we
% install a hook into |\everypar|. This may fail under some
% circumstances because of the violent way \LaTeX\ handles
% this token register. (We hope for the best.)
%
% Given that the abstract is formatted as part of the 
% title page, it, too, must have the status of a
% data command. However, its syntax as a \LaTeX\ environment
% prevents a simple approach using |\appdef|.
% Instead, we temporarily endow the |abstract|
% environment with the ability to set type.
% 
% The procedure |\everypar@hook| takes advantage
% of a trick for returning to vertical mode, effectively
% aborting the nascent paragraph. The only trace thereof
% in the vertical list being the |\parskip| glue.
%    \begin{macrocode}
\AtBeginDocument{%
 \everypar{%
  \everypar@hook
 }%
}%
\def\everypar@hook{%
 \setbox\z@\lastbox\par\removelastskip
 \everypar{}%
 \maketitle@trigger
 \leavevmode
}%
%    \end{macrocode}
%
% We use |\@startsection@hook| in |\@startsection| (see above).
%
%    \begin{macrocode}
\def\@startsection@hook{\maketitle@trigger}%
%    \end{macrocode}
%
% We headpatch into the |listof|\dots\ procedures:
%    \begin{macrocode}
\prepdef\tableofcontents{\listof@hook}%
\prepdef\listoffigures{\listof@hook}%
\prepdef\listoftables{\listof@hook}%
\def\listof@hook{\maketitle@trigger}%
%    \end{macrocode}
%
%    \begin{macrocode}
\def\maketitle@disarm{%
 \global\let\maketitle@trigger\relax
 \global\let\@startsection@hook\@empty
 \global\let\listof@hook\@empty
 \global\let\everypar@hook\@empty
}%
\def\maketitle@trigger{%
 \maketitle
}%
%    \end{macrocode}
%
%
% \subsection{Document Body}%
% Herein are definitions of formatting procedures that occur in the body 
% of the article.
%
% \subsubsection{Footnotes}%
%
% The following code redefines a few things so that spacing in footnotes 
% is different than \LaTeX\ defaults.  The adjustment to |\skip\footins| 
% increases
% the amount of vertical space inserted before the |\footnoterule|.  
% The commented out definitions of |\@makefntext| give proper 
% indentation.
% \begin{macro}{\footins}
% \begin{macro}{\footnotesep}
% \begin{macro}{\@makefntext}
%
% Style Note: in the two-column preprint style, these values are 
% different.
%    \begin{macrocode}
\setlength{\skip\footins}{4ex\@plus1ex\@minus.5ex}%
\setlength\footnotesep{3ex}%
%\long\def\@makefntext#1{\noindent\hb@xt@\z@{\hss$^{\@thefnmark}$}#1}
%\long\def\@makefntext#1{\indent\hbox{$^{\@thefnmark}$}#1}
\long\def\@makefntext@pptt#1{%
 \noindent
 \hbox\@to\z@{\hss$^{\@thefnmark}$}%
 #1%
}%
%    \end{macrocode}
% \end{macro}
% \end{macro}
% \end{macro}
%
% FIXME: doublecheck the |\xdef|.
% \begin{macro}{editornote}
% \begin{macro}{\theeditornote}
% \begin{macro}{\notetoeditor}
% The |\notetoeditor| command puts text into the footnote area that 
% contains some
% special instruction to the editor.
% It places a superior ``En'' at the point in
% the text where the note appears, then makes a footnote labelled
% ``En NOTE TO EDITOR: <text>''
% \begin{verbatim}
% \notetoeditor{<text>}
% \end{verbatim}
% FIXME: cf. definition of |\footnote|.
%
%    \begin{macrocode}
\newcounter{editornote}%
\def\theeditornote{{\rmfamily E}\arabic{editornote}}%
\newcommand\notetoeditor[1]{%
 \stepcounter{editornote}%
 \begingroup
  \def\protect{\noexpand\protect\noexpand}%
  \xdef\@thefnmark{\theeditornote}%
 \endgroup
 \@footnotemark\@footnotetext{NOTE TO EDITOR: #1}%
}%
%    \end{macrocode}
% \end{macro}
% \end{macro}
% \end{macro}
%
% \subsubsection{Sections}%
%
% Modify the indentation behavior so that {\em all} paragraphs, 
% including those first after section heads, are indented.
% This is consonant with the styles of many journals,
% although it's non-standard typographic practice.
%
% Style Note: In the two-column preprint style, we
% add simple logic to engage |\twocolumn| when the first |\section| 
% command is encountered if this has not been done already.
%
% \begin{macro}{\section}
% \begin{macro}{\thesection}
% Level-one head.
%    \begin{macrocode}
\renewcommand\section{%
 \@startsection{section}{1}%
  {\z@}{5ex\@plus.5ex}{1ex\@plus.2ex}{\normalsize\bfseries}%
}%
\def\thesection{\@arabic{\c@section}}%
\def\clear@section@page{}%
%    \end{macrocode}
% \end{macro}
% \end{macro}
%
%
% \begin{macro}{\subsection}
% \begin{macro}{\thesubsection}
% Level-two head.
%    \begin{macrocode}
\renewcommand\subsection{%
 \@startsection{subsection}{2}%
  {\z@}{5ex\@plus.5ex}{1ex\@plus.2ex}{\normalsize\bfseries}%
}%
\def\thesubsection{\thesection.\@arabic{\c@subsection}}%
%    \end{macrocode}
% \end{macro}
% \end{macro}
%
%
% \begin{macro}{\subsubsection}
% \begin{macro}{\thesubsubsection}
% Level-three head.
%    \begin{macrocode}
\renewcommand\subsubsection{%
 \@startsection{subsubsection}{3}%
  {\z@}{5ex\@plus.5ex}{1ex\@plus.2ex}{\normalsize\itshape}%
}%
\def\thesubsubsection{\thesubsection.\@arabic{\c@subsubsection}}%
%    \end{macrocode}
% \end{macro}
% \end{macro}
%
% \begin{macro}{\subsubsubsection}
% \begin{macro}{\thesubsubsubsection}
% Define sectioning command below |\subsubsection|
%    \begin{macrocode}
\newcommand\subsubsubsection{%
 \@startsection{subsubsection}{4}%
  {\z@}{-15\p@\@plus-5\p@\@minus-2\p@}{5\p@}{\normalfont\normalsize\itshape}%
}%
\def\thesubsubsubsection{\thesubsubsection.\@arabic{\c@subsubsubsection}}%
%    \end{macrocode}
% \end{macro}
% \end{macro}
%
% \begin{macro}{\theparagraph}
% The run-in head.
%    \begin{macrocode}
\def\theparagraph{\thesubsubsection.\@arabic{\c@paragraph}}%
%    \end{macrocode}
% \end{macro}
%
% \begin{macro}{\acknowledgments}
% \begin{macro}{\acknowledgements}
% For ApJ, acknowledgments are set off from main body text simply by
% vertical space (no heading or type size change).  I do the same for
% the AAS manuscript style.  For use in producing more polished output
% for other journals, this macro might generate text saying
% ``Acknowledg[e]ments''.
%
% An alternative spelling of the command is provided.
%
%    \begin{macrocode}
\newcommand\acknowledgments{\vskip 3ex\@plus.8ex\@minus.4ex}%
\let\acknowledgements=\acknowledgments
%    \end{macrocode}
% \end{macro}
% \end{macro}
%
%
% \subsubsection{Equation Numbering}%
%
% \begin{macro}{\appendix}
% Redefine |\appendix| so that equation numbers are prepended with 
% |\theappendix| number.  Reset equation counter so that equation
% numbers in the appendix begin again with 1.  
% The |\section| command is redefined by
% |\appendix| to clear the equation counter for each section, so
% that equation numbers will look like A1, A2,\dots, B1, B2,
%
% Note that the
% |\appendix| command is {\em not} a sectioning command; it is merely
% a delimiting markup between the main body of the paper and the
% appendix.  Sections in the appendix are still started with
% |\section| commands.  
% etc. 
%
% If we are in two-column mode at this point,
% we switch back to one column and set a flag to remind us to
% return to two-column mode upon entering the bibliography.
%
%    \begin{macrocode}
\renewcommand\appendix{%
 \par
 \if@twocolumn\@restonecoltrue\onecolumn\fi
 \setcounter{section}{\z@}%
 \setcounter{subsection}{\z@}%
 \setcounter{equation}{\z@}%
%\setcounter{figure}{\z@}%
%\setcounter{table}{\z@}%
%\def\fnum@figure{Figure \Alph{section}\thefigure:}%
%\def\fnum@table{Table \Alph{section}\thetable:}%
 \def\thesection{\Alph{section}}%
 \def\theequation{%
 %\hbox{\normalsize\Alph{section}\arabic{equation}}%
  \thesection\arabic{equation}%
 }%
 \@addtoreset{equation}{section}%
 \appendix@figtab@defs
 \def\section{%
  \@startsection
   {section}{1}{\z@}%
   {5ex\@plus.5ex}{1ex\@plus.2ex}{\normalsize\bfseries}%
   }%
}%
\let\appendix@figtab@defs\@empty
\def\appendix@figtab@defs@pptt{}%
%    \end{macrocode}
% \end{macro}
%
% \begin{macro}{mathletters}
% Define |mathletters| environment.  This is a wrapper that goes
% \emph{around} (outside of) |equation| or |eqnarray| environments.
%
%    \begin{macrocode}
\newcounter{cureqno}%
\newenvironment{mathletters}{%
 \refstepcounter{equation}%
 \setcounter{cureqno}{\value{equation}}%
 \let\@curtheeqn\theequation
 \edef\@tempa{\theequation}%
 \expandafter\def
 \expandafter\theequation
 \expandafter{\@tempa\alph{equation}}%
 \setcounter{equation}{0}%
}{%
 \let\theequation\@curtheeqn
 \setcounter{equation}{\value{cureqno}}%
}%
%    \end{macrocode}
% \end{macro}
%
% \begin{macro}{\eqnum}
% The |\eqnum| command allows the author to override all automatic
% equation marking schemes.
%
%    \begin{macrocode}
\newcommand\eqnum[1]{%
 \def\theequation{#1}%
 \let\@currentlabel\theequation
 \addtocounter{equation}{\m@ne}%
}%
%    \end{macrocode}
% \end{macro}
%
% Here ends the code for the document class.
% Hereinafter appears code that is common to 
% the document class and document style.
%    \begin{macrocode}
%</class>
%    \end{macrocode}
%
% \subsubsection{Bibliography}%
% For the time being, I don't think
% it would be a good idea to force an implementation based on BibTeX,
% although I'll prepare for that eventuality by using the constructs and
% general ideas of the \LaTeX\ |thebibliography| environment.  
% Adjustments have been made so that the |\cite|--|\bibitem| mechanism 
% can be used to cite references symbolically while maintaining proper 
% citation syntax within the paper; the author must create the citation 
% label for each reference in proper journal format in the |\bibitem| 
% command.
% \begin{verbatim}
%    \begin{thebibliography}
%    \bibitem[label]{<texttag>}
%    \reference <bibliographic data>
%    \end{thebibliography}
% \end{verbatim}
%
% Observe that the bibliographic data supplied by the author must also
% conform to the standards of the journal.  I have elected not to burden
% authors with tedious \LaTeX\ commands which would require many curly 
% braces to delimit the bibliographic fields because many of the 
% journals I have targeted in this project have agreed to reduce 
% typographic overhead (bolding, italicizing, etc.) in reference lists.  
% (See Abt's editorial in ApJ 357.)
% It is the responsibility of the author to get these fields in the 
% proper order with the correct punctuation; the information will be 
% typeset as is, i.e., in roman with no size or style changes.
%
% Macros are provided for many of the oft-referenced journals so that 
% authors may use the \LaTeX\ names rather than having to look up a 
% particular journal's
% specific abbreviation.  In principle, all the journals should be using 
% the same abbreviations, too, but I anticipate some changes in the 
% specific abbreviations before a set is finally settled on.  As long as 
% these macros are kept up to date, authors need not be concerned about 
% such editorial changes.
%
% Corresponding with the non-BibTeX references environment and 
% |\reference| command for the reference list is the |\markcite| 
% command.
% This can be  used to identify in-text citations; in this case, the 
% author provides the actual citation text as opposed to the symbolic 
% tag.
% The use of this is encouraged, even for journals where no special 
% formatting of citation call-outs is required; it is crucial for the 
% future use of on-line browsers.
%
% Note: |\markcite| is deprecated (V5.0) and has been eliminated.
%
% The entire bibliography section of code is incorporated into the
% V4.1 document style.
%    \begin{macrocode}
%<*class|style>
%    \end{macrocode}
%
% \begin{macro}{thebibliography}
% \begin{macro}{\clear@thebibliography@page}
% \begin{macro}{\thebib@list}
% \begin{macro}{\thebib@list@pptt}
%
% Note on |\clear@thebibliography@page|: this procedure is nil except in
% the manuscript substyle, where it does a |\clearpage|, and 
% the two-column preprint substyle, where it returns to two column mode.
%
%    \begin{macrocode}
\renewenvironment{thebibliography}[1]{%
 \clear@thebibliography@page
 \subsection*{REFERENCES}%
 \thebib@list
 \def\refpar{\relax}%
 \def\newblock{\hskip .11em\@plus.33em\@minus.07em}%
 \sloppy
 \clubpenalty4000
 \widowpenalty4000
 \sfcode`\.=1000\relax
}{%
 \endlist
 \revtex@pageid
}%
\def\clear@thebibliography@page{}%
\def\thebib@list{%
 \list{\null}{%
  \leftmargin 3em\labelwidth\z@\labelsep\z@\itemindent-\leftmargin
  \usecounter{enumi}%
 }%
}%
\def\thebib@list@pptt{%
 \list{\null}{%
  \leftmargin 1.2em\labelwidth\z@\labelsep\z@\itemindent-\leftmargin
  \usecounter{enumi}%
 }%
}%
%    \end{macrocode}
% \end{macro}
% \end{macro}
% \end{macro}
% \end{macro}
%
% \begin{macro}{references}
% It is not necessary to use the thebibliography environment.
% This style defines a |references| environment that simply sets off the 
% list of references and adjusts spacing parameters.
% It is not possible to use |\bibitem|s within the |references| 
% environment.
% (The |\reference| command cannot be shortened to |\ref|;
% |\ref| is used in standard \LaTeX\ for cross-referencing.)
% \begin{verbatim}
%    \begin{references}
%    \reference <bibliographic data>
%    \end{references}
% \end{verbatim}
%
%    \begin{macrocode}
\newenvironment{references}{%
 \clear@thebibliography@page
 \subsection*{REFERENCES}%
 \bgroup
  \setlength\parindent\z@
  \setlength\parskip\itemsep
  \let\refpar\references@refpar
}{%
  \refpar
 \egroup
 \revtex@pageid
}%
\def\references@refpar@mss{%
 \par\setlength\hangindent{3em}\hangafter\@ne
}%
\def\references@refpar@pptt{%
 \par\setlength\hangindent{1.2em}\hangafter\@ne
}%
\let\references@refpar\references@refpar@mss
%    \end{macrocode}
% \end{macro}
%
% \begin{macro}{\reference}
% \begin{macro}{\@reference}
% List of references.
%    \begin{macrocode}
\newcommand\reference{%
 \@ifnextchar\bgroup{\@reference}{%
  \@latexerr{Missing key on reference command}{%
   Each reference command should have a key corresponding
   to a \protect\markcite\space somewhere in the text
  }%
 }%
}%
\def\@reference#1{\relax\refpar}%
%    \end{macrocode}
% \end{macro}
% \end{macro}
% 
% \begin{macro}{\markcite}
%
% The |\markcite| command is now deprecated and will give
% a diagnostic message in the \TeX\ log (NB: without halting \TeX).
%    \begin{macrocode}
\newcommand\markcite[1]{\remove@command\markcite}%
%    \end{macrocode}
% \end{macro}
%
% \begin{macro}{\@biblabel}
% \begin{macro}{\@cite}
% Formatting choices for citations.
%
% FIXME: doublecheck compatibility with natbib
%    \begin{macrocode}
 \def\@biblabel#1{\relax}%
%\def\@cite#1#2{({#1\if@tempswa , #2\fi})}%
 \def\@cite#1#2{#1\if@tempswa , #2\fi}%
%    \end{macrocode}
% \end{macro}
% \end{macro}
%
% This completes the section of code incorporated into the
% V4.1 document style.
%    \begin{macrocode}
%</class|style>
%    \end{macrocode}
%
%
% Hereinafter appears code for the document class.
%    \begin{macrocode}
%<*class>
%    \end{macrocode}
%
% \subsubsection{Tables and Figures}%
%
% \begin{macro}{topnumber}
% Explicitly set the counter for the number of floats allowed at the
% top of the page
%
%    \begin{macrocode}
\setcounter{topnumber}{7}%
%    \end{macrocode}
% \end{macro}
%
% \begin{macro}{\tablenotemark}
% \begin{macro}{\tablenotetext}
% Footnotes for tables are to be identified with lower case alphabetics,
% rather than numeric.  |\tablenotemark| and |\tablenotetext| commands 
% are provided analogous to |\footnotemark| and |\footnotetext|, but 
% with the identifier {\em required}, like |\altaffilmark| and 
% |\altaffiltext|.
% The |\tablenotetext| must be specified before the |\end{table}|,
% since |\end{table}| displays it.
%
% \begin{verbatim}
%     \tablenotemark{<tag letter(s)>}
%     \tablenotetext{<alpha tag>}{<text>}
% \end{verbatim}
%
% Note that \m{tag letter} is identified with \m{alpha tag}
% for corresponding text.
% It is the
% responsibility of the author to get the correspondence correct.
%
% The table footnotes are coupled to the table in which they occur, 
% rather than being associated with a particular page, and they are 
% printed with the table (relatively close to the caption) instead of 
% appearing at the extreme bottom of the page.
% This is done to ensure that the notes wind up on the same page as the 
% table, since tables are floats and can migrate
% from one page to another.  There is the additional benefit that the 
% notes are unambiguously attached to the proper table, which is what is 
% required in the event that the editorial requirement of one table only 
% per page is ever changed to permit several tables per page.
%
% Footnotes are \emph{not} supported for figures.
%
% Here we give the formatting procedures for the mark, the note, and 
% the sequence of notes as a whole.
% The code for plano tables provides alternative procedures,
% albeit a common queueing mechanism is used.
%    \begin{macrocode}
\newskip\tnotemarkskip
\tnotemarkskip7pt
\newdimen\@abovenoteskip%
\newcommand\tablenotemark[1]{%
\rlap{$^{\mathrm #1}$}\hskip\tnotemarkskip\ignorespaces% Fixed: the space after notemark
}%
\def\@tablenotetext#1#2{%
 \vspace{.5ex}%
 {\noindent\llap{$^{#1}$}#2\par}%
}%
\def\@tablenotes#1{%
 \par
 \vspace{4.5ex}\footnoterule\vspace{.5ex}%
 {\footnotesize#1}%
}%
\def\@tablenotes@pptt#1{%
 \par
 \vspace{3.2ex}\footnoterule\vspace{.4ex}%
 {\footnotesize#1}%
}%
%    \end{macrocode}
% \end{macro}
% \end{macro}
%
% \begin{macro}{\tablenotetext}
% \begin{macro}{\spew@tblnotes}
%
% The command |\tablenotetext| enqueues the table note text.
%
% The tablenotes are
% dequeued onto the vertical list by |\spew@tblnotes|, which is
% executed at the end of table processing.
% Note that |\spew@tblnotes| also clears the |\tblnote@list|.
%
% Praise and honor are due to Kresten Krab Thorup (krab@iesd.auc.dk) for
% posting ftn.sty to the Usenet, from which I drew the technique.
%
%    \begin{macrocode}
\AtBeginDocument{%
 \let\tblnote@list\@empty
}%
%
\newcommand\tablenotetext[2]{\ifdim\@abovenoteskip=0pt\global\@abovenoteskip=20pt\fi% 
 \appgdef\tblnote@list{\hsize\pt@width\leftskip\z@\rightskip\z@%
 \@tablenotetext{#1}{\parfillskip\z@ plus1fil#2\endgraf}}% 
}% 
%
\def\spew@tblnotes{%
 \@ifx@empty\tblnote@list{}{%
  \@tablenotes{\tblnote@list}%
  \global\let\tblnote@list\@empty
 }%
}%
%    \end{macrocode}
% \end{macro}
% \end{macro}
%
% \begin{macro}{table}
% \begin{macro}{table*}
% Patch |\spew@tblnotes| into the end-processing for tables:
%
%    \begin{macrocode}
\prepdef\endtable{\spew@tblnotes}%
\expandafter\prepdef\csname endtable*\endcsname{\spew@tblnotes}%
%    \end{macrocode}
% \end{macro}
% \end{macro}
%
% \begin{macro}{\tableline}
% Add a |\tableline| command for use in tabular environments.  This 
% command produces the horizontal rule(s) between the column headings 
% and the body of the table.
% Authors are discouraged from using any |\hlines| themselves,
% and are typically forbidden from using vertical rules by editors.
%
%    \begin{macrocode}
\let\tableline=\hline
%    \end{macrocode}
% \end{macro}
%
% \begin{macro}{\@makecaption}
% |\@makecaption| determines formatting for figure and table captions.  
% For manuscripts, all captions should be set as flushleft paragraphs.
%
%    \begin{macrocode}
\long\def\@makecaption#1#2{\vskip 2ex\noindent#1\ #2\par}%
%    \end{macrocode}
% \end{macro}
%
% \begin{macro}{\tablenum}
% \begin{macro}{\figurenum}
% The |\tablenum| and |\figurenum| commands allows the author to
% override all automatic table- and figure marking schemes.
%
%    \begin{macrocode}
\newcommand\tablenum[1]{%
 \def\thetable{#1}%
 \let\@currentlabel\thetable
 \addtocounter{table}{\m@ne}%
}%
\newcommand\figurenum[1]{%
 \def\thefigure{#1}%
 \let\@currentlabel\thefigure
 \addtocounter{figure}{\m@ne}%
}%
%    \end{macrocode}
% \end{macro}
% \end{macro}
%
% \begin{macro}{\placetable}
% \begin{macro}{\placefigure}
% \begin{macro}{\placeplate}
% \begin{macro}{\@place@float}
% Usually, a float (table, figure, or plate) is placed in an article at 
% or near the location that it is first mentioned.
% However, there are times when an author wishes to refer to a float
% {\em before} it is placed, for instance in an introduction.
% The editors and typesetters however, may
% have trouble deciding just where the float should be placed.  The
% following commands allow the author to tell the editor and publisher 
% exactly where they want to put the item.
% It requires that a |\label{|\m{key}|}| command be
% used in the markup for the float
% (note that |\label{}| should appear {\em after} any |\caption| 
% command, so that the numbering will be correct).
%
% \begin{verbatim}
%    \placetable{<key>}
%    \placefigure{<key>}
%    \placeplate{<key>}
% \end{verbatim}
%
% Style Note: In the preprint style, the command emits no text;
% in the manuscript style, it places a note to the editors in the body 
% of the text.
%    \begin{macrocode}
\newcommand\placetable{\@place@float{TABLE}}%
\newcommand\placefigure{\@place@float{FIGURE}}%
\newcommand\placeplate{\@place@float{PLATE}}%
\def\@place@float#1#2{%
 \vspace{0.5ex}%
 \begin{center}EDITOR: PLACE #1 \ref{#2} HERE.\end{center}%
 \vspace{0.5ex}%
}%
%    \end{macrocode}
% \end{macro}
% \end{macro}
% \end{macro}
% \end{macro}
%
% \begin{macro}{\figcaption}
% \begin{macro}{\@figcaption}
% \begin{macro}{\@xfigcaption}
% Sometimes the figures are not included in the text, only the figure
% captions. To facilitate this, a |\figcaption| command is
% provided, which formats the figure captions as needed, and
% increments the figure counter automatically.
% The optional argument allows the markup to include, e.g. the filename
% of the external figure.
%
%    \begin{macrocode}
\newcommand\figcaption{\@testopt{\@xfigcaption}{}}%
\def\@xfigcaption[#1]#2{{\def\@captype{figure}\caption{#2}}}%
%    \end{macrocode}
% \end{macro}
% \end{macro}
% \end{macro}
%
% \begin{macro}{\dummytable}
% Similarly, we sometimes want to put a placeholder in a document 
% corresponding to a plano table (published elsewhere) so that cross  
% references, etc. work properly with autoincremented counters.
% Thus, for example, one can do:
% \begin{verbatim}
% \dummytable
% \label{mytable}
% \end{verbatim}
% in the document, and refer to this table by number.
%
%    \begin{macrocode}
\newcommand\dummytable{\refstepcounter{table}}%
%    \end{macrocode}
% \end{macro}
%
%
% Here ends the code for the document class.
% Hereinafter appears code that is common to 
% the document class and document style.
%    \begin{macrocode}
%</class>
%    \end{macrocode}
%
% \subsubsection{Plano Tables}%
%
% Plano tables are camera-ready pages that are included in a typeset
% article the same way as line drawings, graphs, and other camera-ready
% art, i.e., as ``planographic'' figures.  There are desirable 
% properties for a formatter to have that are somewhat above and beyond 
% \LaTeX\'s |tabular| environment that would facilitate the creation of 
% such tables.
% Among these are breaking long tables across pages, using footnotes
% in a table, specifying comments and references for tables, etc.
%
% Note for V5.0:
% Plano tables are no longer supported, but the capabilities of 
% the |deluxetable| are still part of \aastex.
% What this means for the user is that all tables should be coded
% in \LaTeX, whether using the standard \LaTeX\ mechanisms or
% using the enhancements of |deluxetable|, however, the 
% appearance of the table in the printed journal will be under the control
% of the publisher and not the author.
%
% Style Note:
% ApJ ``standard'' widths for tables are 21 and 42 picas, which are the
% widths of single and double columns in the journal.  For wide tables,
% any width less than 42 is acceptable, although using widths of 30pc
% and up will result in better appearance when centering.  Landscape
% orientation is acceptable, and I have provided the pseudo-option
% |\ptlandscape| that resets the page dimensions; the author is still
% responsible for running the local DVI translator with the proper
% parameters so that landscape pages emerge from the laser printer.
%
% Style Note:
% AJ has no ``standard'' widths for camera-ready tables,
% they simply need to be narrower than the page.
% Landscape orientation is acceptable,
% and I have provided the pseudo-option |\ptlandscape| that resets the
% page dimensions; the author is still responsible for running the
% local DVI translator with the proper parameters so that landscape
% pages actually emerge from the laser printer.
%
% What follows was inspired by supertab.sty by Theo Jurriens, 1988.
%
% We define descriptive markup for table elements.
% Observe that the formatting is embedded in the macro definitions
% so that the author doesn't need to be concerned about stylistic 
% nuances.
%
% \begin{verbatim}
%    \tablecaption{<text>}
%    \tablehead{<text>}
%    \colhead{<heading>}
%    \tabletail{<text>}
%    \tablewidth{<dimen>}
%    \tablecolumns{<num>}
%    \tableheadfrac{<num>}
%    \rotate
% \end{verbatim}
%
% The entire deluxetable section of code is incorporated into the 
% V4.1 document style.
%    \begin{macrocode}
%<*class|style>
%    \end{macrocode}
%
% \begin{macro}{\pt@box}
% \begin{macro}{\pt@width}
% \begin{macro}{\pt@column}
% \begin{macro}{\pt@nlines}
% \begin{macro}{\pt@ncol}
% \begin{macro}{\pt@page}
% Allocate registers for use in |deluxetable| (plano) tables.
%    \begin{macrocode}
\newbox\pt@box
\newdimen\pt@width
\newcount\pt@line
\newcount\pt@column
\newcount\pt@nlines
\newcount\pt@ncol
\newcount\pt@page
%    \end{macrocode}
% \end{macro}
% \end{macro}
% \end{macro}
% \end{macro}
% \end{macro}
% \end{macro}
%
% \begin{macro}{\colhead}
% \begin{macro}{\tablecolumns}
% \begin{macro}{\tablecaption}
% \begin{macro}{\tablehead}
% \begin{macro}{\tabletail}
% \begin{macro}{\tablewidth}
% \begin{macro}{\tableheadfrac}
% See |deluxetable| environment definition for |\pt@addcol|.
% The definition was moved so that |\tablecolumns| only affects the current 
% table
%
% FIXME: document syntax of |\tablehead{\colhead{<text>}}|
% and alternative of using |\tablecolumns|.
% Latter was added to assist Evan's translator, therefore is
% highly recommended.
%    \begin{macrocode}
\newcommand\colhead[1]{\multicolumn{1}{c}{#1}\pt@addcol}%
%
\gdef\pt@footnotesize{\string\footnotesize} 
\gdef\pt@scriptsize{\string\scriptsize}
\newcommand\tabletypesize[1]{\def\pt@typesize{#1}\gdef\@typesize{\string#1}
 \ifx\@typesize\pt@footnotesize
      \def\pt@headfrac{\pt@headfrac@ass@footnotesize}
  \else
      \ifx\@typesize\pt@scriptsize
          \def\pt@headfrac{\pt@headfrac@ass@scriptsize}
      \else
          \def\pt@headfrac{\pt@headfrac@ass@normalsize}% 
      \fi
\fi
}% Access to different type sizes in deluxetable 
\def\pt@typesize{}% 
%
\newcommand\tablecolumns[1]{%
 \pt@column=#1\relax%
 \pt@ncol=#1\relax%
 \global\let\pt@addcol\@empty%
}%
%
\newcommand\tablecaption[1]{\gdef\pt@caption{\tnotemarkskip8pt#1}}%Fixed: space after mark 
\newcommand\tablehead[1]{%
 \gdef\pt@head{%
  \hline\hline
  \relax\\[-1.7ex]%
  #1\hskip\tabcolsep
  \\[.7ex]%
  \hline
  \relax\\[-1.5ex]%
 }%
}%
\def\tablehead@aj#1{%
 \gdef\pt@head{%
  #1\hskip\tabcolsep
  \\[.7ex]%
  \hline
  \relax\\[-1.5ex]%
 }%
}%
\newif\if@pt@rot
\newcommand\rotate{\@pt@rottrue}%
\newcommand\tabletail[1]{\gdef\pt@tail{#1}}%
\newcommand\tablewidth[1]{\pt@width=#1\relax}%
\newcommand\tableheadfrac[1]{\gdef\pt@headfrac{#1}}%
%    \end{macrocode}
% Initialization:
%    \begin{macrocode}
\AtBeginDocument{%
 \let\pt@caption\@empty%
 \let\pt@head\@empty%
 \let\pt@tail\@empty%
 \pt@width\textwidth%
 %\def\pt@headfrac{.1}%
 \def\pt@headfrac{\pt@headfrac@ass@normalsize}% initialize typesize to consider 
}%                                            % while calculating rows
%    \end{macrocode}
% \end{macro}
% \end{macro}
% \end{macro}
% \end{macro}
% \end{macro}
% \end{macro}
% \end{macro}
%
% \begin{macro}{\pt@calcnlines}
% Calculate the number of tabular rows that will fit in the body of the
% table.  The height of the table body and the number of lines are
% calculated thus:
% \makeatletter
% \begin{eqnarray}
%  \mbox{\texttt{\@backslashchar bodyheight}}&
%  {}={}&
%  \mbox{\texttt{\@backslashchar textheight}}
%  \times (1 - \mbox{\texttt{\@backslashchar pt@headfrac}})\nonumber\\
%  \mbox{\texttt{\@backslashchar pt@nlines}}&
%  {}={}&
%  \mbox{\texttt{\@backslashchar bodyheight}}
%  / (\mbox{\texttt{\@backslashchar arraystretch}}
%  \times \mbox{\texttt{\@backslashchar baselineskip}})\nonumber
% \end{eqnarray}
% \makeatother
%
% The row height can be seen to be valid algebraically if one examines 
% the definitions for |\strutbox|, |\@setstrut|, and |\@array|.
% I could eliminate
% this calculation, as well as its attendant parameters, if it were 
% possible to determine the accumulated height of the |\halign| box as 
% \TeX\ goes along.
%
% Art says: of course this is not possible in \LaTeX\
% without supplementing it with a specialized output routine.
%    \begin{macrocode}
%
\newdimen\tabbaseskip%
%
\def\pt@calcnlines{% 
 \begingroup%
  \if@pt@rot\textheight\textwidth\fi% rotate tables
  \pt@typesize% Type sizes in deluxetable
  \@tempdima\pt@headfrac\textheight 
  \@tempdimb\textheight\advance\@tempdimb by -\@tempdima
  \@tempdima\arraystretch\baselineskip
  \global\tabbaseskip\baselineskip
  \divide\@tempdimb by\@tempdima
  \global\pt@nlines\@tempdimb
 \endgroup
}% 
%    \end{macrocode}
% \end{macro}
%
% \begin{macro}{\pt@tabacol}
% Set |\@acol|
% to |\pt@tabacol|, which is called when building the preamble for the
% |\halign| (assuming the first column format specification is |l|, 
% which it should be for ApJ tables) so that |\tabskip| within the table 
% is set to |\fill|.
%
% Doing so permits us to toggle between natural and fixed-width
% tables without requiring the user to change markup.  Note that by
% forcing |\tabskip| always to be |\fill|, I pretty well nullify any
% |@{\hspace}| commands the user puts in the template, but this 
% environment is for people who don't want to think about that sort of 
% stuff anyway.
%    \begin{macrocode}
\def\pt@tabacol{%
 \edef\@preamble{\@preamble\hskip\tabcolsep\tabskip\fill}%
}%
%    \end{macrocode}
% \end{macro}
% \end{macro}
%
% \begin{macro}{\pt@tmpcapwidth}
% \begin{macro}{\@makecaption@plano}
%
% Permit different identifying strings, one for the first page of the
% table and one for continuation pages.  \LaTeX's |\fnum@table| is set 
% to one or the other automatically in the environment.
%    \begin{macrocode}
\newdimen\pt@tmpcapwidth
\long\def\@makecaption@plano#1#2{%
 \@ifdim{\pt@width>\z@}{%
  \pt@tmpcapwidth\pt@width
 }{%
  \pt@tmpcapwidth\hsize
 }%
%\center
  \parbox{\pt@tmpcapwidth}{%
   \@make@caption@text{#1}{#2}%
  }%
 \endcenter
}%
%    \end{macrocode}
% \end{macro}
% \end{macro}
%
% \begin{macro}{\fnum@ptable}
% \begin{macro}{\fnum@ptablecont}
% \begin{macro}{\@make@caption@text}
% \begin{macro}{\@makecaption@plano@cont}
% 
% The default formatting procedures for 
% plano table captions.
%    \begin{macrocode}
\def\fnum@ptable{Table \thetable}%
\def\fnum@ptablecont{Table \thetable---{\rmfamily Continued}}%
\long\def\@make@caption@text#1#2{%
 \center\rmfamily#1.\quad#2\endcenter
}%
\long\def\@makecaption@plano@cont#1#2{%
 \center\rmfamily#1\endcenter
 \vskip 2.5ex
}%
%    \end{macrocode}
% \end{macro}
% \end{macro}
% \end{macro}
% \end{macro}
%
% \begin{macro}{\fnum@ptablecont@pptt}
% \begin{macro}{\@make@caption@text@pptt}
% \begin{macro}{\@makecaption@plano@cont@pptt}
% Formatting procedures for 
% plano table captions in the 
% preprint style.
%    \begin{macrocode}
\def\fnum@ptablecont@pptt{Table \thetable---{\itshape Continued}}%
\long\def\@make@caption@text@pptt#1#2{%
 \center\sc#1\\[.5ex]#2\endcenter
}%
\long\def\@makecaption@plano@cont@pptt#1#2{%
 \center\sc#1\endcenter
 \vskip 1.5ex
}%
%    \end{macrocode}
% \end{macro}
% \end{macro}
% \end{macro}
%
% \begin{macro}{\fnum@ptable@apj}
% \begin{macro}{\fnum@ptablecont@apj}
% \begin{macro}{\@make@caption@text@apj}
% \begin{macro}{\@makecaption@plano@cont@apj}
% Formatting procedures for 
% plano table captions in the 
% ApJ style.
%    \begin{macrocode}
\def\fnum@ptable@apj{TABLE \thetable}%
\def\fnum@ptablecont@apj{TABLE \thetable---{\itshape Continued}}%
\long\def\@make@caption@text@apj#1#2{%
 \center\Large\sc#1\\[.5ex]#2\endcenter
}%
\long\def\@makecaption@plano@cont@apj#1#2{%
 \center\Large\sc#1\endcenter
 \vskip 1.5ex
}%
%    \end{macrocode}
% \end{macro}
% \end{macro}
% \end{macro}
% \end{macro}
%
% \begin{macro}{\fnum@ptable@aj}
% \begin{macro}{\fnum@ptablecont@aj}
% \begin{macro}{\@make@caption@text@aj}
% \begin{macro}{\@makecaption@plano@cont@aj}
% Formatting procedures for 
% plano table captions in the 
% AJ style.
%    \begin{macrocode}
\def\fnum@ptable@aj{Table \thetable.}%
\def\fnum@ptablecont@aj{Table \thetable.\space{\rmfamily (continued)}}%
\long\def\@make@caption@text@aj#1#2{%
 \center\large{\sc#1}\space#2\endcenter
}%
\long\def\@makecaption@plano@cont@aj#1#2{%
 \center\large{\sc#1}\endcenter
 \vskip 1.5ex
}%
%    \end{macrocode}
% \end{macro}
% \end{macro}
% \end{macro}
% \end{macro}
%
% \begin{macro}{deluxetable}
% \begin{macro}{planotable}
% \begin{macro}{\startdata}
% \begin{macro}{\enddata}
%
% Implementation Note:
% It is desirable for the last portion of
% a multi-page table to be |\topinsert|ed, rather than |\midinsert|ed.
% At the moment, this doesn't happen.  It's not straightforward to
% make that preference known at the {\em end} of table processing.
%
% Note the context switch that provides for alternative formatting for 
% table notes in plano tables.
% 
% In the start processing of |deluxetable|, the first order
% of business is to put away the argument, which is 
% just like the argument of \LaTeX's |tabular|.
%    \begin{macrocode}
\newbox\spew@tblnotesbox%
\newbox\restof@tblnotesbox%
\global\newdimen\noteheight%
\global\newdimen\pt@tabnoteminht%
\global\pt@tabnoteminht=0pt%
\newenvironment{deluxetable}[1]{%
% Fixed: table immediately after the title
 \maketitle
 \def\pt@format{\string#1}%
%    \end{macrocode}
% Next come the custom modifications to |tabular|'s behavior,
% the row terminator that allows pagebreaks and 
% the alignment preamble parser that always sets
% |\tabskip| to |\hfill|.
% Note: |\nl| and |\nextline| are now deprecated.
%    \begin{macrocode}
 \let\@acol\pt@tabacol% 
 \let\@tabularcr\@ptabularcr%
%    \end{macrocode}
% Now come the context switches that control how
% table notes are formatted.
%    \begin{macrocode}
 \let\@tablenotetext\@ptablenotetext%
 \let\@tablenotes\@ptablenotes%
%    \end{macrocode}
% The remainder of the start processing simply initializes
% counters that are used in |deluxetable| processing.
%    \begin{macrocode}
 \global\pt@ncol\z@%
 \global\pt@column\z@%
 \global\pt@page\@ne%
 \gdef\pt@addcol{\global\advance\pt@ncol\@ne}%
 \if@pt@rot\leavevmode\fi%
}{%
%    \end{macrocode}
% End processing for |deluxetable| entails
% formatting the table notes, outputting the diagnostic
% message for the last table page and wrapping up the float.
%
% Finding the available space for tablenotes below the table body
%    \begin{macrocode}
  \ifx\@typesize\pt@footnotesize%%% considering the table typesize in calculation
    \global\pt@tabnoteminht=\pt@notemin@footnote\baselineskip
  \else
    \ifx\@typesize\pt@scriptsize
       \global\pt@tabnoteminht=\pt@notemin@script\baselineskip
    \else
       \global\pt@tabnoteminht=\pt@notemin@normal\baselineskip    
    \fi
  \fi  
 \global\advance\pt@tabnoteminht-\@abovenoteskip
 \global\@abovenoteskip=0pt
 \setbox\spew@tblnotesbox\vbox{\spew@tblnotes}%
 %%
 %% Finding the available space for tablenotes in current page   
 \if@pt@rot\global\noteheight\textwidth\else\global\noteheight\textheight\fi%
 \global\advance\noteheight-\ht\captionbox%
 \global\advance\noteheight-\pt@line\tabbaseskip
 \if@twocolumn  %% tables in preprint2 style
    \ifx\@typesize\pt@scriptsize
    \global\advance\noteheight-6\baselineskip% 
    \else
    \global\advance\noteheight-3\baselineskip% 
    \fi
  \else
    \global\advance\noteheight-2\baselineskip% 
  \fi
%    \end{macrocode}
%%
%% Print the maximum lines of notes below table within current page
%    \begin{macrocode}
   \ifdim\noteheight<\pt@tabnoteminht% if not possible to keep the minimum two lines 
     \global\noteheight=0pt%         % the whole notes move to next page.
     \global\setbox\restof@tblnotesbox\vbox{\unvbox\spew@tblnotesbox}% 
   \else
   %%
   %% if possible two lines, check the length of notes 
     \ifdim\ht\spew@tblnotesbox>\noteheight%
        % Print maximum notes in available space and store the remaining part
        \hbox\@to\hsize{\hfil\vsplit\spew@tblnotesbox to \noteheight\hfil}%
        \global\setbox\restof@tblnotesbox\vbox{\unvbox\spew@tblnotesbox}% 
     \else%
        % Print the whole notes in current page.
        \vbox{\hbox\@to\hsize{\hfil\box\spew@tblnotesbox\hfil}}%
     \fi
   \fi
 \typeout@deluxetable% 
 \endcenter% 
 \end@plano@float%
%    \end{macrocode}
%%
%% Print the remaining tablenotes into new page(s). 
%    \begin{macrocode}
%% print each page until the remaining box empty.
 \loop\ifvoid\restof@tblnotesbox\else%
 \clearpage
 %%
 %% check whether the remaining box is a full/partial page
 \ifdim\ht\restof@tblnotesbox<\if@pt@rot\hsize\else\vsize\fi% 
     \vbox to \textheight{\if@pt@rot\vfill\fi\hbox to \textwidth{\if@pt@rot\else\hfil\fi
     \if@pt@rot\rotatebox{90}{\box\restof@tblnotesbox}
     \else{\box\restof@tblnotesbox}\fi\hfil}\vfill}%
 \else%
    \vbox to \textheight{\vfill\hbox to \textwidth{\hfil%
    \if@pt@rot\rotatebox{90}{\vsplit\restof@tblnotesbox to \textwidth}
    \else{\vsplit\restof@tblnotesbox to \textheight}\fi\hfil}\vfill}%
 \fi%
 \clearpage
 \repeat% 
 \addtocounter{table}{\m@ne}%
 \tabletypesize{\normalsize}% 
}%
\let@environment{planotable}{deluxetable}%
%
\def\@plano@float{% Invoked by \startdata
 \begingroup%
  \if@pt@rot\columnwidth\textheight\fi% Deluxetable table \rotate
  \@plano@float@{table}% 
}% 
\def\end@plano@float{% S/B invoked by \enddata; instead by \enddeluxetable 
  \end@plano@float@\endgroup}% 
\def\@plano@float@{\@float}% Extra layer of abstraction for float processing 
\def\end@plano@float@{% 
\end@float}%
\newdimen\firsttabskip
\firsttabskip-\hsize 
\appdef\@endfloatbox{%
\if@pt@rot
\global\setbox\@currbox\vbox{\hskip\firsttabskip\global\firsttabskip\z@%
\rotatebox{90}{\box\@currbox}}\else\global\firsttabskip\z@\fi% 
% Fix the roate table on first page
}% 
%
\def\typeout@deluxetable@mss{%
 \typeout{%
  Table \thetable\space has been set to width \the\pt@width
 }%
}%
\def\typeout@deluxetable@ppt{%
 \typeout{%
  Page  \the\pt@page \space of table \thetable\space has been set to 
  width \the\pt@width\space with \the\pt@nlines\space lines per page
 }%
}%
% FIXME: default setting in case of unknown substyle
\let\typeout@deluxetable\typeout@deluxetable@mss
\newcommand\startdata{%
 \pt@calcnlines
 \@ifdim{\pt@width>\z@}{%
  \def\@halignto{\@to\pt@width}%
 }{%
  \def\@halignto{}%
 }%
 \let\fnum@table=\fnum@ptable
 \let\@makecaption\@makecaption@plano
 \global\pt@line\z@
 \start@pt@tabular
 \after@startdata
}%
%
\global\newbox\captionbox
\def\start@pt@tabular{\par%
 \@plano@float 
 \center 
 \expandafter\caption\expandafter{\pt@caption}% 
 \global\setbox\captionbox\vbox{\expandafter\caption\expandafter{\pt@caption}}
 \pt@typesize% Type sizes in deluxetable 
 \expandafter\@tabular\expandafter{\pt@format}% 
}%
% 
\def\set@pt@box#1{\setbox\pt@box}%
\def\after@startdata{\pt@head}%
\def\after@startdata@aj{\hline\hline\relax\\[-1.7ex]\pt@head}%
\def\enddata{%
 \crcr
 \noalign{\vskip .7ex}%
 \before@enddata% 
 \endtabular%
 \setbox\pt@box\lastbox%
 \pt@width\wd\pt@box
 \hbox \@to \hsize{\hfil\box\pt@box\hfil}%
 \ignorespaces}%
\def\before@enddata{\hline}%
\def\before@enddata@aj{\hline\hline}%
\def\nl{\substitute@command\nl\\}%
\def\nextline{\substitute@command\nextline\\}%
%    \end{macrocode}
% \end{macro}
% \end{macro}
% \end{macro}
% \end{macro}
%
% \begin{macro}{\@ptabularcr}
% \begin{macro}{\@argptabularcr}
% A different row terminator is needed so that the line counter
% can be checked prior to adding the line to the alignment box.
% This command is syntactically equivalent to |\\|.
%
% User Note: This code is syntactically the same as the \LaTeX\
% code for |\@tabularcr|, however the latter semantically ignores
% its *-form.
% In the case of |\@ptabularcr|, the * suppresses the possibility
% of a page break.
%
% Implementation Note: This code, like |\@tabularcr| relies on the 
% idiom |{\ifnum0=`}\fi|, an obscure construction that is explained
% by the \TeX book on pages 240 and 385--386.
% The need for this construct is the case where a |\\| is immediately
% followed by a |&| or a |\cr|. The construct |\ifnum0=`{\fi}| undoes
% the effect of the first.
%
%    \begin{macrocode}
\def\@ptabularcr{%
 {\ifnum0=`}\fi% A klootch just in case the next token is & or \cr
 \@ifstar{%
  \@testopt{\@argptabularcr\@empty}\z@skip%
 }{%
  \@testopt{\@argptabularcr\ptable@@split}\z@skip%
 }%
}%
%    \end{macrocode}
%
% Recalculating the lines per page if any optional spacing between rows 
%    \begin{macrocode}
\newdimen\pt@reduceline
\newdimen\extra@vspace
\extra@vspace\z@
\def\@xargptarraycr#1{\@tempdima #1%
   \global\advance\extra@vspace\@tempdima%        % adding the optional spaces
   \advance\@tempdima\dp \@arstrutbox%
   \vrule \@height\z@ \@depth\@tempdima \@width\z@%
   \global\pt@reduceline\arraystretch\tabbaseskip% total lines to be reduced
   \@whilenum\pt@reduceline<\extra@vspace%       % reducing the excess lines% 
   \do{% 
   \global\advance\extra@vspace-\pt@reduceline%
   \global\advance\pt@line1
   }% 
\cr}%
%    \end{macrocode}
%
%    \begin{macrocode}
\def\@argptabularcr#1[#2]{%
 \ifnum0=`{\fi}% To undo the effect of the klootch.
 \@ifdim{#2>\z@}{%
  \unskip\@xargptarraycr{#2}%
 }{%
  \@yargarraycr{#2}%
 }\ptable@split#1%
}%
%    \end{macrocode}
%
% \end{macro}
% \end{macro}
%
% \begin{macro}{\ptable@split}
% \begin{macro}{\ptable@@split}
% \begin{macro}{\before@suspendpt}
% \begin{macro}{\before@resumept}
% \begin{macro}{\before@suspendpt@aj}
% \begin{macro}{\before@resumept@aj}
% \begin{macro}{\typeout@pt@nl@ppt}
% \begin{macro}{\typeout@pt@nl@mss}
% \begin{macro}{\typeout@pt@nl@aj}
% The decision to break the table is made by |\ptable@split|,
% which does its calculations inside a |\noalign|,
% a trick I got from William Baxter.
% Note that if the *-form of the |\\| was used, 
% the token thrown by |\aftergroup| is simply |\@empty|,
% regardless of what the line count is;
% this is how we prevent a pagebreak at this line.
%    \begin{macrocode}
\def\ptable@split#1{%
 \noalign{%
  \global\advance\pt@line\@ne%
  \@ifnum{\pt@line<\pt@nlines}{}{%
   \aftergroup#1%
  }%
 }%
}%
\def\ptable@@split{%
  \before@suspendpt%
  \endtabular%
  \setbox\pt@box\lastbox%
  \pt@width\wd\pt@box\box\pt@box%
  \typeout@pt@nl%
  \global\advance\pt@page\@ne%
  \endcenter%
  \end@plano@float%
  \clearpage
  \global\extra@vspace\z@
  \addtocounter{table}{-2}%
  \let\fnum@table=\fnum@ptablecont%
  \let\@makecaption\@makecaption@plano@cont%
  \global\pt@ncol=\pt@column%  Either 0 or value of \tablecolumns
  \global\pt@line\z@%
  \start@pt@tabular%
  \before@resumept%
  \pt@head%
}%
\def\before@suspendpt{}%
\def\before@resumept{}%
\def\before@suspendpt@aj{\@tabularcr\noalign{\vskip .7ex}\hline}%
\def\before@resumept@aj{\hline\relax\@tabularcr[-1.7ex]}%
\def\typeout@pt@nl@ppt{%
 \typeout{%
  Page \the\pt@page\space of table \thetable\space has been set to
  width \the\pt@width
 }%
}%
\def\typeout@pt@nl@mss{%
 \typeout{%
  Page \the\pt@page\space of table \thetable\space has been set to
  width \the\pt@width\space with \the\pt@nlines\space lines per page
 }%
}%
\def\typeout@pt@nl@aj{%
 \typeout{%
  Table \thetable\space has been set to
  width \the\pt@width\space with \the\pt@nlines\space lines per page
 }%
}%
\let\typeout@pt@nl\typeout@pt@nl@mss
%    \end{macrocode}
% \end{macro}
% \end{macro}
% \end{macro}
% \end{macro}
% \end{macro}
% \end{macro}
% \end{macro}
% \end{macro}
% \end{macro}
%
% \begin{macro}{\tablevspace}
% \begin{macro}{\tablebreak}
% Author markup commands for planotable environment line breaks.
%
% In version 4.0, it was necessary to use |\nl| or |\nextline| instead 
% of |\\| and to employ an alternative mechanism for introducing extra 
% vertical space:
% \begin{verbatim}
%    \nl == \nextline   Used in place of LaTeX's \\
%    \tablevspace{<dimen>}
%    \tablebreak
% \end{verbatim}
%
% In this version of \aastex\, the user can simply employ \LaTeX's |\\| 
% as per usual.
% The user can supply an optional argument to govern the vertical 
% spacing as usual, and the *-form prevents a pagebreak.
%
% Consequently, |\nl|, |\nextline|, and |\tablevspace| are now deprecated.
%
% Note that there is a significant problem with the |\tablevspace| command that 
% recommends against its use:
% if the table splits at that point, the extra vertical space 
% will fall at the top of the next page's table.
%    \begin{macrocode}
\newcommand\tablevspace[1]{\substitute@command\tablevspace\\[#1]}%
%    \end{macrocode}
% In order to force a |deluxetable| to split (over pages), the user
% gives the |\tablebreak| command.
%
% Art says: As a means of controlling page breaks,
% I prefer to see the user employ the *-form of the 
% row terminator |\\| to prevent the indicated line
% from splitting from the next line.
% Besides, there is a drawback in using |\tablebreak|, namely
% that it is syntactically different from |\\|, and
% unlike any extant \LaTeX\ capability.
% The ``\LaTeX\ way'' of forcing a page break would be 
% to use the |\pagebreak| or |\newpage| commands.
% So I deprecate |\tablebreak|.
%    \begin{macrocode}
\newcommand\tablebreak{\cr\ptable@@split}%\\{\cr\ptable@@split}%
%    \end{macrocode}
% \end{macro}
% \end{macro}
%
% \begin{macro}{\cutinhead}
% \begin{macro}{\cutinhead@ppt}
% \begin{macro}{\sidehead}
% We have some specialty heads that are sometimes wanted in tables.
% Art says: Why do these use |\@ptabularcr|?
% \begin{verbatim}
%    \cutinhead{<text>}
%    \sidehead{<text>}
% \end{verbatim}
%    \begin{macrocode}
\newcommand\cutinhead[1]{%
 \noalign{\vskip 1.5ex}%
 \hline
 \@ptabularcr
 \noalign{\vskip -4ex}%
 \multicolumn{\pt@ncol}{c}{#1}%
 \@ptabularcr
 \noalign{\vskip .8ex}%
 \hline
 \@ptabularcr
 \noalign{\vskip -2ex}%
}%
\def\cutinhead@ppt#1{%
 \noalign{\vskip 1.5ex}%
 \hline
 \@ptabularcr
 \noalign{\vskip -2ex}% Style Note: in apj, it is -1.5ex
 \multicolumn{\pt@ncol}{c}{#1}%
 \@ptabularcr
 \noalign{\vskip .8ex}%
 \hline
 \@ptabularcr
 \noalign{\vskip -2ex}%
}%
\newcommand\sidehead[1]{%
 \noalign{\vskip 1.5ex}%
 \multicolumn{\pt@ncol}{@{\hskip\z@}l}{#1}%
 \@ptabularcr
 \noalign{\vskip .5ex}%
}%
%    \end{macrocode}
% \end{macro}
% \end{macro}
% \end{macro}
%
% \begin{macro}{\@ptablenotetext}
% \begin{macro}{\@ptablenotes}
% 
% Here we provide for the alternative formatting procedures for
% table notes in plano tables. 
% Note that the |\tablenotemark| does not change.
%    \begin{macrocode}
\def\@ptablenotetext#1#2{%
 \vspace{0ex}% Style Note: in ppt, it is gone
 {%\parbox{\pt@width}%
 {\hskip1em$^{\mathrm#1}$#2}\par}% 
}%
\def\@ptablenotes#1{%
 \par
 \vspace{2ex}%
 {\setlength\parskip{1.5ex}#1}%
}%
\def\@ptablenotes@apj#1{%
 \par
 \vspace{2ex}%
 {#1}%
}%
%    \end{macrocode}
% \end{macro}
% \end{macro}
%
% \begin{macro}{\tablerefs}
% \begin{macro}{\@tableref}
% \begin{macro}{\tablecomments}
% \begin{macro}{\@tablecom}
% Special kinds of table footnotes.  Sometimes authors tabulate things
% which have corresponding references, and it may be desirable to 
% associate these references with the table rather than
% (or in addition to) the formal reference list.
% Occasionally, authors wish to append a short
% paragraph of explanatory notes that pertain to the entire table, but
% which are different than the caption.
% \begin{verbatim}
%    \tablerefs{<reference list>}
%    \tablecomments{<text>}
% \end{verbatim}
%
% In implementation, we enqueue and dequeue using the same
% mechanism as the table notes.
%    \begin{macrocode}
%
\newcommand\tablerefs[1]{\ifdim\@abovenoteskip=0pt\global\@abovenoteskip=10pt\fi
\appgdef\tblnote@list{\hsize\pt@width\leftskip\z@\rightskip\z@%
\@tableref{\parfillskip\z@ plus1fil#1\endgraf}}}% 
\def\@tableref#1{% 
 \par 
 \vspace*{3ex}% 
 {%\parbox{\pt@width} %%%% 
 {\hskip1em\rmfamily References. --- #1}\par}% 
}% 
\newcommand\tablecomments[1]{\ifdim\@abovenoteskip=0pt\global\@abovenoteskip=10pt\fi
\appgdef\tblnote@list{\hsize\pt@width\leftskip\z@\rightskip\z@%
\@tablecom{\parfillskip\z@ plus1fil#1\endgraf}}}% 
\def\@tablecom#1{% 
 \par 
 \vspace*{3ex}% 
 {%\parbox{\pt@width} %%% 
{\hskip1em\rmfamily Note. --- #1}\par}% 
}% 
%    \end{macrocode}
% \end{macro}
% \end{macro}
% \end{macro}
% \end{macro}
%
% \begin{macro}{\@tableref@pptt}
% \begin{macro}{\@tablecom@pptt}
% Table references and comments in the two-column preprint style:
%    \begin{macrocode}
\def\@tableref@pptt#1{% 
 \par 
 \vspace*{3ex}%
 {%\parbox{\pt@width} %%% 
 {\hskip1em{\sc References.---}#1}\par}% 
}% 
\def\@tablecom@pptt#1{% 
 \vspace*{3ex}{%
  %\parbox{\pt@width} %%% 
  {\hskip1em{\sc Note.---}#1}\par}% 
}% 
%    \end{macrocode}
% \end{macro}
% \end{macro}
%
% Here ends the section of code incorporated into the 
% V4.1 document style.
%    \begin{macrocode}
%</class|style>
%    \end{macrocode}
%
% Hereinafter appears code for the document class.
%    \begin{macrocode}
%<*class>
%    \end{macrocode}
%
% \subsubsection{Float Parameters}%
%
% \begin{macro}{plate}
% Create a new type of float for ``plates''.
% Note that the |\ftype@plate| must be a power of 2, 
% and that plates are put into the LOF (via |\ext@plate|).
%
%    \begin{macrocode}
\newcounter{plate}%
\def\ftype@plate{4}%
\def\ext@plate{lof}%
\newenvironment{plate}{\@float{plate}}{\end@float}%
\newenvironment{plate*}{\@dblfloat{plate}}{\end@dblfloat}%
\let\platewidth=\tablewidth
\newcommand\platenum[1]{%
 \def\theplate{#1}%
 \let\@currentlabel\theplate
 \addtocounter{plate}{\m@ne}%
}%
%    \end{macrocode}
% \end{macro}
%
% \begin{macro}{\thefigure}
% \begin{macro}{\thetable}
% \begin{macro}{\theplate}
% Float numbering is arabic.
%
%    \begin{macrocode}
\def\thefigure{\@arabic\c@figure}%
\def\thetable{\@arabic\c@table}%
\def\theplate{\@arabic\c@plate}%
%    \end{macrocode}
% \end{macro}
% \end{macro}
% \end{macro}
%
% \begin{macro}{\fnum@figure}
% \begin{macro}{\fnum@table}
% \begin{macro}{\fnum@plate}
%
% Formatting of the float number:
%    \begin{macrocode}
\def\fnum@figure{{\rmfamily Fig.\space\thefigure.---}}%
\def\fnum@table{{\rmfamily Table \thetable:}}%
\def\fnum@plate{{\bfseries Plate \theplate.}}%
%    \end{macrocode}
% \end{macro}
% \end{macro}
% \end{macro}
%
% \begin{macro}{\fps@figure}
% \begin{macro}{\fps@table}
% \begin{macro}{\fps@plate}
%
% Float positions:
%    \begin{macrocode}
%
\def\fps@figure{bp}%
\def\fps@table{bp}%
\def\fps@plate{bp}%
%    \end{macrocode}
% \end{macro}
% \end{macro}
% \end{macro}
%
%
% \subsubsection{EPS Inclusions}%
% FIXME: requires package epsf.sty
%
% FIXME: prefer |\includegraphics|
%
% FIXME: implement autoscaling in |\plotone| and |\plottwo|.
%
% \begin{macro}{\eps@scaling}
% \begin{macro}{\epsscale}
% \begin{macro}{\plotone}
% \begin{macro}{\plottwo}
% \begin{macro}{\plotfiddle}
% Support EPS inclusion with the native \LaTeX\ |\includegraphics| 
% command.
% A large number of DVI translators are supported by this mechanism,
% although certain other features of \aastex, notably turnpage tables,
% are supported only by PostScript output devices.
% \begin{verbatim}
%    \plotone{<epsfile>}
%    \plottwo{<epsfile>}{<epsfile>}
%    \plotfiddle{<epsfile>}{<vsize>}{<rot>}{<hsf>}
%     {<vsf>}{<htrans>}{<vtrans>}
%    <epsfile>    name of EPS file
% \end{verbatim}
% |\plotone| inserts the plot in a space that is |\columnwidth| wide;
% |\plottwo| inserts two plots side by side on the same line.
%
% The |\plotfiddle| command now 
% uses only its first six arguments;
% only |voffset| is ignored.
% The parameters |hsf| and |vsf| serves as a scale factor for both
% horizontal and vertical.
%
%    \begin{macrocode}
%\epsfverbosetrue
\def\eps@scaling{1.0}%
\newcommand\epsscale[1]{\gdef\eps@scaling{#1}}%
% 28 Oct 2003 SZ
% added typeout commands to \plotone, \plottwo, \plotfiddle
% this was a change made to the .cls for WPR
% that was not picked up in the .dtx
\newcommand\plotone[1]{%
 \typeout{Plotone included the file #1}
 \centering
 \leavevmode
 \includegraphics[width={\eps@scaling\columnwidth}]{#1}%
}%
\newcommand\plottwo[2]{{%
 \typeout{Plottwo included the files #1 #2}
 \centering
 \leavevmode
 \columnwidth=.45\columnwidth
 \includegraphics[width={\eps@scaling\columnwidth}]{#1}%
 \hfil
 \includegraphics[width={\eps@scaling\columnwidth}]{#2}%
}}%
%
\def\plotfiddle#1#2#3#4#5#6#7{{
\typeout{Plotfiddle included the file #1}
\centering\leavevmode%% Re-implement from v4.0
\vbox to #2{\rule{0pt}{#2}}
\hspace*{#6pt}\includegraphics[width=#4pt, height=#5pt, angle=#3, origin=c]{#1}}%
}
%
%    \end{macrocode}
% \end{macro}
% \end{macro}
% \end{macro}
% \end{macro}
% \end{macro}
%
%
% \subsubsection{Abbreviations for Journals}%
%
% The object here is to provide authors
% with convenient shorthands for the most ``popular'' (often-cited)
% journals; the author can use these markup tags without being concerned
% about the exact form of the journal abbreviation, or its formatting.
% It is up to the keeper of the macros to make sure the macros expand
% to the proper text.  If macro package writers agree to all use the
% same TeX command name, authors only have to remember one thing, and
% the style file will take care of editorial preferences.  This also
% applies when a single journal decides to revamp its abbreviating
% scheme, as happened with the ApJ (Abt 1991).
%
%    \begin{macrocode}
\let\jnl@style=\rmfamily
\def\ref@jnl#1{{\jnl@style#1}}%
\newcommand\aj{\ref@jnl{AJ}}%
          % Astronomical Journal
\newcommand\actaa{\ref@jnl{Acta Astron.}}%
	  % Acta Astronomica
\newcommand\araa{\ref@jnl{ARA\&A}}%
          % Annual Review of Astron and Astrophys
\newcommand\apj{\ref@jnl{ApJ}}%
          % Astrophysical Journal
\newcommand\apjl{\ref@jnl{ApJ}}%
          % Astrophysical Journal, Letters
\newcommand\apjs{\ref@jnl{ApJS}}%
          % Astrophysical Journal, Supplement
\newcommand\ao{\ref@jnl{Appl.~Opt.}}%
          % Applied Optics
\newcommand\apss{\ref@jnl{Ap\&SS}}%
          % Astrophysics and Space Science
\newcommand\aap{\ref@jnl{A\&A}}%
          % Astronomy and Astrophysics
\newcommand\aapr{\ref@jnl{A\&A~Rev.}}%
          % Astronomy and Astrophysics Reviews
\newcommand\aaps{\ref@jnl{A\&AS}}%
          % Astronomy and Astrophysics, Supplement
\newcommand\azh{\ref@jnl{AZh}}%
          % Astronomicheskii Zhurnal
\newcommand\baas{\ref@jnl{BAAS}}%
          % Bulletin of the AAS
\newcommand\caa{\ref@jnl{Chinese Astron. Astrophys.}}%
	  % Chinese Astronomy and Astrophysics	  
\newcommand\cjaa{\ref@jnl{Chinese J. Astron. Astrophys.}}%
	  % Chinese Journal of Astronomy and Astrophysics	  
\newcommand\icarus{\ref@jnl{Icarus}}%
	  % Icarus
\newcommand\jcap{\ref@jnl{J. Cosmology Astropart. Phys.}}%
	  % Journal of Cosmology and Astroparticle Physics
\newcommand\jrasc{\ref@jnl{JRASC}}%
          % Journal of the RAS of Canada
\newcommand\memras{\ref@jnl{MmRAS}}%
          % Memoirs of the RAS
\newcommand\mnras{\ref@jnl{MNRAS}}%
          % Monthly Notices of the RAS
\newcommand\na{\ref@jnl{New A}}%
	  % New Astronomy
\newcommand\nar{\ref@jnl{New A Rev.}}%
	  % New Astronomy Review
\newcommand\pra{\ref@jnl{Phys.~Rev.~A}}%
          % Physical Review A: General Physics
\newcommand\prb{\ref@jnl{Phys.~Rev.~B}}%
          % Physical Review B: Solid State
\newcommand\prc{\ref@jnl{Phys.~Rev.~C}}%
          % Physical Review C
\newcommand\prd{\ref@jnl{Phys.~Rev.~D}}%
          % Physical Review D
\newcommand\pre{\ref@jnl{Phys.~Rev.~E}}%
          % Physical Review E
\newcommand\prl{\ref@jnl{Phys.~Rev.~Lett.}}%
          % Physical Review Letters
\newcommand\pasa{\ref@jnl{PASA}}%
	  % Publications of the Astron. Soc. of Australia
\newcommand\pasp{\ref@jnl{PASP}}%
          % Publications of the ASP
\newcommand\pasj{\ref@jnl{PASJ}}%
          % Publications of the ASJ
\newcommand\qjras{\ref@jnl{QJRAS}}%
          % Quarterly Journal of the RAS
\newcommand\rmxaa{\ref@jnl{Rev. Mexicana Astron. Astrofis.}}%
	  % Revista Mexicana de Astronomia y Astrofisica
\newcommand\skytel{\ref@jnl{S\&T}}%
          % Sky and Telescope
\newcommand\solphys{\ref@jnl{Sol.~Phys.}}%
          % Solar Physics
\newcommand\sovast{\ref@jnl{Soviet~Ast.}}%
          % Soviet Astronomy
\newcommand\ssr{\ref@jnl{Space~Sci.~Rev.}}%
          % Space Science Reviews
\newcommand\zap{\ref@jnl{ZAp}}%
          % Zeitschrift fuer Astrophysik
\newcommand\nat{\ref@jnl{Nature}}%
          % Nature
\newcommand\iaucirc{\ref@jnl{IAU~Circ.}}%
          % IAU Cirulars
\newcommand\aplett{\ref@jnl{Astrophys.~Lett.}}%
          % Astrophysics Letters and Communications
\newcommand\apspr{\ref@jnl{Astrophys.~Space~Phys.~Res.}}%
          % Astrophysics Space Physics Research
\newcommand\bain{\ref@jnl{Bull.~Astron.~Inst.~Netherlands}}%
          % Bulletin Astronomical Institute of the Netherlands
\newcommand\fcp{\ref@jnl{Fund.~Cosmic~Phys.}}%
          % Fundamental Cosmic Physics
\newcommand\gca{\ref@jnl{Geochim.~Cosmochim.~Acta}}%
          % Geochimica Cosmochimica Acta
\newcommand\grl{\ref@jnl{Geophys.~Res.~Lett.}}%
          % Geophysics Research Letters
\newcommand\jcp{\ref@jnl{J.~Chem.~Phys.}}%
          % Journal of Chemical Physics
\newcommand\jgr{\ref@jnl{J.~Geophys.~Res.}}%
          % Journal of Geophysical Research
\newcommand\jqsrt{\ref@jnl{J.~Quant.~Spec.~Radiat.~Transf.}}%
          % Journal of Quantitiative Spectroscopy and Radiative Trasfer
\newcommand\memsai{\ref@jnl{Mem.~Soc.~Astron.~Italiana}}%
          % Mem. Societa Astronomica Italiana
\newcommand\nphysa{\ref@jnl{Nucl.~Phys.~A}}%
          % Nuclear Physics A
\newcommand\physrep{\ref@jnl{Phys.~Rep.}}%
          % Physics Reports
\newcommand\physscr{\ref@jnl{Phys.~Scr}}%
          % Physica Scripta
\newcommand\planss{\ref@jnl{Planet.~Space~Sci.}}%
          % Planetary Space Science
\newcommand\procspie{\ref@jnl{Proc.~SPIE}}%
          % Proceedings of the SPIE
\let\astap=\aap
\let\apjlett=\apjl
\let\apjsupp=\apjs
\let\applopt=\ao
%    \end{macrocode}
%
%
% Here ends the code for the document class.
% Hereinafter appears code that is common to 
% the document class and document style.
%    \begin{macrocode}
%</class>
%    \end{macrocode}
%
% \subsubsection{Text Entities}%
%
% The entire text entities section of code is incorporated into the 
% V4.1 document style.
%    \begin{macrocode}
%<*class|style>
%    \end{macrocode}
%
% Handy little things---that everybody works out for themselves anyway!
% Many of these come from Springer's A\&A package, and some were 
% contributed by Francois Schweizer at DTM.
% You can go berserk making these up,
% especially when you start getting into the composite ones.  We have
% tried to select a tractable number that were useful, and somewhat
% difficult to get right because fussy kerning or some such is required.
% Most can be used in or out of math mode with impunity; |\alt| and 
% |\agt| are relations and can only be used in math mode.
%
% The first group consists of
% various kinds of phantoms used for aligning columns in tables;
% |\phn| is a phantom numeral,
% |\phd|    a phantom decimal,
% |\phs|    a phantom minus sign, and
% |\phm| is a generic phantom.
%
%    \begin{macrocode}
\newcommand\phn{\phantom{0}}%
\newcommand\phd{\phantom{.}}%
\newcommand\phs{\phantom{$-$}}%
\newcommand\phm[1]{\phantom{#1}}%
%    \end{macrocode}
% 
% Wait! These two macros are defined in \file{amssymb.sty}!
%    \begin{macrocode}
%\newcommand\lesssim{%
% \mathrel{\mbox{\rlap{\mbox{\lower4pt\hbox{$\sim$}}}\mbox{$<$}}}%
%}%
%\newcommand\gtrsim{%
% \mathrel{\mbox{\rlap{\mbox{\lower4pt\hbox{$\sim$}}}\mbox{$>$}}}%
%}%
%
\let\la=\lesssim            % For Springer A&A compliance...
\let\ga=\gtrsim
%    \end{macrocode}
% 
% Various and sundry shorthand for text entities:
%    \begin{macrocode}
\newcommand\sq{\mbox{\rlap{$\sqcap$}$\sqcup$}}%
\newcommand\arcdeg{\mbox{$^\circ$}}%
\newcommand\arcmin{\mbox{$^\prime$}}%
\newcommand\arcsec{\mbox{$^{\prime\prime}$}}%
\newcommand\fd{\mbox{$.\!\!^{\mathrm d}$}}%
\newcommand\fh{\mbox{$.\!\!^{\mathrm h}$}}%
\newcommand\fm{\mbox{$.\!\!^{\mathrm m}$}}%
\newcommand\fs{\mbox{$.\!\!^{\mathrm s}$}}%
\newcommand\fdg{\mbox{$.\!\!^\circ$}}%
\newcommand\farcm@mss{\mbox{$.\mkern-4mu^\prime$}}%
\let\farcm\farcm@mss
\newcommand\farcs@mss{\mbox{$.\!\!^{\prime\prime}$}}%
\let\farcs\farcs@mss
\newcommand\fp{\mbox{$.\!\!^{\scriptscriptstyle\mathrm p}$}}%
\newcommand\micron{\mbox{$\mu$m}}%
%    \end{macrocode}
%
% Alternative definitions for ApJ
%    \begin{macrocode}
\def\farcm@apj{%
 \mbox{.\kern -0.7ex\raisebox{.9ex}{\scriptsize$\prime$}}%
}%
\def\farcs@apj{%
 \mbox{%
  \kern  0.13ex.%
  \kern -0.95ex\raisebox{.9ex}{\scriptsize$\prime\prime$}%
  \kern -0.1ex%
 }%
}%
%    \end{macrocode}
%
%
% \begin{macro}{\case}
% Permit author to typeset ``case'' fractions.  This is sometimes
% wanted in displayed equations, during which \LaTeX\ will set fractions
% specified as |\frac{x}{y}| as ``built-up'' fractions (numerator and
% denominator at body text size).
%
%    \begin{macrocode}
\newcommand\case[2]{\mbox{$\frac{#1}{#2}$}}%
%    \end{macrocode}
% \end{macro}
%
% \begin{macro}{\slantfrac}
% \begin{macro}{\onehalf}
% \begin{macro}{\onethird}
% \begin{macro}{\twothirds}
% \begin{macro}{\onequarter}
% \begin{macro}{\threequarters}
% \begin{macro}{\ubvr}
% \begin{macro}{\ub}
% \begin{macro}{\bv}
% \begin{macro}{\vr}
% \begin{macro}{\ur}
% Permit author to typeset ``text fractions''---the |\slantfrac|
% command should be used in text mode, \emph{not} in math mode.  
% Note that this is different from a ``shilling'' fraction, which the 
% author can produce without any special formatting markup.
% Define markup shorthands for several common fractions using |\slantfrac|.
%
% The |\slantfrac| command is reimplemented in this version of \aastex---
% it used to produce a ``special fraction'', now a case fraction.
% The old definition appears below for reference.
%
% I do not understand the |\ubvr| commands.
%    \begin{macrocode}
%\newcommand\slantfrac[2]{\mbox{$\,^#1\!/_#2$}}%
\newcommand\slantfrac{\case}%
\newcommand\onehalf{\slantfrac{1}{2}}%
\newcommand\onethird{\slantfrac{1}{3}}%
\newcommand\twothirds{\slantfrac{2}{3}}%
\newcommand\onequarter{\slantfrac{1}{4}}%
\newcommand\threequarters{\slantfrac{3}{4}}%
%
\newcommand\ubvr{\mbox{$U\!BV\!R$}}%% UBVR system
\newcommand\ub{\mbox{$U\!-\!B$}}%   % U-B
\newcommand\bv{\mbox{$B\!-\!V$}}%   % B-V
\newcommand\vr{\mbox{$V\!-\!R$}}%   % V-R
\newcommand\ur{\mbox{$U\!-\!R$}}%   % U-R
%    \end{macrocode}
% \end{macro}
% \end{macro}
% \end{macro}
% \end{macro}
% \end{macro}
% \end{macro}
% \end{macro}
% \end{macro}
% \end{macro}
% \end{macro}
% \end{macro}
%
% \begin{macro}{\ion}
% Notation for atomic species (ionization levels).  The ionization state
% is specified as the second argument, and should be given as a numeral.
% The macro has to expand the numeric state into the proper notation for
% the publication (roman, numeric, plus signs, etc.)
% \begin{verbatim}
%    \ion{<element>}{<ionization state>}
% \end{verbatim}
%    \begin{macrocode}
\newcommand\ion[2]{#1$\;${\small\rmfamily\@Roman{#2}}\relax}%
%    \end{macrocode}
% \end{macro}
%
% \begin{macro}{\nodata}
% Centered ellipsis for use in tables (for unknown values).  |\nodata|
% is generalized markup for this notion: format of null-valued entries
% in tables is style-specific.
%
%    \begin{macrocode}
%\newcommand\cellipsis{\hfill$\cdots$\hfill}%
%\newcommand\nodata{\multicolumn{1}{c}{$\cdots$}}%
\newcommand\nodata{ ~$\cdots$~ }%
%    \end{macrocode}
% \end{macro}
%
% \begin{macro}{\diameter}
% The |\diameter| command produces a circle with a slash through it;
% added for compatibility with A\&A and MNRAS packages.
%    \begin{macrocode}
\newcommand\diameter{\ooalign{\hfil/\hfil\crcr\mathhexbox20D}}%
%    \end{macrocode}
% \end{macro}
%
% \begin{macro}{\degr}
% The |\degr| command is a synonym for |\arcdeg|;
% added for compatibility with A\&A and MNRAS packages.
%    \begin{macrocode}
\newcommand\degr{\arcdeg}%
%    \end{macrocode}
% \end{macro}
%
% \begin{macro}{\Sun}%
% \begin{macro}{\Sol}%
% \begin{macro}{\sun}%
% \begin{macro}{\Mercury}%
% \begin{macro}{\Venus}%
% \begin{macro}{\Earth}%
% \begin{macro}{\Mars}%
% \begin{macro}{\Jupiter}%
% \begin{macro}{\Saturn}%
% \begin{macro}{\Uranus}%
% \begin{macro}{\Neptune}%
% \begin{macro}{\Pluto}%
% \begin{macro}{\Moon}%
% \begin{macro}{\ascnode}%
% \begin{macro}{\descnode}%
%  Names for the planets and ascending node and descending node.
%    \begin{macrocode}
\newcommand\Sun{\sun}% Sun symbol, "S"
\newcommand\Sol{\sun}%
\newcommand\sun{\odot}%
\newcommand\Mercury{\astro{\char1}}% Mercury symbol, "1"		
\newcommand\Venus{\astro{\char2}}% Venus symbol, "2"		
\newcommand\Earth{\earth}% Earth symbol, "3"		
\newcommand\Terra{\earth}%
\newcommand\earth{\oplus}%
\newcommand\Mars{\astro{\char4}}% Mars symbol, "4"		
\newcommand\Jupiter{\astro{\char5}}% Jupiter symbol, "5"		
\newcommand\Saturn{\astro{\char6}}% Saturn symbol, "6"		
\newcommand\Uranus{\astro{\char7}}% Uranus symbol, "7"		
\newcommand\Neptune{\astro{\char8}}% Neptune symbol, "8"		
\newcommand\Pluto{\astro{\char9}}% Pluo symbol, "9"		
\newcommand\Moon{\astro{\char10}}% Moon symbol, "M"		
\newcommand\Luna{\Moon}%
%\def\ascnode{\astro{N}}% ascending node (dragon's head), (only in cmastro)
%\def\descnode{\astro{D}}% descending node (dragon's tail), (only in cmastro)
%    \end{macrocode}
% \end{macro}
% \end{macro}
% \end{macro}
% \end{macro}
% \end{macro}
% \end{macro}
% \end{macro}
% \end{macro}
% \end{macro}
% \end{macro}
% \end{macro}
% \end{macro}
% \end{macro}
% \end{macro}
% \end{macro}
%
% \begin{macro}{\Aries}%
% \begin{macro}{\VEq}%
% \begin{macro}{\Taurus}%
% \begin{macro}{\Gemini}%
% \begin{macro}{\Cancer}%
% \begin{macro}{\Leo}%
% \begin{macro}{\Virgo}%
% \begin{macro}{\Libra}%
% \begin{macro}{\AEq}%
% \begin{macro}{\Scorpius}%
% \begin{macro}{\Sagittarius}%
% \begin{macro}{\Capricornus}%
% \begin{macro}{\Aquarius}%
% \begin{macro}{\Pisces}%
% Zodiacal symbols, and autumnal and vernal equinox.
%    \begin{macrocode}
\newcommand\Aries{\astro{\char11}}%
\newcommand\VEq{\Aries}% vernal equinox (Aries)
\newcommand\Taurus{\astro{\char12}}%
\newcommand\Gemini{\astro{\char13}}%
\newcommand\Cancer{\astro{\char14}}%
\newcommand\Leo{\astro{\char15}}%
\newcommand\Virgo{\astro{\char16}}%
\newcommand\Libra{\astro{\char17}}%
\newcommand\AEq{\Libra}% autumnal equinox (Libra)
\newcommand\Scorpius{\astro{\char18}}%
\newcommand\Sagittarius{\astro{\char19}}%
\newcommand\Capricornus{\astro{\char20}}%
\newcommand\Aquarius{\astro{\char21}}%
\newcommand\Pisces{\astro{\char22}}%
%    \end{macrocode}
% \end{macro}
% \end{macro}
% \end{macro}
% \end{macro}
% \end{macro}
% \end{macro}
% \end{macro}
% \end{macro}
% \end{macro}
% \end{macro}
% \end{macro}
% \end{macro}
% \end{macro}
% \end{macro}
%
% \begin{macro}{\astro}
% Font for astronomical symbols.
% If the document lacks a |astro| class option, then the
% |\astro@font| font switch gives a warning message suggesting how to correct things.
%
% The font is loaded at a size commensurate with that of basal text.
% We may need to fine-tune things by scaling the value; the |astro| font's metrics
% are, weirdly, about one quarter the em square. 
%
% FIXME: scale the astro font?
%
% FIXME: provide a better default.
%    \begin{macrocode}
\def\load@astro{%
 \dimen@=1\aas@ptsize\p@
 \font\astro@font=Astrosym at\dimen@
}%
\def\astro#1{\leavevmode\hbox{\astro@font#1}}%
\def\astro@font{%
 \ClassWarning{aastex}{%
  Please use class option `astro', since you are using the astro font.%
 }%
}%
%    \end{macrocode}
% \end{macro}
%
% \begin{macro}{\sbond}
% \begin{macro}{\dbond}
% \begin{macro}{\tbond}
% Chemical bonds in text
%    \begin{macrocode}
\newcommand\sbond{\chem@bnd{\@sbnd}}%
\newcommand\dbond{\chem@bnd{\@dbnd}}%
\newcommand\tbond{\chem@bnd{\@tbnd}}%
%    \end{macrocode}
% \end{macro}
% \end{macro}
% \end{macro}
%
% \begin{macro}{\chem@bnd}
% \begin{macro}{\@sbnd}
% \begin{macro}{\@dbnd}
% \begin{macro}{\@tbnd}
% Low-level macros for chemical bonds in text
%    \begin{macrocode}
\def\chem@bnd#1{%
 {%
  \kern.1em\relax
  \setbox\z@\hbox{M}%
  \dimen@ii.8em\relax
  \p@=.1em\relax
  \dimen@.5\ht\z@\dimen@i-\dimen@
  \advance\dimen@1.5\p@\advance\dimen@i-1.0\p@
  #1%
  \kern.1em\relax
  }%
 }%
%
\def\@sbnd{%
 \advance\dimen@-1.5\p@\advance\dimen@i1.5\p@
 \vrule\@height\dimen@\@depth\dimen@i\@width\dimen@ii\nobreak
 }%
%
\def\@dbnd{%
 \advance\dimen@-0.5\p@\advance\dimen@i0.5\p@
 \vrule\@height\dimen@\@depth\dimen@i\@width\dimen@ii\nobreak
 \advance\dimen@-1.5\p@\advance\dimen@i1.5\p@
 \hskip-\dimen@ii
 \vrule\@height\dimen@\@depth\dimen@i\@width\dimen@ii\nobreak
 }%
%
\def\@tbnd{%
 \vrule\@height\dimen@\@depth\dimen@i\@width\dimen@ii\nobreak
 \advance\dimen@-1.5\p@\advance\dimen@i1.5\p@
 \hskip-\dimen@ii
 \vrule\@height\dimen@\@depth\dimen@i\@width\dimen@ii\nobreak
 \advance\dimen@-1.5\p@\advance\dimen@i1.5\p@
 \hskip-\dimen@ii
 \vrule\@height\dimen@\@depth\dimen@i\@width\dimen@ii\nobreak
 }%
%    \end{macrocode}
% \end{macro}
% \end{macro}
% \end{macro}
% \end{macro}
%
% \begin{macro}{\LaTeX}
% This is a generalization of the definition of |\LaTeX| that comes in 
% \LaTeX.
% The change here ensures that \LaTeX\ will be typeset in the current 
% style, rather than always in roman.
% Note that the appearance of the \LaTeX\ logo will not necessarily 
% coincide with that promulgated by Team \LaTeX\ in the Companions,
% but I claim that, in the absence of a standard set by Lamport himself,
% we need have no concerns about being taken to task by the purity police.
%
% Implementation note: we use herein the |\count@| scratch register.
% if |\everymath| begins using that register, this macro is hosed.
%    \begin{macrocode}
\renewcommand\LaTeX{%
 \leavevmode
 L%
 \raise.42ex\hbox{%
  \count@=\the\fam
  $\fam\count@\scriptstyle\kern-.3em A$%
 }%
 \kern-.15em\TeX
}%
%    \end{macrocode}
% \end{macro}
%
% \begin{macro}{\anchor}
% Syntax: |\anchor|\arg{href}\arg{text}.
%
% A general-purpose hypertext link tag, using the parlance of
% the HTML element |<<A>>| (anchor).
% |<href>| is a \emph{full} URI, including the |scheme:| designator
% and may be active in some sort of hyperref-savvy system;
% for now it has no semantic.
% The |<text>| argument is the copy to set in the document,
% much like the content of the HTML |<<A>>| element.
% 
%    \begin{macrocode}
\newcommand\anchor[2]{#2}%
%    \end{macrocode}
% \end{macro}
%
% \begin{macro}{\url}
% Syntax: |\url|\oarg{href}\arg{text}.
%
% We provide |\url| with the |<href>| as an optional argument,
% to support the degenerate case where an author wishes to express a URL in the text.
% In that case, |<href>|\({}={}\)|<text>|.
% If the author supplies both arguments, this command looks like an alias for |\anchor|.
%    \begin{macrocode}
\newcommand\url{\@dblarg\@url}%
\def\@url[#1]{\anchor{#1}}%
%    \end{macrocode}
% \end{macro}
%
% \begin{macro}{\email}
% Syntax: |\email|\arg{address}.
%
% Identify email addresses anywhere in text (not just in the author group).
%
% This markup will also permit production software at the Press
% to make |mailto:| links, etc., in online products.
% |<address>| should \emph{not} contain the |mailto:| scheme designator.
%
% In general text, |\email{foo@bar}| produces |\textt{mailto:foo@bar}|.
% The |\email| command is also used in the author lists;
% |\authoremail| is deprecated.
%
% FIXME: what happens if there is no title page?
%    \begin{macrocode}
\def\@text@email#1{#1}%
\def\authoremail{\substitute@command\authoremail\email}%
%    \end{macrocode}
% \end{macro}
%
% \begin{macro}{\objectname}
% Syntax: |\objectname|\oarg{catalog ID}\arg{text}.
%
% Celestial objects can and should be identified by the author to facilitate
% post-processing of articles by the astronomical catalog databases.
% At the moment, this is done by eye, although when articles contain these
% marks, the Press could pass along a list of `objects in article'
% to database personnel.  Obviously, more sophisticated technical solutions
% are possible.
% 
% In the future, software could construct queries to those databases in the
% online editions, hence it may be useful for the author to supply a catalog identifier.
% We may choose to suppress the |<catalog ID>| argument in the
% documentation in the near term, since we don't have an architecture
% and a nomenclature for linking to them at this point.
%    \begin{macrocode}
\newcommand\objectname{\@testopt\@objectname{[}}
\def\@objectname[#1]#2{#2}%
%    \end{macrocode}
% macro for object
%    \begin{macrocode}
\newcommand\object{\@testopt\@object{[}}% 
\def\@object[#1]#2{#2}%
%    \end{macrocode}
%% macro for facility
%    \begin{macrocode}
\newcommand\facility{\@testopt\@facility{[}}% 
\def\@facility[#1]#2{#2}%
%    \end{macrocode}
%% macro for supportfrom
%    \begin{macrocode}
\newcommand\supportfrom{\@testopt\@supportfrom{[}}% 
\def\@supportfrom[#1]#2{#2}%
%    \end{macrocode}
%% macro for dataset
%    \begin{macrocode}
\newcommand\dataset{\@testopt\@dataset{[}}% 
\def\@dataset[#1]#2{#2}%
%    \end{macrocode}
% \end{macro}
%
% Here ends the section of code incorporated into the 
% V4.1 document style.
%    \begin{macrocode}
%</class|style>
%    \end{macrocode}
%
%
% Hereinafter appears code for the document class.
%    \begin{macrocode}
%<*class>
%    \end{macrocode}
%
% \subsubsection{Debugging}%
%
% \begin{macro}{\sizrpt}
% Reveal the em-width and ex-height of the current font:
%    \begin{macrocode}
\newcommand\sizrpt{%
 (\fontname\the\font): em=\the\fontdimen6\font, ex=\the\fontdimen5\font
 \typeout{%
  (\fontname\the\font): em=\the\fontdimen6\font, ex=\the\fontdimen5\font
 }%
}%
%    \end{macrocode}
% \end{macro}
%
%
%
% \subsection{Substyles for Manuscript and Preprint}%
%
% This portion defines the class options that account for
% the differences between the various types of document formats
% (substyles) covered by this class.
%
% \subsubsection{Preprint class option}%
%
% \begin{macro}{\aas@preprint}
% The user should invoke the |preprint| document class option
% if wishing the document formatted in preprint format.
%    \begin{macrocode}
\def\aas@preprint{%
 \def\revtex@genre{preprint}%
 \let\ptlandscape\@empty
 \setlength\@slugcmmntwidth{.67\textwidth}%
 \let\@makeslugcmmnt\@makeslugcmmnt@ppt
 \let\@abstract\@abstract@ppt
 \let\make@title\make@title@ppt
 \let\make@affil\make@affil@ppt
 \let\notetoeditor\@gobble
 \let\clear@section@page\@empty
 \let\placetable\@gobble
 \let\placefigure\@gobble
 \let\placeplate\@gobble
 \let\typeout@deluxetable\typeout@deluxetable@ppt
 \let\typeout@pt@nl\typeout@pt@nl@mss
 \ps@plaintop
 \let\references@refpar\references@refpar@mss
 \let\cutinhead\cutinhead@ppt
  \def\pt@headfrac@ass@normalsize{.1}
  \def\pt@headfrac@ass@footnotesize{.12}
  \def\pt@headfrac@ass@scriptsize{.15}
  \def\pt@notemin@normal{4.5}
  \def\pt@notemin@footnote{4.5}
  \def\pt@notemin@script{5}
}%
%    \end{macrocode}
% \end{macro}
%
% \subsubsection{Two-column preprint class option}%
%
% \begin{macro}{\aas@preprint2}
% The user should invoke the |preprint2| document class option
% if wishing the document formatted in two-column preprint format.
%
%    \begin{macrocode}
\@namedef{aas@preprint2}{%
 \aas@preprint
 \def\revtex@genre{2-column preprint}%
 \ps@plain
 \@twocolumntrue
 \@rightskip=\z@\@plus4em
 \rightskip\@rightskip
 \opt@just 
 \just@just % -SZ restored preprint2 to full justification
 \setlength\parindent{1.2em}%
 \setlength\parskip{0.5ex}%
 \setlength\columnsep{0.5in}%
 \setlength\hoffset{-0.2in}%
 \tolerance=600
 \setlength\emergencystretch{6\p@}%
 \def\baselinestretch{1.0}%
 \def\@tightleading{1.0}%
 \let\singlespace\@empty
 \let\doublespace\@empty
 \let\@dates\@dates@pptt
 \let\@abstract\@abstract@pptt
 \let\make@title\make@title@pptt
 \let\make@author\make@author@pptt
 \let\make@affil\make@affil@pptt
 \let\@keywords\@keywords@pptt
 \setlength{\skip\footins}{3ex\@plus1ex\@minus.5ex}%
 \setlength\footnotesep{2ex}%
 \let\@makefntext\@makefntext@pptt
 \let\section\section@pptt
 \let\subsection\subsection@pptt
 \let\subsubsection\subsubsection@pptt
 \let\section@centering\@empty
 \let\appendix@figtab@defs\appendix@figtab@defs@pptt
 \def\clear@thebibliography@page{%
  \if@restonecol\if@twocolumn\else\twocolumn\fi\fi
 }%
 \let\references@refpar\references@refpar@pptt
 \let\thebib@list\thebib@list@pptt
 \let\@tablenotes\@tablenotes@pptt
 \let\fnum@ptablecont\fnum@ptablecont@pptt
 \let\@make@caption@text\@make@caption@text@pptt
 \let\@makecaption@plano@cont\@makecaption@plano@cont@pptt
 \let\cutinhead\cutinhead@ppt
 \def\@plano@float@{\@dblfloat}% Extra layer of abstraction for float processing
 \def\end@plano@float@{\end@dblfloat}%
 %\let\@tableref\@tableref@pptt
 \let\@tablecom\@tablecom@pptt
 \def\pt@headfrac@ass@normalsize{.12}
 \def\pt@headfrac@ass@footnotesize{.15}
 \def\pt@headfrac@ass@scriptsize{.18}
 \def\pt@notemin@normal{5}
 \def\pt@notemin@footnote{5}
 \def\pt@notemin@script{5}
}%
%    \end{macrocode}
% \end{macro}
%
% \begin{macro}{\section@pptt}
% \begin{macro}{\subsection@pptt}
% \begin{macro}{\subsubsection@pptt}
% Level one-, two-, and three heads for two-column preprint substyle.
%    \begin{macrocode}
\def\section@pptt{%
 \@startsection{section}{1}%
  {\z@}{2.3ex\@plus1ex\@minus.2ex}%
  {1.5ex\@plus.2ex}{\normalsize\bfseries}%
}%
\def\subsection@pptt{%
 \@startsection{subsection}{2}%
  {\z@}{2ex\@plus1ex\@minus.2ex}{1ex\@plus.2ex}{\normalsize\bfseries}%
}%
\def\subsubsection@pptt{%
 \@startsection{subsubsection}{3}%
  {\z@}{2ex\@plus1ex\@minus.2ex}{1ex\@plus.2ex}{\normalsize\itshape}%
}%
%    \end{macrocode}
% \end{macro}
% \end{macro}
% \end{macro}
%
% \subsubsection{Manuscript class option}%
%
% \begin{macro}{\aas@manuscript}
% The user should invoke the |manuscript| document class option
% if wishing the document formatted in manuscript format.
%
%    \begin{macrocode}
\def\aas@manuscript{%
 \def\revtex@genre{manuscript}%
 \let\ptlandscape\@empty
 \def\baselinestretch{\@doubleleading}%
 \let\footnote@size\@empty
 \def\tighten{\def\baselinestretch{\@doubleleading}}%
 \def\singlespace{}%
 \def\doublespace{}%
 \opt@tighten
 \setlength\@slugcmmntwidth{\textwidth}%
 \let\@makeslugcmmnt\@makeslugcmmnt@mss
%    \end{macrocode}
%
% Page layout dimensions:
%    \begin{macrocode}
 \@rightskip=\z@\@plus4em
 \rightskip\@rightskip
 \opt@just
%    \end{macrocode}
%
% Other switches:
%    \begin{macrocode}
 \let\sluginfo\sluginfo@mss
 \let\@abstract\@abstract@mss
 \let\@keywords\@keywords@mss % pagebreak after keywords
 \let\make@title\make@title@mss
 \def\clear@section@page{\clearpage}%
 \def\clear@thebibliography@page{\clearpage}%
 \let\typeout@deluxetable\typeout@deluxetable@mss
 \let\typeout@pt@nl\typeout@pt@nl@mss
 \ps@plaintop
 \let\references@refpar\references@refpar@mss
  \def\pt@headfrac@ass@normalsize{.08}
  \def\pt@headfrac@ass@footnotesize{.1}
  \def\pt@headfrac@ass@scriptsize{.12}
  \def\pt@notemin@normal{3}
  \def\pt@notemin@footnote{3}
  \def\pt@notemin@script{4}
}%
%    \end{macrocode}
% \end{macro}
%
%
% Here ends the code for the document class.
% Hereinafter appears code that is common to 
% the document class and document style.
%    \begin{macrocode}
%</class>
%    \end{macrocode}
%
% \subsection{For Editorial Use Only}%
% The following commands permit administrative details or production
% instructions to be passed through to production software at the publisher.
% These commands are not to be used by the author.
%
% The entire editorial entities section of code is incorporated into the 
% V4.1 document style.
%    \begin{macrocode}
%<*class|style>
%    \end{macrocode}
%
% \begin{macro}{\journal}
% \begin{macro}{\volume}
% \begin{macro}{\issue}
% |\journal|\arg{identifier}, |\volume|\arg{text}, |\issue|\arg{text}:
% The identification token for the target journal, and the volume
% and issue numbers, are specified in these commands.
% In the \LaTeX\ style, they produce no output; they are simply a |\@gobble|.
%    \begin{macrocode}
\def\journal{\@gobble}%
\def\volume{\@gobble}%
\def\issue{\@gobble}%
%    \end{macrocode}
% \end{macro}
% \end{macro}
% \end{macro}
%
% \begin{macro}{\SGMLbi}
% \begin{macro}{\SGMLbsc}
% \begin{macro}{\SGMLclc}
% |\SGMLbi|\arg{text}, |\SGMLbsc|\arg{text}, |\SGMLclc|\arg{text}:
% These commands mark text that should be in bold italic, bold script (calligraphic),
% or \emph{always} in mixed case.  The first two occur principally in math mode,
% and they permit editors and keyboardists a more direct approach for specifying
% these changes.  The latter item is used around text, notably chemical symbols,
% that must remain in mixed case, even when presented in an all-caps environment.
% In the \LaTeX\ style, these commands just pass the \arg{text} through.
%    \begin{macrocode}
\def\SGMLbi#1{#1}%
\def\SGMLbsc#1{#1}%
\def\SGMLclc#1{#1}%
%    \end{macrocode}
% \end{macro}
% \end{macro}
% \end{macro}
%
% \begin{macro}{\SGMLentity}
% |\SGMLentity|\arg{text}:
% This command permits editors and keyboardists to introduce special symbols
% with their SGML entity names.  \arg{text} does \emph{not} contain the leading
% ampersand (|&|) or the trailing semicolon (|;|).
% In the \LaTeX\ style, this command should add the ampersand and the
% semicolon to \arg{text}.
%    \begin{macrocode}
\def\SGMLentity#1{\&#1;}%
%    \end{macrocode}
% \end{macro}
%
% \begin{macro}{\SGML}
% |\SGML|\arg{text}:
% This command permits SGML instructions and markup to passed through the
% translation process.
% In the \LaTeX\ style, it produces no output; it is simply a |\@gobble|.
%    \begin{macrocode}
\def\SGML{\@gobble}%
%    \end{macrocode}
% \end{macro}
%
%%%%%%%%%%%%%%%%%%%%%%%%%%%%%%%%%%%%%%%%%
%%
% \begin{macro}{\@makecol}
%% Modify the |\@makecol| macro for one-column abstract.
%    \begin{macrocode}
\newif\if@printfoot%
\global\@printfoottrue%
\global\newbox\my@footins
\gdef \@makecol@pptt {%
   \ifvoid\footins
       \setbox\@outputbox \box\@cclv
   \else
       \if@printfoot
       \else
           \global\setbox\my@footins\vbox{\box\footins}
           %\global\advance\c@colht by -\ht\my@footins
       \fi
       \setbox\@outputbox \vbox{%
       \global\setbox\footins\vbox{\unvbox\my@footins\unvbox\footins}
       \boxmaxdepth \@maxdepth
       \unvbox\@cclv
       \if@printfoot
       \vskip\skip\footins
       \color@begingroup
         \normalcolor
         \footnoterule
	 \unvbox \footins
       \color@endgroup
       \else
       \global\@printfoottrue
       \fi
       }%
   \fi
   \xdef\@freelist{\@freelist\@midlist}%
   \global \let \@midlist \@empty
	   \@combinefloats
   \ifvbox\@kludgeins
     \@makespecialcolbox
   \else
       \setbox\@outputbox \vbox to \@colht{%
       \@texttop
       \dimen@ \dp\@outputbox
       \unvbox \@outputbox
       \ifbotragg%
       \vskip -\dimen@
       \@textbottom
       \fi
       }%
   \fi
   \global \maxdepth \@maxdepth
}
%    \end{macrocode}
% \end{macro}
%%
%%
%%%%%%%%%%%%%%%%%%%%%%%%%%%%%%%%%%%%%%%%%
% \begin{macro}{\strip}
%% macro to switch from one-column long abstract to two-column normal text
%    \begin{macrocode}
\iflong@abstract
\newtoks\juo@sar \juo@sar={}
\def\@ca#1#2#3#4{}
\@ifundefined{stripsep}{\newskip\stripsep\stripsep 15pt}{}
\newskip\m@addvipersep
\m@addvipersep\z@
\newskip\c@addvipersep
\c@addvipersep\z@
\newdimen\ht@strip
\newdimen\right@cor
\newdimen\cor@height \cor@height=0pt
\newdimen\juo@pr
\newdimen\juo@ht
\newbox\@viper
\newcount\juo@sk
\newdimen\c@colht
\newbox\my@outputbox
%% macro for adding strip
\def\add@strip#1#2#3#4{\begingroup%
  \xdef\ex@{\global\noexpand\juo@sar{\the\juo@sar\noexpand\@ca{#1}{#2}{#3}{#4}}}\ex@%
\endgroup}
%% macro for remove strip
\def\remove@strip#1{\ifx#1\@empty\global\juo@sk=0\else
   \global\advance\juo@sk by-1\expandafter\next@item\the #1\@@#1\fi}
\def\next@item \@ca#1#2#3#4#5\@@#6{\global #6={#5}\global\juo@pr=#1\global\juo@ht=#2%
\global\cor@height=#3\global\m@addvipersep=#4}
%% macro for strip command
\def\strip{\@ifnextchar[{\@strip}{\@strip[0pt/0pt]}}
\def\@strip[#1/#2]{\global\@tempdima=#1\global\@tempdimb=#2%
        \global \setbox\@viper\vbox\bgroup%
        \hsize\textwidth
        \@parboxrestore
        \col@number \@ne
        \vrule height\topskip width\z@ depth\z@}
\def\endstrip{%
    \egroup
    \if@firstcolumn
    \ifdim\pagetotal>\z@
       \global\ht@strip\pagegoal
       \global\advance\ht@strip by-\pagetotal
       \global\advance\ht@strip by-\stripsep
       \global\advance\ht@strip by-\ht\@viper
       \global\advance\ht@strip by-\stripsep
       \global\juo@ht=\ht\@viper
    \else
       \global\ht@strip\z@
    \fi
    \ifdim\ht@strip>\z@
       \add@strip{\the\pagetotal}{\the\juo@ht}{\the\@tempdima}{\the\@tempdimb}%
       \global\advance\juo@sk by1
    \fi
  \fi
}
%    \end{macrocode}
% \end{macro}
% \begin{macro}{\@outputdblcol}
% Modification in |\@outputdblcol| command for long abstract
%    \begin{macrocode}
\newif\ifbotragg\botraggfalse
\def\@outputdblcol@pptt{%
\if@firstcolumn
    \global\@firstcolumnfalse
    \global\setbox\@leftcolumn\box\@outputbox
    \global\c@colht\@colht\global\c@addvipersep\z@
    \ifnum\juo@sk>0
        \remove@strip\juo@sar
        \global\@colht\the\juo@pr
        \global\@printfootfalse
        \ifdim\right@cor>\z@
            \global\advance\@colht by\right@cor
        \fi
        \global\advance\@colht by\cor@height
        \global\cor@height=0pt
    \fi
\else
    \ifdim\c@colht>\@colht%
        \@tempdima\@colht%
        \ifdim\@colht>\z@%% 
            \ifvoid\footins%
            \else%
                \vskip-\skip\footins%
            \fi%
            \splittopskip0pt%
            \setbox\my@outputbox\vbox{%
            \vspace*{\@colht}%% 
            \global\botraggfalse%
            \vskip.7\baselineskip%
            \vsplit\@outputbox to \baselineskip%
            \unvbox\@outputbox}%
            \unvbox\my@outputbox%
            \ifvoid\my@outputbox%
               \global\botraggtrue
            \fi%
            \global \@printfoottrue
            \global\c@addvipersep\m@addvipersep%
        \else%
        \fi%
        \ifnum\juo@sk>0%
        \else%
            \global\@colht\c@colht%
        \fi%
    \else%
        \global \@firstcolumntrue%
        \global\setbox\@outputbox\vbox to \ht\@leftcolumn{\unvbox\@outputbox}%
        \setbox\@outputbox \vbox{\hb@xt@\textwidth {%
        \hb@xt@\columnwidth {\box\@leftcolumn \hss}%
        \hfil
        \vrule \@width\columnseprule
        \hfil
        \hb@xt@\columnwidth{\box\@outputbox \hss}}%
        \vss}%
        \@combinedblfloats
        \@outputpage
        \begingroup
        \@dblfloatplacement
        \@startdblcolumn
        \@whilesw\if@fcolmade \fi
        {\@outputpage\@startdblcolumn}%
        \endgroup
        \global\juo@sk=0\global\juo@sar={}%
        \global\cor@height\z@\global\c@addvipersep\z@
    \fi
\fi
}
\fi
%    \end{macrocode}
% \end{macro}
%
% Here ends the code for both document class and document style.
%    \begin{macrocode}
%</class|style>
%    \end{macrocode}
%
%
%
% \section{Tail of Document Style Module}%
%    \begin{macrocode}
%<*style>
%    \end{macrocode}
%
% This is not good: we now define a command for the
% document style that is defined elsewhere for the class.
% FIXME: only one definition of |\email|
%    \begin{macrocode}
%%\def\email#1{\appdef\@author{\make@authoremail{#1}}}%
%    \end{macrocode}
%
% Here is another klootch:
% we have to provide a meaning to |\footnote@size|;
% the following code duplicates code in the class file.
% FIXME: only one code for |\footnote@size|!
%    \begin{macrocode}
\AtBeginDocument{%
 \ifx\footnote@size\undefined
  \let\footnote@size\footnotesize
 \fi
}%
%    \end{macrocode}
%
% Finally, execute the commands queued by |\AtEndOfClass|.
%    \begin{macrocode}
\@styleh@@k
%    \end{macrocode}
%
% Here ends the document style.
%    \begin{macrocode}
%</style>
%    \end{macrocode}
%
% \section{Programming Conventions}%
% In writing the above code, I cleave to certain conventions, noted 
% here.
% My goal in explaining them is to encourage others maintaining this 
% body of code to consider following them as well.
% The benefits are twofold: 
% Some of the coding conventions aim to avoid programming pitfalls;
% following them reduces maintenance costs.
% Others make the code easier to read; following these eases the 
% process of understanding how the code works.
%
% And, for my part, I prefer to read code of this type.
%
% \subsection{Whitespace Conventions}%
% Exactly where code lines break and indent, and where additional
% whitespace is inserted is explained here.
% \begin{itemize}
% \item
% Each new macro definition or register assignment begins a new line.
% Therefore, |\def|, |\newcommand|, and their ilk will start in column 
% 1.
% \item
% Code is indented one space for each level of nesting within braces 
% |{}|.
% \item
% Likewise, if possible, for |\if...| and matching |\fi|.
% \item
% However, the closing brace or |\fi| is outdented by one so that it 
% falls
% at the same level of indentation as its matching brace or |\if|,
% and it appears alone on its line.
% \item
% Use of the |tab| character is deprecated
% (tabs are not standardized across all applications and operating 
% systems).
% \item
% Lines of code are limited to 72 characters.
% I follow this convention mostly to ease the transmission of files
% via email (a deprecated practice) and to accomodate people with
% small monitors.
% But |ltxdoc| output looks better with the shorter lines, too.
% \item
% Extraneous whitespace is avoided by using the comment character |%|.
% In most cases, if falling at the end of a line of code,
% a brace will be immediately followed by a comment character,
% as will the macro parameter |#1| and any one-letter control sequence, 
% like |\\|.
% \end{itemize}
% These conventions taken together are illustrated by the following:
% \begin{verbatim}
%\def\prepdef#1#2{%
% \@ifxundefined#1{\toks@{}}{\toks@\expandafter{#1}}%
% \toks@ii{#2}%
% \edef#1{\the\toks@ii\the\toks@}%
%}%
% \end{verbatim}
% In the above, the definition of |\prepdef| would not fit on a single 
% line,
% and required breaking. The first and last lines have matching braces, 
% and are a the same level of indentation, with the last line containing 
% a single brace.
%
% Each of the three intervening lines has balanced braces and is 
% indented by one space. Each line that would otherwise end in a single 
% brace character is terminated by a comment character.
%
% Some coders rely on the fact that a space character seen by \TeX's 
% scanner while in vertical mode is a no-op.
% Be that as it may, I eliminate them unless actually intentional.
%
% \subsection{Conventions For Procedures}%
% Here are some of my preferences when writing procedures:
% \begin{itemize}
%
% \item
% I dislike defining a macro within another macro, especially when the 
% pattern part is non-nil. 
% You know, you are not saving much space in |mem| when you do this, 
% right?
% You do save space in the hash table and the string pool, though.
% On the other hand, we are not dealing with small \TeX\ engines 
% anymore; Team \LaTeX\ has made sure of this.
%
% \item
% If two or more macros have very similar replacement parts, consider
% layering.
%
% \item
% Macros may perform parsing, may maintain tokens or registers, or may 
% set type (produce marks). I try to avoid combining the three functions 
% in a single macro.
%
% \item
% When a procedure both does assignments and sets type, I try to have a 
% clean separations between the two activities. Try to avoid:
% \begin{verbatim}
% \vskip10pt
% \parindent=0pt
% \leavevmode
% \end{verbatim}
%
% \item
% The Boolean calculus (cf. |\@ifx|)
% is very useful and produces code that executes nicely.
% Using it also helps avoid your having to debug problems where 
% |\if...| and |\fi| are not properly balanced
% (a nightmare in case you didn't already experience it).
%
% \end{itemize}
%
% \subsection{Conventions For \LaTeX}%
% Team \LaTeX\ make certain recommendations in \file{clsguide.tex}.
% Ones that I particularly pay attention to are:
% \begin{itemize}
%
% \item
% For the sake of ``color safety'',
% use |\sbox| rather than |\setbox|, |\mbox| rather than |\hbox|, and 
% |\parbox| or |minipage| rather than |\vbox|.
%
% \item
% Use |\newcommand| and |\newenvironment| to declare user-level commands
% and environments. Avoid the idiom |\def\foo|, |\def\endfoo| to define 
% an environment.
% Ideally, all user-level markup could be extracted from the 
% document class by grepping on |\newcommand| and |\newenvironment|.
%
% \item
% Prefer to use |\setlength| to assign registers.
%
% \end{itemize}
% I cannot help but notice that much of the code of \LaTeX\ itself fails 
% to comply with many of the coding recommendations of Team \LaTeX.
%
%
% \section{Yet to do}%
% The following work is considered for the release of the \aastex\ 5.0 
% package:
% \begin{description}
%
% \item[dot accents]
% Augment |\dot| and |\ddot| with |\dddot| and |\ddddot|.
%
% \item[Posters]
% Andrea wants to do posters. Ref best practices list. Sample document?
%
% \item[BibTeX]
% Support for BibTeX-generated bibliographies.
% Not a problem if using |thebibliography| and |\bibitem|.
% Hoguvain and Berend have written a \file{astrobib.bst}.
%
% \item[Landscape tables]
% Automatically rotate landscape tables if supported by output device.
% User may specify \cmd\rotate in the preamble of |deluxetable|.
% Works for any DVI translator targetting a PostScript printer.
%
% This command will be used in the pseudopreamble of deluxetable to create
% a turnpage (broadside) table in the midst of an upright (portrait) article.
% Implementation will be dependent on authors having a PostScript print engine.
%
% The |\rotate| command is implemented syntactically, but not semantically.
%
% \item[Two-column preprint formatting]
% Automatically put text in 2-col, appendices in 1-col, references in 
% 2-col, obviating the user's need to insert markup. Done.
%
% If using native \LaTeX, all float types are supported, columns are not 
% balanced, and the page ends.
%
% If using |multicol|, only full-page floats are supported, columns 
% balance, and pages do not break.
%
% \item[Table autoplacement]
% Possibly support automatically placing tables as follows:
% less than 20pc width---single column,
% up to 44pc width---full width,
% over 44pc width---landscape.
%
% \item[tighten]
% Class option |tighten| and commands |\doublespace|, |\singlespace| are noops in
% |manuscript| substyle. In others, what should be done?
%
% \item[documentstyle]
% What should happen to a compatibility mode document?
% Wouldn't we have to support the exact same document style options?
%
% \item[substyle names]
% Should we keep the old names or use the new ones?
%
% \item[FIXME]
% Resolve all the remaining notes.
%
% \item[Line lengths]
% Source lines should be no loonger than 72 characters. Done.
% But re-check.
%
% \item[Testing]
% In particular, the |\@ptabularcr| code is new,
% but I have done some testing.
%
% \item[Distribution]
% Why not make this class a supported \LaTeX\ oackage?
% It would go in: |CTAN:/macros/latex/contrib/supported|.
%
% \item[Macro names verbatim]
% Use |\cmd{\foo}| when the control sequence name
% appears in a moving argument.
% Use |\cs{foo}| when the name involves non-alpha characters.
%
% \item[Changes from 4.0]
% It would be good to have a section in the users guide 
% called ``Changes from 4.0''.
% Summarizing
% the new markup commands, 
% the removed markup commands, 
% any alteration in formatting,
% etc.
%
% \item[Deprecated features list]
% To add to user documentation.
%
% \item[ltxdoc stuff]
% |\begin{environment}|, |\generate|
%
% \item[Best practices list]
% To add to user documentation.
% \begin{itemize}
% \item How to code celestial coordinates. 
% \item How to specify object names (+,-). 
% \item How to include figures. 
% \item How to rotate tables. 
% \item |\ge|, |\le| exist---use them! 
% \item Don't include gigantic macro lists---please include only the ones you actually use. 
% \item |\newcommand| is recommended; please avoid |\def|.
% \item Ionization potentials? 
% \end{itemize}
%
% \item
% Natbib commands in user guide.
%
% \end{description}
%
% \section{Recently Done}%
% 
% \begin{description}
%
% \item
% We have to say proactively that plano tables aren't accepted anymore,
% and not just remove those words from the AASTeX user guide.
% That's why the |apjpt| and |aj_pt| substyles are removed from v4.1.
% Done 9113.
%
% \item
% \cmd\SGMLcite. (CB): Deprecated in favor of <TT>\cite*. Unimplemented.
% \item
% \cmd\slantfrac reimplemented with \cmd\case. Done 9113.
%
% \item[Deprecate commands]
% Deprecate |\nl|, |\nextline|, |\tablevspace|, and |\tablebreak|. Done 9113.
%
% \item[URL commands]
% Functionality covered by |hyperref| package.
% Syntax: |\anchor|\arg{HREF}\arg{text}.
% Syntax: |\url|\oarg{URL-specification}\arg{text}.
% Syntax: |\email|\arg{address}.
% Ref: \file{url.sty} in |/contrib/other/misc|. \file{hyperref.sty} in |/contrib/supported/hyperref|.
% Done 9113.
%
% \item[Celestial objects]
% Syntax: |\objectname|\oarg{catalogID}\arg{text}. Done 9113.
%
% \item[\texttt{slantfrac}]
% \texttt{slantfrac} to be deprecated. |\onehalf| and friends to be 
% implemented with |\frac|. Done 9113.
%
% \item[Extensions to \texttt{cite}]
% Support multiple citations in natural language. Ref email from Anton. Done with |natbib|.
%
% \item[\texttt{natbib}]
% I incorporated |natbib| into \aastex\ as a vehicle for the expanded
% citation functionality. |natbib| seems the right choice because
% it is a part of |latex/contrib/supported|, is well thought out and implemented
% and because it provided more than enough functionality, albeit
% via different syntax.
%
% \item[\texttt{amsmath}]
% The |amsmath| package is the successor to the \LaTeX2.09 |amssymb| package.
% I have incorporated it into \aastex\ (that is, it is automatically loaded),
% making it a constant part of \aastex\ rather than an option.
%
% The |amsmath| package is a part of |latex/packages/|, so it is an
% ``annointed'' package---much more blessed than even the 
% |latex/contrib/supported| stuff. I have no hesitation about 
% mandating its use.
%
% \item[Unknown Symbols]
% The following control sequence names cause \LaTeX\ to complain:
% |\lhd|, |\rhd|, |\unlhd|, |\unrhd|, |\sqsubset|, |\sqsupset|,
% |\Join|, |\leadsto|, |\Box|, |\Diamond|, |\mho|.
% The message seen is:
% \begin{verbatim}
% ! LaTeX Error: Command \lhd not provided in base LaTeX2e.
% \end{verbatim}
% This problem is now fixed (9113).
%
% \item[UCPress administrative markup]
% Syntax: \cmd\journal\arg{ID}, \cmd\SGML\arg{text}, \cmd\SGMLentity\arg{text}.
% No production. Not documented in user guide. Done 9113.
%
% \item[Evan's Markup]
% Done 9113.
% \begin{itemize}
% \item
% \cmd\journal\arg{ApJ}, 
% \cmd\articleID\arg{number}, 
% \cmd\volno\arg{number}\arg{issue}.
% \item
% \cmd\journalID\arg{number}
% \item
% \cmd\url\oarg[URL]\arg{text or URL}
% \item
% \cmd\corpauth\arg{name} appears in arg of |\author|.
% \item
% \cmd\SGML\arg{text}
% \item
% \cmd\clc\arg{text} goes to \cmd\SGMLclc\arg{text}
% \item
% \cmd\SGMLentity\arg{text}
% \item
% \cmd\SGMLbi\arg{text}
% \item
% \cmd\SGMLbsc\arg{text}
% \end{itemize}
%
% \item[Email]
% Syntax: \cmd\email\arg{address}.
%
% \item[\cs{objectname}]
% \cmd\objectname\oarg{catalog id}\arg{text}
%
% \item[Evan's Points]
% \begin{itemize}
% \item
% \cmd\shorttitle\arg{text} and \cmd\shortauthors\arg{text}
% instead of |\lefthead| and |\righthead|. Done.
% \item
% \cmd\anchor is \cmd\url
% \item
% No plates anymore.
% \item
% Turnpage figures are a requirement.
% \end{itemize}
%
% \end{description}
%
% \section{User Documentation and Sample Documents}%
% User documentation for \aastex\ is the file \file{aasguide.tex},
% generated by the {\sc docstrip} module following this section.
%
% Note: 
% This file requires the \file{aastex.cls} document class itself,
% thereby lengthening the bootstrap process.
% The requirement enters in the section documenting
% the text entities defined by the class (see key |spec.symbols|).
%
% Sample files for this distribution of \aastex\ include
% \begin{itemize}
% \item
% |sample.tex| is an \aastex\ tutorial paper with a fairly complete 
% set of markup.
% \item
% \item
% |table.tex| is a sample \aastex\ table with complex column headings.
% \end{itemize}%
%
% Since the user documentation and sample files are not heavily commented,
% there is little utility in displaying them in the programmer documentation,
% therefore, I bring the class doc to an end here.
%
% \Finale
% %Here ends the programmer's documentation.
% \endinput
%
% User documentation module
%    \begin{macrocode}
%<*guide>
\documentclass[preprint2]{aastex}%
\makeindex
%
\hyphenation{com-pu-scripts}
\newcommand\aastex{AAS\TeX}%
\newcommand\aastexhome{http://www.journals.uchicago.edu/AAS/AASTeX}%
% 
\makeatletter
\renewenvironment{theindex}
               {\if@twocolumn
                  \@restonecolfalse
                \else
                  \@restonecoltrue
                \fi
                \columnseprule \z@
                \columnsep 35\p@
                \twocolumn[\section{\indexname}]%
%               \@mkboth{\MakeUppercase\indexname}%
%                       {\MakeUppercase\indexname}%
                \thispagestyle{plain}\parindent\z@
                \parskip\z@ \@plus .3\p@\relax
                \let\item\@idxitem}
               {\if@restonecol\onecolumn\else\clearpage\fi}
\makeatother
%
\makeatletter % This code stolen from ltxguide.cls:
% Some hacks with verbatim... NB: this would be better done with the
% verbatim package, but this document has to run on any LaTeX
% installation.

\let\o@verbatim\verbatim
\def\verbatim{%
  \ifhmode\unskip\par\fi
% \nopagebreak              % Overridden by list penalty
  \ifx\@currsize\normalsize
     \small
  \fi
  \o@verbatim
}

% Here we extend the font-setting command to include making <> active
% (ie adjusting the input encoding).
\renewcommand \verbatim@font {%
  \normalfont \ttfamily
  \catcode`\<=\active
  \catcode`\>=\active
}

% Make |...| a synonym for \verb|...|.
\RequirePackage{shortvrb}
\MakeShortVerb{\|}


% Make active <...> produce italics surrounded by angle brackets
% (used in verbatim and \verb).
% << produces a less-than, and >> produces a greater-than.

\begingroup
  \catcode`\<=\active
  \catcode`\>=\active
  \gdef<{\@ifnextchar<\@lt\@meta}
  \gdef>{\@ifnextchar>\@gt\@gtr@err}
  \gdef\@meta#1>{\m{#1}}
  \gdef\@lt<{\char`\<}
  \gdef\@gt>{\char`\>}
\endgroup
\def\@gtr@err{%
   \ClassError{ltxguide}{%
      Isolated \protect>%
   }{%
      In this document class, \protect<...\protect>
      is used to indicate a parameter.\MessageBreak
      I've just found a \protect> on its own.
      Perhaps you meant to type \protect>\protect>?
   }%
}
\def\verbatim@nolig@list{\do\`\do\,\do\'\do\-}

% Various forms of argument:

\newcommand{\m}[1]{\mbox{$\langle$\it #1\/$\rangle$}}
\renewcommand{\arg}[1]{{\tt\string{}\m{#1}{\tt\string}}}
\newcommand{\oarg}[1]{{\tt[%]
}\m{#1}{\tt%[
]}}
\def\cmd#1{\cs{\expandafter\cmd@to@cs\string#1}}
\def\cmd@to@cs#1#2{\char\number`#2\relax}
\DeclareRobustCommand\cs[1]{\texttt{\char`\\#1}}
\newcommand*{\file}[1]{\texttt{#1}}%

 % end of code stolen from ltxguide.cls. Thanks, Alan.

 % code stolen from doc.dtx:
\def\GetFileInfo#1{%
  \def\filename{#1}%
  \def\@tempb##1 ##2 ##3\relax##4\relax{%
    \def\filedate{##1}%
    \def\fileversion{##2}%
    \def\fileinfo{##3}}%
  \edef\@tempa{\csname ver@#1\endcsname}%
  \expandafter\@tempb\@tempa\relax? ? \relax\relax}

 % end of code stolen from doc.dtx. Thanks FMi.
\makeatother
%    \end{macrocode}
%
% Title block of user guide.
%    \begin{macrocode}
\title{%
 The \aastex\ \LaTeXe\ Macros\\for Manuscript Preparation
}%

%\author{Chris Biemesderfer}
%\affil{Seagoat Consulting}

%\author{Revised by Jeannette Barnes, May 1995}
%\affil{American Astronomical Society}
%\affil{NOAO, P.O. Box 26732, Tucson, AZ 85726}

%\author{Revised for \fileversion\ by Arthur Ogawa, \filedate}
%\affil{\TeX\ Consultants}
%\email{ogawa@teleport.com}

\expandafter\GetFileInfo\expandafter{\jobname.tex}%

%\date{\filedate, \fileversion}

\begin{document}

\maketitle
\tableofcontents
%    \end{macrocode}
%
% Introduction
%    \begin{macrocode}
\section{Introduction}

The American Astronomical Society (AAS) has developed a markup package to assist authors in preparing manuscripts intended for submission to AAS-affiliated journals, as well as to other journals that will accept \aastex\ manuscripts.

The most important aspect of the \aastex\
package is that it defines the set of commands, or markup, 
that can be used to identify the structural elements of papers.
When articles are marked up using this set of standard commands, 
they may then be submitted electronically to the editorial
offices and fed into the electronic production of the journals.

This guide contains basic instructions for creating
manuscripts using the \aastex\ markup package. 
Authors are expected to be familiar with the editorial
requirements of the journals so that they can make
appropriate submissions; they should also have a basic 
knowledge of \LaTeX---for instance, knowing
how to set up equations using \LaTeX\ commands.
A number of useful publications about \LaTeX\ (and \TeX) are listed in
the reference section of this guide.

{\em Authors who wish to submit papers electronically to the ApJ,  AJ, or PASP are strongly encouraged use the \aastex\ markup package as described in this guide.}

%    \end{macrocode}
%
% Package context
%    \begin{macrocode}

\section{\aastex\ Article Markup}

This section describes the commands in the \aastex\
package that an author might enter in a manuscript that is being
prepared for electronic submission to one of the journals.
The commands will be described in roughly the same order as they
would appear in a manuscript.
The reader will also find it helpful to examine the
sample files that are distributed with the package and to consult the 
information available on the AASTeX Web site located at \aastexhome.
Authors are also reminded to review the instructions to authors and electronic 
submissions guidelines for the specific journals to which they submit 
their papers.

\subsection{Preamble}\label{sec-preamble}

In \LaTeX\ manuscripts, the preamble is that portion of the file before
the \verb"\begin{document}" command.

\subsubsection{Getting Started}

The first piece of markup in the manuscript declares the
overall style of the document.  Any commands that appear before this
markup will be ignored.
\begin{quote}\index{documentclass@\verb`\documentclass`}
\begin{verbatim}
\documentclass{aastex}
\end{verbatim}
\end{quote}
This specifies the document class as \texttt{aastex} with the
default style (|manuscript|).
The paper copy produced by this style file will be double spaced.
Any tables included in the main body of the manuscript will also be
double spaced.

Other substyles are available. They are discussed in \S~\ref{styles}.

\subsubsection{Defining New Commands}

\aastex\ allows authors to define their own commands with 
\LaTeX's \verb"\newcommand". 
(Authors should not use the plain \TeX\ \verb"\def" command in 
AAS journal submissions.) 
Authors' \verb"\newcommand" definitions must be placed in the 
document preamble. 
 
In general, author-defined commands that are 
abbreviations or shorthands, for instance, 
\begin{quote} 
\begin{verbatim} 
\newcommand{\lte}{local thermodynamic 
	equilibrium} 
\end{verbatim} 
\end{quote} 
, are acceptable and can be easily handled 
by journal offices and publishers during data conversion.  
However, abbreviations that attempt to define new symbols by using 
\LaTeX\ commands for repositioning text tend to 
cause problems in the publication process and should be avoided.
 
In particular, author-defined commands that use any of the commands 
listed below are apt to cause problems during data conversion. 
\begin{quote} 
\verb"\hskip", \verb"\vskip", 
\verb"\raise", \verb"\raisebox", 
\verb"\lower", \verb"\rlap", 
\verb"\kern", \verb"\lineskip", 
\verb"\char", \verb"\mathchar", 
\verb"\mathcode", \verb"\buildref", 
\verb"\mathrel", \verb"\baselineskip" 
\end{quote} 
Consequently, authors are strongly discouraged from 
using them. 
 
Extra symbols are defined for \aastex, 
some specifically for an astronomical context, others more 
broadly used in math and physics. 
In particular, the AMS has additional symbol fonts that are 
available in a standard \LaTeX\ package (\verb"amssymb"). 
All of these symbols are depicted in the additional symbosl tables 
supplied with the package and on the \aastex\ Web site. 
 
Before defining a new symbol command, 
authors are advised to consult these tables to see whether the  
symbol they need already exists.  If it does, they should use the 
corresponding markup command. 
Authors should \emph{not} redefine existing command names. 
When one of these commands is encountered in an electronic 
manuscript submitted to a journal, an author's redefinition will 
be ignored and the originally-defined command used. 

\subsubsection{Editorial Information}

A number of markup commands are available for 
editorial office use in recording the publication history 
for each manuscript. 
These commands should be used by 
the editorial and production offices only. 
\begin{quote} 
\begin{verbatim} 
\received{<receipt date>} 
\revised{<revision date>} 
\accepted{<acceptance date>} 
 
\ccc{<code>} 
\cpright{<type>}{<year>} 
\end{verbatim} 
\end{quote} 
 
Copyright information should be specified with the commands \verb"\cpright"\ 
and \verb"\ccc".  The type of copyright and the corresponding year 
are given in \verb"\cpright". Valid copyright types are as follows. 
\begin{quote} 
\begin{tabular}{l@{\quad}p{1.8in}} 
\tt AAS & Copyright has been assigned to the AAS\\
\tt ASP & Copyright has been assigned to the ASP\\
\tt PD & The article is in the public domain\\
\tt none & No copyright is claimed for the article 
\end{tabular} 
\end{quote} 
The copyright type is case sensitive, so the type string must be 
entered exactly as given above. 

The Copyright Clearing Center code may be given in the \verb"\ccc"\ 
command. The code is taken as regular text, so any special characters, 
notably ``\$,'' must be escaped as appropriate. 

\subsubsection{Short Comment on Title Page}

Authors who wish to include a short remark on the title page,
such as the name and date of the journal to which an article
has been submitted, may do so with the following command.
\begin{quote}
\begin{verbatim}
\slugcomment{<text>}
\end{verbatim}
\end{quote}
In the
\texttt{manuscript} style, 
these comments appear on the title page after the title and authors;
in the
\texttt{preprint} style,
they are placed at the upper right corner
of the title page.

\subsubsection{Running Heads}\label{shorthead}
 
Authors are invited to supply running head information using the 
following commands. 
\begin{quote} 
\begin{verbatim} 
\shorttitle{<text>} 
\shortauthors{<text>} 
\end{verbatim} 
\end{quote} 
Two different kinds of data are generally supplied in running heads.  
The left head contains an author list, (last 
names, possibly truncated as ``et al.''), while the right head 
is an abbreviated form of the paper title.  This running head 
information will not appear on the \LaTeX-printed page but will be 
passed through to  copy editing staff 
for inclusion in the published version. 

Editors and publishers impose varying requirements 
on the brevity of these data.  A good rule of thumb is to limit the list 
of authors to three or else use ``et al.,'' and to limit the 
short form of the title to 40--45 characters. 
The editors may choose to modify the author-supplied 
running heads. 

\subsection{Starting the Main Body}

The preamble is a control section. 
None of the markup that appears in the preamble 
actually typesets anything. The author must include a
\begin{quote}
\begin{verbatim}
\begin{document}
\end{verbatim}
\end{quote}
command to identify the beginning of the main, typeset 
portion of the manuscript.

\subsection{Title and Author Information}
 
Authors should use the 
\verb"\title"\ 
 and \verb"\author"\  commands to specify title 
and author information and the \verb"\affil" command 
to indicate the author's primary affiliation. 
Each \verb"\author"\ 
 command 
should be followed by a corresponding \verb"\affil"\ 
and optional \verb"\email"\  command. 
\begin{quote} 
\begin{verbatim} 
\title{<text>} 
\author{<name(s)>} 
\affil{<affiliation>} 
\affil{<address>} 
\email{<e-mail address>} 
\and 
\end{verbatim} 
\end{quote} 
 Line breaks may be inserted in the title 
with the \verb"\\" command.  (Long titles will 
be broken automatically, so the \verb"\\" markup is not required.) 
If the title is explicitly broken over several lines, the 
preferred style for titles in AAS and ASP journals is the so-called 
``inverted pyramid'' style.  In this style, the longest line 
is the first (or top) line, and each succeeding line is shorter. 
The text of the title should be entered in mixed case; 
it will be printed in upper case or mixed case according to 
the style of the publication.  
Footnotes are permissible in titles. Be careful to ensure that 
alternate affiliations (see below) are properly numbered if a 
footnote to the title is specified. 
 
Authors' names should be entered in mixed case. 
Names that appear together in the author list for authors who 
have the same primary affiliation should be specified in a single 
\verb"\author"\  command. 
Each author group should be followed by 
an \verb"\affil"\  command giving the principle 
affiliation of those authors.  Physical and postal address information 
for the specified institution  may be included with \verb"\affil". 
The address can be broken over several lines using the 
\verb"\\" command to indicate 
the line breaks. 
Usually, however, postal information will fit on one line. 
When there is more than one \verb"\author"\  command, the final  
one should be preceded by the \verb"\and"\ command. 
 
Authors often have affiliations in addition to their principle employer. 
These alternate affiliations may be specified with the \verb"\altaffilmark"\ 
 and \verb"\altaffiltext"\ 
 commands. 
These behave like the \verb"\footnotemark"\ 
and |\footnotetext|  commands of \LaTeX\ except that they do not take  
optional arguments. \verb"\altaffilmark"\ 
 is appended to authors' names in the \verb"\author"\ 
 list and generates superscript identification numbers. 
The text for the individual alternate affiliations is generated by the 
\verb"\altaffiltext"\ 
 command. 
\begin{quote} 
\begin{verbatim} 
\altaffilmark{<key number(s)>} 
\altaffiltext{<numerical key>}{<text>} 
\end{verbatim} 
\end{quote} 
It is up to the author to make sure that 
 each \m{key number}  in his or her \verb"\altaffilmark"\ 
 matches the \m{numerical key} for 
the corresponding \verb"\altaffiltext". 

When there is a lengthy author list, all author names may be 
specified in a single \verb"\author"\ 
 command with affiliations 
specified using the \verb"\altaffilmark"\ me\-chan\-ism. 
In these cases, no \verb"\affil"\ 
 commands are used, and in print, the 
affiliations would all be listed in a footnote block at the bottom 
of the title page. 


\subsection{Abstract}

A paper's abstract should be marked with the
\texttt{abstract} environment.
\begin{quote}
\begin{verbatim}
\begin{abstract}
<abstract text>
\end{abstract}
\end{verbatim}
\end{quote}

If an author is using the \texttt{preprint2} style, and the abstract is too long to fit on the title page, the \texttt{longabstract} option may be used in the \texttt{documentclass} to break the abstract to a new page.

\noindent 
%\begin{quote}
\begin{verbatim}
\documentclass[preprint2,longabstract]{aastex}
\end{verbatim}
%\end{quote}

\noindent Note that authors should use the \texttt{longabstract} option only in \texttt{preprint2} mode and only when the abstract is too long to fit on the title page. 

\subsection{Keywords}

Keywords, or subject headings, are accommodated 
as a single piece of text. 
\begin{quote} 
\begin{verbatim} 
\keywords{<text>} 
\end{verbatim} 
\end{quote} 
If authors supply keywords, they must be delimited by whatever 
punctuation is required by the journal. 
They should be specified in alphabetical order. 
The \verb"\keywords"\ 
 command will print the proper leading text---``Keywords:,'' ``Subject headings:,'' etc.---according to journal style. 

\subsection{Comments to Editors}

Authors may make notes or comments to the copy editor with the
\verb`\notetoeditor` command.
\begin{quote}
\begin{verbatim}
\notetoeditor{<text>}
\end{verbatim}
\end{quote}
This command behaves like a
footnote. Output to the printed page is produced only in the
\texttt{manuscript} style.

\subsection{Sections}

\aastex\ supports four levels of section headings. 
\begin{quote} 
\begin{verbatim} 
\section{<heading>} 
\subsection{<heading>} 
\subsubsection{<heading>} 
\paragraph{<heading>} 
\end{verbatim} 
\end{quote} 
Section headings should be given in upper case or mixed case, depending 
on the style of the journal. 
Note that these commands delimit sections by marking the 
\emph{beginning} of each section; 
there are no separate commands to mark the ends. 

\subsection{Figure and Table Placement}   \label{place}

When preparing a manuscript for submission to an AAS journal, figures and tables do not generally need to be ``placed'' in the text of the
document where an author would like them to appear but
rather may simply follow the main body of the text. 
However, authors may indicate to the editors the preferred placement of
these items by use of the \verb"\place*" commands.
\begin{quote}
\begin{verbatim}
\placetable{<key>}
\placefigure{<key>}
\end{verbatim}
\end{quote}
The \verb"\place*{"\m{key}\verb"}" commands are similar to the 
\verb"\ref" 
 command in La\TeX\ 
and require corresponding \verb"\label"\ 
 commands to link them to the 
proper elements. 
 
When used in the 
\texttt{manuscript} style, the \verb"\place*"\  commands will print a short 
message to the editor about figure or table placement.  
In the other styles, nothing is printed. 

\subsection{Acknowledgments}

\aastex\ supports an
\verb"\acknowledgments"%
\index{acknowledgments@\verb`\acknowledgments`} section.
\begin{quote}
\begin{verbatim}
\acknowledgments
 <acknowledgments text>
\end{verbatim}
\end{quote}
In the \aastex\ styles, acknowledgments are set off from the 
conclusion of the body with vertical space. Note the acknowledgments 
command takes no arguments. 

\subsection{Facilities}

To help organizations obtain information on the effectiveness of their
telescopes, the AAS has created a group of keywords for telescope
facilities.  Using a common set of keywords will make searches for 
this information significantly easier and more accurate.  
In addition, the use of facility
keywords will be useful for linking papers that utilize the 
same telescopes together within the framework of the Virtual Observatory. 
A facilities keyword list is available through a
link at the AASTeX Web page, \aastexhome.

The facilities list should appear after the 
acknowledgments section.
\begin{quote}
\begin{verbatim}
Facilities: \facility{<facility ID>}, 
   \facility{<facility ID>}, 
   \facility{<facility ID>}, ...
\end{verbatim}
\end{quote}
As part of the {\it facility ID} argument, 
the author may also include 
the name of the instrument in parentheses, e.g. 
|\facility{HST(WFPC2)}| or |\facility{MMT(Blue channel spectrograph)}|. 
There is no limit to the number of facility keywords 
that may be included in a paper.

\subsection{Appendices}
 
When one or more appendices are needed in a paper, the point where the 
main body text ends and the appendix begins should be marked 
with the \verb"\appendix"\ 
 command. 
\begin{quote} 
\begin{verbatim} 
\section{<body section>} 
\appendix 
\section{<appendix section>} 
\end{verbatim} 
\end{quote} 
The \verb"\appendix"\ 
 command takes care of a number of internal 
housekeeping concerns, such as identifying sections with letters 
instead of numerals, and resetting the equation counter. 
Note that the \verb"\appendix"\ 
 command takes no arguments. 
Sections in the appendix should be headed with \verb"\section"\ 
 commands. 

\subsection{Equations}

Display equations can be typeset in \LaTeX\ in a number of ways.  
The following three are probably of greatest use in AASTeX. 
\begin{quote} 
\begin{verbatim} 
\begin{displaymath} 
\end{displaymath} 
 
\begin{equation} 
\end{equation} 
 
\begin{eqnarray} 
\end{eqnarray} 
\end{verbatim} 
\end{quote} 
The 
\texttt{displaymath}  environment 
will break out a single, 
unnumbered formula.  The \texttt{equation} environment does the same 
thing except that the equation is 
autonumbered by \LaTeX. 
To set several formul\ae\ in which vertical alignment 
is required, or to display a long equation across multiple lines, use the 
\texttt{eqnarray} environment. Each line of the \texttt{eqnarray} 
will be numbered 
unless a \verb"\nonumber"\ 
 command is inserted 
before the equation line delimiter 
(\verb"\\"). 
\LaTeX's equation counter is \emph{not} incremented when 
\verb"\nonumber"\ 
 is used. 
 
Authors may occasionally wish to group related equations together and 
identify them with letters appended to the equation number. 
When this is desired, such related equations should still be set 
in \texttt{equation} or \texttt{eqnarray} environments, whichever is 
appropriate, and then grouped  within 
the \texttt{mathletters} environment. 
\begin{quote} 
\begin{verbatim} 
\begin{mathletters} 
<equation> or <eqnarray> 
\end{mathletters} 
\end{verbatim} 
\end{quote} 
 
It is possible to override \LaTeX's automatic numbering within the 
\texttt{equation} or \texttt{eqnarray} environments using 
\begin{quote} 
\begin{verbatim} 
\eqnum{<text>} 
\end{verbatim} 
\end{quote} 
When \verb"\eqnum"\ 
 is specified inside an \texttt{equation} environment  
or on a particular equation within an \texttt{eqnarray}, the text 
supplied as an argument to \verb"\eqnum"\ 
 is used as the equation 
identifier. 
\LaTeX's equation counter is \emph{not} incremented when \verb"\eqnum"\ 
 is used. 
\verb"\eqnum"\ 
 must be used \emph{inside} the math environment. 
 
If, as a consequence of the use of \verb"\eqnum"\ 
 or \verb"\nonumber", 
\LaTeX's equation counter gets out of synch with the author's 
intended sequence, 
the counter may be reset to a particular value. 
\begin{quote} 
\begin{verbatim} 
\setcounter{equation}{<number>} 
\end{verbatim} 
\end{quote} 
The equation counter should be set to the number corresponding to the 
last equation that was formatted; therefore, it is most appropriate for this 
command to appear immediately after an \texttt{equation} or 
\texttt{eqnarray} environment.  
The command must be used 
\emph{outside} the math environments. 
 
The \texttt{eqsecnum} style file can also be used to modify the way equations 
are numbered. See \S~\ref{styles} for details. 

\subsection{Citations and Bibliography}   \label{bibliography} 
 
Two options are available for marking citations 
and formatting reference lists: the standard \LaTeX\ \verb"thebibliography"\ 
environment, and the \aastex\ \verb"references"\  environment. 
Authors are strongly encouraged to use \verb"thebibliography"\ 
in their electronic submissions. 
 
Please note that the bibliographic data supplied by the author in 
the reference list must conform to the standards of the journal. 
Many of the journals that 
accept \aastex\ agreed to reduce typographic overhead, bold, italic, etc., 
in reference lists \citep{Abt90}, and the AAS 
has elected not to burden authors with tedious markup commands 
to delimit the bibliographic fields.  Instead, citations will be 
typeset in roman with no size or style changes. 
It is the responsibility of the author to 
arrange the required bibliographic fields in the 
proper order with the correct punctuation, according to journal 
style. 
  
\subsubsection{The {\tt thebibliography} Environment}   \label{bib} 
 
The preferred method for reference management is to use \LaTeX's 
\texttt{thebibliography} environment, marking citations in the body 
of the paper with 
\verb"\citep"\ or \verb"\citet"\ 
and associating references with them using \verb"\bibitem". 
The \verb"\cite"-\verb"\bibitem"\ 
mechanism associates citations and references symbolically 
while maintaining proper citation syntax within the paper. 
In the \verb"\bibitem"\ 
command, the author should specify citation data inside 
square brackets and a citation key in curly braces 
for each reference. (The |\bibitem| command is described 
in detail in the next section.)

\noindent
%\begin{quote} 
\begin{verbatim} 
   \begin{thebibliography}{<dummy>} 
   \bibitem[<cite data>]{<key>} <bibliographic data> 
      . 
      . 
   \end{thebibliography} 
\end{verbatim} 
%\end{quote} 

\noindent Note that the argument \m{dummy} to the start command of the 
environment is not 
used in the \aastex\ package, but it is included to be consistent with the 
syntax of standard \LaTeX. It is acceptable to simply insert an 
empty pair of curly braces at the end of the \verb"\begin{thebibliography}"\ 
command. 
 
\subsubsection{Specifying Bibliographic and Citation Information}\label{spec-bib-data} 
 
\aastex\ uses the |natbib| package \citep{Daly98} for 
citation management.  
The |natbib| package re-implements \LaTeX's |\cite| command, 
offering greater flexibility for managing citations in the 
author-year form. 
 
When using |natbib|, bibliographic data are defined 
in |\bibitem| commands. 

\noindent
%\begin{quote} 
\begin{verbatim} 
   \bibitem[<author>(<year>)]{<key>} 
         <bibliographic data> 
\end{verbatim} 
%\end{quote} 

\noindent
The square-bracketed argument of the |\bibitem| 
contains the \m{author} portion of the citation  
followed by the \m{year} set off in parentheses. 
The parentheses are important---|natbib| uses them to 
determine the year portion of the citation---so be sure to 
include them. 
The argument \m{key} in curly braces is the code name 
by which the citation is referenced in the text. 
 
When placing citations in the text, the author should use 
either a |\citep| or a |\citet| command. 
\begin{quote} 
\begin{verbatim} 
\citep{<key(s)>}
\citet{<key(s)>} 
\end{verbatim} 
\end{quote} 
The |\citep| command produces a citation that is entirely 
set off by parentheses, e.g.\ ``(Cox 1995),'' while |\citet| permits 
the author's name to form part of the text, e.g.\ ``Cox (1995).''   
The plain \LaTeX\ |\cite| command behaves like |\citet|. 
 
The citation \m{key} must correspond to the \m{key} in a |\bibitem| command. 
During processing, information from the square-bracketed argument of the 
key's |\bibitem| is inserted in 
the text at the location of the \verb"\cite"\ command. 
Multiple citation keys are separated by commas, e.g., 
|\citep|\allowbreak 
|{knuth84,|\allowbreak 
|cox95,|\allowbreak 
|lamport94}|. 

|\citep| and |\citet| each take optional arguments that specify 
extra text to be appended to the citation label. Text in the first 
set of square brackets will appear before the cite while
text in the second set will appear after it. 
For instance, |\citep[chap. 2]{jon90}| 
would produce the citation ``(Jones et al., 1990, chap. 2),''  
|\citep[see][]{jon90}| would produce ``(see Jones et al, 1990),'' and 
|\citep[see][chap. 2]|\allowbreak|{jon90}| would print 
``(see Jones et al., 1990, chap. 2).'' 
In addition, the * form of the 
cite commands will print the full author list instead of the 
abbreviated form. 

The syntax discussed above should be sufficient for the 
vast majority of cases; however, \aastex\ 
does use the full |natbib| implementation, so many more 
syntax options are available. For details on 
the full range of |natbib| citation options, see the section 
on using |natbib| on the \aastex\ Web 
site and the |natbib| package documentation. 
 
It is not possible to use \verb"\bibitem" 
within \aastex's \texttt{references} environment 
(\S~\ref{refenv}), 
nor will \verb"\cite"\ commands work properly in the main body 
if \verb"\bibitem" commands are absent. 

\subsubsection{The {\tt references} Environment}\label{refenv} 
 
Some authors might prefer to enter citations directly 
into the body of an article. If so, the |references| environment may be 
used to format the reference list. The 
\texttt{references}  environment 
simply sets off 
the list of references and adjusts spacing parameters. 
\begin{quote} 
\begin{verbatim} 
\begin{references} 
\reference{<key>} <bibliographic data> 
   . 
   . 
\end{references} 
\end{verbatim} 
\end{quote} 
 
While the 
\texttt{references}  environment 
remains supported in \aastex, 
we anticipate that authors will prefer the stronger capabilities of 
the standard \LaTeX\ |thebibliography| commands as extended by |natbib|. 

\subsubsection{Abbreviations for Journal Names} 
 
\aastex\ commands are available for many of the most--frequently-referenced 
journals so that authors may use the markup rather than having 
to look up a particular journal's abbreviation. 
In principle, all the journals should be using the 
same abbreviations, but it is fair to anticipate some changes in 
the specific abbreviations before a system is finally settled on. 
If authors use the journal macros, they need not worry about 
changes to journal style governing abbreviations in citations. 
A listing of the current abbreviation macros appears 
in Table~\ref{journame}. 
\begin{table} 
\footnotesize
\begin{center} 
\caption{Abbreviations for Journal Names}\label{journame} 
\begin{tabular}{l@{\quad}p{2.5in}} 
\verb"\araa" & Annual Review of Astronomy and Astrophysics\\ 
\verb"\actaa" & Acta Astronomica\\ 
\verb"\ao" & Applied Optics\\ 
\verb"\aj" & Astronomical Journal\\ 
\verb"\azh" & Astronomicheskii Zhurnal\\ 
\verb"\aap" & Astronomy and Astrophysics\\ 
\verb"\aapr" & Astronomy and Astrophysics Reviews\\ 
\verb"\apj" & Astrophysical Journal\\ 
\verb"\apjl" & \rule[.5ex]{2em}{.4pt}, Letters to the Editor\\ 
\verb"\apjs" & \rule[.5ex]{2em}{.4pt}, Supplement Series\\ 
\verb"\aplett" & Astrophysics Letters and Communications\\ 
\verb"\apspr" & Astrophysics Space Physics Research\\ 
\verb"\apss" & Astrophysics and Space Science\\ 
\verb"\aaps" & \rule[.5ex]{2em}{.4pt}, Supplement Series\\ 
\verb"\baas" & Bulletin of the AAS\\ 
\verb"\bain" & Bulletin Astronomical Inst. Netherlands\\ 
\verb"\caa" & Chinese Astronomy and Astrophysics\\ 
\verb"\cjaa" & Chinese Journal of Astronomy and Astrophysics\\ 
\verb"\fcp" & Fundamental Cosmic Physics\\
\verb"\gca" & Geochimica Cosmochimica Acta\\ 
\verb"\grl" & Geophysics Research Letters\\ 
\verb"\iaucirc" & IAU Circular\\
\verb"\icarus" & Icarus\\ 
\verb"\jcp" & Journal of Chemical Physics\\ 
\verb"\jgr" & Journal of Geophysical Research\\ 
\verb"\jcap" & Journal of Cosmology and Astroparticle Physics\\
\verb"\jqsrt" & Journal of Quantitative Specstroscopy and Radiative Transfer\\
\verb"\jrasc" & Journal of the RAS of Canada\\ 
\verb"\memras" & Memoirs of the RAS\\ 
\verb"\memsai" & Mem. Societa Astronomica Italiana\\
\verb"\mnras" & Monthly Notices of the RAS\\ 
\verb"\na" & New Astronomy\\ 
\verb"\nar" & New Astronomy Review\\
\verb"\nat" & Nature\\ 
\verb"\nphysa" & Nuclear Physics A\\ 
\verb"\physscr" & Physica Scripta\\ 
\verb"\pra" & Physical Review A\\ 
\verb"\prb" & Physical Review B\\ 
\verb"\prc" & Physical Review C\\ 
\verb"\prd" & Physical Review D\\ 
\verb"\pre" & Physical Review E\\ 
\verb"\prl" & Physical Review Letters\\ 
\verb"\physrep" & Physics Reports\\ 
\verb"\planss" & Planetary Space Science\\ 
\verb"\procspie" & Proceedings of the SPIE\\ 
\verb"\pasa" & Publications of the ASA\\ 
\verb"\pasj" & Publications of the ASJ\\ 
\verb"\pasp" & Publications of the ASP\\ 
\verb"\qjras" & Quarterly Journal of the RAS\\ 
\verb"\rmxaa" & Revista Mexicana de Astronomia y Astrofisica \\
\verb"\skytel" & Sky and Telescope\\ 
\verb"\solphys" & Solar Physics\\ 
\verb"\ssr" & Space Science Reviews\\ 
\verb"\sovast" & Soviet Astronomy\\ 
\verb"\zap" & Zeitschrift fuer Astrophysik\\ 
\end{tabular} 
\end{center} 
\end{table} 

\subsection{Figures} 
 
\subsubsection{Electronic Art} 
 
If an author wishes 
to embed graphics in a manuscript, it is necessary that 
the graphics files conform 
to the Encapsulated PostScript (EPS) standard \citep{PLRM}. The 
author must also have an appropriate DVI translator, 
one that targets PostScript output devices. (Detailed information 
on preparing and submitting electronic art is 
available in the submissions instructions for the journals.) 
 
Several commands are available for including EPS files 
in \aastex\ manuscripts. They should be placed within the 
 \texttt{figure} environment. 
\begin{quote} 
\begin{verbatim} 
\begin{figure} 
\figurenum{<text>} 
\epsscale{<num>} 
\plotone{<epsfile>} 
\plottwo{<epsfile>}{<epsfile>} 
\caption{<text>} 
\end{figure} 
\end{verbatim} 
\end{quote} 
When \verb"\figurenum" is specified inside the {\tt figure} environment, 
the text supplied as an argument to \verb"\figurenum" will be used as the 
figure identifier. 
\LaTeX's figure counter is \emph{not} incremented when \verb"\figurenum" 
is used. 
\verb"\figurenum" must be used \emph{inside} the {\tt figure} environment. 
 
\verb"\plotone" inserts the graphic in the named EPS file, 
scaled so that the horizontal 
dimension fits the width of the body text;  
the vertical dimension is scaled to maintain the aspect ratio. 
\verb"\plottwo" inserts two plots next to each other. 
Scale factors are determined automatically from information in the 
EPS file. 
 
The automatic scaling may be overridden with the command 
\verb"\epsscale{<num>}", where \m{num} is a scaling factor 
 in decimal units, e.g., 0.80. 
 
The |\plotone| and |\plottwo| 
macros are invocations of the |graphicx| 
|\includegraphics| command. In most instances, using 
|\plotone| or |\plottwo| should work for placing 
figures in \aastex\ documents. However, if more 
flexibility is needed, the |\includegraphics| command 
may be invoked directly. For instance, to rotate 
an image by 90 degrees, use 
\begin{quote}
\begin{verbatim}
\includegraphics[angle=90]{<epsfile>}
\end{verbatim}
\end{quote}
See the |graphicx| package documentation or the \aastex\ Web site for a 
complete list of the available arguments to |\includegraphics|. 

Authors are strongly encouraged to use |\plotone|, |\plottwo|, or |\includegraphics| to place their EPS figures. However, as of the current release, the old |\plotfiddle| macro from v4.0 has been reintroduced and may be used if the desired figure placement cannot be achieved with one of the methods above. The syntax of the command is 
\begin{quote}
\begin{verbatim}
\plotfiddle{<epsfile>}{<vsize>}{<rot>}{<hsf>}
       {<vsf>}{<htrans>}{<vtrans>}
\end{verbatim}
\end{quote}
where the arguments are
\begin{center}
\begin{tabular}{l@{\quad}p{2in}}
\tt epsfile & name of the EPS file \\
\tt vsize & vertical white space to allow for plot (\LaTeX\ dimension)\\
\tt rot & rotation angle (degrees)\\
\tt hsf & horizontal width of scaled figure (PS points)\\
\tt vsf & vertical height of scaled figure (PS points)\\
\tt htrans & horizontal translation (PS points)\\
\tt vtrans & vertical translation (PS points)\\
\end{tabular}
\end{center}
PostScript points are 1/72 inches, so an \m{htrans} of 72 moves
the graphic 1 inch to the right. Note that the \m{vtrans} argument is 
discarded in the reimplented macro and is included only for 
backward compatibility with \aastex\ v4.0. 

\subsubsection{Figure Captions} \label{legends} 
 
Regardless of whether an author includes electronic art in a manuscript, figure captions, or legends, should be provided. If art is included in the document, use the \verb"\caption"\ command within the \texttt{figure} environment.  If electronic art is not included, figure captions may be grouped together at the end of the document using \verb"\figcaption".
\begin{quote} 
\begin{verbatim} 
\figcaption[<filename>]{<text>\label{<key>}} 
\end{verbatim} 
\end{quote} 
The optional argument, \m{filename}, can be used to iden\-ti\-fy 
the file for the corresponding figure; 
\m{text} refers to the caption for that figure.  The author 
may provide a \verb"\label"\ with a unique \m{key} for cross-referencing 
purposes. 
 
When the \verb"\figcaption"\  command is used, the figure 
identification label, 
e.g., ``Figure 1,'' is generated automatically by the command itself, so 
there is no need to key this information. 
 There is an upper limit of seven figure captions per page. 
Footnotes are \emph{not} supported for figures. 

\subsection{Tables}  \label{tables}
 
There is support in the \aastex\ package for tables via two mechanisms: 
\LaTeX's standard \texttt{table} environment,  
and the 
\texttt{deluxetable} environment, which allows for the formatting 
of lengthy tabular material.  Tables may be 
marked up using either mechanism, although use of 
\texttt{deluxetable} is preferred. 
Authors should \emph{not} use the \LaTeX\ 
\texttt{tabbing} environment when preparing electronic 
submissions. 
 
\subsubsection{The \texttt{deluxetable} Environment}  \label{dte}

Authors are encouraged to use the 
\texttt{deluxetable} environment to format their tables since it automatically 
handles many formatting tasks, including table numbering and insertion 
of horizontal rules. It also provides mechanisms for breaking tables and 
controlling width and vertical spacing that are unavailable in the \LaTeX\ 
\texttt{tabular} environment. 
 
The \texttt{deluxetable} environment 
is delimited by \LaTeX's familiar 
\verb"\begin"\  and \verb"\end"\  constructs. 
The content consists of preamble commands and table data, 
the latter delimited by \verb"\startdata"\ 
 and \verb"\enddata". 
\begin{quote} 
\begin{verbatim} 
\begin{deluxetable}{<cols>} 
<preamble commands> 
\startdata 
<table data> 
\enddata 
\end{deluxetable} 
\end{verbatim} 
\end{quote} 
The argument \m{cols} specifies the justification for each col\-umn. 
An alignment token, ``l,'' ``c,'' or ``r,'' is given for each column, 
indicating flush left, centered, or flush right. 

\subsubsection{Preamble to the \texttt{deluxetable}}%

There are several commands in the 
\texttt{deluxetable} environment 
that 
must be given in the preamble. 
\begin{quote} 
\begin{verbatim} 
\tabletypesize{<font size command>}
\rotate 
\tablewidth{<dimen>} 
\tablenum{<text>} 
\tablecolumns{<num>} 
\tablecaption{<text>\label{<key>}} 
\tablehead{<text>} 
\end{verbatim} 
\end{quote} 
 
If a table is too wide for the printed page, the font size of the table 
can be cahnged with the \verb"\tabletypesize" command, which takes as 
an argument one of the the font size change commands: 
\verb"\small"\ 
 (11pt), \verb"\footnotesize"\ 
 (10pt), or  
\verb"\scriptsize"\ 
 (8pt). 
 
To force a table to be set in landscape orientation,  
use the \verb"\rotate"\ \label{cmd-rotate} command. Note that 
most DVI previewers will not properly render rotated |deluxetable| 
output, so in order to see what the table looks like, 
it must be output to PostScript and viewed there.
 
The width of a deluxetable is defined by \verb"\tablewidth". 
If this command is omitted, the default width is the width of the page. 
The table can be set to its natural width by specifying 
a dimension of \verb"0pt". 
Long tables may have a natural width that is 
different for each page.  The natural width for each page will be 
printed to the log file during processing.  
Authors may then use this log information to 
define a fixed table width in order to give the 
table a more uniform appearance across pages. 
 
It is possible to override \LaTeX's automatic numbering within the 
\texttt{deluxetable} environment. 
When \verb"\tablenum"\ 
 is specified inside a \texttt{deluxetable} 
preamble, 
the text supplied as an argument to \verb"\tablenum"\ 
 will be used as the 
table identifier. 
\LaTeX's equation counter is \emph{not} incremented when 
\verb"\tablenum"\ 
 is used. 

The caption (actually, the title) of the table is specified 
in \verb"\tablecaption". 
The text of \verb"\tablecaption"\ 
 should be brief; 
explanatory notes should be specified in the end notes to the table 
(see \S \ref{tabendnotes} 
 below).  If the caption 
does not appear 
centered above the table after processing, then specify the width of 
the table explicitly with the \verb"\tablewidth"\ 
 command and rerun 
\LaTeX\ on the file.   If an author supplies a \verb"\label"\ for cross-referencing purposes, this, too, should be included in the \verb"\tablecaption". 
 
Column headings are specified with \verb"\tablehead". 
Within the \verb"\tablehead", each column heading should be given 
in a \verb"\colhead", which will ensure that the heading 
is centered 
on the natural width of the column.  
There should be a heading for each column so that there are as 
many \verb"\colhead"\ 
 commands in the \verb"\tablehead"\ 
 as there 
are data columns. 
 
\begin{quote} 
\begin{verbatim} 
\tablehead{ 
\colhead{<heading>} & \colhead{<heading>}} 
\end{verbatim} 
\end{quote} 
 
If more complicated column headings are required, 
any valid \texttt{tabular} command that constitutes a proper 
head line in a \LaTeX\ table may be used. For example, the  
|\multicolumn| command below would create a table head with 
text centered over five columns.
\begin{quote}
\begin{verbatim}
\multicolumn{5}{c}{<text>}
\end{verbatim}
\end{quote}
Consult \citet{Lamport} or \citet{Kopka99} for further details 
on the available table commands. 
 
The \verb"\tablecolumns{"\m{num}\verb"}" command 
is necessary if an author has 
multi-line column headings produced by \verb"\tablehead"\ 
 or other \LaTeX\ 
commands and is using either the \verb"\cutinhead"\ 
 or \verb"\sidehead"\ 
markup (see below).  The \m{num} argument should be  
set to the true number of columns in the 
table.  The command must come before the \verb"\startdata"\ 
 command. 
 
%It is possible that a complicated table heading will overflow 
%the vertical space allotted for the table heading. 
%The fraction of the page allocated 
%for the table heading may be changed with \verb"\tableheadfrac". 
%The \m{num} argument to \verb"\tableheadfrac"\ 
% should be the 
%decimal fraction of the page used for heading information. 
%The default value is 0.1, meaning that 10\% of the page height 
%is reserved for the table heading.  It should rarely be necessary 
%to change this value. 

\subsubsection{Content of the \texttt{deluxetable}}%
 
After the table title and column headings have been specified, 
data rows can be entered. 
Data rows are delimited with the \verb"\startdata"\ 
 and \verb"\enddata"\  commands. 
The end of each row is indicated with the standard \LaTeX\ \verb"\\" command. 
Data cells within a row are separated with \& (ampersand) characters. 
\begin{quote} 
\begin{verbatim} 
\startdata 
<data line>\\ 
<cell>&<cell>&<cell>\\ 
<more data lines>\\
\enddata 
\end{verbatim} 
\end{quote} 
 
Column alignment within the data columns can be adjusted with the \TeX\ 
\verb"\phantom{"\m{string}\verb"}" command, 
where \m{string} can be any character, e.g., \verb"\phantom{$\arcmin$}". 
A blank character of width \m{string} is then inserted in the table. 
Four commands have been predefined for this purpose. 
\begin{quote} 
\begin{tabular}{l@{\quad}p{2in}} 
{\verb"\phn"% 
} & {phantom numeral 0-9}\\ 
{\verb"\phd"% 
} & {phantom decimal point}\\ 
{\verb"\phs"% 
} & {phantom $\pm$ sign}\\ 
{\verb"\phm{"\m{string}\verb"}"} & {generic phantom}\\ 
\end{tabular} 
\end{quote} 
 
Extra vertical space can be inserted between rows with an optional argument to the 
\verb"\\" command. 
\begin{quote} 
\begin{verbatim} 
\\[<dimen>] 
\end{verbatim} 
\end{quote} 
The argument is a dimension 
and may be specified in any units that are legitimate in \LaTeX. 

In a table, it may happen that several rows of data are 
associated with a single object or item. 
Such logical groupings should not be broken across pages. 
In these cases, the |\tablebreak| command may be used to force 
a page break at the desired point. 
\begin{quote} 
\begin{verbatim} 
<table row>\\
\tablebreak
<next table row>\\
\end{verbatim} 
\end{quote} 
This command can be used any time that the default |deluxetable| page 
breaks need to be overridden.
 
Journals often require that table cells that contain no data 
be explicitly marked.  This is to differentiate such cells from 
blank cells, which are frequently interpreted as implicitly 
repeating the entry in the corresponding cell in the row preceding. 
Table cells for which there are no data should contain 
a \verb"\nodata"\ 
 command. 
\begin{quote} 
\begin{verbatim} 
\nodata 
\end{verbatim} 
\end{quote} 
 
Within the |deluxetable| body, two kinds of special heads are 
allowed, |\cutinhead| and |\sidehead|. 
A cut-in head is a piece of text centered across the width 
of the table. It is spaced above and below 
from the data rows that precede and 
follow it and will appear set off by rules in the \LaTeX\ output.   
Similarly, the command for a side head produces a row spanning 
the width of the table but with the text left justified.  
\begin{quote} 
\begin{verbatim} 
\cutinhead{<text>} 
\sidehead{<text>} 
\end{verbatim} 
\end{quote} 
 
Table footnotes (more properly, table \emph{end notes}) 
may also be used in the \texttt{deluxetable} environment.  
Their use is described in detail in \S \ref{tabendnotes}. 

\subsubsection{The \texttt{table} Environment}

Authors may also compose tables using the \texttt{table} environment. 
\begin{quote} 
\begin{verbatim} 
\begin{table} 
\end{table} 
\end{verbatim} 
\end{quote} 
The \texttt{table}\ environment 
encloses not only the tabular 
material but also any title or footnote information 
associated with the table. 
 
Titles, or captions, for tables are indicated with a |\caption| 
command 
\begin{quote} 
\begin{verbatim} 
\caption{<text>\label{<key>}} 
\end{verbatim} 
\end{quote} 
A table label, e.g. ``Table 2,'' is generated automatically 
by \verb"\caption". 
 The author 
may provide a \verb"\label"\ in the caption with a unique \m{key} 
for cross-referencing 
purposes. 
 
The table body should appear within the \texttt{tabular} environment. 
\begin{quote} 
\begin{verbatim} 
\begin{tabular}{<cols>} 
\end{tabular} 
\end{verbatim} 
\end{quote} 
The alignment tokens in \m{cols} specify the justification for each column. The letters ``l,'' ``c,'' or ``r'' is given for each column, 
indicating left, center, or right justification. 
Consult \citet{Lamport} for details about using 
the 
\texttt{tabular} environment 
to prepare tables. 

Each \texttt{tabular} table must appear within a \texttt{table} 
environment. There should be only one \texttt{tabular} table per 
\texttt{table} 
environment. 
If the journal requests manuscripts with only one table per page, 
the author may need to insert a \verb"\clearpage"\ 
 command after especially short tables. 
 
Use the |\tableline| command to insert horizontal rules in the 
|tabular| environment.
\begin{quote} 
\verb"\tableline" 
 \end{quote} 
The use of vertical rules should be avoided. 

As with the \texttt{deluxetable} environment, it 
 is possible to override \LaTeX's automatic numbering within the 
\texttt{table} environment using \verb"\tablenum". 
\verb"\tablenum"\ 
 must be used \emph{inside} the \texttt{table} 
environment. 
	
\subsubsection{Table End Notes} \label{tabendnotes}

\aastex\ supports footnotes and end notes within tables;
this support applies to both the
\texttt{deluxetable} environment
and the standard \LaTeX\ \texttt{table} environment.

Footnotes for tables are usually identified by lowercase letters
rather than numbers. Use the |\tablenotemark| and |\tablenotetext| 
commands to supply table footnotes. 
As with \verb"\altaffilmark" and \verb"\altaffiltext", a note label, 
usually a letter, is required. 
\begin{quote}
\begin{verbatim}
\tablenotemark{<key letter(s)>}
\tablenotetext{<alpha key>}{<text>}
\end{verbatim}
\end{quote}
The \m{key letter} of the |\tablenotemark| should be the same as the
\m{alpha key} for the corresponding \m{text}.
It is the responsibility of the author to make the correspondence
correct.

Sometimes authors tabulate materials that have corresponding
references and may want to associate these references with the table.  
Authors may also wish to append a short paragraph of explanatory
notes that pertain to the entire table. These elements should be 
specified with the commands below.
\begin{quote}
\begin{verbatim}
\tablerefs{<reference list>}
\tablecomments{<text>}
\end{verbatim}
\end{quote}
The \verb"\tablenotetext", \verb"\tablecomments", and 
\verb"\tablerefs" commands \emph{must}
 be specified after the 
\verb"\end{tabular}" or \verb"\enddata" and before 
the closing \verb"\end{table}" or \verb"\end{deluxetable}". 

\subsection{Supplemental Materials}

For many years now, authors have been taking advantage of the AAS journals'
ability to post supplemental materials with their papers in the electronic 
editions.  Even though each paper must stand on its own 
scientifically without the supplements, these materials are reviewed in the 
peer review process and should be included in initial manuscript 
submissions. As with regular figures and tables, papers with online only 
data must reference each electronic object in the main text 
and include an explanation of what the reader will 
find in the electronic edition.  

The four most popular types of supplemental materials are 
machine-readable tables, online color figures, online-only figures, and
animations.  Please see each journal's Web site for details on 
what types of supplemental materials are acceptable, how to
submit these materials, author tools for preparing them, 
and their associated costs.

\subsubsection{Machine-readable Tables}

Online-only tables submitted to AAS journals are converted 
to a machine-readable 
format for presentation in the electronic edition.  
Machine-readable tables have two
parts: the formatted ASCII data and a metadata header that provides
format, units, and short explanations of each column of data.  This
structure is designed to provide maximum flexibility and ease of use
for readers who wish to further manipulate the data with their own 
computer programs or with software like Excel. 

For each machine-readable table, the author should include a short sample 
version of the table in his or her \LaTeX\ submission. This sample version 
will appear in the print edition as well as 
in HTML in the electronic edition.  The sample table should be
5 to 15 lines long.  It should include a table note at the end 
 with text indicating that a machine-readable 
version will be available in the electronic edition. For instance, 
\begin{quote}
\begin{verbatim}
\tablecomments{Table 1 is 
published in its entirety in 
the electronic edition of 
the Astrophysical Journal. 
A portion is shown here 
for guidance regarding its 
form and content.}
\end{verbatim}
\end{quote}
Each example table must 
be cited and numbered as if it were a fully printed table.

\subsubsection{Online Color Figures}

For some journals, authors may submit figures 
that will appear in black and white in the 
print edition of a journal but in color in the electronic edition. 
In these cases, authors must submit separate black and white and RGB color 
versions of each figure labelled according to the 
file-naming conventions required for the publication.  
For the benefit of print journal readers, figure captions should be written 
with the color information placed inside parentheses, 
for instance, ``The dotted 
line (colored blue in electronic edition) 
is the H$_o$ = 75 km s-1 Mpc-1 model,'' and should 
include a note directing the reader to 
``See the electronic edition of the Journal 
for a color version of this figure.''

\subsubsection{Online-only Figures}

This feature is mainly useful for articles that contain large compendia 
of identification charts and other supplemental graphic material that need not
be printed in full in the paper journal. Online-only figures are intended to 
provide supplemental information that is not critical to the scientific 
content of the article but that might nonetheless be of interest to the reader.

Online-only figures must be mentioned explicitly by number and appear 
in correct numerical order in the body of the text.  At least one figure 
in a series must be displayed as an example figure for the print edition. 
The caption should carry a message indicating that 
more figures are available in the electronic edition---for instance, 
``Plots for all sources are available in the electronic edition of the journal.'' 
Enough information must be included in the figure caption for readers of the 
print edition to determine what is contained in the online-only figures.

\subsubsection{Animations}

Currently, only animations in the MPEG format are accepted by the AAS journals.  
Authors must supply a still frame from the animation in EPS format marked 
up like a regular figure that will serve as an example for the reader. 
They should include text in the caption for the 
still frame indicating that an animation is 
available electronically. For instance, 
``This figure is also available as an mpeg 
animation in the electronic edition of the {\it Astrophysical Journal}.'' 
As with online-only figures, authors must include enough information in the 
figure caption for readers of the print edition to determine what the 
animation illustrates.  

\subsection{Miscellaneous}

\subsubsection{Celestial Objects and Data Sets} 

Authors who wish to have the most important objects in their paper
linked to a data center in the electronic edition may do so using the 
\begin{quote}
\begin{verbatim}
\objectname[<catalog ID>]{<text>}
\end{verbatim}
\end{quote}
macro, or its alias |\object|.  
The text contained in the required argument 
will be printed in the paper and will serve as a link anchor 
in the electronic edition. The catalog ID 
given as an optional argument will be carried through as the identification 
key in the link to a data center. Note that links will only 
be activated if the name provided in the argument is recognized by a 
participating data center.  It is the author's responsibility to use the
correct identifier.   

Similar markup is available for linking to data sets hosted
at participating data centers.
\begin{quote}
\begin{verbatim}
\dataset[<catalog ID>]{<text>}
\end{verbatim}
\end{quote}
In the paper, the text in the required argument will be 
printed while the the catalog ID value will be passed through to form 
links to data centers.  

When an article contains |\object| or |\dataset| commands, the 
publisher will be able to use the markup to 
pass along a list of objects and data sets used in the 
paper to database personnel. Software can then be used to construct links to 
those databases.
Please check with each journal's Web site for instructions on how to 
determine the object and data set identifiers, the location of 
verifications tools, and information on where in the paper these 
macros can best be used.

\subsubsection{Ionic Species and Chemical Bonds} 
 
When discussing atomic species, ionization levels can be indicated 
with the following command. 
\begin{quote} 
\begin{verbatim} 
\ion{<element>}{<level>} 
\end{verbatim} 
\end{quote} 
The ionization state is specified as the second argument  
and should be given as a numeral. 
For example, ``\ion{Ca}{3}'' would be marked up as \verb"\ion{Ca}{3}". 
 
For single, double, and triple chemical bonds, use the following macros. 
\begin{quote} 
\begin{verbatim} 
\sbond 
\dbond 
\tbond 
\end{verbatim} 
\end{quote} 
 
\subsubsection{Fractions} 
 
\aastex\ contains commands that permit authors to specify alternate 
forms for fractions. 
\LaTeX\ will set fractions in displayed math as built-up fractions; 
however, it is sometimes desirable to use case fractions in 
displayed equations. 
In such instances, one should use \verb"\case"\ 
rather than \verb"\frac". Note, however, that authors submitting 
manuscripts electronically to AAS journals 
should generally find it unnecessary to use any markup other than the 
standard \LaTeX\ \verb"\frac". 
 
\begin{center} 
\renewcommand{\arraystretch}{1.4} 
\begin{tabular}{@{}llc@{}} 
Built-up & \verb"\frac{1}{2}" & $\displaystyle\frac{1}{2}$ \\
Case     & \verb"\case{1}{2}" & $\case{1}{2}$ \\ 
Shilled  & \verb"1/2" & $1/2$ \\ 
\end{tabular} 
\end{center} 
 
\subsubsection{Astronomical Symbols}

As mentioned earlier, the \aastex\ package 
contains a collection of assorted macros 
for symbols and abbreviations specific to an astronomical context. 
These are commonly useful and also somewhat difficult for authors 
to produce themselves because fussy kerning is required. 
See the symbols pages provided with the package distribution. 
Most of these commands can be used in both running text and math. However, 
\verb"\lesssim"\  and \verb"\gtrsim"\  can only be used in math mode. 
 
\subsubsection{Hypertext Constructs} 
 
The \verb"\anchor"\ \label{cmd-anchor} command is a general-purpose 
hypertext link tag, associating \verb"<text>" in the manuscript with 
the specified resource (\verb"<href>"). 
\begin{quote} 
\begin{verbatim} 
\anchor{<href>}{<text>} 
\end{verbatim} 
\end{quote} 
\verb"<href>"\ should be specified as a \emph{full} URI, including the 
\verb"scheme:"\  designator (|http:|, |ftp:|, etc.). 
 
The \verb"\url" command supports the special case where an author 
wishes to express a URL in the text. 
\begin{quote} 
\begin{verbatim} 
\url{<text>} 
\end{verbatim} 
\end{quote} 
 
The \verb"\email"\  command is used to identify e-mail addresses 
anywhere in the manuscript. 
The text of the argument is the e-mail address. 
Please do \emph{not} prepend the \verb|mailto:| part. 
\begin{quote} 
\begin{verbatim} 
\email{<address>} 
\end{verbatim} 
\end{quote} 
This command should be used to indicate authors' e-mail addresses 
in author lists at the beginning of manuscripts. 
 
\subsection{Concluding the File} 
 
The last command in the electronic manuscript file should be the 
\begin{quote} 
\begin{verbatim} 
\end{document} 
\end{verbatim} 
\end{quote} 
command, which must appear after all the material in the paper. 
This command directs the formatter to finish processing the manuscript. 
 
\section{Style Options}   \label{styles} 

\subsection{Manuscript Style} 

The default style option is the |manuscript| style. This style will 
produce double-spaced pages printed in a single column at the width 
of the page. 
 
\subsection{Preprints} 
 
Two preprint styles are available. 
The \texttt{preprint} style is similar to the |manuscript| style, 
but it is single-spaced and fully justified.   
The \texttt{preprint2} style produces a two-column preprint. 

\subsubsection{Single-column Preprint}

A single-column preprint format may be specified with the 
\texttt{preprint} style option. 
\begin{quote} 
\begin{verbatim} 
\documentclass[preprint]{aastex} 
\end{verbatim} 
\end{quote} 
The size of the typeface used is under author control by way of 
\LaTeX's \m{nn}\texttt{pt} class options 
(where \m{nn} is 10, 11, or 12). 
Use of 10 point type is not recommended with the 
\texttt{preprint} style. 
 
\subsubsection{Two-column Preprint} 
 
The 
\texttt{preprint2} style 
has the principle function of providing two-column formatting. 
\begin{quote} 
\begin{verbatim} 
\documentclass[preprint2]{aastex} 
\end{verbatim} 
\end{quote} 
It is important to remember 
that text lines are considerably shorter when two columns are typeset 
side by side on a page.  Long equations, wide tables and figures, and 
the like, may not typeset in this format without some adjustments. 
The expenditure of great effort to adapt copy and markup for 
two-column pages is counterproductive. 
Remember that the main 
goal of this package is to produce draft-, or referee-, format pages.  
It is the responsibility of the editors and publishers to 
produce publication-format papers for the journals. 
 
The \texttt{preprint2} style 
sets the article's front matter---the title, author, abstract, 
and keyword material---on a separate page at full text width. 
The body of the article is set in a two-column page grid, 
the appendices in a one-column page grid, 
and the bibliography in a two-column page grid. 
(This manual was prepared using the \texttt{preprint2} style.)
 
The author may supply \LaTeX's \verb"\twocolumn"\  or |\onecolumn| 
commands whenever desired. 
Be aware, however, that using explicit column-switching commands can 
cause formatting problems. 


 
\subsection{The {\tt eqsecnum} Style}

The \texttt{eqsecnum} style file can be used to modify the way equations 
are numbered. 
\begin{quote} 
\begin{verbatim} 
\documentclass[eqsecnum]{aastex} 
\end{verbatim} 
\end{quote} 
Normally, equations are numbered sequentially through the 
entire paper, starting over at ``(A1)'' if there is an appendix. 
If \texttt{eqsecnum} appears in the \texttt{documentclass} command, 
equation num\-bers will be sequential through each section and will be 
formatted ``({\it sec-eqn}),'' where ``{\it sec}'' is the current section number and 
``{\it eqn}'' is the number of the equation within that section. 
 
\subsection{The {\tt flushrt} Style }

A \verb"flushrt"\ style option is available for authors that prefer to 
have their margins left and right justified. 
\begin{quote} 
\begin{verbatim} 
\documentclass[eqsecnum,flushrt]{aastex} 
\end{verbatim} 
\end{quote} 
Note the \texttt{preprint} and \texttt{preprint2} styles
are already flush right by default.

 
\section{Additional Documentation}  \label{docs} 
 
The preceding explanation of the markup commands in the 
\aastex\ package has merit for defining syntax, but many 
authors will prefer to examine the sample papers that are 
included with the style files. 
The files of interest are described below. 
 
A comprehensive example employing nearly all of the capabilities 
of the package (in terms of markup as well as formatting) 
is in \verb"sample.tex". 
This file is annotated with comments that describe 
the purpose of most of the markup. 
\verb"sample.tex"\ includes three tables: two marked up using the 
\texttt{deluxetable} environment 
and another table using the \LaTeX\ \texttt{table} 
environment. 
 
In \verb"table.tex", a complex but short example of the 
\texttt{deluxetable} environment 
demonstrates some of the techniques 
that can be used to generate complex column headings and to align 
variable-width columns. 
Here the \LaTeX\ \verb"\multicolumn" command is used to span a heading 
over several columns.  When \verb"\multicolumn" is used along with the 
\verb"\cutinhead" command, the \verb"\tablecolumns" command must be used to specify 
the number of columns in the table---otherwise the \verb"\cutinhead" command 
will not work properly. 
This table also makes use of the \verb"\phn" command to better align some of the 
columns. 
 
This user guide (\verb"aasguide.tex") 
is also marked up with the \aastex\ package, 
although it is not exemplary as a scientific paper. 
 
Many of the markup commands described in the preceding 
sections are standard \LaTeX\ commands. The reader who is 
unfamiliar with their syntax is referred to the \LaTeX\ 
references works cited in the bibliography, in particular 
\citet{Kopka99} and \citet{Lamport}.
 
Authors who wish to know the ins and outs of \TeX\ itself 
should read the \emph{\TeX book} \citep{Knuth}.  
This resource contains a good deal of information about 
typography in general.  Many details of mathematical typography are 
discussed in \emph{Mathematics into Type} by \citet{Swanson}. 
 
\section{Acknowledgments} 
 
\aastex\ was designed and written by Chris Biemesderfer in 1988. 
Substantial revisions were made by Lee Brotzman and Pierre 
Landau when the package was updated to v4.0. 
\aastex\ was rewritten as a \LaTeXe\ class by Arthur Ogawa 
for the  v5.0 release. It was updated to v5.2 by 
SR Nova Private Ltd. 
The documentation has benefited from revisions by 
Jeannette Barnes, Sara Zimmerman, and Greg Schwartz. 
 
\begin{thebibliography}{} 
 
\bibitem[Abt(1990)]{Abt90} 
  Abt, H.  1990, \apj, 357, 1 (editorial) 
 
\bibitem[Adobe(1999)]{PLRM} 
  Adobe Systems, Inc.  1999, 
  \anchor{http://www.adobe.com/prodindex/postscript/} 
  {PostScript Language Reference Manual} (Reading, MA: Addison-Wesley) 
 
\bibitem[Daly(1998)]{Daly98} 
  Daly, P. 1998, \emph{Natural Sciences Citations and References} 
  (\texttt{natbib} package documentation) 
 
\bibitem[Goossens, Mittelbach, and Samarin(1994)]{Goossens94} 
  Goossens, M., Mittelbach, F., and Samarin, A. 1994, 
  \emph{The \LaTeX\ Companion} (Reading, MA: Addison-Wesley) 
 
\bibitem[Hahn(1993)]{Hahn93} 
  Hahn, J. 1993,  \emph{\LaTeX\ for Everyone} 
  (Englewood Cliffs, NJ: Prentice-Hall) 
 
\bibitem[Knuth(1984)]{Knuth} 
  Knuth, D. 1984, \emph{The \TeX book} (Reading, MA: Addison-Wesley). 
  \newblock Revised to cover \TeX3, 1991. 
 
\bibitem[Kopka and Daly(1999)]{Kopka99} 
  Kopka, H. and Daly, P. 1999,  \emph{A Guide to \LaTeX}, 
  3rd edition, (Reading, MA: Addison-Wesley) 
 
\bibitem[Lamport(1994)]{Lamport} 
  Lamport, L. 1994, \emph{\LaTeX: A Document Preparation System\/}, 
  2nd edition, (Reading, MA: Addison-Wesley) 
 
\bibitem[Swanson(1979)]{Swanson} 
  Swanson, E. 1979, \emph{Mathematics into Type} 
  (Providence, RI: American Mathematical Society) 
 
\end{thebibliography} 
\end{document}
%
%</guide>
%    \end{macrocode}
%
% Symbols module
%    \begin{macrocode}
%<*symbols>
\documentclass[preprint2]{aastex}

\newcommand\aastex{AAS\TeX}%

\pagestyle{myheadings}
\markright{\aastex, AMS math, and \LaTeX\ symbols}%

\begin{document}

\begin{table}
\caption{Additional \aastex\ symbols}
\begin{center}
\tabcolsep5pt
\begin{tabular}{@{}ll@{\hspace*{1em}}ll@{}}
$\lesssim$ & \verb"\lesssim", \verb"\la" &
$\gtrsim$ & \verb"\gtrsim", \verb"\ga" \\
\micron & \verb"\micron" &
\sbond & \verb"\sbond" \\
\dbond & \verb"\dbond" &
\tbond & \verb"\tbond" \\
$\sun$ & \verb"\sun" &
$\earth$ & \verb"\earth" \\
\diameter & \verb"\diameter" &
\\
\arcdeg & \verb"\arcdeg", \verb"\degr" &
\sq & \verb"\sq" \\
\arcmin & \verb"\arcmin" &
\arcsec & \verb"\arcsec" \\
\fd & \verb"\fd" &
\fh & \verb"\fh" \\
\fm & \verb"\fm" &
\fs & \verb"\fs" \\
\fdg & \verb"\fdg" &
\farcm & \verb"\farcm" \\
\farcs & \verb"\farcs" &
\fp & \verb"\fp" \\
\onehalf & \verb"\onehalf" &
\ubvr & \verb"\ubvr" \\
\onethird & \verb"\onethird" &
\ub & \verb"\ub" \\
\twothirds & \verb"\twothirds" &
\bv & \verb"\bv" \\
\onequarter & \verb"\onequarter" &
\vr & \verb"\vr" \\
\threequarters & \verb"\threequarters" &
\ur & \verb"\ur" \\
\end{tabular}
\end{center}
\end{table}

\begin{table}
\caption{Text-mode accents}
\begin{center}
\begin{tabular}{*{2}{ll@{\hspace{4em}}}ll}
\`{o} & \verb"\`{o}" & \={o} & \verb"\={o}" & \t{oo} & \verb"\t{oo}" \\
\'{o} & \verb"\'{o}" & \.{o} & \verb"\.{o}" & \c{o}  & \verb"\c{o}" \\
\^{o} & \verb"\^{o}" & \u{o} & \verb"\u{o}" & \d{o}  & \verb"\d{o}" \\
\"{o} & \verb#\"{o}# & \v{o} & \verb"\v{o}" & \b{o}  & \verb"\b{o}" \\
\~{o} & \verb"\~{o}" & \H{o} & \verb"\H{o}" & & \\
\end{tabular}
\end{center}
\end{table}

\begin{table}
\caption{National symbols}
\begin{center}
\begin{tabular}{*{2}{ll@{\hspace{4em}}}ll}
\oe & \verb"\oe" & \aa & \verb"\aa" & \l  & \verb"\l" \\
\OE & \verb"\OE" & \AA & \verb"\AA" & \L  & \verb"\L" \\
\ae & \verb"\ae" & \o  & \verb"\o"  & \ss & \verb"\ss" \\
\AE & \verb"\AE" & \O  & \verb"\O"  & & \\
\end{tabular}
\end{center}
\end{table}

\begin{table}
\caption{Math-mode accents}
\begin{center}
\begin{tabular}{ll@{\hspace{4em}}ll}
$\hat{a}$   & \verb"\hat{a}"   & $\dot{a}$   & \verb"\dot{a}"   \\
$\check{a}$ & \verb"\check{a}" & $\ddot{a}$  & \verb"\ddot{a}"  \\
$\tilde{a}$ & \verb"\tilde{a}" & $\breve{a}$ & \verb"\breve{a}" \\
$\acute{a}$ & \verb"\acute{a}" & $\bar{a}$   & \verb"\bar{a}"   \\
$\grave{a}$ & \verb"\grave{a}" & $\vec{a}$   & \verb"\vec{a}"   \\
\end{tabular}
\end{center}
\end{table}

\begin{table}
\caption{Greek and Hebrew letters (math mode)}
\begin{center}
\begin{tabular}{cl@{\hspace{3em}}cl}
$\alpha$   & \verb"\alpha"   & $\nu$      & \verb"\nu"      \\
$\beta$    & \verb"\beta"    & $\xi$      & \verb"\xi"      \\
$\gamma$   & \verb"\gamma"   & $o$        & \verb"o"        \\
$\delta$   & \verb"\delta"   & $\pi$      & \verb"\pi"      \\
$\epsilon$ & \verb"\epsilon" & $\rho$     & \verb"\rho"     \\
$\zeta$    & \verb"\zeta"    & $\sigma$   & \verb"\sigma"   \\
$\eta$     & \verb"\eta"     & $\tau$     & \verb"\tau"     \\
$\theta$   & \verb"\theta"   & $\upsilon$ & \verb"\upsilon" \\
$\iota$    & \verb"\iota"    & $\phi$     & \verb"\phi"     \\
$\kappa$   & \verb"\kappa"   & $\chi$     & \verb"\chi"     \\
$\lambda$  & \verb"\lambda"  & $\psi$     & \verb"\psi"     \\
$\mu$      & \verb"\mu"      & $\omega$   & \verb"\omega"   \\

$\digamma$ & \verb"\digamma" &
$\varkappa$ & \verb"\varkappa" \\
$\varepsilon$ & \verb"\varepsilon" & $\varsigma$ & \verb"\varsigma" \\
$\vartheta$   & \verb"\vartheta"   & $\varphi$   & \verb"\varphi"   \\
$\varrho$     & \verb"\varrho"     & & \\

$\Gamma$  & \verb"\Gamma"  & $\Sigma$   & \verb"\Sigma"   \\
$\Delta$  & \verb"\Delta"  & $\Upsilon$ & \verb"\Upsilon" \\
$\Theta$  & \verb"\Theta"  & $\Phi$     & \verb"\Phi"     \\
$\Lambda$ & \verb"\Lambda" & $\Psi$     & \verb"\Psi"     \\
$\Xi$     & \verb"\Xi"     & $\Omega$   & \verb"\Omega"   \\
$\Pi$     & \verb"\Pi"     & & \\

$\aleph$ & \verb"\aleph" &
$\beth$ & \verb"\beth" \\
$\gimel$ & \verb"\gimel" &
$\daleth$ & \verb"\daleth" \\
\end{tabular}
\end{center}
\end{table}

\begin{table}
\caption{Binary operators (math mode)}
\begin{center}
\begin{tabular}{cl@{\hspace{3em}}cl}
$\pm$       & \verb"\pm"       & $\cap$             & \verb"\cap" \\
$\mp$       & \verb"\mp"       & $\cup$             & \verb"\cup" \\
$\setminus$ & \verb"\setminus" & $\uplus$           & \verb"\uplus" \\
$\cdot$     & \verb"\cdot"     & $\sqcap$           & \verb"\sqcap" \\
$\times$    & \verb"\times"    & $\sqcup$           & \verb"\sqcup" \\
$\ast$      & \verb"\ast"      & $\triangleleft$    & \verb"\triangleleft" \\
$\star$     & \verb"\star"     & $\triangleright$   & \verb"\triangleright" \\
$\diamond$  & \verb"\diamond"  & $\wr$              & \verb"\wr" \\
$\circ$     & \verb"\circ"     & $\bigcirc$         & \verb"\bigcirc" \\
$\bullet$   & \verb"\bullet"   & $\bigtriangleup$   & \verb"\bigtriangleup" \\
$\div$      & \verb"\div"      & $\bigtriangledown$ & \verb"\bigtriangledown" \\
$\lhd$      & \verb"\lhd"      & $\rhd$             & \verb"\rhd" \\%FIXME
$\vee$      & \verb"\vee"      & $\odot$            & \verb"\odot" \\
$\wedge$    & \verb"\wedge"    & $\dagger$          & \verb"\dagger" \\
$\oplus$    & \verb"\oplus"    & $\ddagger$         & \verb"\ddagger" \\
$\ominus$   & \verb"\ominus"   & $\amalg$           & \verb"\amalg" \\
$\otimes$   & \verb"\otimes"   & $\unlhd$           & \verb"\unlhd" \\
$\oslash$   & \verb"\oslash"   & $\unrhd$           & \verb"\unrhd"
\end{tabular}
\end{center}
\end{table}

\clearpage

\begin{table}
\caption{AMS binary operators (math mode)}
\begin{center}
\begin{tabular}{ll@{\hspace{2em}}ll}
$\dotplus$ & \verb"\dotplus" &
$\ltimes$ & \verb"\ltimes" \\
$\smallsetminus$ & \verb"\smallsetminus" &
$\rtimes$ & \verb"\rtimes" \\
$\Cap$ & \verb"\Cap", \verb"\doublecap" &
$\leftthreetimes$ & \verb"\leftthreetimes" \\
$\Cup$ & \verb"\Cup", \verb"\doublecup" &
$\rightthreetimes$ & \verb"\rightthreetimes" \\
$\barwedge$ & \verb"\barwedge" &
$\curlywedge$ & \verb"\curlywedge" \\
$\veebar$ & \verb"\veebar" &
$\curlyvee$ & \verb"\curlyvee" \\
$\doublebarwedge$ & \verb"\doublebarwedge" &
\\
$\boxminus$ & \verb"\boxminus" &
$\circleddash$ & \verb"\circleddash" \\
$\boxtimes$ & \verb"\boxtimes" &
$\circledast$ & \verb"\circledast" \\
$\boxdot$ & \verb"\boxdot" &
$\circledcirc$ & \verb"\circledcirc" \\
$\boxplus$ & \verb"\boxplus" &
$\centerdot$ & \verb"\centerdot" \\
$\divideontimes$ & \verb"\divideontimes" &
$\intercal$ & \verb"\intercal" \\
\end{tabular}
\end{center}
\end{table}

\begin{table}
\caption{Miscellaneous symbols}
\begin{center}
\begin{tabular}{ll@{\hspace{4em}}ll}
\dag  & \verb"\dag"  & \S & \verb"\S" \\
\copyright & \verb"\copyright"& \ddag & \verb"\ddag"\\
 \P & \verb"\P" & \pounds    & \verb"\pounds" \\
\#    & \verb"\#"    & \$ & \verb"\$"\\
 \% & \verb"\%" & \&    & \verb"\&"  \\
  \_ & \verb"\_" & \{    & \verb"\{"   \\
 \} & \verb"\}" & & \\
\end{tabular}
\end{center}
\end{table}

\begin{table}
\caption{Miscellaneous symbols (math mode)}
\begin{center}
\begin{tabular}{cl@{\hspace{3em}}cl}
$\aleph$   & \verb"\aleph"   & $\prime$       & \verb"\prime"       \\
$\hbar$    & \verb"\hbar"    & $\emptyset$    & \verb"\emptyset"    \\
$\imath$   & \verb"\imath"   & $\nabla$       & \verb"\nabla"       \\
$\jmath$   & \verb"\jmath"   & $\surd$        & \verb"\surd"        \\
$\ell$     & \verb"\ell"     & $\top$         & \verb"\top"         \\
$\wp$      & \verb"\wp"      & $\bot$         & \verb"\bot"         \\
$\Re$      & \verb"\Re"      & $\|$           & \verb"\|"           \\
$\Im$      & \verb"\Im"      & $\angle$       & \verb"\angle"       \\
$\partial$ & \verb"\partial" & $\triangle$    & \verb"\triangle"    \\
$\infty$   & \verb"\infty"   & $\backslash$   & \verb"\backslash"   \\
$\Box$     & \verb"\Box"     & $\Diamond$     & \verb"\Diamond"     \\
$\forall$  & \verb"\forall"  & $\sharp$       & \verb"\sharp"       \\
$\exists$  & \verb"\exists"  & $\clubsuit$    & \verb"\clubsuit"    \\
$\neg$     & \verb"\neg"     & $\diamondsuit$ & \verb"\diamondsuit" \\
$\flat$    & \verb"\flat"    & $\heartsuit$   & \verb"\heartsuit"   \\
$\natural$ & \verb"\natural" & $\spadesuit$   & \verb"\spadesuit"   \\
$\mho$     & \verb"\mho"     &                &
\end{tabular}
\end{center}
\end{table}

\begin{table}
\caption{AMS miscellaneous symbols (math mode)}
\begin{center}
\begin{tabular}{ll@{\hspace{2em}}ll}
$\hbar$ & \verb"\hbar" &
$\backprime$ & \verb"\backprime" \\
$\hslash$ & \verb"\hslash" &
$\varnothing$ & \verb"\varnothing" \\
$\vartriangle$ & \verb"\vartriangle" &
$\blacktriangle$ & \verb"\blacktriangle" \\
$\triangledown$ & \verb"\triangledown" &
$\blacktriangledown$ & \verb"\blacktriangledown" \\
$\square$ & \verb"\square" &
$\blacksquare$ & \verb"\blacksquare" \\
$\lozenge$ & \verb"\lozenge" &
$\blacklozenge$ & \verb"\blacklozenge" \\
$\circledS$ & \verb"\circledS" &
$\bigstar$ & \verb"\bigstar" \\
$\angle$ & \verb"\angle" &
$\sphericalangle$ & \verb"\sphericalangle" \\
$\measuredangle$ & \verb"\measuredangle" &
\\
$\nexists$ & \verb"\nexists" &
$\complement$ & \verb"\complement" \\
$\mho$ & \verb"\mho" &
$\eth$ & \verb"\eth" \\
$\Finv$ & \verb"\Finv" &
$\diagup$ & \verb"\diagup" \\
$\Game$ & \verb"\Game" &
$\diagdown$ & \verb"\diagdown" \\
$\Bbbk$ & \verb"\Bbbk" &
$\restriction$ & \verb"\restriction" \\
\end{tabular}
\end{center}
\end{table}

\begin{table}
\tabcolsep1pt
\caption{Arrows (math mode)}
\begin{center}
\begin{tabular}{clcl}
$\leftarrow$          & \verb"\leftarrow" &
$\longleftarrow$      & \verb"\longleftarrow" \\
$\Leftarrow$          & \verb"\Leftarrow" &
$\Longleftarrow$      & \verb"\Longleftarrow" \\
$\rightarrow$         & \verb"\rightarrow" &
$\longrightarrow$     & \verb"\longrightarrow" \\
$\Rightarrow$         & \verb"\Rightarrow" &
$\Longrightarrow$     & \verb"\Longrightarrow" \\
$\leftrightarrow$     & \verb"\leftrightarrow" &
$\longleftrightarrow$ & \verb"\longleftrightarrow" \\
$\Leftrightarrow$     & \verb"\Leftrightarrow" &
$\Longleftrightarrow$ & \verb"\Longleftrightarrow" \\
$\mapsto$             & \verb"\mapsto" &
$\longmapsto$         & \verb"\longmapsto" \\
$\hookleftarrow$      & \verb"\hookleftarrow" &
$\hookrightarrow$     & \verb"\hookrightarrow" \\
$\leftharpoonup$      & \verb"\leftharpoonup" &
$\rightharpoonup$     & \verb"\rightharpoonup" \\
$\leftharpoondown$    & \verb"\leftharpoondown" &
$\rightharpoondown$   & \verb"\rightharpoondown" \\
$\rightleftharpoons$  & \verb"\rightleftharpoons" &
$\leadsto$            & \verb"\leadsto" \\
$\uparrow$            & \verb"\uparrow" &
$\Updownarrow$        & \verb"\Updownarrow" \\
$\Uparrow$            & \verb"\Uparrow" &
$\nearrow$            & \verb"\nearrow" \\
$\downarrow$          & \verb"\downarrow" &
$\searrow$            & \verb"\searrow" \\
$\Downarrow$          & \verb"\Downarrow" &
$\swarrow$            & \verb"\swarrow" \\
$\updownarrow$        & \verb"\updownarrow" &
$\nwarrow$            & \verb"\nwarrow"
\end{tabular}
\end{center}
\end{table}

\begin{table}
\tabcolsep1pt
\caption{AMS arrows (math mode)}
\begin{center}
\begin{tabular}{llll}
$\dashleftarrow$ & \verb"\dashleftarrow" &
$\dashrightarrow$ & \verb"\dashrightarrow" \\
$\leftleftarrows$ & \verb"\leftleftarrows" &
$\rightrightarrows$ & \verb"\rightrightarrows" \\
$\leftrightarrows$ & \verb"\leftrightarrows" &
$\rightleftarrows$ & \verb"\rightleftarrows" \\
$\Lleftarrow$ & \verb"\Lleftarrow" &
$\Rrightarrow$ & \verb"\Rrightarrow" \\
$\twoheadleftarrow$ & \verb"\twoheadleftarrow" &
$\twoheadrightarrow$ & \verb"\twoheadrightarrow" \\
$\leftarrowtail$ & \verb"\leftarrowtail" &
$\rightarrowtail$ & \verb"\rightarrowtail" \\
$\looparrowleft$ & \verb"\looparrowleft" &
$\looparrowright$ & \verb"\looparrowright" \\
$\leftrightharpoons$ & \verb"\leftrightharpoons" &
$\rightleftharpoons$ & \verb"\rightleftharpoons" \\
$\curvearrowleft$ & \verb"\curvearrowleft" &
$\curvearrowright$ & \verb"\curvearrowright" \\
$\circlearrowleft$ & \verb"\circlearrowleft" &
$\circlearrowright$ & \verb"\circlearrowright" \\
$\Lsh$ & \verb"\Lsh" &
$\Rsh$ & \verb"\Rsh" \\
$\upuparrows$ & \verb"\upuparrows" &
$\downdownarrows$ & \verb"\downdownarrows" \\
$\upharpoonleft$ & \verb"\upharpoonleft" &
$\upharpoonright$ & \verb"\upharpoonright" \\
$\downharpoonleft$ & \verb"\downharpoonleft" &
$\downharpoonright$ & \verb"\downharpoonright" \\
$\multimap$ & \verb"\multimap" &
$\rightsquigarrow$ & \verb"\rightsquigarrow" \\
$\leftrightsquigarrow$ & \verb"\leftrightsquigarrow" &
\\
$\nleftarrow$ & \verb"\nleftarrow" &
$\nrightarrow$ & \verb"\nrightarrow" \\
$\nLeftarrow$ & \verb"\nLeftarrow" &
$\nRightarrow$ & \verb"\nRightarrow" \\
$\nleftrightarrow$ & \verb"\nleftrightarrow" &
$\nLeftrightarrow$ & \verb"\nLeftrightarrow" \\
\end{tabular}
\end{center}
\end{table}

\begin{table}
\caption{Relations (math mode)}
\begin{center}
\begin{tabular}{cl@{\hspace{4em}}cl}

$\leq$        & \verb"\leq"        & $\geq$        & \verb"\geq" \\
$\prec$       & \verb"\prec"       & $\succ$       & \verb"\succ" \\
$\preceq$     & \verb"\preceq"     & $\succeq$     & \verb"\succeq" \\
$\ll$         & \verb"\ll"         & $\gg$         & \verb"\gg" \\
$\subset$     & \verb"\subset"     & $\supset$     & \verb"\supset" \\
$\subseteq$   & \verb"\subseteq"   & $\supseteq$   & \verb"\supseteq" \\
$\sqsubset$   & \verb"\sqsubset"   & $\sqsupset$   & \verb"\sqsupset" \\
$\sqsubseteq$ & \verb"\sqsubseteq" & $\sqsupseteq$ & \verb"\sqsupseteq" \\
$\in$         & \verb"\in"         & $\ni$         & \verb"\ni" \\
$\vdash$      & \verb"\vdash"      & $\dashv$      & \verb"\dashv" \\
$\smile$      & \verb"\smile"      & $\mid$        & \verb"\mid" \\
$\frown$      & \verb"\frown"      & $\parallel$   & \verb"\parallel" \\
$\neq$        & \verb"\neq"        & $\perp$       & \verb"\perp" \\
$\equiv$      & \verb"\equiv"      & $\cong$       & \verb"\cong" \\
$\sim$        & \verb"\sim"        & $\bowtie$     & \verb"\bowtie" \\
$\simeq$      & \verb"\simeq"      & $\propto$     & \verb"\propto" \\
$\asymp$      & \verb"\asymp"      & $\models$     & \verb"\models" \\
$\approx$     & \verb"\approx"     & $\doteq$      & \verb"\doteq" \\
              &                    & $\Join$       & \verb"\Join"
\end{tabular}
\end{center}
\end{table}

\begin{table}
\tabcolsep1pt
\caption{AMS binary relations (math mode)}
\begin{center}
\begin{tabular}{ll@{\hspace{1em}}ll}
$\leqq$ & \verb"\leqq" &
$\geqq$ & \verb"\geqq" \\
$\leqslant$ & \verb"\leqslant" &
$\geqslant$ & \verb"\geqslant" \\
$\eqslantless$ & \verb"\eqslantless" &
$\eqslantgtr$ & \verb"\eqslantgtr" \\
$\lesssim$ & \verb"\lesssim" &
$\gtrsim$ & \verb"\gtrsim" \\
$\lessapprox$ & \verb"\lessapprox" &
$\gtrapprox$ & \verb"\gtrapprox" \\
$\approxeq$ & \verb"\approxeq" &
$\eqsim$ & \verb"\eqsim" \\
$\lessdot$ & \verb"\lessdot" &
$\gtrdot$ & \verb"\gtrdot" \\
$\lll$ & \verb"\lll", \verb"\llless" &
$\ggg$ & \verb"\ggg", \verb"\gggtr" \\
$\lessgtr$ & \verb"\lessgtr" &
$\gtrless$ & \verb"\gtrless" \\
$\lesseqgtr$ & \verb"\lesseqgtr" &
$\gtreqless$ & \verb"\gtreqless" \\
$\lesseqqgtr$ & \verb"\lesseqqgtr" &
$\gtreqqless$ & \verb"\gtreqqless" \\
$\doteqdot$ & \verb"\doteqdot", \verb"\Doteq" &
$\eqcirc$ & \verb"\eqcirc" \\
$\risingdotseq$ & \verb"\risingdotseq" &
$\circeq$ & \verb"\circeq" \\
$\fallingdotseq$ & \verb"\fallingdotseq" &
$\triangleq$ & \verb"\triangleq" \\
$\backsim$ & \verb"\backsim" &
$\thicksim$ & \verb"\thicksim" \\
$\backsimeq$ & \verb"\backsimeq" &
$\thickapprox$ & \verb"\thickapprox" \\
$\subseteqq$ & \verb"\subseteqq" &
$\supseteqq$ & \verb"\supseteqq" \\
$\Subset$ & \verb"\Subset" &
$\Supset$ & \verb"\Supset" \\
$\sqsubset$ & \verb"\sqsubset" &
$\sqsupset$ & \verb"\sqsupset" \\
$\preccurlyeq$ & \verb"\preccurlyeq" &
$\succcurlyeq$ & \verb"\succcurlyeq" \\
$\curlyeqprec$ & \verb"\curlyeqprec" &
$\curlyeqsucc$ & \verb"\curlyeqsucc" \\
$\precsim$ & \verb"\precsim" &
$\succsim$ & \verb"\succsim" \\
$\precapprox$ & \verb"\precapprox" &
$\succapprox$ & \verb"\succapprox" \\
$\vartriangleleft$ & \verb"\vartriangleleft" &
$\vartriangleright$ & \verb"\vartriangleright" \\
$\trianglelefteq$ & \verb"\trianglelefteq" &
$\trianglerighteq$ & \verb"\trianglerighteq" \\
$\vDash$ & \verb"\vDash" &
$\Vdash$ & \verb"\Vdash" \\
$\Vvdash$ & \verb"\Vvdash" &
\\
$\smallsmile$ & \verb"\smallsmile" &
$\shortmid$ & \verb"\shortmid" \\
$\smallfrown$ & \verb"\smallfrown" &
$\shortparallel$ & \verb"\shortparallel" \\
$\bumpeq$ & \verb"\bumpeq" &
$\between$ & \verb"\between" \\
$\Bumpeq$ & \verb"\Bumpeq" &
$\pitchfork$ & \verb"\pitchfork" \\
$\varpropto$ & \verb"\varpropto" &
$\backepsilon$ & \verb"\backepsilon" \\
$\blacktriangleleft$ & \verb"\blacktriangleleft" &
$\blacktriangleright$ & \verb"\blacktriangleright" \\
$\therefore$ & \verb"\therefore" &
$\because$ & \verb"\because" \\
\end{tabular}
\end{center}
\end{table}

\begin{table}
\caption{AMS negated relations (math mode)}
\begin{center}
\begin{tabular}{ll@{\hspace{2em}}ll}
$\nless$ & \verb"\nless" &
$\ngtr$ & \verb"\ngtr" \\
$\nleq$ & \verb"\nleq" &
$\ngeq$ & \verb"\ngeq" \\
$\nleqslant$ & \verb"\nleqslant" &
$\ngeqslant$ & \verb"\ngeqslant" \\
$\nleqq$ & \verb"\nleqq" &
$\ngeqq$ & \verb"\ngeqq" \\
$\lneq$ & \verb"\lneq" &
$\gneq$ & \verb"\gneq" \\
$\lneqq$ & \verb"\lneqq" &
$\gneqq$ & \verb"\gneqq" \\
$\lvertneqq$ & \verb"\lvertneqq" &
$\gvertneqq$ & \verb"\gvertneqq" \\
$\lnsim$ & \verb"\lnsim" &
$\gnsim$ & \verb"\gnsim" \\
$\lnapprox$ & \verb"\lnapprox" &
$\gnapprox$ & \verb"\gnapprox" \\
$\nprec$ & \verb"\nprec" &
$\nsucc$ & \verb"\nsucc" \\
$\npreceq$ & \verb"\npreceq" &
$\nsucceq$ & \verb"\nsucceq" \\
$\precneqq$ & \verb"\precneqq" &
$\succneqq$ & \verb"\succneqq" \\
$\precnsim$ & \verb"\precnsim" &
$\succnsim$ & \verb"\succnsim" \\
$\precnapprox$ & \verb"\precnapprox" &
$\succnapprox$ & \verb"\succnapprox" \\
$\nsim$ & \verb"\nsim" &
$\ncong$ & \verb"\ncong" \\
$\nshortmid$ & \verb"\nshortmid" &
$\nshortparallel$ & \verb"\nshortparallel" \\
$\nmid$ & \verb"\nmid" &
$\nparallel$ & \verb"\nparallel" \\
$\nvdash$ & \verb"\nvdash" &
$\nvDash$ & \verb"\nvDash" \\
$\nVdash$ & \verb"\nVdash" &
$\nVDash$ & \verb"\nVDash" \\
$\ntriangleleft$ & \verb"\ntriangleleft" &
$\ntriangleright$ & \verb"\ntriangleright" \\
$\ntrianglelefteq$ & \verb"\ntrianglelefteq" &
$\ntrianglerighteq$ & \verb"\ntrianglerighteq" \\
$\nsubseteq$ & \verb"\nsubseteq" &
$\nsupseteq$ & \verb"\nsupseteq" \\
$\nsubseteqq$ & \verb"\nsubseteqq" &
$\nsupseteqq$ & \verb"\nsupseteqq" \\
$\subsetneq$ & \verb"\subsetneq" &
$\supsetneq$ & \verb"\supsetneq" \\
$\varsubsetneq$ & \verb"\varsubsetneq" &
$\varsupsetneq$ & \verb"\varsupsetneq" \\
$\subsetneqq$ & \verb"\subsetneqq" &
$\supsetneqq$ & \verb"\supsetneqq" \\
$\varsubsetneqq$ & \verb"\varsubsetneqq" &
$\varsupsetneqq$ & \verb"\varsupsetneqq" \\
\end{tabular}
\end{center}
\end{table}

\begin{table}
\caption{Variable-sized symbols (math mode)}
\begin{center}
\begin{displaymath}
\begin{array}{ccl@{\hspace{2em}}ccl}
\sum & \displaystyle \sum & \hbox{\verb"\sum"} &
\bigcap & \displaystyle \bigcap & \hbox{\verb"\bigcap"} \\
\prod & \displaystyle \prod & \hbox{\verb"\prod"} &
\bigcup & \displaystyle \bigcup & \hbox{\verb"\bigcup"} \\
\coprod & \displaystyle \coprod & \hbox{\verb"\coprod"} &
\bigsqcup & \displaystyle \bigsqcup & \hbox{\verb"\bigsqcup"} \\
\int & \displaystyle \int & \hbox{\verb"\int"} &
\bigvee & \displaystyle \bigvee & \hbox{\verb"\bigvee"} \\
\oint & \displaystyle \oint & \hbox{\verb"\oint"} &
\bigwedge & \displaystyle \bigwedge & \hbox{\verb"\bigwedge"} \\
\bigodot & \displaystyle \bigodot & \hbox{\verb"\bigodot"} &
\bigotimes & \displaystyle \bigotimes & \hbox{\verb"\bigotimes"} \\
\bigoplus & \displaystyle \bigoplus & \hbox{\verb"\bigoplus"} &
\biguplus & \displaystyle \biguplus & \hbox{\verb"\biguplus"} \\
\end{array}
\end{displaymath}
\end{center}
\end{table}

\begin{table}
\caption{Delimiters (math mode)}
\begin{center}
\begin{tabular}{lc@{\hspace{2em}}lc}
$($            & \verb"("            & $)$            & \verb")" \\
$[$            & \verb"["            & $]$            & \verb"]" \\
$\{$           & \verb"\{"           & $\}$           & \verb"\}" \\
$\lfloor$      & \verb"\lfloor"      & $\rfloor$      & \verb"\rfloor" \\
$\lceil$       & \verb"\lceil"       & $\rceil$       & \verb"\rceil" \\
$\langle$      & \verb"\langle"      & $\rangle$      & \verb"\rangle" \\
$/$            & \verb"/"            & $\backslash$   & \verb"\backslash" \\
$\vert$        & \verb"\vert"        & $\Vert$        & \verb"\Vert" \\
$\uparrow$     & \verb"\uparrow"     & $\Uparrow$     & \verb"\Uparrow" \\
$\downarrow$   & \verb"\downarrow"   & $\Downarrow$   & \verb"\Downarrow" \\
$\updownarrow$ & \verb"\updownarrow" &
$\Updownarrow$ & \verb"\Updownarrow" \\
$\ulcorner$ & \verb"\ulcorner" &
$\urcorner$ & \verb"\urcorner" \\
$\llcorner$ & \verb"\llcorner" &
$\lrcorner$ & \verb"\lrcorner" \\
\end{tabular}
\end{center}
\end{table}

\begin{table}
\caption{Function names (math mode)}
\begin{verbatim}
  \arccos   \csc    \ker      \min
  \arcsin   \deg    \lg       \Pr
  \arctan   \det    \lim      \sec
  \arg      \dim    \liminf   \sin
  \cos      \exp    \limsup   \sinh
  \cosh     \gcd    \ln       \sup
  \cot      \hom    \log      \tan
  \coth     \inf    \max      \tanh
\end{verbatim}
\end{table}
\end{document}
%</symbols>
% \end{macrocode}
%
% Sample document modules
%
% Here begins the module `sample`.
%    \begin{macrocode}
%<*sample>
%%
%% Beginning of file 'sample.tex'
%%
%% Modified 2005 June 21
%%
%% This is a sample manuscript marked up using the
%% AASTeX v5.x LaTeX 2e macros.

%% The first piece of markup in an AASTeX v5.x document
%% is the \documentclass command. LaTeX will ignore
%% any data that comes before this command.

%% The command below calls the preprint style
%% which will produce a one-column, single-spaced document.
%% Examples of commands for other substyles follow. Use
%% whichever is most appropriate for your purposes.
%%
%%\documentclass[12pt,preprint]{aastex}

%% manuscript produces a one-column, double-spaced document:

\documentclass[manuscript]{aastex}

%% preprint2 produces a double-column, single-spaced document:

%% \documentclass[preprint2]{aastex}

%% Sometimes a paper's abstract is too long to fit on the
%% title page in preprint2 mode. When that is the case,
%% use the longabstract style option.

%% \documentclass[preprint2,longabstract]{aastex}

%% If you want to create your own macros, you can do so
%% using \newcommand. Your macros should appear before
%% the \begin{document} command.
%%
%% If you are submitting to a journal that translates manuscripts
%% into SGML, you need to follow certain guidelines when preparing
%% your macros. See the AASTeX v5.x Author Guide
%% for information.

\newcommand{\vdag}{(v)^\dagger}
\newcommand{\myemail}{skywalker@galaxy.far.far.away}

%% You can insert a short comment on the title page using the command below.

\slugcomment{Not to appear in Nonlearned J., 45.}

%% If you wish, you may supply running head information, although
%% this information may be modified by the editorial offices.
%% The left head contains a list of authors,
%% usually a maximum of three (otherwise use et al.).  The right
%% head is a modified title of up to roughly 44 characters.
%% Running heads will not print in the manuscript style.

\shorttitle{Collapsed Cores in Globular Clusters}
\shortauthors{Djorgovski et al.}

%% This is the end of the preamble.  Indicate the beginning of the
%% paper itself with \begin{document}.

\begin{document}

%% LaTeX will automatically break titles if they run longer than
%% one line. However, you may use \\ to force a line break if
%% you desire.

\title{Collapsed Cores in Globular Clusters, \\
    Gauge-Boson Couplings, and AAS\TeX\ Examples}

%% Use \author, \affil, and the \and command to format
%% author and affiliation information.
%% Note that \email has replaced the old \authoremail command
%% from AASTeX v4.0. You can use \email to mark an email address
%% anywhere in the paper, not just in the front matter.
%% As in the title, use \\ to force line breaks.

\author{S. Djorgovski\altaffilmark{1,2,3} and Ivan R. King\altaffilmark{1}}
\affil{Astronomy Department, University of California,
    Berkeley, CA 94720}

\author{C. D. Biemesderfer\altaffilmark{4,5}}
\affil{National Optical Astronomy Observatories, Tucson, AZ 85719}
\email{aastex-help@aas.org}

\and

\author{R. J. Hanisch\altaffilmark{5}}
\affil{Space Telescope Science Institute, Baltimore, MD 21218}

%% Notice that each of these authors has alternate affiliations, which
%% are identified by the \altaffilmark after each name.  Specify alternate
%% affiliation information with \altaffiltext, with one command per each
%% affiliation.

\altaffiltext{1}{Visiting Astronomer, Cerro Tololo Inter-American Observatory.
CTIO is operated by AURA, Inc.\ under contract to the National Science
Foundation.}
\altaffiltext{2}{Society of Fellows, Harvard University.}
\altaffiltext{3}{present address: Center for Astrophysics,
    60 Garden Street, Cambridge, MA 02138}
\altaffiltext{4}{Visiting Programmer, Space Telescope Science Institute}
\altaffiltext{5}{Patron, Alonso's Bar and Grill}

%% Mark off your abstract in the ``abstract'' environment. In the manuscript
%% style, abstract will output a Received/Accepted line after the
%% title and affiliation information. No date will appear since the author
%% does not have this information. The dates will be filled in by the
%% editorial office after submission.

\begin{abstract}
This is a preliminary report on surface photometry of the major
fraction of known globular clusters, to see which of them show the signs
of a collapsed core.
We also explore some diversionary mathematics and recreational tables.
\end{abstract}

%% Keywords should appear after the \end{abstract} command. The uncommented
%% example has been keyed in ApJ style. See the instructions to authors
%% for the journal to which you are submitting your paper to determine
%% what keyword punctuation is appropriate.

\keywords{globular clusters: general --- globular clusters: individual(NGC 6397,
NGC 6624, NGC 7078, Terzan 8}

%% From the front matter, we move on to the body of the paper.
%% In the first two sections, notice the use of the natbib \citep
%% and \citet commands to identify citations.  The citations are
%% tied to the reference list via symbolic KEYs. The KEY corresponds
%% to the KEY in the \bibitem in the reference list below. We have
%% chosen the first three characters of the first author's name plus
%% the last two numeral of the year of publication as our KEY for
%% each reference.


%% Authors who wish to have the most important objects in their paper
%% linked in the electronic edition to a data center may do so by tagging
%% their objects with \objectname{} or \object{}.  Each macro takes the
%% object name as its required argument. The optional, square-bracket 
%% argument should be used in cases where the data center identification
%% differs from what is to be printed in the paper.  The text appearing 
%% in curly braces is what will appear in print in the published paper. 
%% If the object name is recognized by the data centers, it will be linked
%% in the electronic edition to the object data available at the data centers  

\section{Introduction}

A focal problem today in the dynamics of globular clusters is core collapse.  
It has been predicted by theory for decades \citep{hen61,lyn68,spi85}, but
observation has been less alert to the phenomenon. For many years the
central brightness peak in M15 \citep{kin75,new78} seemed a unique anomaly.  
Then \citet{aur82} suggested a central peak in \object{NGC 6397}, and a 
limited photographic survey of ours \citep[Paper I]{djo84} found three more 
cases, \objectname{NGC 6624}, \objectname[M 15]{NGC 7078}, and 
\object[Cl 1938-341]{Terzan 8}), whose sharp center had often been 
remarked on \citep{can78}.

\section{Observations}

%% In a manner similar to \objectname authors can provide links to dataset
%% hosted at participating data centers via the \dataset{} command.  The
%% second curly bracket argument is printed in the text while the first
%% parentheses argument serves as the valid data set identifier.  Large
%% lists of data set are best provided in a table (see Table 3 for an example).
%% Valid data set identifiers should be obtained from the data center that
%% is currently hosting the data.

All our observations were short direct exposures with CCD's.  We also have
 a random {\it Chandra} data set \dataset{ADS/Sa.ASCA\#X/86008020} and a
neat \dataset[ADS/Sa.HST#Y0Q70101T]{HST FOS spectrum} that readers can access
via the links in the electronic edition.  Unfortunately this has nothing 
whatsoever to do with this research.  At Lick Observatory we used a TI 500$\times$500 
chip and a GEC 575$\times$385, on the 1-m Nickel reflector.  The only
filter available at Lick was red.  At CTIO we used a GEC 575$\times$385, with
$B, V,$ and $R$ filters, and an RCA 512$\times$320, with $U, B, V, R,$ and $I$
filters, on the 1.5-m reflector. In the CTIO observations we tried to
concentrate on the shortest practicable wavelengths; but faintness, reddening,
and poor short-wavelength sensitivity often kept us from observing in $U$ or
even in $B$. All four cameras had scales of the order of 0.4 arcsec/pixel, and
our field sizes were around 3 arcmin.

The CCD images are unfortunately not always suitable, for very poor
clusters or for clusters with large cores.  Since the latter are easily
studied by other means, we augmented our own CCD profiles by collecting
from the literature a number of star-count profiles 
\citep{kin68,pet76,har84,ort85}, as well as photoelectric profiles 
\citep{kin66,kin75} and electronographic profiles \citep{kro84}.
In a few cases we judged normality by eye estimates on one of the Sky
Surveys.

%% In this section, we use  the \subsection command to set off
%% a subsection.  \footnote is used to insert a footnote to the text.

%% Observe the use of the LaTeX \label
%% command after the \subsection to give a symbolic KEY to the
%% subsection for cross-referencing in a \ref command.
%% You can use LaTeX's \ref and \label commands to keep track of
%% cross-references to sections, equations, tables, and figures.
%% That way, if you change the order of any elements, LaTeX will
%% automatically renumber them.

%% This section also includes several of the displayed math environments
%% mentioned in the Author Guide.

\section{Helicity Amplitudes}

It has been realized that helicity amplitudes provide a convenient means
for Feynman diagram\footnote{Footnotes can be inserted like this.}
evaluations.  These amplitude-level techniques
are particularly convenient for calculations involving many Feynman
diagrams, where the usual trace techniques for the amplitude
squared becomes unwieldy.  Our calculations use the helicity techniques
developed by other authors \cite[]{hag86}; we briefly summarize below.

\subsection{Formalism} \label{bozomath}

%% The equation environment wil produce a numbered display equation.

A tree-level amplitude in $e^+e^-$ collisions can be expressed in
terms of fermion strings of the form
\begin{equation}
\bar v(p_2,\sigma_2)P_{-\tau}\hat a_1\hat a_2\cdots
\hat a_nu(p_1,\sigma_1) ,
\end{equation}
where $p$ and $\sigma$ label the initial $e^{\pm}$ four-momenta
and helicities $(\sigma = \pm 1)$, $\hat a_i=a^\mu_i\gamma_\nu$
and $P_\tau=\frac{1}{2}(1+\tau\gamma_5)$ is a chirality projection
operator $(\tau = \pm1)$.  The $a^\mu_i$ may be formed from particle
four-momenta, gauge-boson polarization vectors or fermion strings with
an uncontracted Lorentz index associated with final-state fermions.

%% The \notetoeditor{TEXT} command allows the author to communicate
%% information to the copy editor.  This information will appear as a
%% footnote on the printed copy for the manuscript style file.  Nothing will
%% appear on the printed copy if the preprint or
%% preprint2 style files are used.

%% The eqnarray environment produces multi-line display math. The end of
%% each line is marked with a \\. Lines will be numbered unless the \\
%% is preceded by a \nonumber command.
%% Alignment points are marked by ampersands (&). There should be two
%% ampersands (&) per line.

In the chiral \notetoeditor{Figures 1 and 2 should appear side-by-side in
print} representation the $\gamma$ matrices are expressed
in terms of $2\times 2$ Pauli matrices $\sigma$ and the unit matrix 1 as
\begin{eqnarray}
\gamma^\mu  & = &
 \left(
\begin{array}{cc}
0 & \sigma^\mu_+ \\
\sigma^\mu_- & 0
\end{array}     \right) ,
 \gamma^5= \left(
\begin{array}{cc}
-1 &   0\\
0 &   1
\end{array}     \right)  , \nonumber \\
\sigma^\mu_{\pm}  & = &   ({\bf 1} ,\pm \sigma) , \nonumber
\end{eqnarray}
giving
\begin{equation}
\hat a= \left(
\begin{array}{cc}
0 & (\hat a)_+\\
(\hat a)_- & 0
\end{array}\right), (\hat a)_\pm=a_\mu\sigma^\mu_\pm ,
\end{equation}
The spinors are expressed in terms of two-component Weyl spinors as
\begin{equation}
u=\left(
\begin{array}{c}
(u)_-\\
(u)_+
\end{array}\right), v={\bf (}\vdag_+{\bf ,}   \vdag_-{\bf )} .
\end{equation}

%% Putting eqnarrays or equations inside the mathletters environment groups
%% the enclosed equations by letter. For instance, the eqnarray below, instead
%% of being numbered, say, (4) and (5), would be numbered (4a) and (4b).
%% LaTeX the paper and look at the output to see the results.

The Weyl spinors are given in terms of helicity eigenstates
$\chi_\lambda(p)$ with $\lambda=\pm1$ by
\begin{mathletters}
\begin{eqnarray}
u(p,\lambda)_\pm & = & (E\pm\lambda|{\bf p}|)^{1/2}\chi_\lambda(p) , \\
v(p,\lambda)_\pm & = & \pm\lambda(E\mp\lambda|{\bf p}|)^{1/2}\chi
_{-\lambda}(p)
\end{eqnarray}
\end{mathletters}

%% This section contains more display math examples, including unnumbered
%% equations (displaymath environment). The last paragraph includes some
%% examples of in-line math featuring a couple of the AASTeX symbol macros.

\section{Floating material and so forth}

%% The displaymath environment will produce the same sort of equation as
%% the equation environment, except that the equation will not be numbered
%% by LaTeX.

Consider a task that computes profile parameters for a modified
Lorentzian of the form
\begin{equation}
I = \frac{1}{1 + d_{1}^{P (1 + d_{2} )}}
\end{equation}
where
\begin{displaymath}
d_{1} = \sqrt{ \left( \begin{array}{c} \frac{x_{1}}{R_{maj}}
\end{array} \right) ^{2} +
\left( \begin{array}{c} \frac{y_{1}}{R_{min}} \end{array} \right) ^{2} }
\end{displaymath}
\begin{displaymath}
d_{2} = \sqrt{ \left( \begin{array}{c} \frac{x_{1}}{P R_{maj}}
\end{array} \right) ^{2} +
\left( \begin{array}{c} \case{y_{1}}{P R_{min}} \end{array} \right) ^{2} }
\end{displaymath}
\begin{displaymath}
x_{1} = (x - x_{0}) \cos \Theta + (y - y_{0}) \sin \Theta
\end{displaymath}
\begin{displaymath}
y_{1} = -(x - x_{0}) \sin \Theta + (y - y_{0}) \cos \Theta
\end{displaymath}

In these expressions $x_{0}$,$y_{0}$ is the star center, and $\Theta$ is the
angle with the $x$ axis.  Results of this task are shown in table~\ref{tbl-1}.
It is not clear how these sorts of analyses may affect determination of
 $M_{\sun}$, but the assumption is that the alternate results
should be less than 90\degr\ out of phase with previous values.
We have no observations of \ion{Ca}{2}.
Roughly \slantfrac{4}{5} of the electronically submitted abstracts
for AAS meetings are error-free.

%% If you wish to include an acknowledgments section in your paper,
%% separate it off from the body of the text using the \acknowledgments
%% command.

%% Included in this acknowledgments section are examples of the
%% AASTeX hypertext markup commands. Use \url without the optional [HREF]
%% argument when you want to print the url directly in the text. Otherwise,
%% use either \url or \anchor, with the HREF as the first argument and the
%% text to be printed in the second.

\acknowledgments

We are grateful to V. Barger, T. Han, and R. J. N. Phillips for
doing the math in section~\ref{bozomath}.
More information on the AASTeX macros package is available \\ at
\url{http://www.aas.org/publications/aastex}.
For technical support, please write to
\email{aastex-help@aas.org}.

%% To help institutions obtain information on the effectiveness of their
%% telescopes, the AAS Journals has created a group of keywords for telescope
%% facilities. A common set of keywords will make these types of searches
%% significantly easier and more accurate. In addition, they will also be
%% useful in linking papers together which utilize the same telescopes
%% within the framework of the National Virtual Observatory.
%% See the AASTeX Web site at http://www.journals.uchicago.edu/AAS/AASTeX
%% for information on obtaining the facility keywords.

%% After the acknowledgments section, use the following syntax and the
%% \facility{} macro to list the keywords of facilities used in the research
%% for the paper.  Each keyword will be checked against the master list during
%% copy editing.  Individual instruments or configurations can be provided 
%% in parentheses, after the keyword, but they will not be verified.

{\it Facilities:} \facility{Nickel}, \facility{HST (STIS)}, \facility{CXO (ASIS)}.

%% Appendix material should be preceded with a single \appendix command.
%% There should be a \section command for each appendix. Mark appendix
%% subsections with the same markup you use in the main body of the paper.

%% Each Appendix (indicated with \section) will be lettered A, B, C, etc.
%% The equation counter will reset when it encounters the \appendix
%% command and will number appendix equations (A1), (A2), etc.

\appendix

\section{Appendix material}

Consider once again a task that computes profile parameters for a modified
Lorentzian of the form
\begin{equation}
I = \frac{1}{1 + d_{1}^{P (1 + d_{2} )}}
\end{equation}
where
\begin{mathletters}
\begin{displaymath}
d_{1} = \frac{3}{4} \sqrt{ \left( \begin{array}{c} \frac{x_{1}}{R_{maj}}
\end{array} \right) ^{2} +
\left( \begin{array}{c} \frac{y_{1}}{R_{min}} \end{array} \right) ^{2} }
\end{displaymath}
\begin{equation}
d_{2} = \case{3}{4} \sqrt{ \left( \begin{array}{c} \frac{x_{1}}{P R_{maj}}
\end{array} \right) ^{2} +
\left( \begin{array}{c} \case{y_{1}}{P R_{min}} \end{array} \right) ^{2} }
\end{equation}
\begin{eqnarray}
x_{1} & = & (x - x_{0}) \cos \Theta + (y - y_{0}) \sin \Theta \\
y_{1} & = & -(x - x_{0}) \sin \Theta + (y - y_{0}) \cos \Theta
\end{eqnarray}
\end{mathletters}

For completeness, here is one last equation.
\begin{equation}
e = mc^2
\end{equation}

%% The reference list follows the main body and any appendices.
%% Use LaTeX's thebibliography environment to mark up your reference list.
%% Note \begin{thebibliography} is followed by an empty set of
%% curly braces.  If you forget this, LaTeX will generate the error
%% "Perhaps a missing \item?".
%%
%% thebibliography produces citations in the text using \bibitem-\cite
%% cross-referencing. Each reference is preceded by a
%% \bibitem command that defines in curly braces the KEY that corresponds
%% to the KEY in the \cite commands (see the first section above).
%% Make sure that you provide a unique KEY for every \bibitem or else the
%% paper will not LaTeX. The square brackets should contain
%% the citation text that LaTeX will insert in
%% place of the \cite commands.

%% We have used macros to produce journal name abbreviations.
%% AASTeX provides a number of these for the more frequently-cited journals.
%% See the Author Guide for a list of them.

%% Note that the style of the \bibitem labels (in []) is slightly
%% different from previous examples.  The natbib system solves a host
%% of citation expression problems, but it is necessary to clearly
%% delimit the year from the author name used in the citation.
%% See the natbib documentation for more details and options.

\begin{thebibliography}{}
\bibitem[Auri\`ere(1982)]{aur82} Auri\`ere, M.  1982, \aap,
    109, 301
\bibitem[Canizares et al.(1978)]{can78} Canizares, C. R.,
    Grindlay, J. E., Hiltner, W. A., Liller, W., \&
    McClintock, J. E.  1978, \apj, 224, 39
\bibitem[Djorgovski \& King(1984)]{djo84} Djorgovski, S.,
    \& King, I. R.  1984, \apjl, 277, L49
\bibitem[Hagiwara \& Zeppenfeld(1986)]{hag86} Hagiwara, K., \&
    Zeppenfeld, D.  1986, Nucl.Phys., 274, 1
\bibitem[Harris \& van den Bergh(1984)]{har84} Harris, W. E.,
    \& van den Bergh, S.  1984, \aj, 89, 1816
\bibitem[H\`enon(1961)]{hen61} H\'enon, M.  1961, Ann.d'Ap., 24, 369
\bibitem[Heiles \& Troland(2003)]{heiles03} Heiles, C. \& Troland, T. H., 2003, \apjs, preprint doi:10.1086/381753
\bibitem[Kim, Ostricker, \& Stone(2003)]{kim03} Kim, W.-T.,  Ostriker, E., \& Stone, J. M., 2003, \apj, 599, 1157
\bibitem[King(1966)]{kin66}  King, I. R.  1966, \aj, 71, 276
\bibitem[King(1975)]{kin75}  King, I. R.  1975, Dynamics of
    Stellar Systems, A. Hayli, Dordrecht: Reidel, 1975, 99
\bibitem[King et al.(1968)]{kin68}  King, I. R., Hedemann, E.,
    Hodge, S. M., \& White, R. E.  1968, \aj, 73, 456
\bibitem[Kron et al.(1984)]{kro84} Kron, G. E., Hewitt, A. V.,
    \& Wasserman, L. H.  1984, \pasp, 96, 198
\bibitem[Lynden-Bell \& Wood(1968)]{lyn68} Lynden-Bell, D.,
    \& Wood, R.  1968, \mnras, 138, 495
\bibitem[Newell \& O'Neil(1978)]{new78} Newell, E. B.,
    \& O'Neil, E. J.  1978, \apjs, 37, 27
\bibitem[Ortolani et al.(1985)]{ort85} Ortolani, S., Rosino, L.,
    \& Sandage, A.  1985, \aj, 90, 473
\bibitem[Peterson(1976)]{pet76} Peterson, C. J.  1976, \aj, 81, 617
\bibitem[Rudnick et al.(2003)]{rudnick03} Rudnick, G. et al., 2003, \apj, 599, 847
\bibitem[Spitzer(1985)]{spi85} Spitzer, L.  1985, Dynamics of
    Star Clusters, J. Goodman \& P. Hut, Dordrecht: Reidel, 109
\bibitem[Treu et al.(2003)]{treu03} Treu, T. et al., 2003, \apj, 591, 53
\end{thebibliography}

\clearpage

%% Use the figure environment and \plotone or \plottwo to include
%% figures and captions in your electronic submission.
%% To embed the sample graphics in
%% the file, uncomment the \plotone, \plottwo, and
%% \includegraphics commands
%%
%% If you need a layout that cannot be achieved with \plotone or
%% \plottwo, you can invoke the graphicx package directly with the
%% \includegraphics command or use \plotfiddle. For more information,
%% please see the tutorial on "Using Electronic Art with AASTeX" in the
%% documentation section at the AASTeX Web site,
%% http://www.journals.uchicago.edu/AAS/AASTeX.
%%
%% The examples below also include sample markup for submission of
%% supplemental electronic materials. As always, be sure to check
%% the instructions to authors for the journal you are submitting to
%% for specific submissions guidelines as they vary from
%% journal to journal.

%% This example uses \plotone to include an EPS file scaled to
%% 80% of its natural size with \epsscale. Its caption
%% has been written to indicate that additional figure parts will be
%% available in the electronic journal.

\begin{figure}
\epsscale{.80}
\plotone{f1.eps}
\caption{Derived spectra for 3C138 \citep[see][]{heiles03}. Plots for all sources are available
in the electronic edition of {\it The Astrophysical Journal}.\label{fig1}}
\end{figure}

\clearpage

%% Here we use \plottwo to present two versions of the same figure,
%% one in black and white for print the other in RGB color
%% for online presentation. Note that the caption indicates
%% that a color version of the figure will be available online.
%%

\begin{figure}
\plottwo{f2.eps}{f2_color.eps}
\caption{A panel taken from Figure 2 of \citet{rudnick03}. 
See the electronic edition of the Journal for a color version 
of this figure.\label{fig2}}
\end{figure}

%% This figure uses \includegraphics to scale and rotate the still frame
%% for an mpeg animation.

\begin{figure}
\includegraphics[angle=90,scale=.50]{f3.eps}
\caption{Animation still frame taken from \citet{kim03}.
This figure is also available as an mpeg
animation in the electronic edition of the
{\it Astrophysical Journal}.}
\end{figure}

%% If you are not including electonic art with your submission, you may
%% mark up your captions using the \figcaption command. See the
%% User Guide for details.
%%
%% No more than seven \figcaption commands are allowed per page,
%% so if you have more than seven captions, insert a \clearpage
%% after every seventh one.

%% Tables should be submitted one per page, so put a \clearpage before
%% each one.

%% Two options are available to the author for producing tables:  the
%% deluxetable environment provided by the AASTeX package or the LaTeX
%% table environment.  Use of deluxetable is preferred.
%%

%% Three table samples follow, two marked up in the deluxetable environment,
%% one marked up as a LaTeX table.

%% In this first example, note that the \tabletypesize{}
%% command has been used to reduce the font size of the table.
%% We also use the \rotate command to rotate the table to
%% landscape orientation since it is very wide even at the
%% reduced font size.
%%
%% Note also that the \label command needs to be placed
%% inside the \tablecaption.

%% This table also includes a table comment indicating that the full
%% version will be available in machine-readable format in the electronic
%% edition.

\clearpage

\begin{deluxetable}{ccrrrrrrrrcrl}
\tabletypesize{\scriptsize}
\rotate
\tablecaption{Sample table taken from \citet{treu03}\label{tbl-1}}
\tablewidth{0pt}
\tablehead{
\colhead{POS} & \colhead{chip} & \colhead{ID} & \colhead{X} & \colhead{Y} &
\colhead{RA} & \colhead{DEC} & \colhead{IAU$\pm$ $\delta$ IAU} &
\colhead{IAP1$\pm$ $\delta$ IAP1} & \colhead{IAP2 $\pm$ $\delta$ IAP2} &
\colhead{star} & \colhead{E} & \colhead{Comment}
}
\startdata
0 & 2 & 1 & 1370.99 & 57.35    &   6.651120 &  17.131149 & 21.344$\pm$0.006  & 2
4.385$\pm$0.016 & 23.528$\pm$0.013 & 0.0 & 9 & -    \\
0 & 2 & 2 & 1476.62 & 8.03     &   6.651480 &  17.129572 & 21.641$\pm$0.005  & 2
3.141$\pm$0.007 & 22.007$\pm$0.004 & 0.0 & 9 & -    \\
0 & 2 & 3 & 1079.62 & 28.92    &   6.652430 &  17.135000 & 23.953$\pm$0.030  & 2
4.890$\pm$0.023 & 24.240$\pm$0.023 & 0.0 & - & -    \\
0 & 2 & 4 & 114.58  & 21.22    &   6.655560 &  17.148020 & 23.801$\pm$0.025  & 2
5.039$\pm$0.026 & 24.112$\pm$0.021 & 0.0 & - & -    \\
0 & 2 & 5 & 46.78   & 19.46    &   6.655800 &  17.148932 & 23.012$\pm$0.012  & 2
3.924$\pm$0.012 & 23.282$\pm$0.011 & 0.0 & - & -    \\
0 & 2 & 6 & 1441.84 & 16.16    &   6.651480 &  17.130072 & 24.393$\pm$0.045  & 2
6.099$\pm$0.062 & 25.119$\pm$0.049 & 0.0 & - & -    \\
0 & 2 & 7 & 205.43  & 3.96     &   6.655520 &  17.146742 & 24.424$\pm$0.032  & 2
5.028$\pm$0.025 & 24.597$\pm$0.027 & 0.0 & - & -    \\
0 & 2 & 8 & 1321.63 & 9.76     &   6.651950 &  17.131672 & 22.189$\pm$0.011  & 2
4.743$\pm$0.021 & 23.298$\pm$0.011 & 0.0 & 4 & edge \\
\enddata
%% Text for table notes should follow after the \enddata but before
%% the \end{deluxetable}. Make sure there is at least one \tablenotemark
%% in the table for each \tablenotetext.
\tablecomments{Table \ref{tbl-1} is published in its entirety in the 
electronic edition of the {\it Astrophysical Journal}.  A portion is 
shown here for guidance regarding its form and content.}
\tablenotetext{a}{Sample footnote for table~\ref{tbl-1} that was generated
with the deluxetable environment}
\tablenotetext{b}{Another sample footnote for table~\ref{tbl-1}}
\end{deluxetable}

%% If you use the table environment, please indicate horizontal rules using
%% \tableline, not \hline.
%% Do not put multiple tabular environments within a single table.
%% The optional \label should appear inside the \caption command.

\clearpage

\begin{table}
\begin{center}
\caption{More terribly relevant tabular information.\label{tbl-2}}
\begin{tabular}{crrrrrrrrrrr}
\tableline\tableline
Star & Height & $d_{x}$ & $d_{y}$ & $n$ & $\chi^2$ & $R_{maj}$ & $R_{min}$ &
\multicolumn{1}{c}{$P$\tablenotemark{a}} & $P R_{maj}$ & $P R_{min}$ &
\multicolumn{1}{c}{$\Theta$\tablenotemark{b}} \\
\tableline
1 &33472.5 &-0.1 &0.4  &53 &27.4 &2.065  &1.940 &3.900 &68.3 &116.2 &-27.639\\
2 &27802.4 &-0.3 &-0.2 &60 &3.7  &1.628  &1.510 &2.156 &6.8  &7.5 &-26.764\\
3 &29210.6 &0.9  &0.3  &60 &3.4  &1.622  &1.551 &2.159 &6.7  &7.3 &-40.272\\
4 &32733.8 &-1.2\tablenotemark{c} &-0.5 &41 &54.8 &2.282  &2.156 &4.313 &117.4 &78.2 &-35.847\\
5 & 9607.4 &-0.4 &-0.4 &60 &1.4  &1.669\tablenotemark{c}  &1.574 &2.343 &8.0  &8.9 &-33.417\\
6 &31638.6 &1.6  &0.1  &39 &315.2 & 3.433 &3.075 &7.488 &92.1 &25.3 &-12.052\\
\tableline
\end{tabular}
%% Any table notes must follow the \end{tabular} command.
\tablenotetext{a}{Sample footnote for table~\ref{tbl-2} that was
generated with the \LaTeX\ table environment}
\tablenotetext{b}{Yet another sample footnote for table~\ref{tbl-2}}
\tablenotetext{c}{Another sample footnote for table~\ref{tbl-2}}
\tablecomments{We can also attach a long-ish paragraph of explanatory
material to a table.}
\end{center}
\end{table}

%% If the table is more than one page long, the width of the table can vary
%% from page to page when the default \tablewidth is used, as below.  The
%% individual table widths for each page will be written to the log file; a
%% maximum tablewidth for the table can be computed from these values.
%% The \tablewidth argument can then be reset and the file reprocessed, so
%% that the table is of uniform width throughout. Try getting the widths
%% from the log file and changing the \tablewidth parameter to see how
%% adjusting this value affects table formatting.

%% The \dataset{} macro has also been applied to a few of the objects to
%% show how many observations can be tagged in a table.

\clearpage

\begin{deluxetable}{lrrrrcrrrrr}
\tablewidth{0pt}
\tablecaption{Literature Data for Program Stars}
\tablehead{
\colhead{Star}           & \colhead{V}      &
\colhead{b$-$y}          & \colhead{m$_1$}  &
\colhead{c$_1$}          & \colhead{ref}    &
\colhead{T$_{\rm eff}$}  & \colhead{log g}  &
\colhead{v$_{\rm turb}$} & \colhead{[Fe/H]} &
\colhead{ref}}
\startdata
HD 97 & 9.7& 0.51& 0.15& 0.35& 2 & \nodata & \nodata & \nodata & $-1.50$ & 2 \\
& & & & & & 5015 & \nodata & \nodata & $-1.50$ & 10 \\
\dataset[ADS/Sa.HST#O6H04VAXQ]{HD 2665} & 7.7& 0.54& 0.09& 0.34& 2 & \nodata & \nodata & \nodata & $-2.30$ & 2 \\
& & & & & & 5000 & 2.50 & 2.4 & $-1.99$ & 5 \\
& & & & & & 5120 & 3.00 & 2.0 & $-1.69$ & 7 \\
& & & & & & 4980 & \nodata & \nodata & $-2.05$ & 10 \\
HD 4306 & 9.0& 0.52& 0.05& 0.35& 20, 2& \nodata & \nodata & \nodata & $-2.70$ & 2 \\
& & & & & & 5000 & 1.75 & 2.0 & $-2.70$ & 13 \\
& & & & & & 5000 & 1.50 & 1.8 & $-2.65$ & 14 \\
& & & & & & 4950 & 2.10 & 2.0 & $-2.92$ & 8 \\
& & & & & & 5000 & 2.25 & 2.0 & $-2.83$ & 18 \\
& & & & & & \nodata & \nodata & \nodata & $-2.80$ & 21 \\
& & & & & & 4930 & \nodata & \nodata & $-2.45$ & 10 \\
HD 5426 & 9.6& 0.50& 0.08& 0.34& 2 & \nodata & \nodata & \nodata & $-2.30$ & 2 \\
\dataset[ADS/Sa.HST#O5F654010]{HD 6755} & 7.7& 0.49& 0.12& 0.28& 20, 2& \nodata & \nodata & \nodata & $-1.70$ & 2 \\
& & & & & & 5200 & 2.50 & 2.4 & $-1.56$ & 5 \\
& & & & & & 5260 & 3.00 & 2.7 & $-1.67$ & 7 \\
& & & & & & \nodata & \nodata & \nodata & $-1.58$ & 21 \\
& & & & & & 5200 & \nodata & \nodata & $-1.80$ & 10 \\
& & & & & & 4600 & \nodata & \nodata & $-2.75$ & 10 \\
\dataset[ADS/Sa.HST#O56D06010]{HD 94028} & 8.2& 0.34& 0.08& 0.25& 20 & 5795 & 4.00 & \nodata & $-1.70$ & 22 \\
& & & & & & 5860 & \nodata & \nodata & $-1.70$ & 4 \\
& & & & & & 5910 & 3.80 & \nodata & $-1.76$ & 15 \\
& & & & & & 5800 & \nodata & \nodata & $-1.67$ & 17 \\
& & & & & & 5902 & \nodata & \nodata & $-1.50$ & 11 \\
& & & & & & 5900 & \nodata & \nodata & $-1.57$ & 3 \\
& & & & & & \nodata & \nodata & \nodata & $-1.32$ & 21 \\
HD 97916 & 9.2& 0.29& 0.10& 0.41& 20 & 6125 & 4.00 & \nodata & $-1.10$ & 22 \\
& & & & & & 6160 & \nodata & \nodata & $-1.39$ & 3 \\
& & & & & & 6240 & 3.70 & \nodata & $-1.28$ & 15 \\
& & & & & & 5950 & \nodata & \nodata & $-1.50$ & 17 \\
& & & & & & 6204 & \nodata & \nodata & $-1.36$ & 11 \\
\cutinhead{This is a cut-in head}
+26\arcdeg2606& 9.7&0.34&0.05&0.28&20,11& 5980 & \nodata & \nodata &$<-2.20$ & 19 \\
& & & & & & 5950 & \nodata & \nodata & $-2.89$ & 24 \\
+26\arcdeg3578& 9.4&0.31&0.05&0.37&20,11& 5830 & \nodata & \nodata & $-2.60$ & 4 \\
& & & & & & 5800 & \nodata & \nodata & $-2.62$ & 17 \\
& & & & & & 6177 & \nodata & \nodata & $-2.51$ & 11 \\
& & & & & & 6000 & 3.25 & \nodata & $-2.20$ & 22 \\
& & & & & & 6140 & 3.50 & \nodata & $-2.57$ & 15 \\
+30\arcdeg2611& 9.2&0.82&0.33&0.55& 2 & \nodata & \nodata & \nodata & $-1.70$ & 2 \\
& & & & & & 4400 & 1.80 & \nodata & $-1.70$ & 12 \\
& & & & & & 4400 & 0.90 & 1.7 & $-1.20$ & 14 \\
& & & & & & 4260 & \nodata & \nodata & $-1.55$ & 10 \\
+37\arcdeg1458& 8.9&0.44&0.07&0.22&20,11& 5296 & \nodata & \nodata & $-2.39$ & 11 \\
& & & & & & 5420 & \nodata & \nodata & $-2.43$ & 3 \\
+58\arcdeg1218&10.0&0.51&0.03&0.36& 2 & \nodata & \nodata & \nodata & $-2.80$ & 2 \\
& & & & & & 5000 & 1.10 & 2.2 & $-2.71$ & 14 \\
& & & & & & 5000 & 2.20 & 1.8 & $-2.46$ & 5 \\
& & & & & & 4980 & \nodata & \nodata & $-2.55$ & 10 \\
+72\arcdeg0094&10.2&0.31&0.09&0.26&12 & 6160 & \nodata & \nodata & $-1.80$ & 19 \\
\sidehead{I'm a side head:}
G5--36 & 10.8& 0.40& 0.07& 0.28& 20 & \nodata & \nodata & \nodata & $-1.19$ & 21 \\
G18--54 & 10.7& 0.37& 0.08& 0.28& 20 & \nodata & \nodata & \nodata & $-1.34$ & 21 \\
G20--08 & 9.9& 0.36& 0.05& 0.25& 20,11& 5849 & \nodata & \nodata & $-2.59$ & 11 \\
& & & & & & \nodata & \nodata & \nodata & $-2.03$ & 21 \\
G20--15 & 10.6& 0.45& 0.03& 0.27& 20,11& 5657 & \nodata & \nodata & $-2.00$ & 11 \\
& & & & & & 6020 & \nodata & \nodata & $-1.56$ & 3 \\
& & & & & & \nodata & \nodata & \nodata & $-1.58$ & 21 \\
G21--22 & 10.7& 0.38& 0.07& 0.27& 20,11& \nodata & \nodata & \nodata & $-1.23$ & 21 \\
G24--03 & 10.5& 0.36& 0.06& 0.27& 20,11& 5866 & \nodata & \nodata & $-1.78$ & 11 \\
& & & & & & \nodata & \nodata & \nodata & $-1.70$ & 21 \\
G30--52 & 8.6& 0.50& 0.25& 0.27& 11 & 4757 & \nodata & \nodata & $-2.12$ & 11 \\
& & & & & & 4880 & \nodata & \nodata & $-2.14$ & 3 \\
G33--09 & 10.6& 0.41& 0.10& 0.28& 20 & 5575 & \nodata & \nodata & $-1.48$ & 11 \\
G66--22 & 10.5& 0.46& 0.16& 0.28& 11 & 5060 & \nodata & \nodata & $-1.77$ & 3 \\
& & & & & & \nodata & \nodata & \nodata & $-1.04$ & 21 \\
G90--03 & 10.4& 0.37& 0.04& 0.29& 20 & \nodata & \nodata & \nodata & $-2.01$ & 21 \\
LP 608--62\tablenotemark{a} & 10.5& 0.30& 0.07& 0.35& 11 & 6250 & \nodata &
\nodata & $-2.70$ & 4 \\
\enddata
\tablenotetext{a}{Star LP 608--62 is also known as BD+1\arcdeg 2341p.  We will
make this footnote extra long so that it extends over two lines.}
%% You can append references to a table using the \tablerefs command.
\tablerefs{
(1) Barbuy, Spite, \& Spite 1985; (2) Bond 1980; (3) Carbon et al. 1987;
(4) Hobbs \& Duncan 1987; (5) Gilroy et al. 1988: (6) Gratton \& Ortolani 1986;
(7) Gratton \& Sneden 1987; (8) Gratton \& Sneden (1988); (9) Gratton \& Sneden 1991;
(10) Kraft et al. 1982; (11) LCL, or Laird, 1990; (12) Leep \& Wallerstein 1981;
(13) Luck \& Bond 1981; (14) Luck \& Bond 1985; (15) Magain 1987;
(16) Magain 1989; (17) Peterson 1981; (18) Peterson, Kurucz, \& Carney 1990;
(19) RMB; (20) Schuster \& Nissen 1988; (21) Schuster \& Nissen 1989b;
(22) Spite et al. 1984; (23) Spite \& Spite 1986; (24) Hobbs \& Thorburn 1991;
(25) Hobbs et al. 1991; (26) Olsen 1983.}
\end{deluxetable}

%% Tables may also be prepared as separate files. See the accompanying
%% sample file table.tex for an example of an external table file.
%% To include an external file in your main document, use the \input
%% command. Uncomment the line below to include table.tex in this
%% sample file. (Note that you will need to comment out the \documentclass,
%% \begin{document}, and \end{document} commands from table.tex if you want
%% to include it in this document.)

%% \input{table}

%% The following command ends your manuscript. LaTeX will ignore any text
%% that appears after it.

\end{document}

%%
%% End of file `sample.tex'.
%</sample>
%    \end{macrocode}
%
% Here begins the module `table`.
%    \begin{macrocode}
%<*table>
%% 
%% Begining of file `table.tex'

%% This complex but short example prepared in the deluxetable environment 
%% demonstrates some of the techniques 
%% that can be used to generate complex column headings and to align 
%% variable-width columns. See the manuscript sample file, sample.tex, 
%% for more table examples.

%% Note this file has its own \documentclass, \begin{document}, and
%% \end{document} commands. If you want to insert this table in another 
%% LaTeX document using the \input command, comment out these lines.

\documentclass{aastex}
\begin{document}

%% Note that the table will print double-spaced since we are using the
%% manuscript style. Change the style to preprint or preprint2 to see 
%% how LaTeX formats the table in those styles.

%% In this example the LaTeX \multicolumn command is used to span a heading 
%% over several columns.  When \multicolumn is used along with the 
%% \cutinhead or \sidehead commands, the \tablecolumns command must 
%% be used to specify the number of columns in the table - 
%% otherwise \cutinhead and \sidehead will not work properly. 
 
%% \cline has been used to produce straddle rules below the spanning heads, 
%% \cutinhead to produce a centered head in the body of the table, and 
%% \sidehead to produce a flush-left head in the body.

%% This table also makes use of the \phn command to better align some of the 
%% columns.  Also see \phd, \phs, and \phm{} - other commands useful for 
%% column alignment.  All of these commands insert a blank space 
%% whose width is  equal to that of a number (\phn), 
%% a decimal point (\phd), a minus sign (\phs), or any 
%% character you wish to use (\phm{text}). 
%% Keep in mind that if you are preparing a table for electronic submission 
%% to one of the journals, you need not worry too much about column 
%% alignment. The editors will fix table alignment as appropriate.

%% If a table is more than one page long, the width of the table can vary 
%% from page to page when the default \tablewidth is used, as below.  The 
%% individual table widths for each page will be written to the log file; a 
%% maximum tablewidth for the table can be computed from these values. 
%% The \tablewidth argument can then be reset and the file reprocessed, so 
%% that the table is uniform throughout the pages. Try getting the widths
%% from the log file and changing the \tablewidth parameter to see how 
%% adjusting this value affects table formatting.

%% The * option to the \\ command has been used in the lines after 
%% the \sidehead to keep them together on the same page. Try taking 
%% the *'s out and LaTeXing the manuscript again to see the difference 
%% in the page breaks. You can group together as many lines as 
%% you like using this command. 


\begin{deluxetable}{rrrrrrrr} 
\tablecolumns{8} 
\tablewidth{0pc} 
\tablecaption{Percentage of Fake Stars Lost} 
\tablehead{ 
\colhead{}    &  \multicolumn{3}{c}{Non-shell Stars} &   \colhead{}   & 
\multicolumn{3}{c}{Shell Stars} \\ 
\cline{2-4} \cline{6-8} \\ 
\colhead{Mag} & \colhead{F336W}   & \colhead{F555W}    & \colhead{F814W} & 
\colhead{}    & \colhead{F336W}   & \colhead{F555W}    & \colhead{F814W}}
\startdata 
20.25 & 2.2$\pm$7.4\phn & \nodata & \nodata & 
& 0.9$\pm$6.8 & \nodata & 0.0$\pm$44.7 \\ 
20.75 & 2.4$\pm$7.8\phn & \nodata & 2.8$\pm$7.4 & 
& 1.7$\pm$6.6 & \nodata & 1.4$\pm$6.7\phn \\ 
21.25 & 0.1$\pm$7.7\phn & \nodata & 1.7$\pm$7.6 & 
& 2.6$\pm$6.5 & \nodata & 0.9$\pm$6.6\phn \\ 
21.75 & 2.4$\pm$4.5\phn & 2.2$\pm$7.4 & 0.1$\pm$7.6 & 
& 7.1$\pm$4.5 & 0.9$\pm$6.8 & 3.3$\pm$6.5\phn \\ 
22.25 & 3.4$\pm$3.1\phn & 1.8$\pm$7.7 & 2.9$\pm$4.4 & 
& 11.8$\pm$3.6 & 0.4$\pm$6.6 & 5.7$\pm$4.4\phn \\ 
22.75 & 4.5$\pm$2.9\phn & 1.8$\pm$7.7 & 4.7$\pm$3.1 & 
& 26.2$\pm$3.6 & 3.4$\pm$6.5 & 10.9$\pm$3.6\phn \\ 
23.25 & 7.0$\pm$2.4\phn & 3.4$\pm$4.5 & 3.7$\pm$2.9 & 
& 44.2$\pm$3.3 & 10.7$\pm$4.5 & 20.6$\pm$3.5\phn \\ 
\cutinhead{More Data} 
23.75 & 12.4$\pm$2.7\phn & 4.1$\pm$3.1 & 6.7$\pm$2.5 & 
& 59.8$\pm$4.0 & 20.1$\pm$3.6 & 32.6$\pm$3.4\phn \\ 
24.25 & 30.2$\pm$3.1\phn & 5.3$\pm$2.9 & 10.0$\pm$2.7 & 
& 74.9$\pm$5.1 & 35.8$\pm$3.6 & 43.1$\pm$4.0\phn \\ 
24.75 & 66.8$\pm$5.5\phn & 10.4$\pm$2.4 & 16.5$\pm$3.2 & 
& 83.7$\pm$6.1 & 56.3$\pm$3.3 & 57.0$\pm$5.2\phn \\ 
25.25 & 87.5$\pm$35.4 & 20.0$\pm$2.7 & 28.0$\pm$5.6 & 
& \nodata & 71.5$\pm$4.0 & 71.8$\pm$6.2\phn \\ 
25.75 & \nodata\phn & 55.3$\pm$3.1 & \nodata & 
& \nodata & 81.2$\pm$5.1 & \nodata\phn \\
26.25 & \nodata\phn & 85.1$\pm$5.5 & \nodata & 
& \nodata & 85.6$\pm$6.1 & \nodata\phn \\
\sidehead{More Data}
27.75 & 12.4$\pm$2.7\phn & 4.1$\pm$3.1 & 6.7$\pm$2.5 & 
& 59.8$\pm$4.0 & 20.1$\pm$3.6 & 32.6$\pm$3.4\phn \\
28.25 & 30.2$\pm$3.1\phn & 5.3$\pm$2.9 & 10.0$\pm$2.7 & 
& 74.9$\pm$5.1 & 35.8$\pm$3.6 & 43.1$\pm$4.0\phn \\ 
29.75 & 66.8$\pm$5.5\phn & 10.4$\pm$2.4 & 16.5$\pm$3.2 & 
& 83.7$\pm$6.1 & 56.3$\pm$3.3 & 57.0$\pm$5.2\phn \\ 
30.25 & 87.5$\pm$35.4 & 20.0$\pm$2.7 & 28.0$\pm$5.6 & 
& \nodata & 71.5$\pm$4.0 & 71.8$\pm$6.2\phn \\ 
31.75 & \nodata\phn & 55.3$\pm$3.1 & \nodata & 
& \nodata & 81.2$\pm$5.1 & \nodata\phn \\ 
32.25 & \nodata\phn & 85.1$\pm$5.5 & \nodata & 
& \nodata & 85.6$\pm$6.1 & \nodata\phn \\ 
33.75 & 12.4$\pm$2.7\phn & 4.1$\pm$3.1 & 6.7$\pm$2.5 & 
& 59.8$\pm$4.0 & 20.1$\pm$3.6 & 32.6$\pm$3.4\phn \\ 
34.25 & 30.2$\pm$3.1\phn & 5.3$\pm$2.9 & 10.0$\pm$2.7 & 
& 74.9$\pm$5.1 & 35.8$\pm$3.6 & 43.1$\pm$4.0\phn \\ 
35.75 & 66.8$\pm$5.5\phn & 10.4$\pm$2.4 & 16.5$\pm$3.2 & 
& 83.7$\pm$6.1 & 56.3$\pm$3.3 & 57.0$\pm$5.2\phn \\ 
36.25 & 87.5$\pm$35.4 & 20.0$\pm$2.7 & 28.0$\pm$5.6 & 
& \nodata & 71.5$\pm$4.0 & 71.8$\pm$6.2\phn \\ 
37.75 & \nodata\phn & 55.3$\pm$3.1 & \nodata & 
& \nodata & 81.2$\pm$5.1 & \nodata\phn \\
38.25 & \nodata\phn & 85.1$\pm$5.5 & \nodata & 
& \nodata & 85.6$\pm$6.1 & \nodata\phn \\
\enddata 
\end{deluxetable} 

 \end{document}

%% 
%% End of file `table.tex'. 


%</table>
%    \end{macrocode}
%
%
% docstrip will complete its |\generate| command at the following line:
\endinput

Textures 2.0.2 (preloaded format=LaTeX R9P0 98.10.16)  20 JAN 1999 12:29
(aasclass
LaTeX2e <1998/06/01>
(aasclass.dtx
File: aasclass.dtx 1999/01/22 5.0rc2 AAS markup document class

LaTeX Warning: File `ltxdoc.cfg' already exists on the system.
               Not generating it from this source.

(ltxdoc.cls
Document Class: ltxdoc 1996/01/11 v2.0r Standard LaTeX documentation class
*************************************
* Local config file ltxdoc.cfg used
*************************************
(ltxdoc.cfg) (article.cls
Document Class: article 1998/05/05 v1.3y Standard LaTeX document class
(size10.clo
File: size10.clo 1998/05/05 v1.3y Standard LaTeX file (size option)
)
\c@part=\count79
\c@section=\count80
\c@subsection=\count81
\c@subsubsection=\count82
\c@paragraph=\count83
\c@subparagraph=\count84
\c@figure=\count85
\c@table=\count86
\abovecaptionskip=\skip41
\belowcaptionskip=\skip42
\bibindent=\dimen102
) (doc.sty
Package: doc 1998/05/19 v2.0b Standard LaTeX documentation package (FMi)
\c@StandardModuleDepth=\count87
\guard@level=\count88
\MacrocodeTopsep=\skip43
\MacroIndent=\dimen103
\macro@cnt=\count89
\MacroTopsep=\skip44
\index@excludelist=\toks14
(multicol.sty
Package: multicol 1998/01/19 v1.5q multicolumn formatting (FMi)
\c@tracingmulticols=\count90
\mult@box=\box26
\multicol@leftmargin=\dimen104
\c@unbalance=\count91
\c@collectmore=\count92
\doublecol@number=\count93
\multicoltolerance=\count94
\multicolpretolerance=\count95
\full@width=\dimen105
\page@free=\dimen106
\premulticols=\dimen107
\postmulticols=\dimen108
\multicolsep=\skip45
\multicolbaselineskip=\skip46
\partial@page=\box27
\last@line=\box28
\mult@rightbox=\box29
\mult@grightbox=\box30
\mult@gfirstbox=\box31
\mult@firstbox=\box32
\@tempa=\box33
\@tempa=\box34
\@tempa=\box35
\@tempa=\box36
\@tempa=\box37
\@tempa=\box38
\@tempa=\box39
\@tempa=\box40
\@tempa=\box41
\@tempa=\box42
\@tempa=\box43
\@tempa=\box44
\@tempa=\box45
\@tempa=\box46
\@tempa=\box47
\@tempa=\box48
\@tempa=\box49
\c@columnbadness=\count96
\c@finalcolumnbadness=\count97
\last@try=\dimen109
\multicolovershoot=\dimen110
\multicolundershoot=\dimen111
\mult@nat@firstbox=\box50
)
\IndexMin=\dimen112
\c@IndexColumns=\count98
\GlossaryMin=\dimen113
\c@GlossaryColumns=\count99
\check@sum=\count100
\bslash@cnt=\count101
\wrong@table=\toks15
\c@CodelineNo=\count102
))
\@indexfile=\write3
Writing index file aasclass.idx
(shortvrb.sty
Package: shortvrb 1998/05/19 v2.0b Standard LaTeX documentation package (FMi)
)
Package shortvrb Info: Made | a short reference for \verb on input line 321.
LaTeX Info: Redefining \cs on input line 358.

LaTeX Warning: Writing file `:aasclass.ins'.



LaTeX Warning: Writing file `:Read.Me'.


(aasclass.aux)
LaTeX Font Info:    Checking defaults for OML/cmm/m/it on input line 585.
LaTeX Font Info:    ... okay on input line 585.
LaTeX Font Info:    Checking defaults for T1/cmr/m/n on input line 585.
LaTeX Font Info:    ... okay on input line 585.
LaTeX Font Info:    Checking defaults for OT1/cmr/m/n on input line 585.
LaTeX Font Info:    ... okay on input line 585.
LaTeX Font Info:    Checking defaults for OMS/cmsy/m/n on input line 585.
LaTeX Font Info:    ... okay on input line 585.
LaTeX Font Info:    Checking defaults for OMX/cmex/m/n on input line 585.
LaTeX Font Info:    ... okay on input line 585.
LaTeX Font Info:    Checking defaults for U/cmr/m/n on input line 585.
LaTeX Font Info:    ... okay on input line 585.
Package shortvrb Info: | already a short reference for \verb on input line 585.
(aasclass.dtx
***************************
* Character table correct *
***************************
(Read.Me
LaTeX Font Info:    External font `cmex10' loaded for size
(Font)              <17.28> on input line 16.
LaTeX Font Info:    External font `cmex10' loaded for size
(Font)              <12> on input line 16.
LaTeX Font Info:    External font `cmex10' loaded for size
(Font)              <8> on input line 16.
LaTeX Font Info:    External font `cmex10' loaded for size
(Font)              <6> on input line 16.
LaTeX Font Info:    External font `cmex10' loaded for size
(Font)              <5> on input line 16.

Overfull \hbox (8.91502pt too wide) in paragraph at lines 24--27
[]\OT1/cmr/m/n/10 The dis-tri-bu-tion point for AAST[]X is \OT1/cmtt/m/n/10 ftp:
//ftp.aas.org/pubs/aastex50.zip\OT1/cmr/m/n/10 ,|
 []

LaTeX Font Info:    External font `cmex10' loaded for size
(Font)              <7> on input line 29.
(aasclass.toc [1

] [2])
\tf@toc=\write4
[3]
Overfull \hbox (15.35983pt too wide) in paragraph at lines 65--68
[]\OT1/cmr/m/n/10 Now type-set \OT1/cmtt/m/n/10 aasclass.ins\OT1/cmr/m/n/10 , th
ereby gen-er-at-ing \OT1/cmtt/m/n/10 aastex.cls \OT1/cmr/m/n/10 and \OT1/cmtt/m/
n/10 aasguide.tex\OT1/cmr/m/n/10 . |
 []

)
LaTeX Font Info:    Font shape `OT1/cmtt/bx/n' in size <10> not available
(Font)              Font shape `OT1/cmtt/m/n' tried instead on input line 79.
[4]
Overfull \hbox (0.5877pt too wide) in paragraph at lines 140--140
 [][]\OT1/cmtt/m/n/9 %<style>  \@obsoletefile{aasclass.cls}{aasclass.sty}\LoadCl
ass{aasclass}% |
 []

[5]
LaTeX Font Info:    Font shape `OT1/cmtt/bx/n' in size <12> not available
(Font)              Font shape `OT1/cmtt/m/n' tried instead on input line 171.
[6] [7] [8] [9] [10]
Overfull \hbox (10.0376pt too wide) in paragraph at lines 579--579
 [][]  \OT1/cmtt/m/n/9 % For version number and date, search on "\fileversion" i
n the .dtx file, |
 []

[11] [12]
Overfull \hbox (5.11002pt too wide) in paragraph at lines 610--613
[]\OT1/cmr/m/n/10 The macro []\OT1/cmtt/m/n/10 \aas@ptsize \OT1/cmr/m/n/10 works
 very much like the []\OT1/cmtt/m/n/10 \@ptsize \OT1/cmr/m/n/10 macro of \OT1/cm
tt/m/n/10 classes.dtx\OT1/cmr/m/n/10 ,|
 []

[13] [14]
Overfull \hbox (47.83716pt too wide) in paragraph at lines 784--784
 [][] \OT1/cmtt/m/n/9 \ifx\undefined\aas@ptsize\expandafter\@firstoftwo\else\exp
andafter\@secondoftwo\fi |
 []

[15]
Overfull \hbox (5.31265pt too wide) in paragraph at lines 784--784
 [][]  \OT1/cmtt/m/n/9 \expandafter\PassOptionsToClass\expandafter{\expandafter1
\aas@ptsizept}% |
 []


Overfull \hbox (47.83716pt too wide) in paragraph at lines 784--784
 [][] \OT1/cmtt/m/n/9 \ifx\undefined\aas@ptsize\expandafter\@firstoftwo\else\exp
andafter\@secondoftwo\fi |
 []


Overfull \hbox (5.31265pt too wide) in paragraph at lines 784--784
 [][]  \OT1/cmtt/m/n/9 \expandafter\PassOptionsToClass\expandafter{\expandafter1
\aas@ptsizept}% |
 []


Overfull \hbox (36.30128pt too wide) in paragraph at lines 792--794
\OT1/cmr/m/n/10 AAST[]X is a vari-ant of the stan-dard L[]T[]X ar-ti-cle class: 
[]\OT1/cmtt/m/n/10 CTAN://macros/latex/unpacked\OT1/cmr/m/n/10 . |
 []


Overfull \hbox (14.63516pt too wide) in paragraph at lines 798--800
[]\OT1/cmr/m/n/10 We al-ways load the stan-dard L[]T[]X sym-bol pack-age: []\OT1
/cmtt/m/n/10 CTAN://macros/latex/base\OT1/cmr/m/n/10 . |
 []


Overfull \hbox (129.57907pt too wide) in paragraph at lines 804--806
[]\OT1/cmr/m/n/10 We al-ways load the stan-dard L[]T[]X ex-tended graph-ics pack
-age: []\OT1/cmtt/m/n/10 CTAN://macros/latex/packages/graphics\OT1/cmr/m/n/10 . 
|
 []


Overfull \hbox (50.91344pt too wide) in paragraph at lines 810--812
[]\OT1/cmr/m/n/10 The AMS []\OT1/cmtt/m/n/10 amsmath \OT1/cmr/m/n/10 pack-age re
-places []\OT1/cmtt/m/n/10 amssym\OT1/cmr/m/n/10 : []\OT1/cmtt/m/n/10 CTAN://mac
ros/latex/packages/amslatex\OT1/cmr/m/n/10 . |
 []


Overfull \hbox (67.91158pt too wide) in paragraph at lines 816--818
[]\OT1/cmr/m/n/10 The con-tributed L[]T[]X []\OT1/cmtt/m/n/10 natbib \OT1/cmr/m/
n/10 pack-age: []\OT1/cmtt/m/n/10 CTAN://macros/latex/contrib/supported/natbib\O
T1/cmr/m/n/10 . |
 []


Overfull \hbox (48.02657pt too wide) in paragraph at lines 822--825
\OT1/cmr/m/n/10 code may serve as a model for load-ing other classes: []\OT1/cmt
t/m/n/10 CTAN://macros/latex/packages/tools\OT1/cmr/m/n/10 . |
 []

LaTeX Font Info:    External font `cmex10' loaded for size
(Font)              <14.4> on input line 830.
[16] [17] [18] [19] [20] [21] [22] [23] [24] [25] [26] [27] [28] [29] [30]
[31] [32] [33] [34]
Overfull \hbox (19.66588pt too wide) in paragraph at lines 2107--2110
\OT1/cmr/m/n/10 The for-mat-ting of the []\OT1/cmtt/m/n/10 \slugcomment \OT1/cmr
/m/n/10 in-for-ma-tion is con-trolled by the []\OT1/cmtt/m/n/10 \@makeslugcmmnt|

 []

[35] [36] [37]
Overfull \hbox (27.72055pt too wide) in paragraph at lines 2338--2338
[]     \OT1/cmtt/m/n/9 \authoraddr{<author's postal address>}% Deprecated (V5.0)
 in favor of \affil[] |
 []

[38] [39] [40] [41] [42] [43] [44] [45] [46] [47] [48] [49]
Overfull \hbox (12.01851pt too wide) in paragraph at lines 3200--3202
[][]\OT1/cmtt/m/n/10 \@makecaption \OT1/cmr/m/n/10 de-ter-mines for-mat-ting for
 fig-ure and ta-ble cap-tions. For manuscripts,|
 []

[50] [51] [52] [53] [54] [55] [56] [57] [58] [59]
Overfull \hbox (19.48749pt too wide) in paragraph at lines 3895--3895
 [][]\OT1/cmtt/m/n/9 \newcommand\tablebreak{\substitute@command\tablebreak\\}%{\
cr\ptable@@split}% |
 []

[60] [61] [62] [63] [64] [65] [66] [67]
Overfull \hbox (10.0376pt too wide) in paragraph at lines 4479--4479
 [][]\OT1/cmtt/m/n/9 %\def\ascnode{\astro{N}}% ascending node (dragon's head), (
only in cmastro) |
 []


Overfull \hbox (19.48749pt too wide) in paragraph at lines 4479--4479
 [][]\OT1/cmtt/m/n/9 %\def\descnode{\astro{D}}% descending node (dragon's tail),
 (only in cmastro) |
 []

[68] [69] [70] [71] [72] [73] [74]
LaTeX Font Info:    Try loading font information for OMS+cmr on input line 4980.

(omscmr.fd
File: omscmr.fd 1998/03/27 v2.5g Standard LaTeX font definitions
)
LaTeX Font Info:    Font shape `OMS/cmr/m/n' in size <10> not available
(Font)              Font shape `OMS/cmsy/m/n' tried instead on input line 4980.
[75] [76]
Overfull \hbox (10.59657pt too wide) in paragraph at lines 5145--5147
[]\OT1/cmr/m/n/10 What should hap-pen to a com-pat-i-bil-ity mode doc-u-ment? Wo
uldn't|
 []

[77]
Overfull \hbox (24.22153pt too wide) in paragraph at lines 5220--5221
[]\OT1/cmr/m/n/10 Deprecate []\OT1/cmtt/m/n/10 \nl\OT1/cmr/m/n/10 , []\OT1/cmtt/
m/n/10 \nextline\OT1/cmr/m/n/10 , []\OT1/cmtt/m/n/10 \tablevspace\OT1/cmr/m/n/10
 , and []\OT1/cmtt/m/n/10 \tablebreak\OT1/cmr/m/n/10 .|
 []


Overfull \hbox (83.01173pt too wide) in paragraph at lines 5223--5229
[]\OT1/cmr/m/n/10 Functionality cov-ered by []\OT1/cmtt/m/n/10 hyperref \OT1/cmr
/m/n/10 pack-age. Syn-tax: []\OT1/cmtt/m/n/10 \anchor{[]}{[]}\OT1/cmr/m/n/10 .|
 []


Overfull \hbox (67.96008pt too wide) in paragraph at lines 5223--5229
\OT1/cmr/m/n/10 Ref: \OT1/cmtt/m/n/10 url.sty \OT1/cmr/m/n/10 in []\OT1/cmtt/m/n
/10 /contrib/other/misc\OT1/cmr/m/n/10 . \OT1/cmtt/m/n/10 hyperref.sty \OT1/cmr/
m/n/10 in []\OT1/cmtt/m/n/10 /contrib/supported/hyperref\OT1/cmr/m/n/10 .|
 []

[78]
Overfull \hbox (11.78822pt too wide) in paragraph at lines 5275--5279
[]\OT1/cmtt/m/n/10 \journal{[]}\OT1/cmr/m/n/10 , \OT1/cmtt/m/n/10 \articleID{[]}
\OT1/cmr/m/n/10 , \OT1/cmtt/m/n/10 \volno{[]}{[]}\OT1/cmr/m/n/10 . |
 []

[79]
*******************
* Checksum passed *
*******************
) (aasclass.ind
Underfull \hbox (badness 5726) in paragraph at lines 6--7
\OT1/cmr/m/n/9 1971, 2000, 2055  
 []

LaTeX Font Info:    External font `cmex10' loaded for size
(Font)              <9> on input line 14.

Underfull \hbox (badness 10000) in paragraph at lines 15--16
\OT1/cmr/m/n/9 1000, 1002-- 
 []


Underfull \hbox (badness 5726) in paragraph at lines 15--16
\OT1/cmr/m/n/9 1006, 1008, 1019  
 []


Underfull \hbox (badness 5726) in paragraph at lines 26--28
\OT1/cmr/m/n/9 1342, 1347, 1671  
 []


Underfull \hbox (badness 10000) in paragraph at lines 36--37
\OT1/cmr/m/n/9 1930, 1933, 
 []


Underfull \hbox (badness 5726) in paragraph at lines 36--37
\OT1/cmr/m/n/9 1937, 1940, 1943  
 []


Underfull \hbox (badness 10000) in paragraph at lines 58--60
\OT1/cmr/m/n/9 1930, 1933, 
 []


Underfull \hbox (badness 5726) in paragraph at lines 58--60
\OT1/cmr/m/n/9 1937, 1940, 1943  
 []


Underfull \hbox (badness 5726) in paragraph at lines 61--62
\OT1/cmr/m/n/9 1417, 1498, 1540  
 []


Underfull \hbox (badness 10000) in paragraph at lines 63--65
\OT1/cmr/m/n/9 1073, 1088, 
 []


Underfull \hbox (badness 5726) in paragraph at lines 63--65
\OT1/cmr/m/n/9 1102, 1112, 1549  
 []

[80]
Underfull \hbox (badness 5726) in paragraph at lines 81--82
\OT1/cmr/m/n/9 1107, 1118, $[]$  
 []


Overfull \hbox (14.06262pt too wide) in paragraph at lines 84--85
[][]\OT1/cmtt/m/n/9 \@latex@warning@no@line   |
 []


Overfull \hbox (9.33768pt too wide) in paragraph at lines 90--91
[][]\OT1/cmtt/m/n/9 \@make@caption@text@aj   |
 []


Overfull \hbox (14.06262pt too wide) in paragraph at lines 91--92
[][]\OT1/cmtt/m/n/9 \@make@caption@text@apj   |
 []


Overfull \hbox (18.78757pt too wide) in paragraph at lines 92--93
[][]\OT1/cmtt/m/n/9 \@make@caption@text@pptt   |
 []


Overfull \hbox (18.78757pt too wide) in paragraph at lines 95--96
[][]\OT1/cmtt/m/n/9 \@makecaption@plano@cont   |
 []


Overfull \hbox (32.9624pt too wide) in paragraph at lines 96--97
[][]\OT1/cmtt/m/n/9 \@makecaption@plano@cont@aj   |
 []


Overfull \hbox (37.68735pt too wide) in paragraph at lines 97--98
[][]\OT1/cmtt/m/n/9 \@makecaption@plano@cont@apj   |
 []


Overfull \hbox (42.4123pt too wide) in paragraph at lines 98--99
[][]\OT1/cmtt/m/n/9 \@makecaption@plano@cont@pptt   |
 []


Underfull \hbox (badness 5726) in paragraph at lines 102--103
\OT1/cmr/m/n/9 1037, 1970, 2050  
 []


Underfull \hbox (badness 10000) in paragraph at lines 126--128
\OT1/cmr/m/n/9 1600, 1603, 
 []


Underfull \hbox (badness 10000) in paragraph at lines 126--128
\OT1/cmr/m/n/9 1606, 1612, 
 []


Underfull \hbox (badness 5726) in paragraph at lines 126--128
\OT1/cmr/m/n/9 1615, 1618, 1624  
 []


Underfull \hbox (badness 10000) in paragraph at lines 148--150
\OT1/cmr/m/n/9 1191, 1216, 
 []


Underfull \hbox (badness 5726) in paragraph at lines 148--150
\OT1/cmr/m/n/9 2028, 2033, 2037  
 []


Underfull \hbox (badness 5726) in paragraph at lines 156--157
\OT1/cmr/m/n/9 1332, 1468, 2017  
 []


Underfull \hbox (badness 10000) in paragraph at lines 167--169
\OT1/cmr/m/n/9 434, 453--455, 
 []

[81]
Underfull \hbox (badness 5726) in paragraph at lines 173--175
\OT1/cmr/m/n/9 1163, 1166, 1175  
 []


Underfull \hbox (badness 5726) in paragraph at lines 178--180
\OT1/cmr/m/n/9 1166, 1499, 1703  
 []


Underfull \hbox (badness 10000) in paragraph at lines 187--189
\OT1/cmr/m/n/9 1930, 1933, 
 []


Underfull \hbox (badness 5726) in paragraph at lines 187--189
\OT1/cmr/m/n/9 1937, 1940, 1943  
 []


Overfull \hbox (9.33768pt too wide) in paragraph at lines 212--213
[][]\OT1/cmtt/m/n/9 \abovedisplayshortskip   |
 []


Underfull \hbox (badness 5726) in paragraph at lines 224--225
\OT1/cmr/m/n/9 1348, 1564, 1672  
 []


Underfull \hbox (badness 10000) in paragraph at lines 247--249
\OT1/cmr/m/n/9 1002--1006, 1084  
 []


Overfull \hbox (4.61273pt too wide) in paragraph at lines 250--251
[][]\OT1/cmtt/m/n/9 \appendix@figtab@defs   |
 []


Overfull \hbox (28.23746pt too wide) in paragraph at lines 251--252
[][]\OT1/cmtt/m/n/9 \appendix@figtab@defs@pptt   |
 []


Underfull \hbox (badness 5726) in paragraph at lines 252--253
\OT1/cmr/m/n/9 1328, 1641, 1647  
 []


Underfull \hbox (badness 5726) in paragraph at lines 258--259
\OT1/cmr/m/n/9 1170, 1210, 1211  
 []


Underfull \hbox (badness 10000) in paragraph at lines 273--275
\OT1/cmr/m/n/9 1869--1871, 
 []


Underfull \hbox (badness 10000) in paragraph at lines 273--275
\OT1/cmr/m/n/9 1874--1880, 
 []


Underfull \hbox (badness 10000) in paragraph at lines 273--275
\OT1/cmr/m/n/9 1882--1884, 
 []


Underfull \hbox (badness 10000) in paragraph at lines 273--275
\OT1/cmr/m/n/9 1886--1891, 
 []


Underfull \hbox (badness 10000) in paragraph at lines 273--275
\OT1/cmr/m/n/9 1893--1897, $[]$  
 []


Underfull \hbox (badness 5726) in paragraph at lines 276--278
\OT1/cmr/m/n/9 1135, 1324, 1398  
 []

[82]
Underfull \hbox (badness 5726) in paragraph at lines 293--295
\OT1/cmr/m/n/9 1995, 2043, 2045  
 []


Overfull \hbox (9.33768pt too wide) in paragraph at lines 302--303
[][]\OT1/cmtt/m/n/9 \belowdisplayshortskip   |
 []


Underfull \hbox (badness 10000) in paragraph at lines 332--335
\OT1/cmr/m/n/9 1038, 1043, 
 []


Underfull \hbox (badness 10000) in paragraph at lines 332--335
\OT1/cmr/m/n/9 1051, 1055, 
 []


Underfull \hbox (badness 10000) in paragraph at lines 332--335
\OT1/cmr/m/n/9 1066, 1072, 
 []


Underfull \hbox (badness 10000) in paragraph at lines 332--335
\OT1/cmr/m/n/9 1083, 1422, 
 []


Underfull \hbox (badness 10000) in paragraph at lines 332--335
\OT1/cmr/m/n/9 1431, 1434, 
 []


Underfull \hbox (badness 10000) in paragraph at lines 332--335
\OT1/cmr/m/n/9 1439, 1442, 
 []


Underfull \hbox (badness 10000) in paragraph at lines 332--335
\OT1/cmr/m/n/9 1448, 1451, 
 []


Underfull \hbox (badness 5726) in paragraph at lines 332--335
\OT1/cmr/m/n/9 1457, 1460, 1511  
 []


Underfull \hbox (badness 5726) in paragraph at lines 335--336
\OT1/cmr/m/n/9 1688, 1693, 1701  
 []


Underfull \hbox (badness 10000) in paragraph at lines 336--338
\OT1/cmr/m/n/9 1869--1871, 
 []


Underfull \hbox (badness 10000) in paragraph at lines 336--338
\OT1/cmr/m/n/9 1874--1880, 
 []


Underfull \hbox (badness 10000) in paragraph at lines 336--338
\OT1/cmr/m/n/9 1884, 1886-- 
 []


Underfull \hbox (badness 10000) in paragraph at lines 336--338
\OT1/cmr/m/n/9 1891, 1893--1897  
 []


Overfull \hbox (32.9624pt too wide) in paragraph at lines 348--350
[][]\OT1/cmtt/m/n/9 \clear@thebibliography@page   |
 []


Underfull \hbox (badness 5726) in paragraph at lines 348--350
\OT1/cmr/m/n/9 1275, 2012, 2058  
 []


Underfull \hbox (badness 5726) in paragraph at lines 350--351
\OT1/cmr/m/n/9 1563, 2057, 2058  
 []


Underfull \hbox (badness 10000) in paragraph at lines 355--357
\OT1/cmr/m/n/9 1056, 1073, 
 []


Underfull \hbox (badness 5726) in paragraph at lines 355--357
\OT1/cmr/m/n/9 1088, 1102, 1112  
 []


Underfull \hbox (badness 5726) in paragraph at lines 361--362
\OT1/cmr/m/n/9 1695, 1696, 1698  
 []


Overfull \hbox (4.61273pt too wide) in paragraph at lines 387--388
[][]\OT1/cmtt/m/n/9 \DeclareRobustCommand   |
 []


Underfull \hbox (badness 10000) in paragraph at lines 391--393
\OT1/cmr/m/n/9 1917, 1924, 
 []


Underfull \hbox (badness 10000) in paragraph at lines 391--393
\OT1/cmr/m/n/9 1925, 1929-- 
 []


Underfull \hbox (badness 10000) in paragraph at lines 391--393
\OT1/cmr/m/n/9 1931, 1933, 
 []


Underfull \hbox (badness 10000) in paragraph at lines 391--393
\OT1/cmr/m/n/9 1937, 1938, 
 []


Underfull \hbox (badness 5726) in paragraph at lines 391--393
\OT1/cmr/m/n/9 1940, 1941, 1943  
 []


Underfull \hbox (badness 10000) in paragraph at lines 393--395
\OT1/cmr/m/n/9 1917, 1924, 
 []


Underfull \hbox (badness 10000) in paragraph at lines 393--395
\OT1/cmr/m/n/9 1925, 1929-- 
 []


Underfull \hbox (badness 10000) in paragraph at lines 393--395
\OT1/cmr/m/n/9 1931, 1933, 
 []


Underfull \hbox (badness 10000) in paragraph at lines 393--395
\OT1/cmr/m/n/9 1937, 1938, 
 []


Underfull \hbox (badness 5726) in paragraph at lines 393--395
\OT1/cmr/m/n/9 1940, 1941, 1943  
 []


Underfull \hbox (badness 10000) in paragraph at lines 395--397
\OT1/cmr/m/n/9 1925, 1930, 
 []


Underfull \hbox (badness 10000) in paragraph at lines 395--397
\OT1/cmr/m/n/9 1932, 1933, 
 []


Underfull \hbox (badness 10000) in paragraph at lines 395--397
\OT1/cmr/m/n/9 1937, 1939, 
 []


Underfull \hbox (badness 5726) in paragraph at lines 395--397
\OT1/cmr/m/n/9 1940, 1942, 1943  
 []

[83]
Underfull \hbox (badness 5726) in paragraph at lines 417--418
\OT1/cmr/m/n/9 1483, 1562, 2023  
 []


Underfull \hbox (badness 10000) in paragraph at lines 419--422
\OT1/cmr/m/n/9 1047, 1051, 
 []


Underfull \hbox (badness 10000) in paragraph at lines 419--422
\OT1/cmr/m/n/9 1058, 1066, 
 []


Underfull \hbox (badness 10000) in paragraph at lines 419--422
\OT1/cmr/m/n/9 1076, 1083, 
 []


Underfull \hbox (badness 10000) in paragraph at lines 419--422
\OT1/cmr/m/n/9 1426, 1431, 
 []


Underfull \hbox (badness 10000) in paragraph at lines 419--422
\OT1/cmr/m/n/9 1434, 1439, 
 []


Underfull \hbox (badness 10000) in paragraph at lines 419--422
\OT1/cmr/m/n/9 1442, 1448, 
 []


Underfull \hbox (badness 10000) in paragraph at lines 419--422
\OT1/cmr/m/n/9 1451, 1457, 
 []


Underfull \hbox (badness 5726) in paragraph at lines 419--422
\OT1/cmr/m/n/9 1460, 1477, 1561  
 []


Overfull \hbox (4.61273pt too wide) in paragraph at lines 437--438
[][]\OT1/cmtt/m/n/9 \eqnarray@fleqn@fixed   |
 []


Overfull \hbox (4.61273pt too wide) in paragraph at lines 475--476
[][]\OT1/cmtt/m/n/9 \fnum@ptablecont@pptt   |
 []


Underfull \hbox (badness 5726) in paragraph at lines 476--477
\OT1/cmr/m/n/9 1503, 1565, $[]$  
 []


Underfull \hbox (badness 5726) in paragraph at lines 485--486
\OT1/cmr/m/n/9 1317, 1322, 2059  
 []


Underfull \hbox (badness 10000) in paragraph at lines 491--492
\OT1/cmr/m/n/9 1847, 1849--1853  
 []


Underfull \hbox (badness 10000) in paragraph at lines 516--518
\OT1/cmr/m/n/9 1377, 1381, 
 []


Underfull \hbox (badness 10000) in paragraph at lines 516--518
\OT1/cmr/m/n/9 1389, 1474, 
 []


Underfull \hbox (badness 10000) in paragraph at lines 516--518
\OT1/cmr/m/n/9 1517, 1526, 
 []


Underfull \hbox (badness 10000) in paragraph at lines 516--518
\OT1/cmr/m/n/9 1527, 1575, 
 []


Underfull \hbox (badness 10000) in paragraph at lines 516--518
\OT1/cmr/m/n/9 1576, 1599, 
 []


Underfull \hbox (badness 5726) in paragraph at lines 516--518
\OT1/cmr/m/n/9 1605, 1611, 1617  
 []

[84]
Underfull \hbox (badness 5726) in paragraph at lines 541--542
\OT1/cmr/m/n/9 1696, 1698, 1705  
 []


Overfull \hbox (23.51251pt too wide) in paragraph at lines 544--545
[][]\OT1/cmtt/m/n/9 \interfootnotelinepenalty   |
 []


Underfull \hbox (badness 10000) in paragraph at lines 552--554
\OT1/cmr/m/n/9 1094, 1106, 
 []


Underfull \hbox (badness 10000) in paragraph at lines 552--554
\OT1/cmr/m/n/9 1117, 1192, 
 []


Underfull \hbox (badness 5726) in paragraph at lines 552--554
\OT1/cmr/m/n/9 1437, 1446, 2038  
 []


Underfull \hbox (badness 10000) in paragraph at lines 585--587
\OT1/cmr/m/n/9 1044, 1056, 
 []


Underfull \hbox (badness 5726) in paragraph at lines 585--587
\OT1/cmr/m/n/9 1115, 1457, 1460  
 []


Underfull \hbox (badness 5726) in paragraph at lines 617--618
\OT1/cmr/m/n/9 1014, $[]$, 1973  
 []


Underfull \hbox (badness 5726) in paragraph at lines 626--627
\OT1/cmr/m/n/9 1972, 2001, 2056  
 []


Underfull \hbox (badness 5726) in paragraph at lines 637--638
\OT1/cmr/m/n/9 1150, 1152, 1157  
 []

[85]
Underfull \hbox (badness 10000) in paragraph at lines 646--647
\OT1/cmr/m/n/9 1822--1825, 1831  
 []


Underfull \hbox (badness 10000) in paragraph at lines 648--650
\OT1/cmr/m/n/9 1818--1827, 
 []


Underfull \hbox (badness 10000) in paragraph at lines 648--650
\OT1/cmr/m/n/9 1829, 1831, 
 []


Underfull \hbox (badness 10000) in paragraph at lines 648--650
\OT1/cmr/m/n/9 1832, 1837, 
 []


Underfull \hbox (badness 10000) in paragraph at lines 648--650
\OT1/cmr/m/n/9 1840, 1846, 
 []


Underfull \hbox (badness 10000) in paragraph at lines 648--650
\OT1/cmr/m/n/9 1848, 1855--1859  
 []


Underfull \hbox (badness 5726) in paragraph at lines 663--664
\OT1/cmr/m/n/9 1614, 1623, 1862  
 []


Underfull \hbox (badness 10000) in paragraph at lines 676--678
\OT1/cmr/m/n/9 1520, 1547, 
 []


Underfull \hbox (badness 10000) in paragraph at lines 676--678
\OT1/cmr/m/n/9 1575, 1598, 
 []


Underfull \hbox (badness 10000) in paragraph at lines 676--678
\OT1/cmr/m/n/9 1601, 1604, 
 []


Underfull \hbox (badness 10000) in paragraph at lines 676--678
\OT1/cmr/m/n/9 1607, 1610, 
 []


Underfull \hbox (badness 10000) in paragraph at lines 676--678
\OT1/cmr/m/n/9 1613, 1616, 
 []


Underfull \hbox (badness 5726) in paragraph at lines 676--678
\OT1/cmr/m/n/9 1619, 1622, 1625  
 []


Underfull \hbox (badness 10000) in paragraph at lines 685--688
\OT1/cmr/m/n/9 1056, 1073, 
 []


Underfull \hbox (badness 10000) in paragraph at lines 685--688
\OT1/cmr/m/n/9 1181, 1187, 
 []


Underfull \hbox (badness 10000) in paragraph at lines 685--688
\OT1/cmr/m/n/9 1192, 1210, 
 []


Underfull \hbox (badness 10000) in paragraph at lines 685--688
\OT1/cmr/m/n/9 1218, 2030, 
 []


Underfull \hbox (badness 5726) in paragraph at lines 685--688
\OT1/cmr/m/n/9 2034, 2038, 2059  
 []


Underfull \hbox (badness 10000) in paragraph at lines 714--716
\OT1/cmr/m/n/9 1629, 1645, 
 []


Underfull \hbox (badness 5726) in paragraph at lines 714--716
\OT1/cmr/m/n/9 1651, 1656, 1660  
 []


Underfull \hbox (badness 5726) in paragraph at lines 717--718
\OT1/cmr/m/n/9 1279, 1634, 1990  
 []


Underfull \hbox (badness 5726) in paragraph at lines 748--749
\OT1/cmr/m/n/9 1149, 1336, 1337  
 []


Underfull \hbox (badness 10000) in paragraph at lines 751--752
\OT1/cmr/m/n/9 1821, 1827, 
 []


Underfull \hbox (badness 5726) in paragraph at lines 751--752
\OT1/cmr/m/n/9 1829, 1837, 1842  
 []

[86]
Underfull \hbox (badness 10000) in paragraph at lines 762--763
\OT1/cmr/m/n/9 1515, 1523, 
 []


Underfull \hbox (badness 5726) in paragraph at lines 762--763
\OT1/cmr/m/n/9 1524, 1557, 1558  
 []


Underfull \hbox (badness 5726) in paragraph at lines 765--766
\OT1/cmr/m/n/9 1370, 1470, 1567  
 []


Underfull \hbox (badness 10000) in paragraph at lines 767--768
\OT1/cmr/m/n/9 1386, 1400, 
 []


Underfull \hbox (badness 5726) in paragraph at lines 767--768
\OT1/cmr/m/n/9 1516, 1517, 1571  
 []


Underfull \hbox (badness 5726) in paragraph at lines 769--770
\OT1/cmr/m/n/9 1548, 1549, 1568  
 []


Underfull \hbox (badness 10000) in paragraph at lines 770--772
\OT1/cmr/m/n/9 1371, 1469, 
 []


Underfull \hbox (badness 10000) in paragraph at lines 770--772
\OT1/cmr/m/n/9 1472, 1567, 
 []


Underfull \hbox (badness 5726) in paragraph at lines 770--772
\OT1/cmr/m/n/9 1602, 1614, 1623  
 []


Underfull \hbox (badness 10000) in paragraph at lines 772--774
\OT1/cmr/m/n/9 1410, 1493, 
 []


Underfull \hbox (badness 5726) in paragraph at lines 772--774
\OT1/cmr/m/n/9 1549, 1586, 1592  
 []


Underfull \hbox (badness 10000) in paragraph at lines 774--776
\OT1/cmr/m/n/9 1471, 1492, 
 []


Underfull \hbox (badness 5726) in paragraph at lines 774--776
\OT1/cmr/m/n/9 1560, 1579, 1585  
 []


Underfull \hbox (badness 10000) in paragraph at lines 779--782
[]$\OT1/cmr/m/n/9 , 1396, 
 []


Underfull \hbox (badness 10000) in paragraph at lines 779--782
\OT1/cmr/m/n/9 1402, 1417, 
 []


Underfull \hbox (badness 10000) in paragraph at lines 779--782
\OT1/cmr/m/n/9 1418, 1487, 
 []


Underfull \hbox (badness 10000) in paragraph at lines 779--782
\OT1/cmr/m/n/9 1493, 1498, 
 []


Underfull \hbox (badness 10000) in paragraph at lines 779--782
\OT1/cmr/m/n/9 1499, 1524, 
 []


Underfull \hbox (badness 10000) in paragraph at lines 779--782
\OT1/cmr/m/n/9 1558, 1580, 
 []


Underfull \hbox (badness 10000) in paragraph at lines 779--782
\OT1/cmr/m/n/9 1586, 1592, 
 []


Underfull \hbox (badness 10000) in paragraph at lines 779--782
\OT1/cmr/m/n/9 1629, 1645, 
 []


Underfull \hbox (badness 5726) in paragraph at lines 779--782
\OT1/cmr/m/n/9 1651, 1656, 1660  
 []


Underfull \hbox (badness 10000) in paragraph at lines 802--808
\OT1/cmr/m/n/9 1714, 1716, 
 []


Underfull \hbox (badness 10000) in paragraph at lines 802--808
\OT1/cmr/m/n/9 1718, 1720, 
 []


Underfull \hbox (badness 10000) in paragraph at lines 802--808
\OT1/cmr/m/n/9 1722, 1724, 
 []


Underfull \hbox (badness 10000) in paragraph at lines 802--808
\OT1/cmr/m/n/9 1726, 1728, 
 []


Underfull \hbox (badness 10000) in paragraph at lines 802--808
\OT1/cmr/m/n/9 1730, 1732, 
 []


Underfull \hbox (badness 10000) in paragraph at lines 802--808
\OT1/cmr/m/n/9 1734, 1736, 
 []


Underfull \hbox (badness 10000) in paragraph at lines 802--808
\OT1/cmr/m/n/9 1738, 1740, 
 []


Underfull \hbox (badness 10000) in paragraph at lines 802--808
\OT1/cmr/m/n/9 1742, 1744, 
 []


Underfull \hbox (badness 10000) in paragraph at lines 802--808
\OT1/cmr/m/n/9 1746, 1748, 
 []


Underfull \hbox (badness 10000) in paragraph at lines 802--808
\OT1/cmr/m/n/9 1750, 1752, 
 []


Underfull \hbox (badness 10000) in paragraph at lines 802--808
\OT1/cmr/m/n/9 1754, 1756, 
 []


Underfull \hbox (badness 10000) in paragraph at lines 802--808
\OT1/cmr/m/n/9 1758, 1760, 
 []


Underfull \hbox (badness 10000) in paragraph at lines 802--808
\OT1/cmr/m/n/9 1762, 1764, 
 []


Underfull \hbox (badness 10000) in paragraph at lines 802--808
\OT1/cmr/m/n/9 1766, 1768, 
 []


Underfull \hbox (badness 10000) in paragraph at lines 802--808
\OT1/cmr/m/n/9 1770, 1772, 
 []


Underfull \hbox (badness 10000) in paragraph at lines 802--808
\OT1/cmr/m/n/9 1774, 1776, 
 []


Underfull \hbox (badness 10000) in paragraph at lines 802--808
\OT1/cmr/m/n/9 1778, 1780, 
 []


Underfull \hbox (badness 10000) in paragraph at lines 802--808
\OT1/cmr/m/n/9 1782, 1784, 
 []


Underfull \hbox (badness 10000) in paragraph at lines 802--808
\OT1/cmr/m/n/9 1786, 1788, 
 []


Underfull \hbox (badness 10000) in paragraph at lines 802--808
\OT1/cmr/m/n/9 1790, 1792, 
 []


Underfull \hbox (badness 10000) in paragraph at lines 802--808
\OT1/cmr/m/n/9 1794, 1796, 
 []


Underfull \hbox (badness 5726) in paragraph at lines 802--808
\OT1/cmr/m/n/9 1798, 1800, 1802  
 []


Underfull \hbox (badness 5726) in paragraph at lines 809--810
\OT1/cmr/m/n/9 1982, 2015, 2063  
 []


Overfull \hbox (9.33768pt too wide) in paragraph at lines 810--811
[][]\OT1/cmtt/m/n/9 \references@refpar@mss   |
 []


Underfull \hbox (badness 5726) in paragraph at lines 810--811
\OT1/cmr/m/n/9 1292, 1982, 2063  
 []


Overfull \hbox (14.06262pt too wide) in paragraph at lines 811--812
[][]\OT1/cmtt/m/n/9 \references@refpar@pptt   |
 []


Underfull \hbox (badness 5726) in paragraph at lines 812--813
\OT1/cmr/m/n/9 1280, 1282, 1301  
 []


Underfull \hbox (badness 5726) in paragraph at lines 821--822
\OT1/cmr/m/n/9 1967, 1986, 2041  
 []


Underfull \hbox (badness 10000) in paragraph at lines 840--842
\OT1/cmr/m/n/9 1442, 1448, 
 []


Underfull \hbox (badness 10000) in paragraph at lines 840--842
\OT1/cmr/m/n/9 1451, 1457, 
 []


Underfull \hbox (badness 5726) in paragraph at lines 840--842
\OT1/cmr/m/n/9 1460, 1656, 1660  
 []


Underfull \hbox (badness 5726) in paragraph at lines 847--848
[]$\OT1/cmr/m/n/9 , 1215, 2007  
 []

[87]
Underfull \hbox (badness 5726) in paragraph at lines 873--875
\OT1/cmr/m/n/9 1105, 1113, 1860  
 []


Underfull \hbox (badness 5726) in paragraph at lines 880--881
\OT1/cmr/m/n/9 1336, 1337, 1475  
 []


Underfull \hbox (badness 5726) in paragraph at lines 895--896
\OT1/cmr/m/n/9 1249, 1276, 2008  
 []


Underfull \hbox (badness 10000) in paragraph at lines 897--899
\OT1/cmr/m/n/9 1529, 1595, 
 []


Underfull \hbox (badness 5726) in paragraph at lines 897--899
\OT1/cmr/m/n/9 1596, 1847, 1958  
 []


Underfull \hbox (badness 10000) in paragraph at lines 929--930
\OT1/cmr/m/n/9 1328, 1331-- 
 []


Underfull \hbox (badness 5726) in paragraph at lines 929--930
\OT1/cmr/m/n/9 1333, 1641, 1647  
 []


Underfull \hbox (badness 5726) in paragraph at lines 937--939
\OT1/cmr/m/n/9 1420, 1969, 2049  
 []


Underfull \hbox (badness 10000) in paragraph at lines 943--945
\OT1/cmr/m/n/9 1233, 1235, 
 []


Underfull \hbox (badness 5726) in paragraph at lines 943--945
\OT1/cmr/m/n/9 1239, 1243, 1244  
 []


Underfull \hbox (badness 10000) in paragraph at lines 945--947
\OT1/cmr/m/n/9 1223, 1346, 
 []


Underfull \hbox (badness 5726) in paragraph at lines 945--947
\OT1/cmr/m/n/9 1347, $[]$, 1677  
 []


Underfull \hbox (badness 5726) in paragraph at lines 950--951
\OT1/cmr/m/n/9 1671, $[]$, 1679  
 []


Underfull \hbox (badness 10000) in paragraph at lines 951--953
\OT1/cmr/m/n/9 1189, 1208, 
 []


Underfull \hbox (badness 5726) in paragraph at lines 951--953
\OT1/cmr/m/n/9 1211, 1223, 1225  
 []


Underfull \hbox (badness 10000) in paragraph at lines 955--958
\OT1/cmr/m/n/9 1225, 1341, 
 []


Underfull \hbox (badness 10000) in paragraph at lines 955--958
\OT1/cmr/m/n/9 1342, 1428, 
 []


Underfull \hbox (badness 10000) in paragraph at lines 955--958
\OT1/cmr/m/n/9 1429, 1437, 
 []


Underfull \hbox (badness 10000) in paragraph at lines 955--958
\OT1/cmr/m/n/9 1445, 1446, 
 []


Underfull \hbox (badness 10000) in paragraph at lines 955--958
\OT1/cmr/m/n/9 1454, 1455, 
 []


Underfull \hbox (badness 10000) in paragraph at lines 955--958
\OT1/cmr/m/n/9 1487, 1492, 
 []


Underfull \hbox (badness 10000) in paragraph at lines 955--958
\OT1/cmr/m/n/9 1579, 1585, 
 []


Underfull \hbox (badness 5726) in paragraph at lines 955--958
\OT1/cmr/m/n/9 1591, $[]$, 1678  
 []

[88]
Underfull \hbox (badness 5726) in paragraph at lines 969--970
\OT1/cmr/m/n/9 1066, 1071, 1083  
 []


Overfull \hbox (18.78757pt too wide) in paragraph at lines 987--988
[][]\OT1/cmtt/m/n/9 \typeout@deluxetable@mss   |
 []


Overfull \hbox (18.78757pt too wide) in paragraph at lines 988--989
[][]\OT1/cmtt/m/n/9 \typeout@deluxetable@ppt   |
 []


Underfull \hbox (badness 10000) in paragraph at lines 1022--1024
\OT1/cmr/m/n/9 1930, 1933, 
 []


Underfull \hbox (badness 5726) in paragraph at lines 1022--1024
\OT1/cmr/m/n/9 1937, 1940, 1943  
 []

) [89] (aasclass.aux) )
Here is how much of TeX's memory you used:
 687 strings out of 3926
 7187 string characters out of 65536
 95346 words of memory out of 192512
 3530 multiletter control sequences out of 10000
 14 hyphenation exceptions out of 3001
 31i,7n,21p,286b,1420s stack positions out of 200i,100n,120p,5000b,4000s

Output written on aasclass (89 pages, 323624 bytes).
elapsed seconds = 124.88

Textures 2.0.2 (preloaded format=plain 96.6.7)  20 JAN 1999 21:48
(aasclass (docstrip.tex
\blockLevel=\count26
\emptyLines=\count27
\processedLines=\count28
\commentsRemoved=\count29
\commentsPassed=\count30
\codeLinesPassed=\count31
\TotalprocessedLines=\count32
\TotalcommentsRemoved=\count33
\TotalcommentsPassed=\count34
\TotalcodeLinesPassed=\count35
\NumberOfFiles=\count36
\inFile=\read0
\inputcheck=\read1
\@tempcnta=\count37
\off@0=\count38
\off@1=\count39
\off@2=\count40
\off@3=\count41
\off@4=\count42
\off@5=\count43
\off@6=\count44
\off@7=\count45
\off@8=\count46
\off@9=\count47
\off@10=\count48
\off@11=\count49
\off@12=\count50
\off@13=\count51
\off@14=\count52
\off@15=\count53
\@temptokena=\toks12
\@maxfiles=\count54
\@maxoutfiles=\count55
Utility: `docstrip' 2.5b <1998/04/28>
English documentation    <1998/04/03>

**********************************************************
* This program converts documented macro-files into fast *
* loadable files by stripping off (nearly) all comments! *
**********************************************************

********************************************************
* No Configuration file found, using default settings. *
********************************************************

)

Generating file(s) aastex.cls aasclass.tex aasguide.tex aassymbols.tex sample.tex table.tex

Processing file aasclass.dtx (class) -> aastex.cls
                             (doc) -> aasclass.tex
                             (guide) -> aasguide.tex
			     (symbols) -> aassymbols.tex
                             (sample) -> sample.tex
                             (table) -> table.tex
File aasclass.dtx ended by \endinput.
Lines  processed: 9686
Comments removed: 3461
Comments  passed: 42
Codelines passed: 6142

***********************************************************
*
* To finish the installation, please move the following
* files into a directory searched by TeX:
*
*   aastex.cls
*
* To produce the user documentation, run aasguide.tex through LaTeX.
*
* To produce the AMS and AAS symbols crib sheet, run aassymbols.tex through LaTEX.
*
* To typeset the aastex sample documents, run sample.tex, sample2.tex, and samp
tbls.tex through LaTeX.
*
* To produce the programmer documentation, run aasclass.dtx through LaTeX.
*
* Happy TeXing
***********************************************************
No pages of output.
elapsed seconds = 46.72

Textures 2.0.2 (preloaded format=LaTeX R9P0 98.10.16)  23 JAN 1999 02:30
(aasguide
LaTeX2e <1998/06/01>
(aasguide.tex
File: aasguide.tex 1999/01/22 5.0rc2 AAS markup document class
(aastex.cls
Document Class: aastex 1999/01/22 5.0rc2 AAS markup document class
Class aastex Info: Original \LaTeX2.09 style by Chris Biemesderfer (chris@ferber
ts.com). Adapted to \LaTeXe by A. Ogawa (ogawa@teleport.com) on input line 74.
(article.cls
Document Class: article 1998/05/05 v1.3y Standard LaTeX document class
(size10.clo
File: size10.clo 1998/05/05 v1.3y Standard LaTeX file (size option)
)
\c@part=\count79
\c@section=\count80
\c@subsection=\count81
\c@subsubsection=\count82
\c@paragraph=\count83
\c@subparagraph=\count84
\c@figure=\count85
\c@table=\count86
\abovecaptionskip=\skip41
\belowcaptionskip=\skip42
\bibindent=\dimen102
) (latexsym.sty
Package: latexsym 1996/11/20 v2.2d Standard LaTeX package (lasy symbols)
\symlasy=\mathgroup4
LaTeX Font Info:    Overwriting symbol font `lasy' in version `bold'
(Font)                  U/lasy/m/n --> U/lasy/b/n on input line 85.
) (graphicx.sty
Package: graphicx 1997/06/09 v1.0d Enhanced LaTeX Graphics (DPC,SPQR)
(keyval.sty
Package: keyval 1998/06/05 v1.11 key=value parser (DPC)
\KV@toks@=\toks14
) (graphics.sty
Package: graphics 1998/05/27 v1.0h Standard LaTeX Graphics (DPC,SPQR)
(trig.sty
Package: trig 1994/10/16 v1.08 sin cos tan (DPC)
) (graphics.cfg)
Package graphics Info: Driver file: textures.def on input line 85.
(textures.def
File: textures.def 1997/5/28 v0.4 Graphics extensions for Textures 1.7 and later
 (Arthur Ogawa/ogawa@teleport.com)
))
\Gin@req@height=\dimen103
\Gin@req@width=\dimen104
) (amssymb.sty
Package: amssymb 1996/11/03 v2.2b
(amsfonts.sty
Package: amsfonts 1997/09/17 v2.2e
\@emptytoks=\toks15
\symAMSa=\mathgroup5
\symAMSb=\mathgroup6
LaTeX Font Info:    Overwriting math alphabet `\mathfrak' in version `bold'
(Font)                  U/euf/m/n --> U/euf/b/n on input line 133.
)) (natbib.sty
Package: natbib 1998/07/14 6.8c (PWD)
\bibhang=\skip43
\bibsep=\skip44
LaTeX Info: Redefining \cite on input line 467.
\c@NAT@ctr=\count87
) (verbatim.sty
Package: verbatim 1997/04/30 v1.5k LaTeX2e package for verbatim enhancements
\every@verbatim=\toks16
\verbatim@line=\toks17
\verbatim@in@stream=\read1
)
\@slugcmmntwidth=\dimen105
\abstract@box=\box26
\c@editornote=\count88
\c@cureqno=\count89
\pt@box=\box27
\pt@width=\dimen106
\pt@line=\count90
\pt@column=\count91
\pt@nlines=\count92
\pt@ncol=\count93
\pt@page=\count94
\pt@tmpcapwidth=\dimen107
\c@plate=\count95
)
\@indexfile=\write3
Writing index file aasguide.idx
(shortvrb.sty
Package: shortvrb 1998/05/19 v2.0b Standard LaTeX documentation package (FMi)
)
Package shortvrb Info: Made | a short reference for \verb on input line 102.
(aasguide.aux)
LaTeX Font Info:    Checking defaults for OML/cmm/m/it on input line 170.
LaTeX Font Info:    ... okay on input line 170.
LaTeX Font Info:    Checking defaults for T1/cmr/m/n on input line 170.
LaTeX Font Info:    ... okay on input line 170.
LaTeX Font Info:    Checking defaults for OT1/cmr/m/n on input line 170.
LaTeX Font Info:    ... okay on input line 170.
LaTeX Font Info:    Checking defaults for OMS/cmsy/m/n on input line 170.
LaTeX Font Info:    ... okay on input line 170.
LaTeX Font Info:    Checking defaults for OMX/cmex/m/n on input line 170.
LaTeX Font Info:    ... okay on input line 170.
LaTeX Font Info:    Checking defaults for U/cmr/m/n on input line 170.
LaTeX Font Info:    ... okay on input line 170.
AAS 2-column preprint substyle, v5.0 <25 Sep 98>.
LaTeX Font Info:    Try loading font information for U+lasy on input line 172.
(ulasy.fd
File: ulasy.fd 1996/11/20 v2.2dLaTeX symbol font definitions
)
LaTeX Font Info:    Font shape `U/lasy/b/n' in size <7> not available
(Font)              Font shape `U/lasy/m/n' tried instead on input line 172.
LaTeX Font Info:    Try loading font information for U+msa on input line 172.
(umsa.fd
File: umsa.fd 1995/01/05 v2.2e AMS font definitions
)
LaTeX Font Info:    Try loading font information for U+msb on input line 172.
(umsb.fd
File: umsb.fd 1995/01/05 v2.2e AMS font definitions
) (aasguide.toc
Underfull \hbox (badness 1675) in paragraph at lines 33--33
 [][] []\OT1/cmr/m/n/10 Abbreviations for Jour-nal 
 []


Overfull \hbox (0.41202pt too wide) in paragraph at lines 42--42
 [][] []\OT1/cmr/m/n/10 Using the \OT1/cmtt/m/n/10 \twocolumn \OT1/cmr/m/n/10 Co
m- 
 []

[1


])
\tf@toc=\write4

Underfull \hbox (badness 3179) in paragraph at lines 204--209
[]\OT1/cmr/m/n/10 Information about ac-quir-ing the AAST[]X
 []


Underfull \hbox (badness 1796) in paragraph at lines 204--209
\OT1/cmr/m/n/10 pack-age will be mailed to you elec-tron-i-cally
 []


Underfull \hbox (badness 1297) in paragraph at lines 204--209
\OT1/cmr/m/n/10 if you send an empty e-mail mes-sage (some
 []


Overfull \hbox (57.47958pt too wide) in paragraph at lines 229--235
[][] 
 []


Overfull \hbox (1.29999pt too wide) in paragraph at lines 237--243
[]\OT1/cmr/m/n/10 The AAS pro-vides user sup-port for the AAST[]X
 []


Overfull \hbox (17.18939pt too wide) in paragraph at lines 237--243
\OT1/cmr/m/n/10 tron-i-cally by send-ing e-mail to \OT1/cmtt/m/n/10 aastex-help@
aas.org\OT1/cmr/m/n/10 .[][] 
 []


Underfull \hbox (badness 3536) in paragraph at lines 255--255
[]\OT1/cmr/bx/n/10 Command De-scrip-tions for Jour-nal
 []

[2]
Overfull \hbox (11.4403pt too wide) in paragraph at lines 306--309
[]\OT1/cmr/m/n/10 The re-main-der of this sec-tion refers to manuscripts
 []


Overfull \hbox (34.9874pt too wide) in paragraph at lines 318--318
 [] \OT1/cmtt/m/n/9 \lower 0.5ex\hbox{${}\buildrel\lt\over\sim{}$}%[] 
 []


Overfull \hbox (6.35497pt too wide) in paragraph at lines 338--361
[]\OT1/cmtt/m/n/10 \lower[]\OT1/cmr/m/n/10 , []\OT1/cmtt/m/n/10 \rlap[]\OT1/cmr/
m/n/10 , []\OT1/cmtt/m/n/10 \kern[]\OT1/cmr/m/n/10 , []\OT1/cmtt/m/n/10 \lineski
p[]\OT1/cmr/m/n/10 , []\OT1/cmtt/m/n/10 \baselineskip[]\OT1/cmr/m/n/10 ,
 []

[3]
Overfull \hbox (25.97986pt too wide) in paragraph at lines 399--406
[][] 
 []


LaTeX Warning: Characters dropped after `\end{verbatim}' on input line 487.

[4]
Overfull \hbox (18.55078pt too wide) in paragraph at lines 557--563
\OT1/cmr/m/n/10 It is up to the au-thor to make sure that []\OT1/cmtt/m/n/10 \al
taffilmark[]
 []

[5]
Underfull \hbox (badness 1019) in paragraph at lines 644--664
\OT1/cmr/m/n/10 the proper el-e-ments for as-sist-ing the ed-i-tors.
 []


Underfull \hbox (badness 1127) in paragraph at lines 644--664
\OT1/cmr/m/n/10 en-vi-ron-ment (see Sec-tions 3.15[] and 5[]); in a
 []

[6]
Overfull \hbox (2.85645pt too wide) in paragraph at lines 824--828
[]\OT1/cmr/m/n/10 Authors should be aware that in the \OT1/cmtt/m/n/10 manuscrip
t[]
 []


Overfull \hbox (7.1341pt too wide) in paragraph at lines 831--838
[]\OT1/cmr/m/n/10 This sec-tion de-scribes the use of the \OT1/cmtt/m/n/10 delux
etable
 []


Underfull \hbox (badness 1142) in paragraph at lines 839--847
[]\OT1/cmr/m/n/10 The \OT1/cmtt/m/n/10 deluxetable[] \OT1/cmr/m/n/10 en-vi-ron-m
ent is de-lim-ited
 []


Overfull \hbox (16.66133pt too wide) in paragraph at lines 882--891
\OT1/cmr/m/n/10 (11pt), the []\OT1/cmtt/m/n/10 \footnotesize[] \OT1/cmr/m/n/10 (
10pt) or the []\OT1/cmtt/m/n/10 \scriptsize[]
 []


Underfull \hbox (badness 1394) in paragraph at lines 897--909
[]\OT1/cmr/m/n/10 The width of a delux-etable is de-fined by
 []


Underfull \hbox (badness 3612) in paragraph at lines 910--921
\OT1/cmr/m/n/10 num-ber-ing within the \OT1/cmtt/m/n/10 deluxetable \OT1/cmr/m/n
/10 en-vi-ron-
 []

[7]
Overfull \hbox (3.66202pt too wide) in paragraph at lines 949--970
[]\OT1/cmr/m/n/10 Column head-ings are spec-i-fied with []\OT1/cmtt/m/n/10 \tabl
ehead[]\OT1/cmr/m/n/10 .
 []


Overfull \hbox (1.81458pt too wide) in paragraph at lines 971--982
\OT1/cmr/m/n/10 pro-duced by []\OT1/cmtt/m/n/10 \tablehead[] \OT1/cmr/m/n/10 or 
other L[]T[]Xcommands
 []


Overfull \hbox (64.80882pt too wide) in paragraph at lines 1040--1049
[][] 
 []

[8]
Underfull \hbox (badness 2165) in paragraph at lines 1158--1160
[]\OT1/cmr/m/n/10 Tabular in-for-ma-tion is type-set within the
 []

[9]
Underfull \hbox (badness 1005) in paragraph at lines 1202--1213
\OT1/cmr/m/n/10 to []\OT1/cmtt/m/n/10 \tablenum[] \OT1/cmr/m/n/10 is used as the
 ta-ble iden-ti-fier.
 []


Underfull \hbox (badness 1418) in paragraph at lines 1269--1272
[]\OT1/cmr/m/n/10 Two op-tions are avail-able for sym-bol-i-cally
 []


Underfull \hbox (badness 1072) in paragraph at lines 1269--1272
\OT1/cmr/m/n/10 dard L[]T[]X method and the []\OT1/cmtt/m/n/10 natbib \OT1/cmr/m
/n/10 pack-age
 []


Underfull \hbox (badness 2753) in paragraph at lines 1336--1348
[]\OT1/cmr/m/n/10 Clusters of galax-ies dis-play X-ray
 []

[10]
Underfull \hbox (badness 2376) in paragraph at lines 1428--1436
\OT1/cmr/m/n/10 mul-ti-ple ci-ta-tion keys, e.g., []\OT1/cmtt/m/n/10 \citep[]{sa
ra86,
 []


Underfull \hbox (badness 1314) in paragraph at lines 1428--1436
[]\OT1/cmtt/m/n/10 odea97,[]maoz95} \OT1/cmr/m/n/10 will pro-duce ``(Sarazin 198
6;
 []


Overfull \hbox (17.03613pt too wide) in paragraph at lines 1439--1443
[]\OT1/cmtt/m/n/10 \citet[]*[[]][[]]{[]} 
 []


Overfull \hbox (17.03613pt too wide) in paragraph at lines 1439--1443
\OT1/cmtt/m/n/10 \citep[]*[[]][[]]{[]} 
 []


Underfull \hbox (badness 1565) in paragraph at lines 1467--1480
\OT1/cmr/m/n/10 mech-a-nism as-so-ciates ci-ta-tions and ref-er-ences
 []

[11]
Underfull \hbox (badness 1019) in paragraph at lines 1516--1525
[]\OT1/cmr/m/n/10 The AAST[]X class uses the []\OT1/cmtt/m/n/10 natbib \OT1/cmr/
m/n/10 pack-
 []


Overfull \hbox (30.9122pt too wide) in paragraph at lines 1570--1613
 [] 
 []


Overfull \hbox (1.94704pt too wide) in paragraph at lines 1626--1626
 []\OT1/cmtt/m/n/9 \figcaption[[]]{[]\label{[]}}[] 
 []


Underfull \hbox (badness 3128) in paragraph at lines 1643--1649
\OT1/cmr/m/n/10 manuscript is be-ing sub-mit-ted elec-tron-i-cally.
 []

[12] [13]
Underfull \hbox (badness 1394) in paragraph at lines 1750--1790
\OT1/cmr/m/n/10 bols avail-able in AAST[]X. If you wish to
 []


Underfull \hbox (badness 1728) in paragraph at lines 1750--1790
\OT1/cmr/m/n/10 use any of these sym-bols, you must ob-tain
 []


Underfull \hbox (badness 4144) in paragraph at lines 1750--1790
\OT1/cmr/m/n/10 avail-able from CTAN: ftp://ctan.tug.org/tex-
 []

[14]
Underfull \hbox (badness 1314) in paragraph at lines 1890--1904
\OT1/cmr/m/n/10 tures, OzT[]X, and CMacT[]X for Mac-in-tosh,
 []


Underfull \hbox (badness 1742) in paragraph at lines 1890--1904
\OT1/cmr/m/n/10 Win-dows, and, on UNIX com-put-ers, \OT1/cmtt/m/n/10 dvips\OT1/c
mr/m/n/10 .
 []


Overfull \hbox (22.93831pt too wide) in paragraph at lines 1890--1904
\OT1/cmr/m/n/10 T[]XLive CD, see []\OT1/cmtt/m/n/10 http://www.tug.org/texlive.h
tml\OT1/cmr/m/n/10 .
 []


Overfull \hbox (8.66833pt too wide) in paragraph at lines 1916--1916
 []      \OT1/cmtt/m/n/9 {[]}{[]}{[]}{[]}{[]}[] 
 []

[15]
Overfull \hbox (11.05083pt too wide) in paragraph at lines 1942--1948
[]\OT1/cmr/m/n/10 For the ad-ven-ture-some, there is also []\OT1/cmtt/m/n/10 \pl
otfiddle[]\OT1/cmr/m/n/10 ,
 []


Underfull \hbox (badness 1142) in paragraph at lines 2006--2014
[]\OT1/cmr/m/n/10 A []\OT1/cmtt/m/n/10 flushrt.sty \OT1/cmr/m/n/10 file op-tion 
is avail-able for
 []


Underfull \hbox (badness 1609) in paragraph at lines 2045--2047
[]\OT1/cmtt/m/n/10 \singlespace[] \OT1/cmr/m/n/10 sets the ver-ti-cal spac-ing t
o a
 []

[16]
Underfull \hbox (badness 1838) in paragraph at lines 2095--2099
[]\OT1/cmr/m/n/10 The pre-ced-ing de-tailed ex-pla-na-tion of the
 []


Overfull \hbox (0.93925pt too wide) in paragraph at lines 2105--2114
[]\OT1/cmr/m/n/10 A more com-pre-hen-sive ex-am-ple re-quir-ing nearly
 []


Underfull \hbox (badness 1325) in paragraph at lines 2105--2114
\OT1/cmr/m/n/10 has three ta-bles: two marked up us-ing the
 []


Underfull \hbox (badness 1014) in paragraph at lines 2115--2120
\OT1/cmr/m/n/10 ta-bles are con-tained in the files []\OT1/cmtt/m/n/10 table1.te
x\OT1/cmr/m/n/10 ,
 []


Underfull \hbox (badness 1490) in paragraph at lines 2115--2120
[]\OT1/cmtt/m/n/10 table2.tex\OT1/cmr/m/n/10 , []\OT1/cmtt/m/n/10 table3.tex\OT1
/cmr/m/n/10 , and []\OT1/cmtt/m/n/10 complext.tex\OT1/cmr/m/n/10 .
 []

[17]
Underfull \hbox (badness 1558) in paragraph at lines 2171--2173
\OT1/cmr/m/it/10 Com-mand Sum-mary \OT1/cmr/m/n/10 (Prov-i-dence, RI: T[]X
 []

LaTeX Font Info:    Try loading font information for OMS+cmr on input line 2215.

(omscmr.fd
File: omscmr.fd 1998/03/27 v2.5g Standard LaTeX font definitions
)
LaTeX Font Info:    Font shape `OMS/cmr/m/n' in size <10> not available
(Font)              Font shape `OMS/cmsy/m/n' tried instead on input line 2215.

Underfull \hbox (badness 1087) in paragraph at lines 2219--2224
\OT1/cmr/m/n/10 takes care of hy-phen-ation au-to-mat-i-cally.
 []

! Extra }, or forgotten \endgroup.
l.2242 ...\verb`\}`---end group}
                                .
? 

Underfull \hbox (badness 1394) in paragraph at lines 2259--2268
\OT1/cmr/m/n/10 the equa-tion be-tween []\OT1/cmtt/m/n/10 \[[] \OT1/cmr/m/n/10 a
nd []\OT1/cmtt/m/n/10 \][] \OT1/cmr/m/n/10 com-mands,
 []

[18]
Underfull \hbox (badness 1990) in paragraph at lines 2301--2312
[]\OT1/cmr/m/n/10 The \OT1/cmtt/m/n/10 deluxetable[] \OT1/cmr/m/n/10 en-vi-ron-m
ent is in-tended
 []


Overfull \hbox (27.21739pt too wide) in paragraph at lines 2313--2317
[]\OT1/cmr/m/n/10 Notes in ta-bles should be marked by []\OT1/cmtt/m/n/10 \table
notemark[]\OT1/cmr/m/n/10 ;
 []


Overfull \hbox (19.93968pt too wide) in paragraph at lines 2313--2317
\OT1/cmr/m/n/10 cor-re-spond-ing text should ap-pear in a []\OT1/cmtt/m/n/10 \ta
blenotetext[]
 []


Overfull \hbox (20.81256pt too wide) in paragraph at lines 2334--2334
 []\OT1/cmtt/m/n/9 \tablecaption{This is a caption.\label{tab1}}[] 
 []


Underfull \hbox (badness 10000) in paragraph at lines 2361--2369
\OT1/cmr/m/n/10 type-set-ting []\OT1/cmtt/m/n/10 CTAN://[]macros/[]latex/[]packa
ges/
 []


Overfull \hbox (2.5793pt too wide) in paragraph at lines 2388--2393
\OT1/cmr/m/n/10 manuscripts that will be sub-mit-ted electronically---
 []

[19]
LaTeX Font Info:    Try loading font information for OML+cmr on input line 2410.

(omlcmr.fd
File: omlcmr.fd 1998/03/27 v2.5g Standard LaTeX font definitions
)
LaTeX Font Info:    Font shape `OML/cmr/m/n' in size <10> not available
(Font)              Font shape `OML/cmm/m/it' tried instead on input line 2410.

Overfull \hbox (15.18954pt too wide) in paragraph at lines 2409--2416
 [] 
 []

! Missing { inserted.
<inserted text> 
                {
l.2441 ...erb`\}`---end group} &
                                 &
? 

Overfull \hbox (85.96729pt too wide) in paragraph at lines 2432--2444
 [] 
 []


Overfull \hbox (5.38377pt too wide) detected at line 2629
[]
 []


Overfull \hbox (54.88292pt too wide) in paragraph at lines 2672--2737
 [] 
 []


Underfull \hbox (badness 1210) in paragraph at lines 2794--2798
\OT1/cmr/m/n/10 mented in Sec-tion E.3[]), ex-cept that you do
 []

[20]
Overfull \hbox (3.27318pt too wide) in paragraph at lines 2808--2809
\OT1/cmr/m/n/10 style op-tion, if used, with the []\OT1/cmtt/m/n/10 preprint2
 []

[21] [22]
Underfull \hbox (badness 2680) in paragraph at lines 2826--2828
[]\OT1/cmr/m/n/10 Replace any []\OT1/cmtt/m/n/10 \lefthead[] \OT1/cmr/m/n/10 com
-mand with
 []


Overfull \hbox (7.85611pt too wide) in paragraph at lines 2835--2836
[][]\OT1/cmtt/m/n/10 \markcite[]{[]}\OT1/cmr/m/n/10 (Headroom
 []


Overfull \hbox (2.7778pt too wide) in paragraph at lines 2843--2844
[][]\OT1/cmtt/m/n/10 \reference[]{[]}\OT1/cmr/m/it/10 reference
 []


Overfull \hbox (51.45876pt too wide) in paragraph at lines 2847--2848
[][]\OT1/cmtt/m/n/10 \bibitem[][[]]{[]}\OT1/cmr/m/it/10 reference
 []


Underfull \hbox (badness 1783) in paragraph at lines 2865--2868
\OT1/cmr/m/n/10 means that V5.0 can-not be run un-der
 []


Underfull \hbox (badness 2351) in paragraph at lines 2909--2912
\OT1/cmr/m/n/10 land-scape ori-en-ta-tion by us-ing the
 []


Underfull \hbox (badness 5563) in paragraph at lines 2913--2918
[]\OT1/cmtt/m/n/10 \\[]\OT1/cmr/m/n/10 , the com-mands []\OT1/cmtt/m/n/10 \table
vspace[]\OT1/cmr/m/n/10 ,
 []

[23]
Underfull \hbox (badness 10000) in paragraph at lines 2924--2926
[]\OT1/cmr/m/n/10 Hypertext-oriented com-mands []\OT1/cmtt/m/n/10 \anchor[]\OT1/
cmr/m/n/10 ,
 []


Underfull \hbox (badness 3482) in paragraph at lines 2929--2931
[]\OT1/cmr/m/n/10 The data com-mands []\OT1/cmtt/m/n/10 \shorttitle[] \OT1/cmr/m
/n/10 and
 []


Underfull \hbox (badness 6032) in paragraph at lines 2929--2931
[]\OT1/cmtt/m/n/10 \shortauthors[] \OT1/cmr/m/n/10 re-place []\OT1/cmtt/m/n/10 \
lefthead[] \OT1/cmr/m/n/10 and
 []


Underfull \hbox (badness 3849) in paragraph at lines 2939--2944
\OT1/cmr/m/n/10 com-mands []\OT1/cmtt/m/n/10 \tablevspace[]\OT1/cmr/m/n/10 , []\
OT1/cmtt/m/n/10 \tablebreak[]\OT1/cmr/m/n/10 ,
 []


Underfull \hbox (badness 3482) in paragraph at lines 2945--2948
[]\OT1/cmr/m/n/10 The data com-mands []\OT1/cmtt/m/n/10 \shorttitle[] \OT1/cmr/m
/n/10 and
 []


Underfull \hbox (badness 6032) in paragraph at lines 2945--2948
[]\OT1/cmtt/m/n/10 \shortauthors[] \OT1/cmr/m/n/10 re-place []\OT1/cmtt/m/n/10 \
lefthead[] \OT1/cmr/m/n/10 and
 []


Underfull \hbox (badness 1142) in paragraph at lines 2962--2964
[]\OT1/cmr/m/n/10 The com-mand []\OT1/cmtt/m/n/10 \markcite[] \OT1/cmr/m/n/10 is
 dep-re-cated.
 []


Underfull \hbox (badness 1142) in paragraph at lines 2966--2968
[]\OT1/cmr/m/n/10 The com-mand []\OT1/cmtt/m/n/10 \SGMLcite[] \OT1/cmr/m/n/10 is
 dep-re-cated.
 []


Underfull \hbox (badness 1275) in paragraph at lines 2975--2976
[]\OT1/cmr/m/n/10 Replace your []\OT1/cmtt/m/n/10 \documentstyle[] \OT1/cmr/m/n/
10 state-ment
 []


Overfull \hbox (3.27318pt too wide) in paragraph at lines 2984--2986
\OT1/cmr/m/n/10 style op-tion, if used, with the []\OT1/cmtt/m/n/10 preprint2
 []


Underfull \hbox (badness 2680) in paragraph at lines 3004--3006
[]\OT1/cmr/m/n/10 Replace any []\OT1/cmtt/m/n/10 \lefthead[] \OT1/cmr/m/n/10 com
-mand with
 []

[24]
Overfull \hbox (7.85611pt too wide) in paragraph at lines 3027--3028
[][]\OT1/cmtt/m/n/10 \markcite[]{[]}\OT1/cmr/m/n/10 (Headroom
 []


Overfull \hbox (2.7778pt too wide) in paragraph at lines 3035--3036
[][]\OT1/cmtt/m/n/10 \reference[]{[]}\OT1/cmr/m/it/10 reference
 []


Overfull \hbox (51.45876pt too wide) in paragraph at lines 3039--3040
[][]\OT1/cmtt/m/n/10 \bibitem[][[]]{[]}\OT1/cmr/m/it/10 reference
 []

(aasguide.ind [25]
! Undefined control sequence.
l.95 ...X\ Archive Network, \see
                                 CTAN{2}
? 
 [26

] [27] [28] [29] [30]) (aasguide.aux)

 *File List*
  aastex.cls    1999/01/22 5.0rc2 AAS markup document class
 article.cls    1998/05/05 v1.3y Standard LaTeX document class
  size10.clo    1998/05/05 v1.3y Standard LaTeX file (size option)
latexsym.sty    1996/11/20 v2.2d Standard LaTeX package (lasy symbols)
graphicx.sty    1997/06/09 v1.0d Enhanced LaTeX Graphics (DPC,SPQR)
  keyval.sty    1998/06/05 v1.11 key=value parser (DPC)
graphics.sty    1998/05/27 v1.0h Standard LaTeX Graphics (DPC,SPQR)
    trig.sty    1994/10/16 v1.08 sin cos tan (DPC)
graphics.cfg
textures.def    1997/5/28 v0.4 Graphics extensions for Textures 1.7 and later (A
rthur Ogawa/ogawa@teleport.com)
 amssymb.sty    1996/11/03 v2.2b
amsfonts.sty    1997/09/17 v2.2e
  natbib.sty    1998/07/14 6.8c (PWD)
verbatim.sty    1997/04/30 v1.5k LaTeX2e package for verbatim enhancements
shortvrb.sty    1998/05/19 v2.0b Standard LaTeX documentation package (FMi)
   ulasy.fd    1996/11/20 v2.2dLaTeX symbol font definitions
    umsa.fd    1995/01/05 v2.2e AMS font definitions
    umsb.fd    1995/01/05 v2.2e AMS font definitions
  omscmr.fd    1998/03/27 v2.5g Standard LaTeX font definitions
  omlcmr.fd    1998/03/27 v2.5g Standard LaTeX font definitions
aasguide.ind
 ***********

 )
Here is how much of TeX's memory you used:
 1876 strings out of 3926
 18884 string characters out of 65536
 95730 words of memory out of 200704
 4653 multiletter control sequences out of 10000
 15 hyphenation exceptions out of 3001
 33i,15n,26p,322b,419s stack positions out of 200i,100n,120p,5000b,4000s

Output written on aasguide (30 pages, 146198 bytes).
elapsed seconds = 62.18

Textures 2.0.2 (preloaded format=LaTeX R9P0 98.10.16)  20 JAN 1999 23:13
(sample1
LaTeX2e <1998/06/01>
(sample.tex
File: sample.tex 1999/01/22 5.0rc2 AAS markup document class
(aastex.cls
Document Class: aastex 1999/01/22 5.0rc2 AAS markup document class
Class aastex Info: Original \LaTeX2.09 style by Chris Biemesderfer (chris@ferber
ts.com). Adapted to \LaTeXe by A. Ogawa (ogawa@teleport.com) on input line 68.
(article.cls
Document Class: article 1998/05/05 v1.3y Standard LaTeX document class
(size12.clo
File: size12.clo 1998/05/05 v1.3y Standard LaTeX file (size option)
)
\c@part=\count79
\c@section=\count80
\c@subsection=\count81
\c@subsubsection=\count82
\c@paragraph=\count83
\c@subparagraph=\count84
\c@figure=\count85
\c@table=\count86
\abovecaptionskip=\skip41
\belowcaptionskip=\skip42
\bibindent=\dimen102
) (latexsym.sty
Package: latexsym 1996/11/20 v2.2d Standard LaTeX package (lasy symbols)
\symlasy=\mathgroup4
LaTeX Font Info:    Overwriting symbol font `lasy' in version `bold'
(Font)                  U/lasy/m/n --> U/lasy/b/n on input line 85.
) (graphicx.sty
Package: graphicx 1997/06/09 v1.0d Enhanced LaTeX Graphics (DPC,SPQR)
(keyval.sty
Package: keyval 1998/06/05 v1.11 key=value parser (DPC)
\KV@toks@=\toks14
) (graphics.sty
Package: graphics 1998/05/27 v1.0h Standard LaTeX Graphics (DPC,SPQR)
(trig.sty
Package: trig 1994/10/16 v1.08 sin cos tan (DPC)
) (graphics.cfg)
Package graphics Info: Driver file: textures.def on input line 85.
(textures.def
File: textures.def 1997/5/28 v0.4 Graphics extensions for Textures 1.7 and later
 (Arthur Ogawa/ogawa@teleport.com)
))
\Gin@req@height=\dimen103
\Gin@req@width=\dimen104
) (amsmath.sty
Package: amsmath 1997/03/20 v1.2d AMS math features
\@mathmargin=\skip43
(amstext.sty
Package: amstext 1996/10/28 v1.2b
(amsgen.sty
File: amsgen 1996/10/29 v1.2b
\@emptytoks=\toks15
\ex@=\dimen105
)) (amsbsy.sty
Package: amsbsy 1996/10/28 v1.2b
\pmbraise@=\dimen106
) (amsopn.sty
Package: amsopn 1996/10/28 v1.2b operator names
)
LaTeX Info: Redefining \frac on input line 188.
\uproot@=\count87
\leftroot@=\count88
\classnum@=\count89
\DOTSCASE@=\count90
LaTeX Info: Redefining \dots on input line 335.
LaTeX Info: Redefining \ldots on input line 418.
LaTeX Info: Redefining \cdots on input line 422.
\Mathstrutbox@=\box26
\strutbox@=\box27
\big@size=\dimen107
LaTeX Font Info:    Redeclaring font encoding OML on input line 500.
LaTeX Font Info:    Redeclaring font encoding OMS on input line 501.
\skewcharcount@=\count91
\familycount@=\count92
\pointcount@=\count93
\accentdimen@=\dimen108
\accentmu@=\count94
\minaw@=\dimen109
\c@MaxMatrixCols=\count95
\dotsspace@=\muskip10
\c@parentequation=\count96
\dspbrk@lvl=\count97
\tag@help=\toks16
\row@=\count98
\column@=\count99
\maxfields@=\count100
\andhelp@=\toks17
\eqnshift@=\dimen110
\alignsep@=\dimen111
\tagshift@=\dimen112
\tagwidth@=\dimen113
\totwidth@=\dimen114
\lineht@=\dimen115
\@envbody=\toks18
\multlinegap=\skip44
\multlinetaggap=\skip45
) (natbib.sty
Package: natbib 1998/07/14 6.8c (PWD)
\bibhang=\skip46
\bibsep=\skip47
LaTeX Info: Redefining \cite on input line 467.
\c@NAT@ctr=\count101
) (verbatim.sty
Package: verbatim 1997/04/30 v1.5k LaTeX2e package for verbatim enhancements
\every@verbatim=\toks19
\verbatim@line=\toks20
\verbatim@in@stream=\read1
)
\@slugcmmntwidth=\dimen116
\abstract@box=\box28
\c@editornote=\count102
\c@cureqno=\count103
\pt@box=\box29
\pt@width=\dimen117
\pt@line=\count104
\pt@column=\count105
\pt@nlines=\count106
\pt@ncol=\count107
\pt@page=\count108
\pt@tmpcapwidth=\dimen118
\c@plate=\count109
) (sample1.aux)
LaTeX Font Info:    Checking defaults for OML/cmm/m/it on input line 65.
LaTeX Font Info:    ... okay on input line 65.
LaTeX Font Info:    Checking defaults for T1/cmr/m/n on input line 65.
LaTeX Font Info:    ... okay on input line 65.
LaTeX Font Info:    Checking defaults for OT1/cmr/m/n on input line 65.
LaTeX Font Info:    ... okay on input line 65.
LaTeX Font Info:    Checking defaults for OMS/cmsy/m/n on input line 65.
LaTeX Font Info:    ... okay on input line 65.
LaTeX Font Info:    Checking defaults for OMX/cmex/m/n on input line 65.
LaTeX Font Info:    ... okay on input line 65.
LaTeX Font Info:    Checking defaults for U/cmr/m/n on input line 65.
LaTeX Font Info:    ... okay on input line 65.
AAS manuscript substyle, v5.0 <25 Sep 98>.

Class aastex Warning: Command \authoremail is deprecated in aastex V 5.0. Using 
\email instead (please fix your document).  on input line 70.

[1

]
LaTeX Font Info:    Try loading font information for U+lasy on input line 97.
(ulasy.fd
File: ulasy.fd 1996/11/20 v2.2dLaTeX symbol font definitions
) [2

] [3] [4

] (sample1.aux)

 *File List*
  aastex.cls    1999/01/22 5.0rc2 AAS markup document class
 article.cls    1998/05/05 v1.3y Standard LaTeX document class
  size12.clo    1998/05/05 v1.3y Standard LaTeX file (size option)
latexsym.sty    1996/11/20 v2.2d Standard LaTeX package (lasy symbols)
graphicx.sty    1997/06/09 v1.0d Enhanced LaTeX Graphics (DPC,SPQR)
  keyval.sty    1998/06/05 v1.11 key=value parser (DPC)
graphics.sty    1998/05/27 v1.0h Standard LaTeX Graphics (DPC,SPQR)
    trig.sty    1994/10/16 v1.08 sin cos tan (DPC)
graphics.cfg
textures.def    1997/5/28 v0.4 Graphics extensions for Textures 1.7 and later (A
rthur Ogawa/ogawa@teleport.com)
 amsmath.sty    1997/03/20 v1.2d AMS math features
 amstext.sty    1996/10/28 v1.2b
  amsgen.sty    
  amsbsy.sty    1996/10/28 v1.2b
  amsopn.sty    1996/10/28 v1.2b operator names
  natbib.sty    1998/07/14 6.8c (PWD)
verbatim.sty    1997/04/30 v1.5k LaTeX2e package for verbatim enhancements
   ulasy.fd    1996/11/20 v2.2dLaTeX symbol font definitions
 ***********

 )
Here is how much of TeX's memory you used:
 1909 strings out of 3926
 19413 string characters out of 65536
 72809 words of memory out of 169984
 4744 multiletter control sequences out of 10000
 14 hyphenation exceptions out of 3001
 35i,8n,23p,263b,265s stack positions out of 200i,100n,120p,5000b,4000s

Output written on sample1 (4 pages, 3183 bytes).
elapsed seconds = 13.10

Textures 2.0.2 (preloaded format=LaTeX R9P0 98.10.16)  20 JAN 1999 23:15
(sample2
LaTeX2e <1998/06/01>
(sample2.tex
File: sample2.tex 1999/01/22 5.0rc2 AAS markup document class
(aastex.cls
Document Class: aastex 1999/01/22 5.0rc2 AAS markup document class
Class aastex Info: Original \LaTeX2.09 style by Chris Biemesderfer (chris@ferber
ts.com). Adapted to \LaTeXe by A. Ogawa (ogawa@teleport.com) on input line 68.
(article.cls
Document Class: article 1998/05/05 v1.3y Standard LaTeX document class
(size12.clo
File: size12.clo 1998/05/05 v1.3y Standard LaTeX file (size option)
)
\c@part=\count79
\c@section=\count80
\c@subsection=\count81
\c@subsubsection=\count82
\c@paragraph=\count83
\c@subparagraph=\count84
\c@figure=\count85
\c@table=\count86
\abovecaptionskip=\skip41
\belowcaptionskip=\skip42
\bibindent=\dimen102
) (latexsym.sty
Package: latexsym 1996/11/20 v2.2d Standard LaTeX package (lasy symbols)
\symlasy=\mathgroup4
LaTeX Font Info:    Overwriting symbol font `lasy' in version `bold'
(Font)                  U/lasy/m/n --> U/lasy/b/n on input line 85.
) (graphicx.sty
Package: graphicx 1997/06/09 v1.0d Enhanced LaTeX Graphics (DPC,SPQR)
(keyval.sty
Package: keyval 1998/06/05 v1.11 key=value parser (DPC)
\KV@toks@=\toks14
) (graphics.sty
Package: graphics 1998/05/27 v1.0h Standard LaTeX Graphics (DPC,SPQR)
(trig.sty
Package: trig 1994/10/16 v1.08 sin cos tan (DPC)
) (graphics.cfg)
Package graphics Info: Driver file: textures.def on input line 85.
(textures.def
File: textures.def 1997/5/28 v0.4 Graphics extensions for Textures 1.7 and later
 (Arthur Ogawa/ogawa@teleport.com)
))
\Gin@req@height=\dimen103
\Gin@req@width=\dimen104
) (amsmath.sty
Package: amsmath 1997/03/20 v1.2d AMS math features
\@mathmargin=\skip43
(amstext.sty
Package: amstext 1996/10/28 v1.2b
(amsgen.sty
File: amsgen 1996/10/29 v1.2b
\@emptytoks=\toks15
\ex@=\dimen105
)) (amsbsy.sty
Package: amsbsy 1996/10/28 v1.2b
\pmbraise@=\dimen106
) (amsopn.sty
Package: amsopn 1996/10/28 v1.2b operator names
)
LaTeX Info: Redefining \frac on input line 188.
\uproot@=\count87
\leftroot@=\count88
\classnum@=\count89
\DOTSCASE@=\count90
LaTeX Info: Redefining \dots on input line 335.
LaTeX Info: Redefining \ldots on input line 418.
LaTeX Info: Redefining \cdots on input line 422.
\Mathstrutbox@=\box26
\strutbox@=\box27
\big@size=\dimen107
LaTeX Font Info:    Redeclaring font encoding OML on input line 500.
LaTeX Font Info:    Redeclaring font encoding OMS on input line 501.
\skewcharcount@=\count91
\familycount@=\count92
\pointcount@=\count93
\accentdimen@=\dimen108
\accentmu@=\count94
\minaw@=\dimen109
\c@MaxMatrixCols=\count95
\dotsspace@=\muskip10
\c@parentequation=\count96
\dspbrk@lvl=\count97
\tag@help=\toks16
\row@=\count98
\column@=\count99
\maxfields@=\count100
\andhelp@=\toks17
\eqnshift@=\dimen110
\alignsep@=\dimen111
\tagshift@=\dimen112
\tagwidth@=\dimen113
\totwidth@=\dimen114
\lineht@=\dimen115
\@envbody=\toks18
\multlinegap=\skip44
\multlinetaggap=\skip45
) (natbib.sty
Package: natbib 1998/07/14 6.8c (PWD)
\bibhang=\skip46
\bibsep=\skip47
LaTeX Info: Redefining \cite on input line 467.
\c@NAT@ctr=\count101
) (verbatim.sty
Package: verbatim 1997/04/30 v1.5k LaTeX2e package for verbatim enhancements
\every@verbatim=\toks19
\verbatim@line=\toks20
\verbatim@in@stream=\read1
)
\@slugcmmntwidth=\dimen116
\abstract@box=\box28
\c@editornote=\count102
\c@cureqno=\count103
\pt@box=\box29
\pt@width=\dimen117
\pt@line=\count104
\pt@column=\count105
\pt@nlines=\count106
\pt@ncol=\count107
\pt@page=\count108
\pt@tmpcapwidth=\dimen118
\c@plate=\count109
) (sample2.aux)
LaTeX Font Info:    Checking defaults for OML/cmm/m/it on input line 145.
LaTeX Font Info:    ... okay on input line 145.
LaTeX Font Info:    Checking defaults for T1/cmr/m/n on input line 145.
LaTeX Font Info:    ... okay on input line 145.
LaTeX Font Info:    Checking defaults for OT1/cmr/m/n on input line 145.
LaTeX Font Info:    ... okay on input line 145.
LaTeX Font Info:    Checking defaults for OMS/cmsy/m/n on input line 145.
LaTeX Font Info:    ... okay on input line 145.
LaTeX Font Info:    Checking defaults for OMX/cmex/m/n on input line 145.
LaTeX Font Info:    ... okay on input line 145.
LaTeX Font Info:    Checking defaults for U/cmr/m/n on input line 145.
LaTeX Font Info:    ... okay on input line 145.
AAS manuscript substyle, v5.0 <25 Sep 98>.
LaTeX Font Info:    Try loading font information for U+lasy on input line 150.
(ulasy.fd
File: ulasy.fd 1996/11/20 v2.2dLaTeX symbol font definitions
) [1

] [2

] [3] [4] [5] [6]
Overfull \hbox (29.01564pt too wide) in paragraph at lines 436--445
 [][] 
 []


Overfull \hbox (29.01564pt too wide) in paragraph at lines 457--457
 []$[]$ 
 []


Overfull \hbox (29.01564pt too wide) in paragraph at lines 457--457
 []$[]$ 
 []


Overfull \hbox (29.01564pt too wide) in paragraph at lines 457--457
 []$[]$ 
 []


Overfull \hbox (29.01564pt too wide) in paragraph at lines 457--457
 []$[]$ 
 []


Overfull \hbox (29.01564pt too wide) in paragraph at lines 457--457
 []$[]$ 
 []

Table 1 has been set to width 498.77063pt
[7]
Page 1 of table 2 has been set to width 424.88113pt with 23 lines per page 
[8] [9]
Page 2 of table 2 has been set to width 438.58932pt with 23 lines per page 
[10]
Page 3 of table 2 has been set to width 427.9186pt with 23 lines per page 
[11]
Page 4 of table 2 has been set to width 426.20332pt with 23 lines per page 
[12]
Page 5 of table 2 has been set to width 421.53383pt with 23 lines per page 
[13]
Table 2 has been set to width 438.2089pt
[14] [15








] [16

] (sample2.aux)

 *File List*
  aastex.cls    1999/01/22 5.0rc2 AAS markup document class
 article.cls    1998/05/05 v1.3y Standard LaTeX document class
  size12.clo    1998/05/05 v1.3y Standard LaTeX file (size option)
latexsym.sty    1996/11/20 v2.2d Standard LaTeX package (lasy symbols)
graphicx.sty    1997/06/09 v1.0d Enhanced LaTeX Graphics (DPC,SPQR)
  keyval.sty    1998/06/05 v1.11 key=value parser (DPC)
graphics.sty    1998/05/27 v1.0h Standard LaTeX Graphics (DPC,SPQR)
    trig.sty    1994/10/16 v1.08 sin cos tan (DPC)
graphics.cfg
textures.def    1997/5/28 v0.4 Graphics extensions for Textures 1.7 and later (A
rthur Ogawa/ogawa@teleport.com)
 amsmath.sty    1997/03/20 v1.2d AMS math features
 amstext.sty    1996/10/28 v1.2b
  amsgen.sty    
  amsbsy.sty    1996/10/28 v1.2b
  amsopn.sty    1996/10/28 v1.2b operator names
  natbib.sty    1998/07/14 6.8c (PWD)
verbatim.sty    1997/04/30 v1.5k LaTeX2e package for verbatim enhancements
   ulasy.fd    1996/11/20 v2.2dLaTeX symbol font definitions
 ***********

 )
Here is how much of TeX's memory you used:
 1937 strings out of 3926
 19626 string characters out of 65536
 84103 words of memory out of 169984
 4769 multiletter control sequences out of 10000
 14 hyphenation exceptions out of 3001
 35i,18n,24p,422b,491s stack positions out of 200i,100n,120p,5000b,4000s

Output written on sample2 (16 pages, 30435 bytes).
elapsed seconds = 17.47

Textures 2.0.2 (preloaded format=LaTeX R9P0 98.10.16)  20 JAN 1999 23:17
(samptbls
LaTeX2e <1998/06/01>
(samptbls.tex
File: samptbls.tex 1999/01/22 5.0rc2 AAS markup document class
(aastex.cls
Document Class: aastex 1999/01/22 5.0rc2 AAS markup document class
Class aastex Info: Original \LaTeX2.09 style by Chris Biemesderfer (chris@ferber
ts.com). Adapted to \LaTeXe by A. Ogawa (ogawa@teleport.com) on input line 68.
(article.cls
Document Class: article 1998/05/05 v1.3y Standard LaTeX document class
(size12.clo
File: size12.clo 1998/05/05 v1.3y Standard LaTeX file (size option)
)
\c@part=\count79
\c@section=\count80
\c@subsection=\count81
\c@subsubsection=\count82
\c@paragraph=\count83
\c@subparagraph=\count84
\c@figure=\count85
\c@table=\count86
\abovecaptionskip=\skip41
\belowcaptionskip=\skip42
\bibindent=\dimen102
) (latexsym.sty
Package: latexsym 1996/11/20 v2.2d Standard LaTeX package (lasy symbols)
\symlasy=\mathgroup4
LaTeX Font Info:    Overwriting symbol font `lasy' in version `bold'
(Font)                  U/lasy/m/n --> U/lasy/b/n on input line 85.
) (graphicx.sty
Package: graphicx 1997/06/09 v1.0d Enhanced LaTeX Graphics (DPC,SPQR)
(keyval.sty
Package: keyval 1998/06/05 v1.11 key=value parser (DPC)
\KV@toks@=\toks14
) (graphics.sty
Package: graphics 1998/05/27 v1.0h Standard LaTeX Graphics (DPC,SPQR)
(trig.sty
Package: trig 1994/10/16 v1.08 sin cos tan (DPC)
) (graphics.cfg)
Package graphics Info: Driver file: textures.def on input line 85.
(textures.def
File: textures.def 1997/5/28 v0.4 Graphics extensions for Textures 1.7 and later
 (Arthur Ogawa/ogawa@teleport.com)
))
\Gin@req@height=\dimen103
\Gin@req@width=\dimen104
) (amsmath.sty
Package: amsmath 1997/03/20 v1.2d AMS math features
\@mathmargin=\skip43
(amstext.sty
Package: amstext 1996/10/28 v1.2b
(amsgen.sty
File: amsgen 1996/10/29 v1.2b
\@emptytoks=\toks15
\ex@=\dimen105
)) (amsbsy.sty
Package: amsbsy 1996/10/28 v1.2b
\pmbraise@=\dimen106
) (amsopn.sty
Package: amsopn 1996/10/28 v1.2b operator names
)
LaTeX Info: Redefining \frac on input line 188.
\uproot@=\count87
\leftroot@=\count88
\classnum@=\count89
\DOTSCASE@=\count90
LaTeX Info: Redefining \dots on input line 335.
LaTeX Info: Redefining \ldots on input line 418.
LaTeX Info: Redefining \cdots on input line 422.
\Mathstrutbox@=\box26
\strutbox@=\box27
\big@size=\dimen107
LaTeX Font Info:    Redeclaring font encoding OML on input line 500.
LaTeX Font Info:    Redeclaring font encoding OMS on input line 501.
\skewcharcount@=\count91
\familycount@=\count92
\pointcount@=\count93
\accentdimen@=\dimen108
\accentmu@=\count94
\minaw@=\dimen109
\c@MaxMatrixCols=\count95
\dotsspace@=\muskip10
\c@parentequation=\count96
\dspbrk@lvl=\count97
\tag@help=\toks16
\row@=\count98
\column@=\count99
\maxfields@=\count100
\andhelp@=\toks17
\eqnshift@=\dimen110
\alignsep@=\dimen111
\tagshift@=\dimen112
\tagwidth@=\dimen113
\totwidth@=\dimen114
\lineht@=\dimen115
\@envbody=\toks18
\multlinegap=\skip44
\multlinetaggap=\skip45
) (natbib.sty
Package: natbib 1998/07/14 6.8c (PWD)
\bibhang=\skip46
\bibsep=\skip47
LaTeX Info: Redefining \cite on input line 467.
\c@NAT@ctr=\count101
) (verbatim.sty
Package: verbatim 1997/04/30 v1.5k LaTeX2e package for verbatim enhancements
\every@verbatim=\toks19
\verbatim@line=\toks20
\verbatim@in@stream=\read1
)
\@slugcmmntwidth=\dimen116
\abstract@box=\box28
\c@editornote=\count102
\c@cureqno=\count103
\pt@box=\box29
\pt@width=\dimen117
\pt@line=\count104
\pt@column=\count105
\pt@nlines=\count106
\pt@ncol=\count107
\pt@page=\count108
\pt@tmpcapwidth=\dimen118
\c@plate=\count109
) (samptbls.aux (table1.aux) (table2.aux) (table3.aux) (complext.aux))
LaTeX Font Info:    Checking defaults for OML/cmm/m/it on input line 70.
LaTeX Font Info:    ... okay on input line 70.
LaTeX Font Info:    Checking defaults for T1/cmr/m/n on input line 70.
LaTeX Font Info:    ... okay on input line 70.
LaTeX Font Info:    Checking defaults for OT1/cmr/m/n on input line 70.
LaTeX Font Info:    ... okay on input line 70.
LaTeX Font Info:    Checking defaults for OMS/cmsy/m/n on input line 70.
LaTeX Font Info:    ... okay on input line 70.
LaTeX Font Info:    Checking defaults for OMX/cmex/m/n on input line 70.
LaTeX Font Info:    ... okay on input line 70.
LaTeX Font Info:    Checking defaults for U/cmr/m/n on input line 70.
LaTeX Font Info:    ... okay on input line 70.
AAS manuscript substyle, v5.0 <25 Sep 98>.
(table1.tex
File: table1.tex 1999/01/22 5.0rc2 AAS markup document class
LaTeX Font Info:    Try loading font information for U+lasy on input line 73.
(ulasy.fd
File: ulasy.fd 1996/11/20 v2.2dLaTeX symbol font definitions
)
Overfull \hbox (28.88113pt too wide) in alignment at lines 73--94
[][][][][][][][][][][]
 []

Page 1 of table 1 has been set to width 396.0pt with 23 lines per page 
[1]
Overfull \hbox (42.58932pt too wide) in alignment at lines 94--114
[][][][][][][][][][][]
 []

Page 2 of table 1 has been set to width 396.0pt with 23 lines per page 
[2]

Class aastex Warning: Command \tablebreak is deprecated in aastex V 5.0. Using \
\ instead (please fix your document).  on input line 133.


Overfull \hbox (28.54768pt too wide) in alignment at lines 114--134
[][][][][][][][][][][]
 []

Page 3 of table 1 has been set to width 396.0pt with 23 lines per page 
[3]
Overfull \hbox (22.67274pt too wide) in alignment at lines 134--154
[][][][][][][][][][][]
 []

Page 4 of table 1 has been set to width 396.0pt with 23 lines per page 
[4]
Overfull \hbox (32.13792pt too wide) in alignment at lines 154--174
[][][][][][][][][][][]
 []

Page 5 of table 1 has been set to width 396.0pt with 23 lines per page 
[5]
Overfull \hbox (22.67274pt too wide) in alignment at lines 174--194
[][][][][][][][][][][]
 []

Page 6 of table 1 has been set to width 396.0pt with 23 lines per page 
[6]
Overfull \hbox (22.67274pt too wide) in alignment at lines 194--214
[][][][][][][][][][][]
 []

Page 7 of table 1 has been set to width 396.0pt with 23 lines per page 
[7]
Overfull \hbox (47.96555pt too wide) in alignment at lines 214--234
[][][][][][][][][][][]
 []

Page 8 of table 1 has been set to width 396.0pt with 23 lines per page 
[8]
Overfull \hbox (26.39375pt too wide) in alignment at lines 234--254
[][][][][][][][][][][]
 []

Page 9 of table 1 has been set to width 396.0pt with 23 lines per page 
[9]

Class aastex Warning: Command \tablebreak is deprecated in aastex V 5.0. Using \
\ instead (please fix your document).  on input line 255.


Overfull \hbox (30.20332pt too wide) in alignment at lines 254--274
[][][][][][][][][][][]
 []

Page 10 of table 1 has been set to width 396.0pt with 23 lines per page 
[10]
Overfull \hbox (35.85893pt too wide) in alignment at lines 274--294
[][][][][][][][][][][]
 []

Page 11 of table 1 has been set to width 396.0pt with 23 lines per page 
[11]
Overfull \hbox (30.20332pt too wide) in alignment at lines 294--314
[][][][][][][][][][][]
 []

Page 12 of table 1 has been set to width 396.0pt with 23 lines per page 
[12]
Overfull \hbox (51.67407pt too wide) in alignment at lines 314--345
[][][][][][][][][][][]
 []

Table 1 has been set to width 396.0pt
) [13] (table2.tex
File: table2.tex 1999/01/20 5.0rc2 AAS markup document class
Page 1 of table 2 has been set to width 358.29579pt with 23 lines per page 
[14]
Page 2 of table 2 has been set to width 358.29579pt with 23 lines per page 
[15]
Page 3 of table 2 has been set to width 366.29016pt with 23 lines per page 
[16]
Page 4 of table 2 has been set to width 350.48625pt with 23 lines per page 
[17]
Table 2 has been set to width 355.84999pt
) [18] (table3.tex
File: table3.tex 1999/01/20 5.0rc2 AAS markup document class

Underfull \hbox (badness 10000) in alignment at lines 71--92
[][][][][]
 []

Page 1 of table 3 has been set to width 324.0pt with 23 lines per page 
[19]
Underfull \hbox (badness 10000) in alignment at lines 92--112
[][][][][]
 []

Page 2 of table 3 has been set to width 324.0pt with 23 lines per page 
[20]
Underfull \hbox (badness 10000) in alignment at lines 112--132
[][][][][]
 []

Page 3 of table 3 has been set to width 324.0pt with 23 lines per page 
[21]
Underfull \hbox (badness 10000) in alignment at lines 132--152
[][][][][]
 []

Page 4 of table 3 has been set to width 324.0pt with 23 lines per page 
[22]
Underfull \hbox (badness 10000) in alignment at lines 152--166
[][][][][]
 []

Table 3 has been set to width 324.0pt
) [23] (complext.tex
File: complext.tex 1999/01/20 5.0rc2 AAS markup document class
Table 4 has been set to width 405.92818pt
) [24] (samptbls.aux (table1.aux) (table2.aux) (table3.aux) (complext.aux))

 *File List*
  aastex.cls    1999/01/20 5.0rc2 AAS markup document class
 article.cls    1998/05/05 v1.3y Standard LaTeX document class
  size12.clo    1998/05/05 v1.3y Standard LaTeX file (size option)
latexsym.sty    1996/11/20 v2.2d Standard LaTeX package (lasy symbols)
graphicx.sty    1997/06/09 v1.0d Enhanced LaTeX Graphics (DPC,SPQR)
  keyval.sty    1998/06/05 v1.11 key=value parser (DPC)
graphics.sty    1998/05/27 v1.0h Standard LaTeX Graphics (DPC,SPQR)
    trig.sty    1994/10/16 v1.08 sin cos tan (DPC)
graphics.cfg
textures.def    1997/5/28 v0.4 Graphics extensions for Textures 1.7 and later (A
rthur Ogawa/ogawa@teleport.com)
 amsmath.sty    1997/03/20 v1.2d AMS math features
 amstext.sty    1996/10/28 v1.2b
  amsgen.sty    
  amsbsy.sty    1996/10/28 v1.2b
  amsopn.sty    1996/10/28 v1.2b operator names
  natbib.sty    1998/07/14 6.8c (PWD)
verbatim.sty    1997/04/30 v1.5k LaTeX2e package for verbatim enhancements
  table1.tex    1999/01/20 5.0rc2 AAS markup document class
   ulasy.fd    1996/11/20 v2.2dLaTeX symbol font definitions
  table2.tex    1999/01/20 5.0rc2 AAS markup document class
  table3.tex    1999/01/20 5.0rc2 AAS markup document class
complext.tex    1999/01/20 5.0rc2 AAS markup document class
 ***********

 )
Here is how much of TeX's memory you used:
 1977 strings out of 3926
 19844 string characters out of 65536
 80327 words of memory out of 169984
 4746 multiletter control sequences out of 10000
 14 hyphenation exceptions out of 3001
 35i,12n,24p,264b,455s stack positions out of 200i,100n,120p,5000b,4000s

Output written on samptbls (24 pages, 52470 bytes).
elapsed seconds = 27.12
\par
}
